\renewcommand{\B}{\mathfrak B}
\setcounter{chapter}{0}
\chapter{概形的语言}
\src{79}
\begin{center}\textbf{本章提要}\end{center}
\begin{enumerate}[label=\S\,\arabic*.,itemsep=0.1em]
\item 仿射概形。
\item 预概形与预概形态射。
\item 预概形的积。
\item 子预概形与浸入态射。
\item 约化预概形；分离条件。
\item 有限性条件。
\item 有理映射。
\item Chevalley 概形。
\item 拟凝聚层的补充知识。
\item 形式概形。
\end{enumerate}

第 1—8 节主要是展开一种将在后文始终使用的语言。不过应当指出，依照本书的总体精神，
第 7、8 节的使用将少于其他各节，其作用也较不根本；事实上，讨论 Chevalley 概形，
只是为了与 Chevalley\cite{b1} 和 Nagata\cite{b9} 的语言建立联系。
第 9 节给出拟凝聚层的一些定义与结果，其中某些内容已不再只是把交换代数中已知的概念
翻译成“几何”语言，而具有整体性质；从以后的各章起，它们将成为态射整体研究中不可缺少的工具。
最后，第 10 节引入概形概念的一种推广；在第 III 章中，我们将借助它来方便地表述并证明
固有态射的上同调研究中的基本结果。另外还应指出，形式概形的概念似乎是表述
“模空间理论”（代数簇的分类问题）中某些事实所不可缺少的。\transnote{
此处原文为 \textit{théorie des modules}，作者紧接着用“代数簇的分类问题”作解释。
这里的 \textit{modules} 指分类参数意义上的模问题，而不是环上模的代数理论。}
第 10 节的结果在第 III 章 §3 之前不会使用，建议读者在读到那里以前暂且跳过这一节。
\par


% END chapters/ch1_intro.tex

\par

% BEGIN chapters/ch1_s1_1.tex
\section{仿射概形}\hypertarget{heading:1:1}{}
\subsection{环的素谱}\hypertarget{heading:1:1.1}{}
\src{80—83}

\egapar{1.1.1}
\textbf{记号。}设 $A$ 为一个（交换）环，$M$ 为一个 $A$-模。在本章及以后各章中，
我们将经常使用下列记号：

$\Spec(A)$ 是 $A$ 的素理想的集合，也称为 $A$ 的\emph{素谱}；
对于 $x\in X=\Spec(A)$，常常用 $\jj_x$ 代替 $x$ 更为方便。
$\Spec(A)$ 为空，当且仅当环 $A$ 为零环。

$A_x=A_{\jj_x}$ 是分母取自 $S=A\setminus\jj_x$ 的（局部）分式环 $S^{-1}A$。

$\mm_x=\jj_xA_{\jj_x}$ 是 $A_x$ 的极大理想。

$k(x)=A_x/\mm_x$ 是 $A_x$ 的剩余域，它与整环 $A/\jj_x$ 的分式域典范同构，
我们将两者等同。

对于 $f\in A$ 及 $x\in X$，$f(x)$ 是 $f$ 模 $\jj_x$ 的类，
视为 $A/\jj_x\subset k(x)$ 的元素。也称 $f(x)$ 为 $f$ 在点 $x\in\Spec(A)$ 的\emph{值}；
关系 $f(x)=0$ 与 $f\in\jj_x$ 等价。

$M_x=M\otimes_A A_x$ 是分母取自 $A\setminus\jj_x$ 的分式模。

$\rr(E)$ 是由 $A$ 的子集 $E$ 生成的理想的根。

对于 $E\subset A$，$V(E)$ 是满足 $E\subset\jj_x$ 的点 $x\in X$ 的集合，
也就是使每个 $f\in E$ 都满足 $f(x)=0$ 的点的集合。因此有
\begin{equation}
\rr(E)=\bigcap_{x\in V(E)}\jj_x.
\tag{1.1.1.1}\label{eq:1:1.1.1.1}
\end{equation}
对于 $f\in A$，记 $V(f)=V(\{f\})$。

$D(f)=X\setminus V(f)$ 是满足 $f(x)\ne0$ 的点 $x\in X$ 的集合。

\egastatement{1.1.2}{命题}
有下列性质：
\begin{enumerate}[label=(\roman*)]
\item $V(0)=X$，$V(1)=\varnothing$。
\item $E\subset E'$ 蕴含 $V(E)\supset V(E')$。
\item 对于 $A$ 的任意子集族 $(E_\lambda)$，有
\[
V\Bigl(\bigcup_\lambda E_\lambda\Bigr)
=V\Bigl(\sum_\lambda E_\lambda\Bigr)
=\bigcap_\lambda V(E_\lambda).
\]
\item $V(EE')=V(E)\cup V(E')$。
\item $V(E)=V(\rr(E))$。
\end{enumerate}

性质 (i)、(ii)、(iii) 是显然的，(v) 由 (ii) 及公式（\hyperref[eq:1:1.1.1.1]{1.1.1.1}）得到。
显然 $V(EE')\supset V(E)\cup V(E')$；反之，若 $x\notin V(E)$ 且 $x\notin V(E')$，
则存在 $f\in E$ 和 $f'\in E'$，使 $f(x)$ 和 $f'(x)$ 在 $k(x)$ 中均非零。
于是 $f(x)f'(x)\ne0$，即 $x\notin V(EE')$，这就证明了 (iv)。

命题（\hyperlink{para:1:1.1.2}{1.1.2}）特别表明，当 $E$ 遍历 $A$ 的所有子集时，形如 $V(E)$ 的集合
恰为 $X$ 上某个拓扑的闭集。我们称此拓扑为\emph{谱拓扑}。\origfootnote{
将这一拓扑引入代数几何的是 Zariski。因此它也常被称为 $X$ 的“Zariski 拓扑”。}
除非另作明确说明，始终假定 $X=\Spec(A)$ 赋予谱拓扑。

\egapar{1.1.3}
对于 $X$ 的任意子集 $Y$，以 $\jj(Y)$ 表示满足 $f(y)=0$（对所有 $y\in Y$）的
元素 $f\in A$ 所成的集合；等价地，$\jj(Y)$ 是 $y\in Y$ 对应的素理想 $\jj_y$ 的交。
显然，$Y\subset Y'$ 蕴含 $\jj(Y)\supset\jj(Y')$，且对于 $X$ 的任意子集族 $(Y_\lambda)$，有
\begin{equation}
\jj\Bigl(\bigcup_\lambda Y_\lambda\Bigr)
=\bigcap_\lambda\jj(Y_\lambda).
\tag{1.1.3.1}\label{eq:1:1.1.3.1}
\end{equation}
最后，有
\begin{equation}
\jj(\{x\})=\jj_x.
\tag{1.1.3.2}\label{eq:1:1.1.3.2}
\end{equation}

\egastatement{1.1.4}{命题}
(i) 对于 $A$ 的任意子集 $E$，有 $\jj(V(E))=\rr(E)$。

(ii) 对于 $X$ 的任意子集 $Y$，有 $V(\jj(Y))=\overline Y$，其中 $\overline Y$ 是 $Y$ 在 $X$ 中的闭包。

(i) 是定义和（\hyperref[eq:1:1.1.1.1]{1.1.1.1}）的直接推论。另一方面，$V(\jj(Y))$ 是包含 $Y$ 的闭集；
反之，若 $Y\subset V(E)$，则对于每个 $f\in E$ 及每个 $y\in Y$ 都有 $f(y)=0$，
因而 $E\subset\jj(Y)$，$V(E)\supset V(\jj(Y))$，这就证明了 (ii)。

\egastatement{1.1.5}{推论}
$X=\Spec(A)$ 的闭子集与 $A$ 中等于自身的根的理想（即素理想的交）之间，
通过两个反序映射 $Y\mapsto\jj(Y)$、$\mathfrak a\mapsto V(\mathfrak a)$ 建立一一对应。
两个闭子集之并 $Y_1\cup Y_2$ 对应于 $\jj(Y_1)\cap\jj(Y_2)$；
任意闭子集族 $(Y_\lambda)$ 的交对应于各 $\jj(Y_\lambda)$ 之和的根。

\egastatement{1.1.6}{推论}
若 $A$ 是诺特环，则 $X=\Spec(A)$ 是诺特空间。

应当注意，这一推论的逆命题不成立。例如，一个非诺特整环可以恰好只有一个非零素理想，
秩为 $1$ 的非离散赋值环就是这样的例子。

作为素谱不是诺特空间的环 $A$ 的例子，可以取无限紧空间 $Y$ 上的数值连续函数环
$\mathscr C(Y)$。我们知道，作为集合，$Y$ 可与 $A$ 的极大理想集合等同；
而且容易看出，$X=\Spec(A)$ 的拓扑在 $Y$ 上诱导的拓扑就是 $Y$ 原有的拓扑。
由于 $Y$ 不是诺特空间，$X$ 也不是诺特空间。

\egastatement{1.1.7}{推论}
对于每个 $x\in X$，$\{x\}$ 的闭包是满足 $\jj_x\subset\jj_y$ 的点 $y\in X$ 的集合。
$\{x\}$ 是闭集，当且仅当 $\jj_x$ 是极大理想。

\egastatement{1.1.8}{推论}
空间 $X=\Spec(A)$ 是 Kolmogoroff 空间。

事实上，若 $x,y$ 是 $X$ 的两个不同点，则或者 $\jj_x\not\subset\jj_y$，
或者 $\jj_y\not\subset\jj_x$，因而其中一个点不属于另一个点的闭包。

\egapar{1.1.9}
由命题（\hyperlink{para:1:1.1.2}{1.1.2}，(iv)），对于 $A$ 的两个元素 $f,g$，有
\begin{equation}
D(fg)=D(f)\cap D(g).
\tag{1.1.9.1}\label{eq:1:1.1.9.1}
\end{equation}
还应指出，由命题（\hyperlink{para:1:1.1.4}{1.1.4}，(i)）及命题（\hyperlink{para:1:1.1.2}{1.1.2}，(v)），
关系 $D(f)=D(g)$ 表示 $\rr(f)=\rr(g)$，也等价于包含 $(f)$ 与包含 $(g)$ 的极小素理想相同。
特别地，当 $f=ug$ 且 $u$ 可逆时，此关系成立。

\egastatement{1.1.10}{命题}
(i) 当 $f$ 遍历 $A$ 时，集合 $D(f)$ 构成 $X$ 的拓扑的一组基。

(ii) 对于每个 $f\in A$，$D(f)$ 都是拟紧的。特别地，$X=D(1)$ 是拟紧的。

(i) 设 $U$ 是 $X$ 的开集。依定义，$U=X\setminus V(E)$，其中 $E$ 是 $A$ 的某个子集。
而 $V(E)=\bigcap_{f\in E}V(f)$，故 $U=\bigcup_{f\in E}D(f)$。

(ii) 由 (i)，只需证明：若 $A$ 的元素族 $(f_\lambda)_{\lambda\in L}$ 满足
$D(f)\subset\bigcup_{\lambda\in L}D(f_\lambda)$，
则存在 $L$ 的有限子集 $J$，使 $D(f)\subset\bigcup_{\lambda\in J}D(f_\lambda)$。
设 $\mathfrak a$ 是由各 $f_\lambda$ 生成的理想；由假设有
$V(f)\supset V(\mathfrak a)$，从而 $\rr(f)\subset\rr(\mathfrak a)$。
由于 $f\in\rr(f)$，存在整数 $n>0$ 使 $f^n\in\mathfrak a$。
于是 $f^n$ 已属于某个有限子族 $(f_\lambda)_{\lambda\in J}$ 生成的理想 $\mathfrak b$，且
\[
V(f)=V(f^n)\supset V(\mathfrak b)
=\bigcap_{\lambda\in J}V(f_\lambda),
\]
即 $D(f)\subset\bigcup_{\lambda\in J}D(f_\lambda)$。

\egastatement{1.1.11}{命题}
对于 $A$ 的每个理想 $\mathfrak a$，$\Spec(A/\mathfrak a)$ 典范地等同于
$\Spec(A)$ 的闭子空间 $V(\mathfrak a)$。

事实上，我们知道，$A/\mathfrak a$ 的理想（分别地，素理想）与 $A$ 中包含 $\mathfrak a$ 的理想
（分别地，素理想）之间，存在保持包含序的典范一一对应。

回忆一下，$A$ 的幂零元集合 $\nrad$（即 $A$ 的幂零根）是等于 $\rr(0)$ 的理想，
也就是 $A$ 的所有素理想之交（0，\hyperlink{para:0:1.1.1}{1.1.1}）。

\egastatement{1.1.12}{推论}
拓扑空间 $\Spec(A)$ 与 $\Spec(A/\nrad)$ 典范同构。

\egastatement{1.1.13}{命题}
$X=\Spec(A)$ 不可约（0，\hyperlink{para:0:2.1.1}{2.1.1}），当且仅当环 $A/\nrad$ 是整环；
等价地，当且仅当理想 $\nrad$ 是素理想。

由推论（\hyperlink{para:1:1.1.12}{1.1.12}），可以只考虑 $\nrad=0$ 的情形。
若 $X$ 可约，则存在两个均不等于 $X$ 的闭子集 $Y_1,Y_2$，使 $X=Y_1\cup Y_2$，从而
$\jj(X)=\jj(Y_1)\cap\jj(Y_2)=0$，且 $\jj(Y_1)$ 与 $\jj(Y_2)$ 均不等于 $(0)$（\hyperlink{para:1:1.1.5}{1.1.5}）；
因此 $A$ 不是整环。反之，若 $A$ 中有两个非零元素 $f,g$，使 $fg=0$，则
$V(f)\ne X$，$V(g)\ne X$（因为 $A$ 的所有素理想之交为 $(0)$），而
$X=V(fg)=V(f)\cup V(g)$。

\egastatement{1.1.14}{推论}
(i) 在 $X=\Spec(A)$ 的闭子集与 $A$ 的根理想之间的一一对应中，
$X$ 的不可约闭子集对应于 $A$ 的素理想。特别地，$X$ 的不可约分支对应于 $A$ 的极小素理想。

(ii) 映射 $x\mapsto\overline{\{x\}}$ 在 $X$ 与 $X$ 的不可约闭子集的集合之间建立一一对应；
换言之，$X$ 的每个不可约闭子集恰有一个泛点。

(i) 直接由（\hyperlink{para:1:1.1.13}{1.1.13}）及（\hyperlink{para:1:1.1.11}{1.1.11}）得到。为证明 (ii)，由（\hyperlink{para:1:1.1.11}{1.1.11}），可以只考虑 $X$ 不可约的情形。
这时，由命题（\hyperlink{para:1:1.1.13}{1.1.13}），$A$ 中存在一个最小素理想 $\nrad$，它对应于 $X$ 的一个泛点。
此外，由于 $X$ 是 Kolmogoroff 空间，其泛点唯一（（\hyperlink{para:1:1.1.8}{1.1.8}）及（0，\hyperlink{para:0:2.1.3}{2.1.3}））。

\egastatement{1.1.15}{命题}
若 $\frJ$ 是 $A$ 的一个包含于根 $\frR(A)$ 中的理想，
则 $V(\frJ)$ 在 $X=\Spec(A)$ 中的唯一邻域是整个空间 $X$。

事实上，依定义，$A$ 的每个极大理想都属于 $V(\frJ)$。
又因为 $A$ 的每个真理想 $\mathfrak a$ 都包含于一个极大理想，
故 $V(\mathfrak a)\cap V(\frJ)\ne\varnothing$，从而得到本命题。\transnote{
原文此处写作“每个理想”；这里只涉及非空闭集 $V(\mathfrak a)$，
因而 $\mathfrak a\ne A$。译文将这一限制明确写出。
$\frR(A)$ 是第 0 章（\hyperlink{para:1:1.1.2}{1.1.2}）定义的根，不是幂零根 $\nrad$。}
\par


% END chapters/ch1_s1_1.tex

\par

% BEGIN chapters/ch1_s1_2.tex
\subsection{环的素谱的函子性质}\hypertarget{heading:1:1.2}{}
\src{83—84}

\egapar{1.2.1}
设 $A,A'$ 是两个环，$\varphi\colon A'\to A$ 是环同态。
对于每个素理想 $x=\jj_x\in\Spec(A)=X$，环 $A'/\varphi^{-1}(\jj_x)$
与 $A/\jj_x$ 的一个子环典范同构，因而是整环；也就是说，
$\varphi^{-1}(\jj_x)$ 是 $A'$ 的素理想。以 $\aphi(x)$ 表示这个素理想，便定义了映射
\[
\aphi\colon X=\Spec(A)\longrightarrow X'=\Spec(A'),
\]
也记作 $\Spec(\varphi)$，称为\emph{与同态 $\varphi$ 相伴的映射}。
以 $\varphi^x$ 表示由 $\varphi$ 取商得到的单同态
$A'/\varphi^{-1}(\jj_x)\to A/\jj_x$，也以同一记号表示它到分式域的典范延拓
\[
\varphi^x\colon k(\aphi(x))\longrightarrow k(x).
\]
因此，依定义，对于每个 $f'\in A'$，有
\begin{equation}
\varphi^x\bigl(f'(\aphi(x))\bigr)=\bigl(\varphi(f')\bigr)(x)
\qquad(x\in X).
\tag{1.2.1.1}\label{eq:1:1.2.1.1}
\end{equation}

\egastatement{1.2.2}{命题}
(i) 对于 $A'$ 的每个子集 $E'$，有
\begin{equation}
(\aphi)^{-1}(V(E'))=V(\varphi(E'));
\tag{1.2.2.1}\label{eq:1:1.2.2.1}
\end{equation}
特别地，对于每个 $f'\in A'$，有
\begin{equation}
(\aphi)^{-1}(D(f'))=D(\varphi(f')).
\tag{1.2.2.2}\label{eq:1:1.2.2.2}
\end{equation}
(ii) 对于 $A$ 的每个理想 $\mathfrak a$，有
\begin{equation}
\overline{\aphi(V(\mathfrak a))}=V(\varphi^{-1}(\mathfrak a)).
\tag{1.2.2.3}\label{eq:1:1.2.2.3}
\end{equation}

事实上，关系 $\aphi(x)\in V(E')$ 依定义等价于
$E'\subset\varphi^{-1}(\jj_x)$，也就等价于 $\varphi(E')\subset\jj_x$，
最终等价于 $x\in V(\varphi(E'))$，从而得到 (i)。
为证明 (ii)，可以假定 $\mathfrak a$ 等于自身的根，因为
$V(\rr(\mathfrak a))=V(\mathfrak a)$（\hyperlink{para:1:1.1.2}{1.1.2}，(v)），且
$\varphi^{-1}(\rr(\mathfrak a))=\rr(\varphi^{-1}(\mathfrak a))$。
令 $Y=V(\mathfrak a)$，$\mathfrak a'=\jj(\aphi(Y))$，则
$\overline{\aphi(Y)}=V(\mathfrak a')$（命题（\hyperlink{para:1:1.1.4}{1.1.4}，(ii)））。
关系 $f'\in\mathfrak a'$ 依定义等价于：对于每个 $x'\in\aphi(Y)$，有 $f'(x')=0$。
由公式（\hyperref[eq:1:1.2.1.1]{1.2.1.1}），这又等价于：对于每个 $x\in Y$，有 $\varphi(f')(x)=0$；
也等价于 $\varphi(f')\in\jj(Y)=\mathfrak a$，这里用了 $\mathfrak a$ 为根理想。
由此得到 (ii)。

\egastatement{1.2.3}{推论}
映射 $\aphi$ 连续。

注意，若 $A''$ 是第三个环，$\varphi'\colon A''\to A'$ 是同态，则有
\[
{}^a(\varphi\circ\varphi')=({}^a\varphi')\circ({}^a\varphi).
\]
\noindent 此结果与推论（\hyperlink{para:1:1.2.3}{1.2.3}）说明：$\Spec(A)$ 是关于 $A$ 的反变函子，
从环范畴取值于拓扑空间范畴。\transnote{
原版在此印作 ${}^a(\varphi'\circ\varphi)=({}^a\varphi)\circ({}^a\varphi')$。
按本段给定的 $\varphi\colon A'\to A$、$\varphi'\colon A''\to A'$，
两边的复合次序均应交换；译文采用与这些源、靶一致的式子。
同一关系在原文（\hyperlink{para:1:1.6.10}{1.6.10}）中按正确次序写出。}

\egastatement{1.2.4}{推论}
假设 $\varphi$ 满足：每个 $f\in A$ 都可写成 $f=h\varphi(f')$，其中 $h$ 在 $A$ 中可逆
（特别地，$\varphi$ 满射时满足此条件）。则 $\aphi$ 是 $X$ 到 $\aphi(X)$ 的同胚。

先证明，对于每个 $E\subset A$，存在 $A'$ 的子集 $E'$，使 $V(E)=V(\varphi(E'))$。
由公理 $(\mathrm T_0)$（\hyperlink{para:1:1.1.8}{1.1.8}）及公式（\hyperref[eq:1:1.2.2.1]{1.2.2.1}），这首先蕴含 $\aphi$ 单射；
再由同一公式可知 $\aphi$ 是同胚。
而对于每个 $f\in E$，只需取一个 $f'\in A'$，使 $h\varphi(f')=f$，其中 $h$ 在 $A$ 中可逆；
这些 $f'$ 所成的集合 $E'$ 就满足要求。

\egapar{1.2.5}
特别地，若 $\varphi$ 是由 $A$ 到商环 $A/\mathfrak a$ 的典范同态，
便重新得到（\hyperlink{para:1:1.1.12}{1.1.12}），\transnote{
此处引文编号照录原版。一般理想 $\mathfrak a$ 的相应结论在（\hyperlink{para:1:1.1.11}{1.1.11}）；
（\hyperlink{para:1:1.1.12}{1.1.12}）是 $\mathfrak a$ 为幂零根时的特例。}
而 $\aphi$ 就是把与 $\Spec(A/\mathfrak a)$ 等同的 $V(\mathfrak a)$
嵌入 $X=\Spec(A)$ 的典范包含映射。

（\hyperlink{para:1:1.2.4}{1.2.4}）的另一个特例给出：

\egastatement{1.2.6}{推论}
若 $S$ 是 $A$ 的一个乘法子集，则素谱 $\Spec(S^{-1}A)$ 连同其拓扑，
典范地等同于 $X=\Spec(A)$ 中由满足 $\jj_x\cap S=\varnothing$ 的点 $x$ 所成的子空间。

事实上，我们知道（0，\hyperlink{para:0:1.2.6}{1.2.6}），$S^{-1}A$ 的素理想恰是满足
$\jj_x\cap S=\varnothing$ 的理想 $S^{-1}\jj_x$，且有
$\jj_x=(i_A^S)^{-1}(S^{-1}\jj_x)$。
因此，只需将推论（\hyperlink{para:1:1.2.4}{1.2.4}）应用于 $i_A^S$。

\egastatement{1.2.7}{推论}
$\aphi(X)$ 在 $X'$ 中稠密，当且仅当核 $\Ker\varphi$ 的每个元素均为幂零元。

事实上，将公式（\hyperref[eq:1:1.2.2.3]{1.2.2.3}）应用于理想 $\mathfrak a=(0)$，得到
$\overline{\aphi(X)}=V(\Ker\varphi)$。
而 $V(\Ker\varphi)=X'$ 当且仅当 $\Ker\varphi$ 包含于 $A'$ 的每个素理想，
即包含于 $A'$ 的幂零根 $\nrad$。\transnote{
原版这一证明的末句将 $X'$、$A'$ 印作 $X$、$A$。
由于 $\Ker\varphi$ 是 $\varphi\colon A'\to A$ 的定义域 $A'$ 中的理想，
这里补出撇号，以与命题及（\hyperref[eq:1:1.2.2.3]{1.2.2.3}）一致。}
\par


% END chapters/ch1_s1_2.tex

\par

% BEGIN chapters/ch1_s1_3.tex
\subsection{与模相伴的层}\hypertarget{heading:1:1.3}{}
\src{84—90}

\egapar{1.3.1}
设 $A$ 是交换环，$M$ 是 $A$-模，$f\in A$，$S_f$ 是由各 $f^n$（$n\geq0$）组成的乘法子集。
回忆一下，我们记 $A_f=S_f^{-1}A$，$M_f=S_f^{-1}M$。
若 $S'_f$ 是由能够整除 $S_f$ 的某个元素的所有 $g\in A$ 所成的饱和乘法子集，
则我们知道，$A_f$ 与 $M_f$ 分别典范地等同于 $(S'_f)^{-1}A$ 与 $(S'_f)^{-1}M$（0，\hyperlink{para:0:1.4.3}{1.4.3}）。

\egastatement{1.3.2}{引理}
下列条件等价：
\begin{enumerate}[label=\alph*)]
\item $g\in S'_f$；
\item $S'_g\subset S'_f$；
\item $f\in\rr(g)$；
\item $\rr(f)\subset\rr(g)$；
\item $V(g)\subset V(f)$；
\item $D(f)\subset D(g)$。
\end{enumerate}
这直接由定义及（\hyperlink{para:1:1.1.5}{1.1.5}）得到。

\egapar{1.3.3}
若 $D(f)=D(g)$，则由引理（\hyperlink{para:1:1.3.2}{1.3.2}，b)）可知 $M_f=M_g$。
更一般地，若 $D(f)\supset D(g)$，即 $S'_f\subset S'_g$，
则由（0，\hyperlink{para:0:1.4.1}{1.4.1}）知存在典范的函子性同态
\[
\rho_{g,f}\colon M_f\longrightarrow M_g.
\]
若 $D(f)\supset D(g)\supset D(h)$，则有（0，\hyperlink{para:0:1.4.4}{1.4.4}）
\begin{equation}
\rho_{h,g}\circ\rho_{g,f}=\rho_{h,f}.
\tag{1.3.3.1}\label{eq:1:1.3.3.1}
\end{equation}

对于给定的 $x\in X=\Spec(A)$，当 $f$ 遍历 $A\setminus\jj_x$ 时，
各 $S'_f$ 构成 $A\setminus\jj_x$ 的一个递增滤过子集族：
因为对 $A\setminus\jj_x$ 的任意两个元素 $f,g$，$S'_f$ 与 $S'_g$ 都包含于 $S'_{fg}$。
又因为这些 $S'_f$ 的并为 $A\setminus\jj_x$，由（0，\hyperlink{para:0:1.4.5}{1.4.5}）可知，
$A_x$-模 $M_x$ 与相对于同态族 $(\rho_{g,f})$ 的归纳极限 $\ilim M_f$ 典范地等同。
以
\[
\rho_{x,f}\colon M_f\longrightarrow M_x
\]
表示 $f\in A\setminus\jj_x$（等价地，$x\in D(f)$）时的典范同态。

\egastatement{1.3.4}{定义}
在 $X=\Spec(A)$ 的由各 $D(f)$（$f\in A$）组成的开集基 $\B$ 上，
考虑预层 $D(f)\mapsto A_f$（分别地，$D(f)\mapsto M_f$）。
与此预层相伴的环层（分别地，$\wt A$-模层）称为素谱 $X$ 的\emph{结构层}
（分别地，\emph{与 $A$-模 $M$ 相伴的层}），
记作 $\wt A$ 或 $\OO_X$（分别地，$\wt M$）
（（\hyperlink{para:1:1.1.10}{1.1.10}），（0，\hyperlink{para:0:3.2.1}{3.2.1} 及 \hyperlink{para:0:3.5.6}{3.5.6}））。

由（0，\hyperlink{para:0:3.2.4}{3.2.4}）可知，茎 $\wt A_x$（分别地，$\wt M_x$）
与环 $A_x$（分别地，$A_x$-模 $M_x$）等同。以
\[
\theta_f\colon A_f\longrightarrow\Gamma(D(f),\wt A),
\qquad
\theta_f\colon M_f\longrightarrow\Gamma(D(f),\wt M)
\]
分别表示典范映射。因而，对于每个 $x\in D(f)$ 及每个 $z\in M_f$，有
\begin{equation}
\bigl(\theta_f(z)\bigr)_x=\rho_{x,f}(z).
\tag{1.3.4.1}\label{eq:1:1.3.4.1}
\end{equation}

\egastatement{1.3.5}{命题}
$\wt M$ 是关于 $M$ 的正合协变函子，从 $A$-模范畴取值于 $\wt A$-模层范畴。

事实上，设 $M,N$ 是两个 $A$-模，$u\colon M\to N$ 是同态。
对于每个 $f\in A$，$u$ 典范地给出一个 $A_f$-模同态 $u_f\colon M_f\to N_f$，
且当 $D(g)\subset D(f)$ 时，下图交换（0，\hyperlink{para:0:1.4.1}{1.4.1}）：
\[
\begin{tikzcd}[column sep=large,row sep=large]
M_f\arrow[r,"u_f"]\arrow[d,"\rho_{g,f}"'] & N_f\arrow[d,"\rho_{g,f}"]\\
M_g\arrow[r,"u_g"'] & N_g.
\end{tikzcd}
\]
这些同态因而定义一个 $\wt A$-模层同态 $\wt u\colon\wt M\to\wt N$（0，\hyperlink{para:0:3.2.3}{3.2.3}）。
此外，对于每个 $x\in X$，$\wt u_x$ 是各 $u_f$ 沿满足 $x\in D(f)$ 的 $f\in A$ 的归纳极限。
因而，由（0，\hyperlink{para:0:1.4.5}{1.4.5}），当将 $\wt M_x$、$\wt N_x$ 分别典范地等同于 $M_x$、$N_x$ 时，
$\wt u_x$ 就与由 $u$ 典范得到的同态 $u_x$ 等同。
若 $P$ 是第三个 $A$-模，$v\colon N\to P$ 为同态，且 $w=v\circ u$，
则显然 $\wt w_x=\wt v_x\circ\wt u_x$，故 $\wt w=\wt v\circ\wt u$。
因此，确实定义了所述协变函子。
这一函子正合，因为对于每个 $x\in X$，$M_x$ 是关于 $M$ 的正合函子（0，\hyperlink{para:0:1.3.2}{1.3.2}）。
另外，由两边的定义（0，\hyperlink{para:0:1.7.1}{1.7.1} 及 \hyperlink{para:0:3.1.6}{3.1.6}），有 $\Supp(\wt M)=\Supp(M)$。

\egastatement{1.3.6}{命题}
对于每个 $f\in A$，开集 $D(f)\subset X$ 典范地等同于素谱 $\Spec(A_f)$；
与 $A_f$-模 $M_f$ 相伴的层 $\wt{M_f}$ 典范地等同于限制 $\wt M|D(f)$。

第一项断言是（\hyperlink{para:1:1.2.6}{1.2.6}）的特例。又若 $g\in A$ 满足 $D(g)\subset D(f)$，
则 $M_g$ 典范地等同于 $M_f$ 的分式模，其分母是 $g$ 在 $A_f$ 中的典范像的各次幂（0，\hyperlink{para:0:1.4.6}{1.4.6}）。
由定义即得 $\wt{M_f}$ 与 $\wt M|D(f)$ 的典范等同。

\egastatement{1.3.7}{定理}
对于每个 $A$-模 $M$ 及每个 $f\in A$，同态
\[
\theta_f\colon M_f\longrightarrow\Gamma(D(f),\wt M)
\]
是双射；换言之，预层 $D(f)\mapsto M_f$ 是一个层。
特别地，$M$ 通过 $\theta_1$ 与 $\Gamma(X,\wt M)$ 等同。

注意，当 $M=A$ 时，$\theta_f$ 是环结构的同态。
因此，定理（\hyperlink{para:1:1.3.7}{1.3.7}）将蕴含：当借助 $\theta_f$ 将环 $A_f$ 与
$\Gamma(D(f),\wt A)$ 等同时，同态 $\theta_f\colon M_f\to\Gamma(D(f),\wt M)$ 是模同构。

先证明 $\theta_f$ 单射。若 $z\in M_f$ 满足 $\theta_f(z)=0$，
这表示对于 $A_f$ 的每个素理想 $\mathfrak p$，存在 $h\notin\mathfrak p$，使 $hz=0$。
于是 $z$ 的零化理想不包含于 $A_f$ 的任何素理想，因而是整个 $A_f$，故 $z=0$。

还需证明 $\theta_f$ 满射。可以归结到 $f=1$ 的情形，
一般情形借助（\hyperlink{para:1:1.3.6}{1.3.6}）作“局部化”即可得到。
于是设 $s$ 是 $\wt M$ 在 $X$ 上的截面；由（\hyperlink{para:1:1.3.4}{1.3.4}）及（\hyperlink{para:1:1.1.10}{1.1.10}，(ii)），
存在 $X$ 的一个有限覆盖 $(D(f_i))_{i\in I}$（$f_i\in A$），
使每个限制 $s_i=s|D(f_i)$ 都可写成 $\theta_{f_i}(t_i)$，其中 $t_i\in M_{f_i}$。
若 $i,j\in I$，将 $s_i$ 与 $s_j$ 在 $D(f_i)\cap D(f_j)=D(f_if_j)$ 上的限制相等写出，
由 $\wt M$ 的定义得到
\begin{equation}
\rho_{f_if_j,f_i}(t_i)=\rho_{f_if_j,f_j}(t_j).
\tag{1.3.7.1}\label{eq:1:1.3.7.1}
\end{equation}
依定义，对于每个 $i\in I$ 可写 $t_i=z_i/f_i^{n_i}$，其中 $z_i\in M$。
由于 $I$ 有限，在将各 $z_i$ 乘以 $f_i$ 的适当次幂之后，可以假定所有 $n_i$ 都等于同一个正整数 $n$。
这时，（\hyperref[eq:1:1.3.7.1]{1.3.7.1}）表示存在整数 $m_{ij}\geq0$，使
\[
(f_if_j)^{m_{ij}}(f_j^nz_i-f_i^nz_j)=0.
\]
还可假定所有 $m_{ij}$ 都等于同一个整数 $m$。
再将 $z_i$ 换成 $f_i^mz_i$，便归结到 $m=0$ 的情形，\transnote{
这里分母指数也相应从 $n$ 改为 $n+m$，之后仍以 $n$ 表示新的共同指数。
这是保持 $t_i$ 不变的同一分数表示的调整。}
也就是对于任意 $i,j$ 都有
\begin{equation}
f_j^nz_i=f_i^nz_j.
\tag{1.3.7.2}\label{eq:1:1.3.7.2}
\end{equation}
这里 $D(f_i^n)=D(f_i)$，且各 $D(f_i)$ 覆盖 $X$，
故各 $f_i^n$ 生成的理想为 $A$；换言之，存在 $g_i\in A$，使 $\sum_i g_if_i^n=1$。
考虑 $M$ 中的元素 $z=\sum_i g_iz_i$。由（\hyperref[eq:1:1.3.7.2]{1.3.7.2}），有
\[
f_i^nz=\sum_j g_jf_i^nz_j
=\Bigl(\sum_j g_jf_j^n\Bigr)z_i=z_i,
\]
因而依定义，$t_i=z/1$ 在 $M_{f_i}$ 中成立。
所以 $s_i$ 是 $\theta_1(z)$ 在 $D(f_i)$ 上的限制，从而 $s=\theta_1(z)$，证明完毕。

\egastatement{1.3.8}{推论}
设 $M,N$ 是两个 $A$-模；典范同态
\[
u\longmapsto\wt u\colon
\Hom_A(M,N)\longrightarrow\Hom_{\wt A}(\wt M,\wt N)
\]
是双射。特别地，$\wt M=0$ 与 $M=0$ 等价。

\begin{sloppypar}
事实上，考虑典范同态 $v\mapsto\Gamma(v)$，它将
$\Hom_{\wt A}(\wt M,\wt N)$ 映到
$\Hom_{\Gamma(\wt A)}(\Gamma(\wt M),\Gamma(\wt N))$；
由定理（\hyperlink{para:1:1.3.7}{1.3.7}），后一个模典范地等同于 $\Hom_A(M,N)$。
只需验证 $u\mapsto\wt u$ 与 $v\mapsto\Gamma(v)$ 互逆。
由 $\wt u$ 的定义，显然 $\Gamma(\wt u)=u$。
另一方面，对于 $v\in\Hom_{\wt A}(\wt M,\wt N)$，令 $u=\Gamma(v)$；
由 $v$ 典范得到的映射
$w\colon\Gamma(D(f),\wt M)\to\Gamma(D(f),\wt N)$ 使下图交换：
\end{sloppypar}
\[
\begin{tikzcd}[column sep=large,row sep=large]
M\arrow[r,"u"]\arrow[d,"\rho_{f,1}"'] & N\arrow[d,"\rho_{f,1}"]\\
M_f\arrow[r,"w"'] & N_f.
\end{tikzcd}
\]
因此，对于每个 $f\in A$，必有 $w=u_f$（0，\hyperlink{para:0:1.2.4}{1.2.4}），这就说明 $\wt{\Gamma(v)}=v$。

\egastatement{1.3.9}{推论}
(i) 设 $u\colon M\to N$ 是 $A$-模同态；
则与 $\Ker u$、$\Im u$、$\Coker u$ 相伴的层分别为
$\Ker\wt u$、$\Im\wt u$、$\Coker\wt u$。
特别地，$\wt u$ 单射（分别地，满射、双射），当且仅当 $u$ 具有相应性质。

(ii) 若 $M$ 是某个 $A$-模族 $(M_\lambda)$ 的归纳极限（分别地，直和），
则 $\wt M$ 在一个典范同构的意义下，是层族 $(\wt{M_\lambda})$ 的归纳极限（分别地，直和）。

(i) 只需将 $M\mapsto\wt M$ 的正合性（\hyperlink{para:1:1.3.5}{1.3.5}）应用于两个 $A$-模正合列
\[
0\longrightarrow\Ker u\longrightarrow M\longrightarrow\Im u\longrightarrow0,
\]
\[
0\longrightarrow\Im u\longrightarrow N\longrightarrow\Coker u\longrightarrow0.
\]
第二项断言随后由定理（\hyperlink{para:1:1.3.7}{1.3.7}）得到。

(ii) 设 $(M_\lambda,g_{\mu\lambda})$ 是以 $M$ 为归纳极限的 $A$-模归纳系，
$g_\lambda\colon M_\lambda\to M$ 为典范同态。
由于对 $\lambda\leq\mu\leq\nu$ 有
$\wt g_\mu\circ\wt g_{\mu\lambda}=\wt g_\lambda$ 及
$\wt g_{\nu\lambda}=\wt g_{\nu\mu}\circ\wt g_{\mu\lambda}$，
故 $(\wt{M_\lambda},\wt g_{\mu\lambda})$ 是 $X$ 上的层的归纳系。
若以 $h_\lambda\colon\wt{M_\lambda}\to\ilim\wt{M_\lambda}$ 表示典范同态，
则存在唯一同态 $v\colon\ilim\wt{M_\lambda}\to\wt M$，使 $v\circ h_\lambda=\wt g_\lambda$。
为证明 $v$ 双射，只需验证对每个 $x\in X$，$v_x$ 是
$(\ilim\wt{M_\lambda})_x$ 到 $\wt M_x$ 的双射。
而 $\wt M_x=M_x$，并且
\[
\Bigl(\ilim\wt{M_\lambda}\Bigr)_x
=\ilim(\wt{M_\lambda})_x
=\ilim(M_\lambda)_x=M_x
\qquad\text{（0，\hyperlink{para:0:1.3.3}{1.3.3}）}.
\]
另一方面，依定义，$(\wt g_\lambda)_x$ 与 $(h_\lambda)_x$ 都等于
由 $(M_\lambda)_x$ 到 $M_x$ 的典范映射。
由 $(\wt g_\lambda)_x=v_x\circ(h_\lambda)_x$ 可知 $v_x$ 是恒等映射。

最后，若 $M$ 是两个 $A$-模 $N,P$ 的直和，则显然 $\wt M=\wt N\oplus\wt P$。
任意直和都是有限直和的归纳极限，故 (ii) 的断言均已证明。

注意，（\hyperlink{para:1:1.3.8}{1.3.8}）证明了：与各 $A$-模的相伴层同构的层组成一个阿贝尔范畴（T，I，1.4）。
还应指出，由（\hyperlink{para:1:1.3.9}{1.3.9}），若 $M$ 是有限型 $A$-模，即存在满同态 $A^n\to M$，
则存在满同态 $\wt A^n\to\wt M$；也就是说，$\wt A$-模层 $\wt M$
由其在 $X$ 上的某个有限截面族生成（0，\hyperlink{para:0:5.1.1}{5.1.1}）；反之亦然。

\egapar{1.3.10}
若 $N$ 是 $A$-模 $M$ 的子模，则由（\hyperlink{para:1:1.3.9}{1.3.9}），典范包含 $j\colon N\to M$
给出单同态 $\wt N\to\wt M$，从而可以将 $\wt N$ 典范地等同于 $\wt M$ 的一个
$\wt A$-子模层；我们始终作此等同。
若 $N,P$ 是 $M$ 的两个子模，则
\begin{equation}
\wt{N+P}=\wt N+\wt P,
\tag{1.3.10.1}\label{eq:1:1.3.10.1}
\end{equation}
\begin{equation}
\wt{N\cap P}=\wt N\cap\wt P.
\tag{1.3.10.2}\label{eq:1:1.3.10.2}
\end{equation}
因为 $N+P$ 是典范同态 $N\oplus P\to M$ 的像，
$N\cap P$ 是典范同态 $M\to(M/N)\oplus(M/P)$ 的核，所以只需应用（\hyperlink{para:1:1.3.9}{1.3.9}）。
由（\hyperref[eq:1:1.3.10.1]{1.3.10.1}）与（\hyperref[eq:1:1.3.10.2]{1.3.10.2}）可知，若 $\wt N=\wt P$，则 $N=P$。

\egastatement{1.3.11}{推论}
在与 $A$-模的相伴层同构的层所组成的范畴中，函子 $\Gamma$ 正合。

事实上，设 $\wt M\xrightarrow{\wt u}\wt N\xrightarrow{\wt v}\wt P$ 是正合列，
对应于两个 $A$-模同态 $u\colon M\to N$、$v\colon N\to P$。
若 $Q=\Im u$，$R=\Ker v$，则
$\wt Q=\Im\wt u=\Ker\wt v=\wt R$（推论（\hyperlink{para:1:1.3.9}{1.3.9}）），故 $Q=R$。

\egastatement{1.3.12}{推论}
设 $M,N$ 是两个 $A$-模。

(i) 与 $M\otimes_A N$ 相伴的层典范地等同于 $\wt M\otimes_{\wt A}\wt N$。

(ii) 若 $M$ 还有有限表示，则与 $\Hom_A(M,N)$ 相伴的层和
$\shHom_{\wt A}(\wt M,\wt N)$ 典范同构。

(i) 层 $\G=\wt M\otimes_{\wt A}\wt N$ 与预层
\[
U\longmapsto\G(U)=\Gamma(U,\wt M)\otimes_{\Gamma(U,\wt A)}\Gamma(U,\wt N)
\]
相伴，其中 $U$ 遍历由各 $D(f)$（$f\in A$）组成的 $X$ 的基（\hyperlink{para:1:1.1.10}{1.1.10}，(i)）。
由（\hyperlink{para:1:1.3.7}{1.3.7}）和（\hyperlink{para:1:1.3.6}{1.3.6}），$\G(D(f))$ 典范地等同于 $M_f\otimes_{A_f}N_f$。
另一方面，此 $A_f$-模与 $(M\otimes_A N)_f$ 典范同构（0，\hyperlink{para:0:1.3.4}{1.3.4}），
而后者又典范地等同于 $\Gamma(D(f),\wt{M\otimes_A N})$（\hyperlink{para:1:1.3.7}{1.3.7} 及 \hyperlink{para:1:1.3.6}{1.3.6}）。
此外，容易验证所得典范同构
\[
\G(D(f))\xrightarrow{\sim}\Gamma(D(f),\wt{M\otimes_A N})
\]
满足与限制算子的相容条件（0，\hyperlink{para:0:1.4.2}{1.4.2}），因而定义一个典范的函子性同构
\[
\wt M\otimes_{\wt A}\wt N\xrightarrow{\sim}\wt{M\otimes_A N}.
\]

(ii) 层 $\G=\shHom_{\wt A}(\wt M,\wt N)$ 与预层
\[
U\longmapsto\G(U)=\Hom_{\wt A|U}(\wt M|U,\wt N|U)
\]
\begin{sloppypar}
相伴，其中 $U$ 遍历由各 $D(f)$ 组成的 $X$ 的基。
而 $\G(D(f))$ 与 $\Hom_{A_f}(M_f,N_f)$ 典范地等同（\hyperlink{para:1:1.3.6}{1.3.6} 及 \hyperlink{para:1:1.3.8}{1.3.8}）；
由关于 $M$ 的假设，后者又典范地等同于 $(\Hom_A(M,N))_f$（0，\hyperlink{para:0:1.3.5}{1.3.5}）。
最后，$(\Hom_A(M,N))_f$ 典范地等同于
$\Gamma(D(f),\wt{\Hom_A(M,N)})$（\hyperlink{para:1:1.3.6}{1.3.6} 及 \hyperlink{para:1:1.3.7}{1.3.7}）。
所得典范同构与限制算子相容（0，\hyperlink{para:0:1.4.2}{1.4.2}），因此定义了典范同构
\end{sloppypar}
\[
\shHom_{\wt A}(\wt M,\wt N)\xrightarrow{\sim}\wt{\Hom_A(M,N)}.
\]

\egapar{1.3.13}
现在设 $B$ 是一个（交换）$A$-代数。
这可解释为：$B$ 是一个 $A$-模，并且给定了元素 $e\in B$ 及
$A$-同态 $\varphi\colon B\otimes_A B\to B$，使下列条件成立：
第一，下面两图交换，其中 $\sigma$ 是典范对称同构：
\[
\begin{tikzcd}[column sep=large,row sep=large]
B\otimes_A B\otimes_A B\arrow[r,"\varphi\otimes1"]\arrow[d,"1\otimes\varphi"']
 & B\otimes_A B\arrow[d,"\varphi"]\\
B\otimes_A B\arrow[r,"\varphi"'] & B;
\end{tikzcd}
\qquad
\begin{tikzcd}[column sep=small,row sep=large]
B\otimes_A B\arrow[rr,"\sigma"]\arrow[dr,"\varphi"'] && B\otimes_A B\arrow[dl,"\varphi"]\\
& B &
\end{tikzcd}
\]
第二，$\varphi(e\otimes x)=\varphi(x\otimes e)=x$。
由（\hyperlink{para:1:1.3.12}{1.3.12}），$\wt A$-模层同态
$\wt\varphi\colon\wt B\otimes_{\wt A}\wt B\to\wt B$ 满足相应条件，
因而在 $\wt B$ 上定义一个 $\wt A$-代数层结构。

同样，给定一个 $B$-模 $N$，等价于给定一个 $A$-模 $N$ 及一个
$A$-同态 $\psi\colon B\otimes_A N\to N$，使下图交换，且 $\psi(e\otimes n)=n$：
\[
\begin{tikzcd}[column sep=large,row sep=large]
B\otimes_A B\otimes_A N\arrow[r,"\varphi\otimes1"]\arrow[d,"1\otimes\psi"']
 & B\otimes_A N\arrow[d,"\psi"]\\
B\otimes_A N\arrow[r,"\psi"'] & N.
\end{tikzcd}
\]
同态 $\wt\psi\colon\wt B\otimes_{\wt A}\wt N\to\wt N$ 满足相应条件，
因而在 $\wt N$ 上定义一个 $\wt B$-模层结构。

同理，若 $u\colon B\to B'$（分别地，$v\colon N\to N'$）是
$A$-代数同态（分别地，$B$-模同态），则 $\wt u$（分别地，$\wt v$）是
$\wt A$-代数层同态（分别地，$\wt B$-模层同态），
$\Ker\wt u$ 是 $\wt B$-理想层（分别地，$\Ker\wt v$、$\Coker\wt v$ 和 $\Im\wt v$
都是 $\wt B$-模层）。若 $N$ 是 $B$-模，则 $\wt N$ 是有限型 $\wt B$-模层，
当且仅当 $N$ 是有限型 $B$-模（0，\hyperlink{para:0:5.2.3}{5.2.3}）。

若 $M,N$ 是两个 $B$-模，则 $\wt B$-模层 $\wt M\otimes_{\wt B}\wt N$
典范地等同于 $\wt{M\otimes_B N}$；同样，当 $M$ 有有限表示时，
$\shHom_{\wt B}(\wt M,\wt N)$ 典范地等同于 $\wt{\Hom_B(M,N)}$。
证明与（\hyperlink{para:1:1.3.12}{1.3.12}）相同。

若 $\frJ$ 是 $B$ 的理想，$N$ 是 $B$-模，则 $\wt{\frJ N}=\wt{\frJ}\,\wt N$。

最后，若 $B$ 是由 $A$-子模 $B_n$（$n\in\mathbb Z$）给出分次的 $A$-代数，
则 $\wt A$-代数层 $\wt B$ 是各 $\wt A$-模层 $\wt{B_n}$ 的直和（\hyperlink{para:1:1.3.9}{1.3.9}），
并由这些 $\wt A$-子模层给出分次；分次公理表示同态
$B_m\otimes_A B_n\to B$ 的像包含于 $B_{m+n}$。
同样，若 $M$ 是由子模 $M_n$ 给出分次的 $B$-模，
则 $\wt M$ 是由各 $\wt{M_n}$ 给出分次的 $\wt B$-模层。

\egapar{1.3.14}
若 $B$ 是 $A$-代数，$M$ 是 $B$ 的子模，则由 $\wt M$ 生成的
$\wt B$ 的 $\wt A$-子代数层（0，\hyperlink{para:0:4.1.3}{4.1.3}）就是 $\wt C$，
其中 $C$ 是由 $M$ 生成的 $B$ 的子代数。
事实上，$C$ 是同态 $\bigotimes_A^n M\to B$（$n\geq0$）的像所成的各子模之和，
因此只需应用（\hyperlink{para:1:1.3.9}{1.3.9}）和（\hyperlink{para:1:1.3.12}{1.3.12}）。
\par


% END chapters/ch1_s1_3.tex

\par

% BEGIN chapters/ch1_s1_4.tex
\subsection{素谱上的拟凝聚层}\hypertarget{heading:1:1.4}{}
\src{90—92}

\egastatement{1.4.1}{定理}
设 $X$ 是环 $A$ 的素谱，$V$ 是 $X$ 的拟紧开子集，$\F$ 是一个 $(\OO_X|V)$-模层。
下列四个条件等价：
\begin{enumerate}[label=\alph*)]
\item 存在 $A$-模 $M$，使 $\F$ 同构于 $\wt M|V$。
\item 存在 $V$ 的有限开覆盖 $(V_i)$，其中各 $V_i$ 均为包含于 $V$ 的
$D(f_i)$（$f_i\in A$），并且对每个 $i$，$\F|V_i$ 同构于形如 $\wt{M_{f_i}}$ 的层，
其中 $M_{f_i}$ 是一个 $A_{f_i}$-模。
\item 层 $\F$ 拟凝聚（0，\hyperlink{para:0:5.1.3}{5.1.3}）。
\item 下列两个性质成立：
\begin{enumerate}[label={$d\arabic*)$},leftmargin=2.7em]
\item 对于每个满足 $D(f)\subset V$ 的 $f\in A$，及每个截面
$s\in\Gamma(D(f),\F)$，存在整数 $n\geq0$，使 $f^ns$ 可延拓为 $\F$ 在 $V$ 上的一个截面。
\item 对于每个满足 $D(f)\subset V$ 的 $f\in A$，及每个在 $D(f)$ 上的限制为 $0$ 的截面
$t\in\Gamma(V,\F)$，存在整数 $n\geq0$，使 $f^nt=0$。
\end{enumerate}
\end{enumerate}

在条件 $d1)$ 和 $d2)$ 的写法中，已依（\hyperlink{para:1:1.3.7}{1.3.7}）默认将 $A$ 与 $\Gamma(\wt A)$ 等同。

$a)\Rightarrow b)$ 直接由（\hyperlink{para:1:1.3.6}{1.3.6}）及各 $D(f_i)$ 构成 $X$ 的拓扑的一组基（\hyperlink{para:1:1.1.10}{1.1.10}）得到。
由于每个 $A$-模都同构于某个同态 $A^{(I)}\to A^{(J)}$ 的余核，
（\hyperlink{para:1:1.3.9}{1.3.9}）说明与每个 $A$-模相伴的层都是拟凝聚的；故 $b)\Rightarrow c)$。
反之，若 $\F$ 拟凝聚，则每个 $x\in V$ 都有一个形如 $D(f)\subset V$ 的邻域，
使 $\F|D(f)$ 同构于某个同态
$\wt{A_f}^{(I)}\to\wt{A_f}^{(J)}$ 的余核，
因而同构于 $A_f$-模 $N$ 的相伴层，其中 $N$ 是对应同态
$A_f^{(I)}\to A_f^{(J)}$ 的余核（\hyperlink{para:1:1.3.8}{1.3.8} 及 \hyperlink{para:1:1.3.9}{1.3.9}）。
因为 $V$ 拟紧，显然 $c)\Rightarrow b)$。

为证明 $b)$ 蕴含 $d1)$ 和 $d2)$，先假定 $V=D(g)$，其中 $g\in A$，
并且 $\F$ 同构于某个 $A_g$-模 $N$ 的相伴层 $\wt N$。
将 $X$ 换成 $V$、$A$ 换成 $A_g$（\hyperlink{para:1:1.3.6}{1.3.6}），便可归结到 $g=1$ 的情形。
这时 $\Gamma(D(f),\wt N)$ 与 $N_f$ 典范地等同（\hyperlink{para:1:1.3.6}{1.3.6} 及 \hyperlink{para:1:1.3.7}{1.3.7}），
故截面 $s\in\Gamma(D(f),\wt N)$ 与形如 $z/f^n$ 的元素等同，其中 $z\in N$。
截面 $f^ns$ 与 $N_f$ 中的元素 $z/1$ 等同，因而是与元素 $z\in N$ 等同的
$\wt N$ 在 $X$ 上的截面在 $D(f)$ 上的限制。这就证明了此情形中的 $d1)$。
同样，$t\in\Gamma(X,\wt N)$ 与某个元素 $z\in N$ 等同，
$t$ 在 $D(f)$ 上的限制与 $z$ 在 $N_f$ 中的像 $z/1$ 等同；
这一像为零，表示存在整数 $n\geq0$，使 $f^nz=0$ 在 $N$ 中成立，
也等价于 $f^nt=0$。

为完成 $b)$ 蕴含 $d1)$、$d2)$ 的证明，只需建立下面的引理。

\egastatement{1.4.1.1}{引理}
假定 $V$ 是有限多个形如 $D(g_i)$ 的集合之并，且各层
$\F|D(g_i)$、$\F|(D(g_i)\cap D(g_j))=\F|D(g_ig_j)$ 都满足 $d1)$ 和 $d2)$。
则 $\F$ 具有以下两个性质：
\begin{enumerate}[label={$d'\arabic*)$},leftmargin=2.7em]
\item 对于每个 $f\in A$ 及每个截面 $s\in\Gamma(D(f)\cap V,\F)$，
存在整数 $n\geq0$，使 $f^ns$ 可延拓为 $\F$ 在 $V$ 上的一个截面。
\item 对于每个 $f\in A$ 及每个在 $D(f)\cap V$ 上的限制为 $0$ 的截面
$t\in\Gamma(V,\F)$，存在整数 $n\geq0$，使 $f^nt=0$。
\end{enumerate}

先证明 $d'2)$。因为 $D(f)\cap D(g_i)=D(fg_i)$，对于每个 $i$，
存在整数 $n_i$，使 $(fg_i)^{n_i}t$ 在 $D(g_i)$ 上的限制为零。
由于 $g_i$ 在 $A_{g_i}$ 中的像可逆，$f^{n_i}t$ 在 $D(g_i)$ 上的限制也为零。
取 $n$ 为各 $n_i$ 中的最大者，即得 $d'2)$。

为证明 $d'1)$，将 $d1)$ 应用于层 $\F|D(g_i)$。
存在整数 $n_i\geq0$ 及 $\F$ 在 $D(g_i)$ 上的截面 $s'_i$，
使 $s'_i$ 延拓了 $(fg_i)^{n_i}s$ 在 $D(fg_i)$ 上的限制。
由于 $g_i$ 在 $A_{g_i}$ 中的像可逆，存在 $\F$ 在 $D(g_i)$ 上的截面 $s_i$，
使 $s'_i=g_i^{n_i}s_i$；于是 $s_i$ 延拓了 $f^{n_i}s$ 在 $D(fg_i)$ 上的限制。
还可以假定所有 $n_i$ 都等于同一个整数 $n$。
由构造，$s_i-s_j$ 在
$D(f)\cap D(g_i)\cap D(g_j)=D(fg_ig_j)$ 上的限制为零。
将 $d2)$ 应用于层 $\F|D(g_ig_j)$，可知存在整数 $m_{ij}\geq0$，
使 $(fg_ig_j)^{m_{ij}}(s_i-s_j)$ 在 $D(g_ig_j)$ 上的限制为零。
由于 $g_ig_j$ 在 $A_{g_ig_j}$ 中的像可逆，
$f^{m_{ij}}(s_i-s_j)$ 在 $D(g_ig_j)$ 上的限制为零。
于是可假定所有 $m_{ij}$ 都等于同一个整数 $m$，
从而存在截面 $s'\in\Gamma(V,\F)$，延拓各 $f^ms_i$。
此截面因而延拓了 $f^{n+m}s$，这就证明了 $d'1)$。

现在还需证明，$d1)$ 和 $d2)$ 蕴含 $a)$。
先证明：若 $\F$ 满足 $d1)$ 和 $d2)$，则对每个满足 $D(g)\subset V$ 的 $g\in A$，
层 $\F|D(g)$ 也满足这两个条件。对于 $d1)$，这是显然的。
另一方面，若 $t\in\Gamma(D(g),\F)$ 在 $D(f)\subset D(g)$ 上的限制为零，
则由 $d1)$，存在整数 $m\geq0$，使 $g^mt$ 延拓为 $\F$ 在 $V$ 上的截面 $s$。
应用 $d2)$ 可知，存在整数 $n\geq0$，使 $f^ns=0$；
由于 $g$ 在 $A_g$ 中的像可逆，故 $f^nt=0$。

由此，因为 $V$ 拟紧，引理（\hyperlink{para:1:1.4.1.1}{1.4.1.1}）说明条件 $d'1)$ 与 $d'2)$ 成立。
考虑 $A$-模 $M=\Gamma(V,\F)$，并定义一个 $\wt A$-模层同态
\[
u\colon\wt M\longrightarrow j_*(\F),
\]
其中 $j\colon V\to X$ 是典范包含映射。
由于各 $D(f)$ 构成 $X$ 的拓扑的一组基，只需对每个 $f\in A$ 定义同态
\[
u_f\colon M_f\longrightarrow\Gamma(D(f),j_*(\F))
=\Gamma(D(f)\cap V,\F),
\]
并验证通常的相容条件（0，\hyperlink{para:0:3.2.5}{3.2.5}）。
因为 $f$ 在 $A_f$ 中的典范像可逆，限制同态
$M=\Gamma(V,\F)\to\Gamma(D(f)\cap V,\F)$ 分解为
\[
M\longrightarrow M_f\xrightarrow{u_f}\Gamma(D(f)\cap V,\F)
\qquad\text{（0，\hyperlink{para:0:1.2.4}{1.2.4}）};
\]
当 $D(g)\subset D(f)$ 时，相容条件的验证是直接的。

下面证明，条件 $d'1)$（分别地，$d'2)$）蕴含每个 $u_f$ 满射（分别地，单射）。
这将说明 $u$ 双射，从而 $\F$ 是一个同构于 $\wt M$ 的 $\wt A$-模层在 $V$ 上的限制。
事实上，若 $s\in\Gamma(D(f)\cap V,\F)$，则由 $d'1)$，存在整数 $n\geq0$，
使 $f^ns$ 延拓为某个截面 $z\in M$。于是 $u_f(z/f^n)=s$，故 $u_f$ 满射。
同样，若 $z\in M$ 满足 $u_f(z/1)=0$，这表示截面 $z$ 在 $D(f)\cap V$ 上的限制为零。
由 $d'2)$，存在整数 $n\geq0$，使 $f^nz=0$，于是 $z/1=0$ 在 $M_f$ 中成立，
故 $u_f$ 单射。证毕。

\egastatement{1.4.2}{推论}
$X$ 的一个拟紧开集上的每个拟凝聚层，都由 $X$ 上的某个拟凝聚层诱导。

\egastatement{1.4.3}{推论}
$X=\Spec(A)$ 上的每个拟凝聚 $\OO_X$-代数层，都同构于形如 $\wt B$ 的
$\OO_X$-代数层，其中 $B$ 是一个 $A$-代数；
每个拟凝聚 $\wt B$-模层都同构于形如 $\wt N$ 的 $\wt B$-模层，其中 $N$ 是一个 $B$-模。

事实上，一个拟凝聚 $\OO_X$-代数层是拟凝聚 $\OO_X$-模层，
因而形如 $\wt B$，其中 $B$ 是一个 $A$-模。
由用 $\wt A$-模层同态 $\wt B\otimes_{\wt A}\wt B\to\wt B$
对 $\OO_X$-代数层结构的刻画，以及（\hyperlink{para:1:1.3.12}{1.3.12}），可知 $B$ 是 $A$-代数。
若 $\F$ 是拟凝聚 $\wt B$-模层，只需证明 $\F$ 也是拟凝聚 $\wt A$-模层，即可同样得出结论。
由于问题是局部的，限制到 $X$ 的某个形如 $D(f)$ 的开集上，便可假定 $\F$ 是某个
$\wt B$-模层同态 $\wt B^{(I)}\to\wt B^{(J)}$ 的余核；
此同态当然也是 $\wt A$-模层同态。
于是由（\hyperlink{para:1:1.3.8}{1.3.8}）及（\hyperlink{para:1:1.3.9}{1.3.9}）得到本命题。
\par


% END chapters/ch1_s1_4.tex

\par

% BEGIN chapters/ch1_s1_5.tex
\subsection{素谱上的凝聚层}\hypertarget{heading:1:1.5}{}
\src{92—93}

\egastatement{1.5.1}{定理}
设 $A$ 是诺特环，$X=\Spec(A)$ 是其素谱，$V$ 是 $X$ 的开子集，
$\F$ 是一个 $(\OO_X|V)$-模层。下列条件等价：
\begin{enumerate}[label=\alph*)]
\item $\F$ 凝聚。
\item $\F$ 有限型且拟凝聚。
\item 存在有限型 $A$-模 $M$，使 $\F$ 同构于层 $\wt M|V$。
\end{enumerate}

$a)$ 显然蕴含 $b)$。为证明 $b)$ 蕴含 $c)$，先注意，由于 $V$ 拟紧（0，\hyperlink{para:0:2.2.3}{2.2.3}），
$\F$ 同构于一个形如 $\wt N|V$ 的层，其中 $N$ 是 $A$-模（\hyperlink{para:1:1.4.1}{1.4.1}）。
而 $N=\ilim M_\lambda$，其中 $M_\lambda$ 遍历 $N$ 的有限型 $A$-子模，
因而由（\hyperlink{para:1:1.3.9}{1.3.9}）有
$\F=\wt N|V=\ilim(\wt{M_\lambda}|V)$。
因为 $\F$ 有限型且 $V$ 拟紧，存在指标 $\lambda$，使
$\F=\wt{M_\lambda}|V$（0，\hyperlink{para:0:5.2.3}{5.2.3}）。

最后证明 $c)$ 蕴含 $a)$。显然，这时 $\F$ 有限型（\hyperlink{para:1:1.3.6}{1.3.6} 及 \hyperlink{para:1:1.3.9}{1.3.9}）；
又因为问题是局部的，可以只考虑 $V=D(f)$、$f\in A$ 的情形。
由于 $A_f$ 是诺特环，最终只需证明：形如 $\wt A^n\to\wt M$ 的同态的核有限型，
其中 $M$ 是一个 $A$-模。
而这种同态形如 $\wt u$，其中 $u$ 是同态 $A^n\to M$（\hyperlink{para:1:1.3.8}{1.3.8}）；
若 $P=\Ker u$，则 $\wt P=\Ker\wt u$（\hyperlink{para:1:1.3.9}{1.3.9}）。
因为 $A$ 是诺特环，$P$ 有限型，证明完毕。

\egastatement{1.5.2}{推论}
在（\hyperlink{para:1:1.5.1}{1.5.1}）的假设下，层 $\OO_X$ 是凝聚环层。

\egastatement{1.5.3}{推论}
在（\hyperlink{para:1:1.5.1}{1.5.1}）的假设下，$X$ 的任意开集上的每个凝聚层，
都由 $X$ 上的某个凝聚层诱导。

\egastatement{1.5.4}{推论}
在（\hyperlink{para:1:1.5.1}{1.5.1}）的假设下，每个拟凝聚 $\OO_X$-模层 $\F$
都是其凝聚 $\OO_X$-子模层的归纳极限。

事实上，$\F=\wt M$，其中 $M$ 是某个 $A$-模，
而 $M$ 是其有限型子模的归纳极限；由（\hyperlink{para:1:1.3.9}{1.3.9}）及（\hyperlink{para:1:1.5.1}{1.5.1}）即得结论。
\par


% END chapters/ch1_s1_5.tex

\par

% BEGIN chapters/ch1_s1_6.tex
\subsection{素谱上拟凝聚层的函子性质}\hypertarget{heading:1:1.6}{}
\src{93—96}

\egapar{1.6.1}
设 $A,A'$ 是两个环，$\varphi\colon A'\to A$ 是同态，
\[
\aphi\colon X=\Spec(A)\longrightarrow X'=\Spec(A')
\]
是与 $\varphi$ 相伴的连续映射（\hyperlink{para:1:1.2.1}{1.2.1}）。我们将定义一个典范的环层同态
\[
\wt\varphi\colon\OO_{X'}\longrightarrow(\aphi)_*(\OO_X).
\]
对于每个 $f'\in A'$，令 $f=\varphi(f')$；由（\hyperref[eq:1:1.2.2.2]{1.2.2.2}），有
$(\aphi)^{-1}(D(f'))=D(f)$。
环 $\Gamma(D(f'),\wt{A'})$ 与 $\Gamma(D(f),\wt A)$
分别等同于 $A'_{f'}$ 与 $A_f$（\hyperlink{para:1:1.3.6}{1.3.6} 及 \hyperlink{para:1:1.3.7}{1.3.7}）。
而同态 $\varphi$ 典范地定义同态
$\varphi_{f'}\colon A'_{f'}\to A_f$（0，\hyperlink{para:0:1.5.1}{1.5.1}）；也就是说，有环同态
\[
\Gamma(D(f'),\wt{A'})
\longrightarrow\Gamma((\aphi)^{-1}(D(f')),\wt A)
=\Gamma(D(f'),(\aphi)_*(\wt A)).
\]
此外，这些同态满足通常的相容条件：当 $D(f')\supset D(g')$ 时，下图交换（0，\hyperlink{para:0:1.5.1}{1.5.1}）：
\[
\begin{tikzcd}[column sep=large,row sep=large]
\Gamma(D(f'),\wt{A'})\arrow[r]\arrow[d]
 & \Gamma(D(f'),(\aphi)_*(\wt A))\arrow[d]\\
\Gamma(D(g'),\wt{A'})\arrow[r]
 & \Gamma(D(g'),(\aphi)_*(\wt A)).
\end{tikzcd}
\]
由于各 $D(f')$ 构成 $X'$ 的拓扑的一组基（0，\hyperlink{para:0:3.2.3}{3.2.3}），
便确实定义了一个 $\OO_{X'}$-代数层同态。
因此，二元组 $\Phi=(\aphi,\wt\varphi)$ 是赋环空间态射
\[
\Phi\colon(X,\OO_X)\longrightarrow(X',\OO_{X'})
\qquad\text{（0，\hyperlink{para:0:4.1.1}{4.1.1}）}.
\]

还应指出，若令 $x'=\aphi(x)$，则同态 $\wt\varphi_x^\sharp$（0，\hyperlink{para:0:3.7.1}{3.7.1}）
正是由 $\varphi\colon A'\to A$ 典范得到的同态
\[
\varphi_x\colon A'_{x'}\longrightarrow A_x
\qquad\text{（0，\hyperlink{para:0:1.5.1}{1.5.1}）}.
\]
事实上，每个 $z'\in A'_{x'}$ 均可写成 $g'/f'$，其中 $f',g'\in A'$ 且 $f'\notin\jj_{x'}$。
所以 $D(f')$ 是 $x'$ 在 $X'$ 中的邻域，而由 $\wt\varphi$ 得到的同态
$\Gamma(D(f'),\wt{A'})\to\Gamma((\aphi)^{-1}(D(f')),\wt A)$ 就是 $\varphi_{f'}$。
考虑对应于 $g'/f'\in A'_{f'}$ 的截面 $s'\in\Gamma(D(f'),\wt{A'})$，
便得到 $\wt\varphi_x^\sharp(z')=\varphi(g')/\varphi(f')$ 在 $A_x$ 中成立。

\egapar{1.6.2}
\textbf{例。}设 $S$ 是 $A$ 的乘法子集，$\varphi$ 是典范同态 $A\to S^{-1}A$。
由（\hyperlink{para:1:1.2.6}{1.2.6}）已知，$\aphi$ 将 $Y=\Spec(S^{-1}A)$ 同胚地映到
$X=\Spec(A)$ 中由满足 $\jj_x\cap S=\varnothing$ 的点组成的子空间。
此外，对于属于这个子空间的每个 $x$，写作 $x=\aphi(y)$（$y\in Y$），
同态 $\wt\varphi_y^\sharp\colon\OO_x\to\OO_y$ 是双射（0，\hyperlink{para:0:1.2.6}{1.2.6}）。
换言之，$\OO_Y$ 与 $\OO_X$ 在 $Y$ 上诱导的层等同。

\egastatement{1.6.3}{命题}
对于每个 $A$-模 $M$，存在从 $\OO_{X'}$-模层 $\wt{\restmod M\varphi}$
到直像 $\Phi_*(\wt M)$ 的典范函子性同构。

为简便，令 $M'=\restmod M\varphi$；对于每个 $f'\in A'$，仍令 $f=\varphi(f')$。
截面模 $\Gamma(D(f'),\wt{M'})$ 与 $\Gamma(D(f),\wt M)$ 分别等同于
$M'_{f'}$ 与 $M_f$，它们分别是 $A'_{f'}$-模与 $A_f$-模。
此外，$A'_{f'}$-模 $\restmod{(M_f)}{\varphi_{f'}}$ 与 $M'_{f'}$ 典范同构（0，\hyperlink{para:0:1.5.2}{1.5.2}）。
因此有 $\Gamma(D(f'),\wt{A'})$-模的函子性同构
\[
\Gamma(D(f'),\wt{M'})
\xrightarrow{\sim}
\restmod{\Gamma((\aphi)^{-1}(D(f')),\wt M)}{\varphi_{f'}}.
\]
这些同构满足通常的与限制相容的条件（0，\hyperlink{para:0:1.5.6}{1.5.6}），因而定义了所述函子性同构。
更具体地，若 $u\colon M_1\to M_2$ 是 $A$-模同态，
则也可将它视为 $A'$-模同态
$\restmod{(M_1)}\varphi\to\restmod{(M_2)}\varphi$；
若将后者记为 $u_{[\varphi]}$，则 $\Phi_*(\wt u)$ 与 $\wt{u_{[\varphi]}}$ 等同。

此证明还说明，对于每个 $A$-代数 $B$，典范的函子性同构
$\wt{\restmod B\varphi}\to\Phi_*(\wt B)$ 是 $\OO_{X'}$-代数层同构。
若 $M$ 是 $B$-模，则典范的函子性同构
$\wt{\restmod M\varphi}\to\Phi_*(\wt M)$ 是 $\Phi_*(\wt B)$-模层同构。

\egastatement{1.6.4}{推论}
直像函子 $\Phi_*$ 在拟凝聚 $\OO_X$-模层范畴上正合。

事实上，显然 $\restmod M\varphi$ 是关于 $M$ 的正合函子，
而 $\wt{M'}$ 是关于 $M'$ 的正合函子（\hyperlink{para:1:1.3.5}{1.3.5}）。

\egastatement{1.6.5}{命题}
设 $N'$ 是 $A'$-模，$N$ 是 $A$-模 $N'\otimes_{A'}\restmod A\varphi$。
存在从 $\OO_X$-模层 $\Phi^*(\wt{N'})$ 到 $\wt N$ 的典范函子性同构。

先注意，映射 $j\colon z'\mapsto z'\otimes1$ 是从 $N'$ 到
$\restmod N\varphi$ 的 $A'$-同态。事实上，依定义，对于 $f'\in A'$，有
\[
(f'z')\otimes1=z'\otimes\varphi(f')=\varphi(f')(z'\otimes1).
\]
由（\hyperlink{para:1:1.3.8}{1.3.8}）得到 $\OO_{X'}$-模层同态
$\wt j\colon\wt{N'}\to\wt{\restmod N\varphi}$；
再由（\hyperlink{para:1:1.6.3}{1.6.3}），可以将 $\wt j$ 看成从 $\wt{N'}$ 到 $\Phi_*(\wt N)$ 的同态。
与此同态 $\wt j$ 典范对应的是同态
$h=\wt j^\sharp\colon\Phi^*(\wt{N'})\to\wt N$（0，\hyperlink{para:0:4.4.3}{4.4.3}）。
我们将证明每个茎上的 $h_x$ 都是双射。

令 $x'=\aphi(x)$，取 $f'\in A'$ 使 $x'\in D(f')$，并令 $f=\varphi(f')$。
环 $\Gamma(D(f),\wt A)$ 与 $A_f$ 等同；
模 $\Gamma(D(f),\wt N)$、$\Gamma(D(f'),\wt{N'})$ 分别与 $N_f$、$N'_{f'}$ 等同。
设截面 $s'\in\Gamma(D(f'),\wt{N'})$ 对应于 $n'/f'^p$（$n'\in N'$），
$s$ 是它经 $\wt j$ 在 $\Gamma(D(f),\wt N)$ 中的像；
则 $s$ 对应于 $(n'\otimes1)/f^p$。
另设 $t\in\Gamma(D(f),\wt A)$ 对应于 $g/f^q$（$g\in A$）；依定义，有
\[
h_x(s'_{x'}\otimes t_x)=t_xs_x
\qquad\text{（0，\hyperlink{para:0:4.4.3}{4.4.3}）}.
\]
而 $N_f$ 可以典范地等同于
$N'_{f'}\otimes_{A'_{f'}}\restmod{(A_f)}{\varphi_{f'}}$（0，\hyperlink{para:0:1.5.4}{1.5.4}）；
这时 $s$ 对应于元素 $(n'/f'^p)\otimes1$，而截面 $y\mapsto t_ys_y$
对应于 $(n'/f'^p)\otimes(g/f^q)$。
（0，\hyperlink{para:0:1.5.6}{1.5.6}）中的相容交换图说明，$h_x$ 就是典范同构
\begin{equation}
N'_{x'}\otimes_{A'_{x'}}\restmod{(A_x)}{\varphi_x}
\xrightarrow{\sim}
N_x=(N'\otimes_{A'}\restmod A\varphi)_x.
\tag{1.6.5.1}\label{eq:1:1.6.5.1}
\end{equation}

此外，设 $v\colon N'_1\to N'_2$ 是 $A'$-模同态。
由于每个 $x'\in X'$ 都满足 $\wt v_{x'}=v_{x'}$，
从上述论证立即得到：$\Phi^*(\wt v)$ 典范地等同于 $\wt{v\otimes1}$。
这就完成了（\hyperlink{para:1:1.6.5}{1.6.5}）的证明。

若 $B'$ 是 $A'$-代数，则从 $\Phi^*(\wt{B'})$ 到
$\wt{B'\otimes_{A'}\restmod A\varphi}$ 的典范同构是 $\OO_X$-代数层同构。
若再有一个 $B'$-模 $N'$，则从 $\Phi^*(\wt{N'})$ 到
$\wt{N'\otimes_{A'}\restmod A\varphi}$ 的典范同构是 $\Phi^*(\wt{B'})$-模层同构。

\egastatement{1.6.6}{推论}
当 $s'$ 遍历 $A'$-模 $\Gamma(\wt{N'})$ 时，
这些截面 $s'$ 在 $\Phi^*(\wt{N'})$ 中的典范像生成
$A$-模 $\Gamma(\Phi^*(\wt{N'}))$。\transnote{
原版此处将 $\Gamma(\wt{N'})$ 称为“$A$-模”；按（\hyperlink{para:1:1.6.5}{1.6.5}），$N'$ 是 $A'$-模，
故这里将该处写作“$A'$-模”。结论中的 $\Gamma(\Phi^*(\wt{N'}))$ 则确为 $A$-模。}

事实上，当将 $N'$、$N$ 分别与 $\Gamma(\wt{N'})$、$\Gamma(\wt N)$ 等同（\hyperlink{para:1:1.3.7}{1.3.7}）时，
这些像就是 $N$ 中的元素 $z'\otimes1$，其中 $z'$ 遍历 $N'$。

\egapar{1.6.7}
在（\hyperlink{para:1:1.6.5}{1.6.5}）的证明中，已顺便证明典范映射（0，\hyperref[eq:0:4.4.3.2]{4.4.3.2}）
\[
\rho\colon\wt{N'}\longrightarrow\Phi_*(\Phi^*(\wt{N'}))
\]
就是同态 $\wt j$，其中
$j\colon N'\to N'\otimes_{A'}\restmod A\varphi$ 是映射 $z'\mapsto z'\otimes1$。
同样，典范映射（0，\hyperref[eq:0:4.4.3.3]{4.4.3.3}）
\[
\sigma\colon\Phi^*(\Phi_*(\wt M))\longrightarrow\wt M
\]
就是 $\wt p$，其中
$p\colon\restmod M\varphi\otimes_{A'}\restmod A\varphi\to M$
是将每个纯张量 $z\otimes a$（$z\in M$，$a\in A$）映为 $az$ 的典范同态。
这直接由定义得到（0，\hyperlink{para:0:3.7.1}{3.7.1}；0，\hyperlink{para:0:4.4.3}{4.4.3}；I，\hyperlink{para:1:1.3.7}{1.3.7}）。

由此可知（0，\hyperlink{para:0:4.4.3}{4.4.3} 及 \hyperref[eq:0:3.5.4.4]{3.5.4.4}），若 $v\colon N'\to\restmod M\varphi$
是 $A'$-同态，则有 $\wt v^\sharp=\wt{v\otimes1}$。\transnote{
此式照录原文的简写。若 $v\otimes1$ 按通常的张量积同态理解，
它的靶是 $\restmod M\varphi\otimes_{A'}A$，还须与上一段的 $p$ 复合。
将靶明确写为 $M$ 时，完整的同态为
$z'\otimes a\mapsto a\,v(z')$，亦即
$\wt v^\sharp=\wt{p\circ(v\otimes1)}$。}

\egapar{1.6.8}
设 $N'_1,N'_2$ 是两个 $A'$-模，且 $N'_1$ 有有限表示。
由（\hyperlink{para:1:1.6.7}{1.6.7}）及（\hyperlink{para:1:1.3.12}{1.3.12}，(ii)）可知，典范同态（0，\hyperlink{para:0:4.4.6}{4.4.6}）
\[
\Phi^*\bigl(\shHom_{\wt{A'}}(\wt{N'_1},\wt{N'_2})\bigr)
\longrightarrow
\shHom_{\wt A}\bigl(\Phi^*(\wt{N'_1}),\Phi^*(\wt{N'_2})\bigr)
\]
就是 $\wt\gamma$，其中 $\gamma$ 是典范的 $A$-模同态
\[
\Hom_{A'}(N'_1,N'_2)\otimes_{A'}A
\longrightarrow
\Hom_A(N'_1\otimes_{A'}A,N'_2\otimes_{A'}A).
\]

\egapar{1.6.9}
设 $\frJ'$ 是 $A'$ 的理想，$M$ 是 $A$-模。
依定义，$\wt{\frJ'}\,\wt M$ 是典范同态
\[
\Phi^*(\wt{\frJ'})\otimes_{\wt A}\wt M\longrightarrow\wt M
\]
的像，故由（\hyperlink{para:1:1.6.5}{1.6.5}）及（\hyperlink{para:1:1.3.12}{1.3.12}，(i)），
$\wt{\frJ'}\,\wt M$ 典范地等同于 $\wt{\frJ'M}$。
特别地，$\Phi^*(\wt{\frJ'})\wt A$ 与 $\wt{\frJ'A}$ 等同。
再考虑函子 $\Phi^*$ 的右正合性，$\wt A$-代数层
$\Phi^*(\wt{A'/\frJ'})$ 便与 $\wt{A/\frJ'A}$ 等同。

\egapar{1.6.10}
设 $A''$ 是第三个环，$\varphi'\colon A''\to A'$ 是同态，令 $\varphi''=\varphi\circ\varphi'$。
直接由定义，有
\[
{}^a\varphi''=({}^a\varphi')\circ({}^a\varphi),
\qquad
\wt{\varphi''}=\wt\varphi\circ\wt{\varphi'}
\qquad\text{（0，\hyperlink{para:0:1.5.7}{1.5.7}）}.
\]
因此 $\Phi''=\Phi'\circ\Phi$；换言之，$(\Spec(A),\wt A)$ 是关于 $A$ 的函子，
从环范畴取值于赋环空间范畴。\transnote{
这里 $\wt\varphi\circ\wt{\varphi'}$ 按不同底空间上层态射的复合来理解（0，\hyperlink{para:0:3.5.2}{3.5.2}）；
若均写成直像层之间的同态，则为
$({}^a\varphi')_*(\wt\varphi)\circ\wt{\varphi'}$。
由 $\Phi''=\Phi'\circ\Phi$ 可见，这个环到赋环空间的函子是反变的。}
\par


% END chapters/ch1_s1_6.tex

\par

% BEGIN chapters/ch1_s1_7.tex
\subsection{仿射概形态射的刻画}\hypertarget{heading:1:1.7}{}
\src{96—97}

\egastatement{1.7.1}{定义}
若赋环空间 $(X,\OO_X)$ 同构于某个形如 $(\Spec(A),\wt A)$ 的赋环空间，其中 $A$ 是环，
则称 $(X,\OO_X)$ 为\emph{仿射概形}。
这时，与环 $A$ 典范地等同的环 $\Gamma(X,\OO_X)$（\hyperlink{para:1:1.3.7}{1.3.7}）
称为仿射概形 $(X,\OO_X)$ 的\emph{环}；在不致混淆时，将它记作 $A(X)$。

为简便起见，当说到仿射概形 $\Spec(A)$ 时，始终是指赋环空间 $(\Spec(A),\wt A)$。

\egapar{1.7.2}
设 $A,B$ 是两个环，$(X,\OO_X)$、$(Y,\OO_Y)$ 是对应于素谱
$X=\Spec(A)$、$Y=\Spec(B)$ 的仿射概形。
由（\hyperlink{para:1:1.6.1}{1.6.1}）已知，每个环同态 $\varphi\colon B\to A$ 对应于态射
\[
\Phi=(\aphi,\wt\varphi)=\Spec(\varphi)
\colon(X,\OO_X)\longrightarrow(Y,\OO_Y).
\]
注意，$\varphi$ 由 $\Phi$ 完全确定，因为依定义有
\[
\varphi=\Gamma(\wt\varphi)\colon
\Gamma(\wt B)\longrightarrow\Gamma((\aphi)_*(\wt A))=\Gamma(\wt A).
\]

\egastatement{1.7.3}{定理}
设 $(X,\OO_X)$、$(Y,\OO_Y)$ 是两个仿射概形。
赋环空间态射 $(\psi,\theta)\colon(X,\OO_X)\to(Y,\OO_Y)$
可写成 $(\aphi,\wt\varphi)$，其中 $\varphi$ 是环同态 $A(Y)\to A(X)$，
当且仅当对于每个 $x\in X$，$\theta_x^\sharp$ 是局部同态
$\OO_{\psi(x)}\to\OO_x$。

令 $A=A(X)$，$B=A(Y)$。此条件是必要的，因为由（\hyperlink{para:1:1.6.1}{1.6.1}）已知，
$\wt\varphi_x^\sharp$ 是由 $\varphi$ 典范得到的同态 $B_{\aphi(x)}\to A_x$；
按 $\aphi(x)=\varphi^{-1}(\jj_x)$ 的定义，此同态是局部的。

证明此条件充分。依定义，$\theta$ 是同态 $\OO_Y\to\psi_*(\OO_X)$，
由此典范地得到环同态
\[
\varphi=\Gamma(\theta)\colon
B=\Gamma(Y,\OO_Y)\longrightarrow
\Gamma(Y,\psi_*(\OO_X))=\Gamma(X,\OO_X)=A.
\]
关于 $\theta_x^\sharp$ 的假设使我们可以取商，得到剩余域之间的单同态
$\theta^x\colon k(\psi(x))\to k(x)$，并且对于每个截面 $f\in\Gamma(Y,\OO_Y)=B$，有
\[
\theta^x(f(\psi(x)))=\varphi(f)(x).
\]
因此，关系 $f(\psi(x))=0$ 等价于 $\varphi(f)(x)=0$。
这表示 $\jj_{\psi(x)}=\jj_{\aphi(x)}$，也可写成
$\psi(x)=\aphi(x)$ 对所有 $x\in X$ 成立，即 $\psi=\aphi$。
我们还知道下图交换（0，\hyperlink{para:0:3.7.1}{3.7.1}）：
\[
\begin{tikzcd}[column sep=large,row sep=large]
B=\Gamma(Y,\OO_Y)\arrow[r,"\varphi"]\arrow[d]
 & \Gamma(X,\OO_X)=A\arrow[d]\\
B_{\psi(x)}\arrow[r,"\theta_x^\sharp"'] & A_x.
\end{tikzcd}
\]
这表示 $\theta_x^\sharp$ 等于由 $\varphi$ 典范得到的同态
$\varphi_x\colon B_{\aphi(x)}\to A_x$（0，\hyperlink{para:0:1.5.1}{1.5.1}）。
由于各 $\theta_x^\sharp$ 完全刻画 $\theta^\sharp$，从而也完全刻画 $\theta$（0，\hyperlink{para:0:3.7.1}{3.7.1}），
按 $\wt\varphi$ 的定义（\hyperlink{para:1:1.6.1}{1.6.1}），便得到 $\theta=\wt\varphi$。

满足（\hyperlink{para:1:1.7.3}{1.7.3}）的条件的赋环空间态射 $(\psi,\theta)$，
称为\emph{仿射概形态射}。

\egastatement{1.7.4}{推论}
设 $(X,\OO_X)$、$(Y,\OO_Y)$ 是两个仿射概形，
$A=\Gamma(\OO_X)$，$B=\Gamma(\OO_Y)$。
则仿射概形态射集合 $\Hom((X,\OO_X),(Y,\OO_Y))$
与从 $B$ 到 $A$ 的环同态集合之间存在典范同构。

也可以说，关于 $A$ 的函子 $(\Spec(A),\wt A)$ 与关于 $(X,\OO_X)$ 的函子
$\Gamma(X,\OO_X)$，给出交换环范畴与仿射概形范畴的对偶范畴之间的等价（T，I，1.2）。

\egastatement{1.7.5}{推论}
若 $\varphi\colon B\to A$ 满射，则对应态射 $(\aphi,\wt\varphi)$
是赋环空间的单态射（参见（\hyperlink{para:1:4.1.7}{4.1.7}））。

事实上，我们知道 $\aphi$ 单射（\hyperlink{para:1:1.2.5}{1.2.5}）。又因为 $\varphi$ 满射，
对于每个 $x\in X$，由 $\varphi$ 取分式环得到的同态
$\wt\varphi_x^\sharp\colon B_{\aphi(x)}\to A_x$ 也满射（0，\hyperlink{para:0:1.5.1}{1.5.1}）。
由此得到结论（0，\hyperlink{para:0:4.1.1}{4.1.1}）。
\par


% END chapters/ch1_s1_7.tex

\par

% BEGIN chapters/ch1_s2_1.tex
\section{预概形与预概形态射}\hypertarget{heading:1:2}{}
\subsection{预概形的定义}\hypertarget{heading:1:2.1}{}
\src{97—98}

\egapar{2.1.1}
给定一个赋环空间 $(X,\OO_X)$。若 $X$ 的开子集 $V$ 所对应的赋环空间
$(V,\OO_X|V)$ 是仿射概形（\hyperlink{para:1:1.7.1}{1.7.1}），则称 $V$ 为一个\emph{仿射开集}。

\egastatement{2.1.2}{定义}
所谓\emph{预概形}，是指一个赋环空间 $(X,\OO_X)$，其中 $X$ 的每一点都具有仿射开邻域。

\egastatement{2.1.3}{命题}
若 $(X,\OO_X)$ 是预概形，则仿射开集构成 $X$ 的拓扑的一组基。

事实上，设 $V$ 为 $x\in X$ 的任意开邻域。依假设，$x$ 具有一个开邻域 $W$，
使得 $(W,\OO_X|W)$ 是仿射概形；以 $A$ 表示它的环。
在空间 $W$ 中，$V\cap W$ 是 $x$ 的开邻域，故存在 $f\in A$，使得
$D(f)$ 是包含在 $V\cap W$ 中的 $x$ 的开邻域（\hyperlink{para:1:1.1.10}{1.1.10}，(i)）。
赋环空间 $(D(f),\OO_X|D(f))$ 是环同构于 $A_f$ 的仿射概形（\hyperlink{para:1:1.3.6}{1.3.6}），命题得证。

\egastatement{2.1.4}{命题}
预概形的底空间是 Kolmogoroff 空间。

事实上，若 $x,y$ 是预概形 $X$ 的两个不同的点，且它们不属于同一个仿射开集，
则显然存在其中一点的一个开邻域不含另一点；若它们属于同一个仿射开集，
则由（\hyperlink{para:1:1.1.8}{1.1.8}）得到同一结论。

\egastatement{2.1.5}{命题}
若 $(X,\OO_X)$ 是预概形，则 $X$ 的每个不可约闭子集都有且只有一个泛点；
因此，映射 $x\mapsto\overline{\{x\}}$ 是 $X$ 到其不可约闭子集的集合的双射。

事实上，设 $Y$ 为 $X$ 的不可约闭子集，$y\in Y$，$U$ 为 $y$ 在 $X$ 中的一个仿射开邻域。
则 $U\cap Y$ 在 $Y$ 中稠密且不可约（0，\hyperlink{para:0:2.1.1}{2.1.1} 及 \hyperlink{para:0:2.1.4}{2.1.4}）。
由（\hyperlink{para:1:1.1.14}{1.1.14}），$U\cap Y$ 是某一点 $x$ 在 $U$ 中的闭包，因而 $Y$ 是 $x$ 在 $X$ 中的闭包。
泛点的唯一性由（\hyperlink{para:1:2.1.4}{2.1.4}）及（0，\hyperlink{para:0:2.1.3}{2.1.3}）得到。

\egapar{2.1.6}
若 $Y$ 是 $X$ 的不可约闭子集，$y$ 是它的泛点，则局部环 $\OO_y$ 也记为
$\OO_{X/Y}$，称为\emph{$X$ 沿 $Y$ 的局部环}，或\emph{$Y$ 在 $X$ 中的局部环}。

若 $X$ 本身不可约，且 $x$ 为其泛点，则也称 $\OO_x$ 为\emph{$X$ 上的有理函数环}（参见 §7）。

\egastatement{2.1.7}{命题}
若 $(X,\OO_X)$ 是预概形，则对于 $X$ 的任意开子集 $U$，
赋环空间 $(U,\OO_X|U)$ 是预概形。

这直接由定义（\hyperlink{para:1:2.1.2}{2.1.2}）及命题（\hyperlink{para:1:2.1.3}{2.1.3}）得到。

称 $(U,\OO_X|U)$ 为 $(X,\OO_X)$ \emph{在 $U$ 上诱导的预概形}，
或 $(X,\OO_X)$ \emph{在 $U$ 上的限制}。

\egapar{2.1.8}
若预概形 $(X,\OO_X)$ 的底空间 $X$ 不可约（分别为连通），则称该预概形
\emph{不可约}（分别为\emph{连通}）。称一个预概形为\emph{整的}，
若它不可约且约化（参见（\hyperlink{para:1:5.1.4}{5.1.4}））。称预概形 $(X,\OO_X)$ \emph{局部整}，
若每个 $x\in X$ 都有一个开邻域 $U$，使得 $(X,\OO_X)$ 在 $U$ 上诱导的预概形是整的。


% END chapters/ch1_s2_1.tex

\par

% BEGIN chapters/ch1_s2_2.tex
\subsection{预概形态射}\hypertarget{heading:1:2.2}{}
\src{98—100}

\egastatement{2.2.1}{定义}
给定两个预概形 $(X,\OO_X)$、$(Y,\OO_Y)$。
所谓从 $(X,\OO_X)$ 到 $(Y,\OO_Y)$ 的\emph{（预概形）态射}，
是指一个赋环空间态射 $(\psi,\theta)$，使得对于每个 $x\in X$，
$\theta_x^\sharp$ 都是局部同态 $\OO_{\psi(x)}\to\OO_x$。

取商以后，同态 $\theta_x^\sharp\colon\OO_{\psi(x)}\to\OO_x$
给出单同态 $\theta^x\colon k(\psi(x))\to k(x)$；由此可将 $k(x)$ 看成域 $k(\psi(x))$ 的扩张。

\egapar{2.2.2}
两个预概形态射 $(\psi,\theta)$、$(\psi',\theta')$ 的复合 $(\psi'',\theta'')$
仍为预概形态射，这由公式
$\theta^{\prime\prime\sharp}=\theta^\sharp\circ\psi^*(\theta^{\prime\sharp})$
（0，\hyperlink{para:0:3.5.5}{3.5.5}）得到。因此，预概形构成一个范畴。依照一般的记号，
以 $\Hom(X,Y)$ 表示从预概形 $X$ 到预概形 $Y$ 的态射集合。

\egastatement{2.2.3}{例}
若 $U$ 是 $X$ 的开子集，则诱导预概形 $(U,\OO_X|U)$ 到 $(X,\OO_X)$ 的
典范嵌入（0，\hyperlink{para:0:4.1.2}{4.1.2}）是预概形态射。而且它是赋环空间的单态射，
因而也是预概形的单态射；这直接由（0，\hyperlink{para:0:4.1.1}{4.1.1}）得到。

\egastatement{2.2.4}{命题}
设 $(X,\OO_X)$ 是预概形，$(S,\OO_S)$ 是环为 $A$ 的仿射概形。
则从预概形 $(X,\OO_X)$ 到预概形 $(S,\OO_S)$ 的态射，
与从 $A$ 到环 $\Gamma(X,\OO_X)$ 的同态之间存在典范的一一对应。

首先指出，若 $(X,\OO_X)$、$(Y,\OO_Y)$ 是任意两个赋环空间，
则从前者到后者的态射 $(\psi,\theta)$ 典范地确定环同态
\[
\Gamma(\theta)\colon\Gamma(Y,\OO_Y)
\longrightarrow\Gamma(Y,\psi_*(\OO_X))=\Gamma(X,\OO_X).
\]
在所考虑的情形中，只须证明：任意同态 $\varphi\colon A\to\Gamma(X,\OO_X)$
都能写成 $\Gamma(\theta)$，且相应的态射唯一。
依假设，$X$ 有一个由仿射开集构成的覆盖 $(V_\alpha)$。
将 $\varphi$ 与限制同态
$\Gamma(X,\OO_X)\to\Gamma(V_\alpha,\OO_X|V_\alpha)$ 复合，得到同态
\[
\varphi_\alpha\colon A\longrightarrow\Gamma(V_\alpha,\OO_X|V_\alpha).
\]
由（\hyperlink{para:1:1.7.3}{1.7.3}），它对应于唯一的态射
$(\psi_\alpha,\theta_\alpha)\colon(V_\alpha,\OO_X|V_\alpha)\to(S,\OO_S)$。
此外，对于任意一对指标 $(\alpha,\beta)$，$V_\alpha\cap V_\beta$ 的每一点都有
一个包含在 $V_\alpha\cap V_\beta$ 中的仿射开邻域 $W$（\hyperlink{para:1:2.1.3}{2.1.3}）。
显然，将 $\varphi_\alpha$、$\varphi_\beta$ 分别与到 $W$ 的限制同态复合，
得到的是同一个同态 $\Gamma(S,\OO_S)\to\Gamma(W,\OO_X|W)$。
因而，由对任意 $\alpha$ 及 $x\in V_\alpha$ 都成立的关系
$(\theta_\alpha^\sharp)_x=(\varphi_\alpha)_x$（\hyperlink{para:1:1.6.1}{1.6.1}），
态射 $(\psi_\alpha,\theta_\alpha)$ 与 $(\psi_\beta,\theta_\beta)$ 在 $W$ 上的限制相同。
于是，存在唯一的赋环空间态射
$(\psi,\theta)\colon(X,\OO_X)\to(S,\OO_S)$，
其在每个 $V_\alpha$ 上的限制为 $(\psi_\alpha,\theta_\alpha)$。
显然，它是预概形态射，且 $\Gamma(\theta)=\varphi$。

设 $u\colon A\to\Gamma(X,\OO_X)$ 是环同态，
$\varphi=(\psi,\theta)\colon(X,\OO_X)\to(S,\OO_S)$ 是对应态射。
则对于每个 $f\in A$，有
\begin{equation}
\varphi^{-1}(D(f))=X_{u(f)}.\tag{2.2.4.1}\label{eq:1:2.2.4.1}
\end{equation}
这里使用的是（0，\hyperlink{para:0:5.5.2}{5.5.2}）中关于局部自由层 $\OO_X$ 的记号。
事实上，只须在 $X$ 本身为仿射概形时验证此式，而这时它就是（\hyperref[eq:1:1.2.2.2]{1.2.2.2}）。

\egastatement{2.2.5}{命题}
在（\hyperlink{para:1:2.2.4}{2.2.4}）的假设下，设 $\varphi\colon A\to\Gamma(X,\OO_X)$ 是环同态，
$f\colon(X,\OO_X)\to(S,\OO_S)$ 是对应的预概形态射，
$\F$、$\G$ 分别是拟凝聚的 $\OO_S$-模层与 $\OO_X$-模层，并令 $M=\Gamma(S,\F)$。
则 $f$-态射 $\F\to\G$（0，\hyperlink{para:0:4.4.1}{4.4.1}）与 $A$-同态
$M\to\restmod{\Gamma(X,\G)}{\varphi}$ 之间存在典范的一一对应。

事实上，仿照（\hyperlink{para:1:2.2.4}{2.2.4}）中的论证，即可归结到 $X$ 为仿射概形的情形；
这时命题由（\hyperlink{para:1:1.6.3}{1.6.3}）及（\hyperlink{para:1:1.3.8}{1.3.8}）得到。

\egapar{2.2.6}
称预概形态射 $(\psi,\theta)\colon(X,\OO_X)\to(Y,\OO_Y)$ 为\emph{开的}
（分别为\emph{闭的}），若对于 $X$ 的每个开子集 $U$（分别为闭子集 $F$），
$\psi(U)$ 在 $Y$ 中开（分别为 $\psi(F)$ 在 $Y$ 中闭）。
若 $\psi(X)$ 在 $Y$ 中稠密，则称 $(\psi,\theta)$ 为\emph{支配态射}；
若 $\psi$ 满射，则称它为\emph{满射态射}。注意，这些条件只涉及连续映射 $\psi$。

\egastatement{2.2.7}{命题}
设
\[
f=(\psi,\theta)\colon(X,\OO_X)\to(Y,\OO_Y),\qquad
 g=(\psi',\theta')\colon(Y,\OO_Y)\to(Z,\OO_Z)
\]
是两个预概形态射。
\begin{enumerate}[label=(\roman*)]
\item 若 $f,g$ 都是开的（分别为闭的、支配的、满射的），则 $g\circ f$ 也具有相应性质。
\item 若 $f$ 满射，且 $g\circ f$ 是闭的，则 $g$ 是闭的。
\item 若 $g\circ f$ 满射，则 $g$ 满射。
\end{enumerate}
断言 (i)、(iii) 是显然的。令 $g\circ f=(\psi'',\theta'')$。
若 $F$ 在 $Y$ 中闭，则 $\psi^{-1}(F)$ 在 $X$ 中闭，故 $\psi''(\psi^{-1}(F))$ 在 $Z$ 中闭。
但由于 $\psi$ 满射，有 $\psi(\psi^{-1}(F))=F$，从而
$\psi''(\psi^{-1}(F))=\psi'(F)$，这就证明了 (ii)。

\egastatement{2.2.8}{命题}
设 $f=(\psi,\theta)$ 为态射 $(X,\OO_X)\to(Y,\OO_Y)$，$(U_\alpha)$ 是 $Y$ 的开覆盖。
则 $f$ 是开的（分别为闭的、满射的、支配的），当且仅当它在每个诱导预概形
$(\psi^{-1}(U_\alpha),\OO_X|\psi^{-1}(U_\alpha))$ 上的限制，
视为从此诱导预概形到诱导预概形 $(U_\alpha,\OO_Y|U_\alpha)$ 的态射时，
是开的（分别为闭的、满射的、支配的）。

此命题直接由定义得到，只须注意：$Y$ 的子集 $F$ 在 $Y$ 中闭（分别为开、稠密），
当且仅当每个 $F\cap U_\alpha$ 在 $U_\alpha$ 中闭（分别为开、稠密）。

\egapar{2.2.9}
设 $(X,\OO_X)$、$(Y,\OO_Y)$ 为两个预概形，并假设 $X,Y$ 具有相同的有限个数的
不可约分支，分别记为 $X_i,Y_i$（$1\leq i\leq n$）；
以 $\xi_i,\eta_i$ 分别表示 $X_i,Y_i$ 的泛点（\hyperlink{para:1:2.1.5}{2.1.5}）。
若态射
\[
f=(\psi,\theta)\colon(X,\OO_X)\longrightarrow(Y,\OO_Y)
\]
满足：对每个 $i$，有 $\psi(\xi_i)=\eta_i$，且
$\theta_{\xi_i}^\sharp\colon\OO_{\eta_i}\to\OO_{\xi_i}$ 是同构，
则称 $f$ 为\emph{双有理态射}。显然，双有理态射是支配的（0，\hyperlink{para:0:2.1.8}{2.1.8}），
因而，若它还是闭的，则它满射。

\egastatement{2.2.10}{记号约定}
在本书以下各处，只要不致混淆，我们将在预概形（分别为态射）的记号中省略
结构层（分别为结构层之间的态射）。若 $U$ 是预概形 $X$ 的底空间的开子集，
则当把 $U$ 作为预概形来谈论时，始终是指在 $U$ 上诱导的预概形。


% END chapters/ch1_s2_2.tex

\par

% BEGIN chapters/ch1_s2_3.tex
\subsection{预概形的粘合}\hypertarget{heading:1:2.3}{}
\src{101}

\egapar{2.3.1}
由定义（\hyperlink{para:1:2.1.2}{2.1.2}）可知，由预概形粘合而得的任意赋环空间（0，\hyperlink{para:0:4.1.6}{4.1.6}）仍然是预概形。
\transnote{此处原文引作（0，\hyperlink{para:0:4.1.6}{4.1.6}）；本版关于赋环空间粘合的条目实际为（0，\hyperlink{para:0:4.1.7}{4.1.7}）。}
特别地，由于每个预概形依定义都有由仿射开集构成的覆盖，
可见每个预概形都能由仿射概形粘合而得。

\egastatement{2.3.2}{例}
设 $K$ 是域，$B=K[s]$、$C=K[t]$ 是 $K$ 上的两个一元多项式环；
令 $X_1=\Spec(B)$、$X_2=\Spec(C)$，它们是两个同构的仿射概形。
在 $X_1$（分别在 $X_2$）中，令 $U_{12}$（分别令 $U_{21}$）为仿射开集
$D(s)$（分别为 $D(t)$）。其环 $B_s$（分别为 $C_t$）由形如
$f(s)/s^n$（分别为 $g(t)/t^n$）的有理分式组成，其中 $f\in B$（分别为 $g\in C$）。
设 $u_{12}\colon U_{21}\to U_{12}$ 为与环同构 $B_s\to C_t$ 对应的预概形同构（\hyperlink{para:1:2.2.4}{2.2.4}），
此环同构把 $f(s)/s^n$ 映到有理分式 $f(1/t)/(1/t)^n$。
可以借助 $u_{12}$ 沿 $U_{12}$、$U_{21}$ 将 $X_1$ 与 $X_2$ 粘合，
因为这里显然没有另需满足的粘合条件。

以后，我们将把如此得到的预概形 $X$ 看成一种一般构造方法的特殊情形（II，2.4.3）。
这里仅证明 $X$ 不是仿射概形；这将由环 $\Gamma(X,\OO_X)$ 同构于 $K$、
因而其素谱仅有一点而得到。事实上，$\OO_X$ 的一个全局截面在 $X_1$
（分别在 $X_2$）上的限制，是一个多项式 $f(s)$（分别为 $g(t)$）；
这里已将 $X_1,X_2$ 等同于 $X$ 中相应的仿射开集。
由定义可知，必须有 $g(t)=f(1/t)$，而这只有在 $f=g\in K$ 时才可能。


% END chapters/ch1_s2_3.tex

\par

% BEGIN chapters/ch1_s2_4.tex
\subsection{局部概形}\hypertarget{heading:1:2.4}{}
\src{101—103}

\egapar{2.4.1}
所谓\emph{局部概形}，是指其环 $A$ 为局部环的仿射概形；
这时 $X=\Spec(A)$ 中有且只有一个闭点 $a$，且对于任意其他点 $b\in X$，
有 $a\in\overline{\{b\}}$（\hyperlink{para:1:1.1.7}{1.1.7}）。

对于任意预概形 $Y$ 及点 $y\in Y$，称局部概形 $\Spec(\OO_y)$ 为
\emph{$Y$ 在点 $y$ 的局部概形}。
设 $V$ 为 $Y$ 中含 $y$ 的仿射开集，$B$ 为仿射概形 $V$ 的环。
$\OO_y$ 与 $B_y$ 典范地等同（\hyperlink{para:1:1.3.4}{1.3.4}），故典范同态 $B\to B_y$ 对应于一个
预概形态射 $\Spec(\OO_y)\to V$（\hyperlink{para:1:1.6.1}{1.6.1}）。
将此态射与典范嵌入 $V\to Y$ 复合，便得到态射 $\Spec(\OO_y)\to Y$；
它不依赖于所选的含 $y$ 的仿射开集 $V$。
事实上，若 $V'$ 为第二个含 $y$ 的仿射开集，则存在第三个含 $y$ 的仿射开集 $W$，
满足 $W\subset V\cap V'$（\hyperlink{para:1:2.1.3}{2.1.3}）；因而可归结到 $V\subset V'$ 的情形。
若 $B'$ 为 $V'$ 的环，则只须注意下图交换（0，\hyperlink{para:0:1.5.1}{1.5.1}）：
\[
\begin{tikzcd}[column sep=large,row sep=large]
 B'\arrow[r]\arrow[dr] & B\arrow[d]\\
 & \OO_y.
\end{tikzcd}
\]
如此定义的态射 $\Spec(\OO_y)\to Y$ 称为\emph{典范态射}。

\egastatement{2.4.2}{命题}
设 $(Y,\OO_Y)$ 为预概形。对于每个 $y\in Y$，令 $(\psi,\theta)$ 为典范态射
\[
(\Spec(\OO_y),\wt{\OO_y})\longrightarrow(Y,\OO_Y).
\]
则 $\psi$ 是从 $\Spec(\OO_y)$ 到 $Y$ 的子空间 $S_y$ 的同胚，
其中 $S_y$ 由满足 $y\in\overline{\{z\}}$ 的点 $z$ 组成，
即由 $y$ 的一般化点组成（0，\hyperlink{para:0:2.1.2}{2.1.2}）。
此外，若 $z=\psi(\mathfrak p)$，则相应的茎同态
$\OO_z\to(\OO_y)_{\mathfrak p}$ 是同构；因此，$(\psi,\theta)$ 是赋环空间的单态射。

由于 $\Spec(\OO_y)$ 的唯一闭点 $a$ 属于此空间每一点的闭包，且 $\psi(a)=y$，
连续映射 $\psi$ 下 $\Spec(\OO_y)$ 的像包含于 $S_y$。
由于 $S_y$ 包含于每个含 $y$ 的仿射开集，可以归结到 $Y$ 为仿射概形的情形；
而这时命题由（\hyperlink{para:1:1.6.2}{1.6.2}）得到。

因此可见（\hyperlink{para:1:2.1.5}{2.1.5}），$\Spec(\OO_y)$ 与 $Y$ 中含 $y$ 的不可约闭子集的集合之间
存在一一对应。

\egastatement{2.4.3}{推论}
点 $y\in Y$ 是 $Y$ 的某个不可约分支的泛点，当且仅当局部环 $\OO_y$ 的
唯一素理想是它的极大理想，换言之，当且仅当 $\OO_y$ 的维数为零。

\egastatement{2.4.4}{命题}
设 $(X,\OO_X)$ 是环为 $A$ 的局部概形，$a$ 是它的唯一闭点，$(Y,\OO_Y)$ 为预概形。
每个态射 $u=(\psi,\theta)\colon(X,\OO_X)\to(Y,\OO_Y)$ 都唯一地分解为
\[
 X\longrightarrow\Spec(\OO_{\psi(a)})\longrightarrow Y,
\]
其中第二个箭头为典范态射，第一个箭头对应于一个局部同态 $\OO_{\psi(a)}\to A$。
这就在态射 $(X,\OO_X)\to(Y,\OO_Y)$ 的集合与局部同态
$\OO_y\to A$（$y\in Y$）的集合之间建立了典范的一一对应。

事实上，对于每个 $x\in X$，有 $a\in\overline{\{x\}}$，
因而 $\psi(a)\in\overline{\{\psi(x)\}}$。
这说明 $\psi(X)$ 包含于每个含 $\psi(a)$ 的仿射开集。
因此可以归结到 $(Y,\OO_Y)$ 为环是 $B$ 的仿射概形的情形；
这时 $u=(\aphi,\wt\varphi)$，其中 $\varphi\in\Hom(B,A)$（\hyperlink{para:1:1.7.3}{1.7.3}）。
此外，$\varphi^{-1}(\jj_a)=\jj_{\psi(a)}$，故 $B\setminus\jj_{\psi(a)}$ 中
每个元素在 $\varphi$ 下的像都在局部环 $A$ 中可逆。
命题所述的分解于是由分式环的泛性质（0，\hyperlink{para:0:1.2.4}{1.2.4}）得到。
反过来，每个局部同态 $\OO_y\to A$ 都对应于唯一的态射
$(\psi',\theta')\colon X\to\Spec(\OO_y)$，使得 $\psi'(a)$ 为对应于 $y$ 的闭点（\hyperlink{para:1:1.7.3}{1.7.3}）；
将它与典范态射 $\Spec(\OO_y)\to Y$ 复合，便得到态射 $X\to Y$，命题得证。

\egapar{2.4.5}
环为域 $K$ 的仿射概形，其底空间只有一点。
若 $A$ 为极大理想是 $\mathfrak m$ 的局部环，则每个局部同态 $A\to K$ 的核都等于
$\mathfrak m$，因而分解为 $A\to A/\mathfrak m\to K$，其中第二个箭头是单同态。
因此，态射 $\Spec(K)\to\Spec(A)$ 与域单同态 $A/\mathfrak m\to K$ 一一对应。
\transnote{上一句按原文保留，但其适用范围须与前一句的“局部同态”一致：
它描述的是将 $\Spec(K)$ 的唯一点映到 $\Spec(A)$ 的闭点的态射。
对于所有态射 $\Spec(K)\to\Spec(A)$，还必须允许该点映到 $\Spec(A)$ 的其他点；
完整的对应关系是下述（\hyperlink{para:1:2.4.6}{2.4.6}）中的 $k(y)\to K$，其中 $y$ 遍历目标的各点。
不能仅由 $A$ 是局部环就省去这一选择。}

设 $(Y,\OO_Y)$ 是预概形。对于每个 $y\in Y$ 及 $\OO_y$ 的理想 $\mathfrak a_y$，
典范同态 $\OO_y\to\OO_y/\mathfrak a_y$ 定义了一个态射
$\Spec(\OO_y/\mathfrak a_y)\to\Spec(\OO_y)$；
将其与典范态射 $\Spec(\OO_y)\to Y$ 复合，得到态射
$\Spec(\OO_y/\mathfrak a_y)\to Y$，也称为\emph{典范态射}。
取 $\mathfrak a_y=\mathfrak m_y$，即 $\OO_y/\mathfrak a_y=k(y)$，
命题（\hyperlink{para:1:2.4.4}{2.4.4}）遂给出：

\egastatement{2.4.6}{推论}
设 $(X,\OO_X)$ 是环为域 $K$ 的局部概形，$\xi$ 是 $X$ 的唯一点，$(Y,\OO_Y)$ 是预概形。
每个态射 $u\colon(X,\OO_X)\to(Y,\OO_Y)$ 都唯一地分解为
\[
 X\longrightarrow\Spec(k(u(\xi)))\longrightarrow Y,
\]
其中第二个箭头是典范态射，第一个箭头对应于单同态 $k(u(\xi))\to K$。
这在态射 $(X,\OO_X)\to(Y,\OO_Y)$ 的集合与单同态 $k(y)\to K$（$y\in Y$）
的集合之间建立了典范的一一对应。

\egastatement{2.4.7}{推论}
对于每个 $y\in Y$，典范态射 $\Spec(\OO_y/\mathfrak a_y)\to Y$ 是赋环空间的单态射。

在 $\mathfrak a_y=0$ 时，我们已知此结论（\hyperlink{para:1:2.4.2}{2.4.2}）；再应用（\hyperlink{para:1:1.7.5}{1.7.5}）即可。

\egastatement{2.4.8}{注}
设 $X$ 是局部概形，$a$ 是其唯一闭点。
由于每个含 $a$ 的仿射开集必为整个 $X$，每个可逆的 $\OO_X$-模层（0，\hyperlink{para:0:5.4.1}{5.4.1}）
必同构于 $\OO_X$，也称它为\emph{平凡的}。
对于一般的仿射概形 $\Spec(A)$，此性质并不一定成立；
我们将在第 V 章看到，若 $A$ 是正规环，则当 $A$ 为析因环时，此性质成立。


% END chapters/ch1_s2_4.tex

\par

% BEGIN chapters/ch1_s2_5.tex
\subsection{一个预概形上的预概形}\hypertarget{heading:1:2.5}{}
\src{103—104}

\egastatement{2.5.1}{定义}
给定一个预概形 $S$。给出一个预概形 $X$ 及一个预概形态射
$\varphi\colon X\to S$，就称为给出了一个\emph{预概形 $S$ 上的预概形 $X$}，
或一个\emph{$S$-预概形}；称 $S$ 为 $S$-预概形 $X$ 的\emph{基预概形}，
称态射 $\varphi$ 为 $S$-预概形 $X$ 的\emph{结构态射}。
当 $S$ 是环为 $A$ 的仿射概形时，也称赋予 $\varphi$ 的 $X$ 为
\emph{环 $A$ 上的预概形}，或\emph{$A$-预概形}。

由（\hyperlink{para:1:2.2.4}{2.2.4}）可知，给出一个环 $A$ 上的预概形，等价于给出一个预概形
$(X,\OO_X)$，其结构层 $\OO_X$ 为一个 $A$-代数层。
因此，\emph{任意预概形都能以唯一的方式看成一个 $\mathbf Z$-预概形}。

若 $\varphi\colon X\to S$ 是 $S$-预概形 $X$ 的结构态射，
则当 $\varphi(x)=s$ 时，称点 $x\in X$ \emph{位于点 $s\in S$ 之上}。
若 $\varphi$ 是支配态射（\hyperlink{para:1:2.2.6}{2.2.6}），则称 $X$ \emph{支配} $S$。

\egapar{2.5.2}
设 $X,Y$ 是两个 $S$-预概形。
若预概形态射 $u\colon X\to Y$ 使下图交换：
\[
\begin{tikzcd}[column sep=large,row sep=large]
X\arrow[r,"u"]\arrow[dr] & Y\arrow[d]\\
& S,
\end{tikzcd}
\]
其中斜向及向下的箭头为结构态射，则称 $u$ 为一个\emph{$S$ 上的预概形态射}，
或一个\emph{$S$-态射}。这蕴含：对于每个 $s\in S$ 及位于 $s$ 之上的 $x\in X$，
$u(x)$ 也位于 $s$ 之上。

由此定义立即可见，两个 $S$-态射的复合仍为 $S$-态射，故 $S$-预概形构成一个范畴。
以 $\Hom_S(X,Y)$ 表示从 $S$-预概形 $X$ 到 $S$-预概形 $Y$ 的 $S$-态射集合；
$S$-预概形 $X$ 的恒等态射记为 $1_X$。
当 $S$ 是环为 $A$ 的仿射概形时，也称 $S$-态射为\emph{$A$-态射}。

\egapar{2.5.3}
若 $X$ 是 $S$-预概形，$v\colon X'\to X$ 是预概形态射，
则复合态射 $X'\to X\to S$ 使 $X'$ 成为 $S$-预概形。
特别地，借助典范嵌入，可以把在 $X$ 的任意开集 $U$ 上诱导的预概形看成 $S$-预概形。

因此，若 $u\colon X\to Y$ 是 $S$-预概形之间的 $S$-态射，
则 $u$ 在 $X$ 的任意开集 $U$ 上诱导的预概形上的限制，是一个 $S$-态射 $U\to Y$。
反过来，设 $(U_\alpha)$ 为 $X$ 的开覆盖，并对每个 $\alpha$ 给出一个
$S$-态射 $u_\alpha\colon U_\alpha\to Y$。
若对于每对指标 $(\alpha,\beta)$，$u_\alpha$ 与 $u_\beta$ 在
$U_\alpha\cap U_\beta$ 上的限制相同，则存在唯一的 $S$-态射 $X\to Y$，
其在每个 $U_\alpha$ 上的限制为 $u_\alpha$。

若 $S$-态射 $u\colon X\to Y$ 满足 $u(X)\subset V$，其中 $V$ 是 $Y$ 的开子集，
则把 $u$ 看成从 $X$ 到 $V$ 的态射时，它仍为 $S$-态射。

\egapar{2.5.4}
设 $S'\to S$ 是预概形态射。
对于每个 $S'$-预概形 $X$，复合态射 $X\to S'\to S$ 使 $X$ 成为 $S$-预概形。
反过来，设 $S'$ 是在 $S$ 的一个开集上诱导的预概形；
若 $X$ 是 $S$-预概形，且其结构态射 $f\colon X\to S$ 满足 $f(X)\subset S'$，
则可以把 $X$ 看成 $S'$-预概形。
在后一情形中，若 $Y$ 是一个 $S$-预概形，其结构态射也将底空间映入 $S'$，
则从 $X$ 到 $Y$ 的每个 $S$-态射也都是 $S'$-态射。

\egapar{2.5.5}
若 $X$ 是 $S$-预概形，$\varphi\colon X\to S$ 是结构态射，
则称从 $S$ 到 $X$ 的一个 $S$-态射为 $X$ 的一个\emph{$S$-截面}，
即一个满足 $\varphi\circ\psi=1_S$ 的预概形态射 $\psi\colon S\to X$。
以 $\Gamma(X/S)$ 表示 $X$ 的 $S$-截面的集合。
\transnote{本条首句原文将 $X$ 排作“un $S$-morphisme”；依（\hyperlink{para:1:2.5.1}{2.5.1}）及本句的结构态射，
应为“un $S$-préschéma”，此处明确按后者译出。}


% END chapters/ch1_s2_5.tex

\par

% BEGIN chapters/ch1_s3_1.tex
\section{预概形的积}\hypertarget{heading:1:3}{}
\subsection{预概形的和}\hypertarget{heading:1:3.1}{}
\src{104}

设 $(X_\alpha)$ 是任意一族预概形，$X$ 是其底空间 $X_\alpha$ 的拓扑和。
于是 $X$ 是一族两两不交的开子空间 $X'_\alpha$ 的并，且对于每个 $\alpha$，
存在从 $X_\alpha$ 到 $X'_\alpha$ 的同胚 $\varphi_\alpha$。
在每个 $X'_\alpha$ 上赋予层 $(\varphi_\alpha)_*(\OO_{X_\alpha})$，
显然便使 $X$ 成为预概形；称其为预概形族 $(X_\alpha)$ 的\emph{和}，记为
$\coprod_\alpha X_\alpha$。
若 $Y$ 为预概形，则映射 $f\mapsto(f\circ\varphi_\alpha)_\alpha$ 是从
$\Hom(X,Y)$ 到积集合 $\prod_\alpha\Hom(X_\alpha,Y)$ 的双射。
特别地，若各 $X_\alpha$ 是结构态射为 $\psi_\alpha$ 的 $S$-预概形，
则唯一的态射 $\psi\colon X\to S$，满足对每个 $\alpha$ 有
$\psi\circ\varphi_\alpha=\psi_\alpha$，使 $X$ 成为 $S$-预概形。
两个预概形 $X,Y$ 的和记为 $X\amalg Y$。
若 $X=\Spec(A)$、$Y=\Spec(B)$，则 $X\amalg Y$ 与 $\Spec(A\times B)$ 典范地等同。


% END chapters/ch1_s3_1.tex

\par

% BEGIN chapters/ch1_s3_2.tex
\subsection{预概形的积}\hypertarget{heading:1:3.2}{}
\src{104—108}

\egastatement{3.2.1}{定义}
给定两个 $S$-预概形 $X,Y$。设三元组 $(Z,p_1,p_2)$ 由一个 $S$-预概形 $Z$
及两个 $S$-态射 $p_1\colon Z\to X$、$p_2\colon Z\to Y$ 组成。
若对于每个 $S$-预概形 $T$，映射 $f\mapsto(p_1\circ f,p_2\circ f)$ 是从
$T$ 到 $Z$ 的 $S$-态射集合，到由一个 $S$-态射 $T\to X$ 和一个 $S$-态射 $T\to Y$
组成的有序对的集合的双射，换言之，给出双射
\[
\Hom_S(T,Z)\xrightarrow{\sim}\Hom_S(T,X)\times\Hom_S(T,Y),
\]
则称 $(Z,p_1,p_2)$ 是 $S$-预概形 $X,Y$ 的一个\emph{积}。

这就是将范畴中两个对象的积的一般概念用于 $S$-预概形范畴（T，I，1.1）。
特别地，两个 $S$-预概形的积在唯一的 $S$-同构意义下唯一。
由于这一唯一性，通常用 $X\times_S Y$（不致混淆时简记为 $X\times Y$）
表示两个 $S$-预概形 $X,Y$ 的积，而在记号中省略态射 $p_1,p_2$；
它们分别称为从 $X\times_S Y$ 到 $X,Y$ 的\emph{典范投影}。
若 $g\colon T\to X$、$h\colon T\to Y$ 是两个 $S$-态射，
则以 $(g,h)_S$ 或简记 $(g,h)$ 表示满足 $p_1\circ f=g$、$p_2\circ f=h$ 的
$S$-态射 $f\colon T\to X\times_S Y$。
若 $X',Y'$ 是两个 $S$-预概形，$p'_1,p'_2$ 是假设存在的积 $X'\times_S Y'$ 的典范投影，
$u\colon X'\to X$、$v\colon Y'\to Y$ 为两个 $S$-态射，
则以 $u\times_S v$（或简记 $u\times v$）表示 $S$-态射
$(u\circ p'_1,v\circ p'_2)_S\colon X'\times_S Y'\to X\times_S Y$。

当 $S$ 是环为 $A$ 的仿射概形时，上述记号中的 $S$ 常被换成 $A$。

\egastatement{3.2.2}{命题}
设 $X,Y,S$ 是环分别为 $B,C,A$ 的三个仿射概形。
令 $Z=\Spec(B\otimes_A C)$，并令 $p_1,p_2$ 为与典范 $A$-同态
\[
u\colon B\longrightarrow B\otimes_A C,\quad b\longmapsto b\otimes1,
\qquad
v\colon C\longrightarrow B\otimes_A C,\quad c\longmapsto1\otimes c
\]
对应的 $S$-态射（\hyperlink{para:1:2.2.4}{2.2.4}）。则 $(Z,p_1,p_2)$ 是 $X,Y$ 的一个积。

由（\hyperlink{para:1:2.2.4}{2.2.4}），只须验证：若将每个 $A$-同态 $f\colon B\otimes_A C\to L$
（$L$ 为 $A$-代数）映到有序对 $(f\circ u,f\circ v)$，便定义了双射
\[
\Hom_A(B\otimes_A C,L)\xrightarrow{\sim}\Hom_A(B,L)\times\Hom_A(C,L).
\]
\origfootnote{这里的记号 $\Hom_A$ 表示 $A$-代数同态的集合。}
这直接由定义及关系 $b\otimes c=(b\otimes1)(1\otimes c)$ 得到。

\egastatement{3.2.3}{推论}
设 $T$ 是环为 $D$ 的仿射概形，
$\alpha=({}^a\rho,\wt\rho)$、$\beta=({}^a\sigma,\wt\sigma)$ 分别是
$S$-态射 $T\to X$、$T\to Y$，其中 $\rho\colon B\to D$、$\sigma\colon C\to D$
是 $A$-同态。则 $(\alpha,\beta)_S=({}^a\tau,\wt\tau)$，其中
$\tau\colon B\otimes_A C\to D$ 是满足 $\tau(b\otimes c)=\rho(b)\sigma(c)$ 的同态。

\egastatement{3.2.4}{命题}
设 $f\colon S'\to S$ 是预概形的单态射（T，I，1.1），
$X,Y$ 是两个 $S'$-预概形，并借助 $f$ 同时看成 $S$-预概形。
则 $S$-预概形 $X,Y$ 的每个积也是 $S'$-预概形 $X,Y$ 的积，反之亦然。

设 $\varphi\colon X\to S'$、$\psi\colon Y\to S'$ 为结构态射。
若 $T$ 是 $S$-预概形，$u\colon T\to X$、$v\colon T\to Y$ 是两个 $S$-态射，
则依定义有 $f\circ\varphi\circ u=f\circ\psi\circ v=\theta$，其中 $\theta$ 为 $T$ 的结构态射。
关于 $f$ 的假设蕴含 $\varphi\circ u=\psi\circ v=\theta'$。
可见，可以将 $T$ 看成结构态射为 $\theta'$ 的 $S'$-预概形，
并将 $u,v$ 看成 $S'$-态射。结合（\hyperlink{para:1:3.2.1}{3.2.1}），命题的结论立即得到。

\egastatement{3.2.5}{推论}
设 $X,Y$ 是两个 $S$-预概形，$\varphi\colon X\to S$、$\psi\colon Y\to S$ 是其结构态射。
若 $S'$ 是 $S$ 的开子集，且 $\varphi(X)\subset S'$、$\psi(Y)\subset S'$，
则 $S$-预概形 $X,Y$ 的每个积也是 $S'$-预概形 $X,Y$ 的积，反之亦然。

只须将（\hyperlink{para:1:3.2.4}{3.2.4}）用于典范嵌入 $S'\to S$。

\egastatement{3.2.6}{定理}
给定两个 $S$-预概形 $X,Y$，其积 $X\times_S Y$ 存在。

我们分几个步骤证明。

\egastatement{3.2.6.1}{引理}
设 $(Z,p,q)$ 是 $X,Y$ 的积，$U,V$ 分别为 $X,Y$ 的开子集。
令 $W=p^{-1}(U)\cap q^{-1}(V)$。
则由 $W$ 及 $p,q$ 在 $W$ 上的限制（分别视为态射 $W\to U$、$W\to V$）
组成的三元组，是 $U,V$ 的一个积。

事实上，若 $T$ 为 $S$-预概形，则可以将 $S$-态射 $T\to W$ 与将 $T$ 映入 $W$ 的
$S$-态射 $T\to Z$ 等同。
设 $g\colon T\to U$、$h\colon T\to V$ 是任意两个 $S$-态射；
可分别把它们看成从 $T$ 到 $X,Y$ 的 $S$-态射。
依假设，存在唯一的 $S$-态射 $f\colon T\to Z$，使得 $g=p\circ f$、$h=q\circ f$。
由于 $p(f(T))\subset U$、$q(f(T))\subset V$，有
$f(T)\subset p^{-1}(U)\cap q^{-1}(V)$，断言得证。

\egastatement{3.2.6.2}{引理}
设 $Z$ 为 $S$-预概形，$p\colon Z\to X$、$q\colon Z\to Y$ 是两个 $S$-态射，
$(U_\alpha)$、$(V_\lambda)$ 分别是 $X,Y$ 的开覆盖。
假设对于每对指标 $(\alpha,\lambda)$，$S$-预概形
$W_{\alpha\lambda}=p^{-1}(U_\alpha)\cap q^{-1}(V_\lambda)$
与 $p,q$ 在其上的限制构成 $U_\alpha,V_\lambda$ 的一个积。
则 $(Z,p,q)$ 是 $X,Y$ 的一个积。

先证明：若 $f_1,f_2$ 是两个 $S$-态射 $T\to Z$，则关系
$p\circ f_1=p\circ f_2$、$q\circ f_1=q\circ f_2$ 蕴含 $f_1=f_2$。
事实上，$Z$ 是各 $W_{\alpha\lambda}$ 的并，故各 $f_1^{-1}(W_{\alpha\lambda})$
构成 $T$ 的开覆盖，各 $f_2^{-1}(W_{\alpha\lambda})$ 也如此。
此外，依假设有
\[
\begin{aligned}
f_1^{-1}(W_{\alpha\lambda})
 &=f_1^{-1}(p^{-1}(U_\alpha))\cap f_1^{-1}(q^{-1}(V_\lambda))\\
 &=f_2^{-1}(p^{-1}(U_\alpha))\cap f_2^{-1}(q^{-1}(V_\lambda))
 =f_2^{-1}(W_{\alpha\lambda}).
\end{aligned}
\]
因此，只须证明对于每对指标，$f_1,f_2$ 在此共同开集上的限制相同。
这些限制可看成从 $f_1^{-1}(W_{\alpha\lambda})$ 到 $W_{\alpha\lambda}$ 的 $S$-态射，
故断言由假设与定义（\hyperlink{para:1:3.2.1}{3.2.1}）得到。

现在设已给出两个 $S$-态射 $g\colon T\to X$、$h\colon T\to Y$。
令 $T_{\alpha\lambda}=g^{-1}(U_\alpha)\cap h^{-1}(V_\lambda)$，
它们构成 $T$ 的开覆盖。依假设，存在 $S$-态射 $f_{\alpha\lambda}$，
使 $p\circ f_{\alpha\lambda}$、$q\circ f_{\alpha\lambda}$ 分别为
$g,h$ 在 $T_{\alpha\lambda}$ 上的限制。
再证明 $f_{\alpha\lambda}$、$f_{\beta\mu}$ 在
$T_{\alpha\lambda}\cap T_{\beta\mu}$ 上的限制相同，即可完成证明。
由定义，此交集在这两个态射下的像均包含于
$W_{\alpha\lambda}\cap W_{\beta\mu}$。由于
\[
W_{\alpha\lambda}\cap W_{\beta\mu}
=p^{-1}(U_\alpha\cap U_\beta)\cap q^{-1}(V_\lambda\cap V_\mu),
\]
由（\hyperlink{para:1:3.2.6.1}{3.2.6.1}），此预概形与 $p,q$ 的限制构成
$U_\alpha\cap U_\beta$、$V_\lambda\cap V_\mu$ 的积。
而 $p\circ f_{\alpha\lambda}$ 与 $p\circ f_{\beta\mu}$ 在
$T_{\alpha\lambda}\cap T_{\beta\mu}$ 上相同，对 $q$ 亦然，
故 $f_{\alpha\lambda}$ 与 $f_{\beta\mu}$ 在此交集上相同。证毕。

\egapar{3.2.6.3}
\emph{设 $(U_\alpha)$、$(V_\lambda)$ 分别为 $X,Y$ 的开覆盖，
并假设每对 $(\alpha,\lambda)$ 所对应的积 $U_\alpha\times_S V_\lambda$ 都存在。
则 $X,Y$ 的积存在。}

将引理（\hyperlink{para:1:3.2.6.1}{3.2.6.1}）用于开集 $U_\alpha\cap U_\beta$、$V_\lambda\cap V_\mu$，
可知分别由 $X,Y$ 在它们上面诱导的 $S$-预概形的积存在。
此外，令 $i=(\alpha,\lambda)$、$j=(\beta,\mu)$。
由积的唯一性，此积到 $U_\alpha\times_S V_\lambda$
（分别到 $U_\beta\times_S V_\mu$）的某个开子集上诱导的 $S$-预概形
$W_{ij}$（分别为 $W_{ji}$）存在典范同构 $h_{ij}$（分别为 $h_{ji}$）。
因而 $f_{ij}=h_{ij}\circ h_{ji}^{-1}$ 是 $W_{ji}$ 到 $W_{ij}$ 的同构。
对于第三对指标 $k=(\gamma,\nu)$，在 $W_{ki}\cap W_{kj}$ 上有
$f_{ik}=f_{ij}\circ f_{jk}$；这是将（\hyperlink{para:1:3.2.6.1}{3.2.6.1}）用于 $U_\beta$、$V_\mu$ 中的开集
$U_\alpha\cap U_\beta\cap U_\gamma$、$V_\lambda\cap V_\mu\cap V_\nu$ 得到的。

因此，存在预概形 $Z$、其底空间的开覆盖 $(Z_i)$，及对每个 $i$ 的一个同构
$g_i\colon Z_i\to U_\alpha\times_S V_\lambda$，
使得对每对 $(i,j)$，在相应的交集上有 $f_{ij}=g_i\circ g_j^{-1}$（\hyperlink{para:1:2.3.1}{2.3.1}）；
并且 $g_i(Z_i\cap Z_j)=W_{ij}$。
若 $p_i,q_i,\theta_i$ 分别为 $S$-预概形 $U_\alpha\times_S V_\lambda$ 的
两个投影及结构态射，则立即可见在 $Z_i\cap Z_j$ 上有
$p_i\circ g_i=p_j\circ g_j$，另外两个态射也满足同样的关系。
故可以由在每个 $Z_i$ 上分别等于 $p_i\circ g_i$、$q_i\circ g_i$、$\theta_i\circ g_i$
来定义预概形态射 $p\colon Z\to X$、$q\colon Z\to Y$、$\theta\colon Z\to S$。
赋予 $\theta$ 的 $Z$ 于是成为 $S$-预概形。

现证明 $Z'_i=p^{-1}(U_\alpha)\cap q^{-1}(V_\lambda)$ 等于 $Z_i$。
对于每个 $j=(\beta,\mu)$，有
\[
Z_j\cap Z'_i=g_j^{-1}\bigl(p_j^{-1}(U_\alpha)\cap q_j^{-1}(V_\lambda)\bigr).
\]
又
\[
p_j^{-1}(U_\alpha)\cap q_j^{-1}(V_\lambda)
=p_j^{-1}(U_\alpha\cap U_\beta)\cap q_j^{-1}(V_\lambda\cap V_\mu).
\]
由（\hyperlink{para:1:3.2.6.1}{3.2.6.1}），$p_j,q_j$ 在该预概形上的限制使其成为
$U_\alpha\cap U_\beta$、$V_\lambda\cap V_\mu$ 的积。
积的唯一性于是给出
$p_j^{-1}(U_\alpha)\cap q_j^{-1}(V_\lambda)=W_{ji}$。
故对于每个 $j$，有 $Z_j\cap Z'_i=Z_j\cap Z_i$，从而 $Z'_i=Z_i$。
由（\hyperlink{para:1:3.2.6.2}{3.2.6.2}），$(Z,p,q)$ 是 $X,Y$ 的积。

\egapar{3.2.6.4}
\emph{设 $\varphi\colon X\to S$、$\psi\colon Y\to S$ 为 $X,Y$ 的结构态射，
$(S_i)$ 为 $S$ 的开覆盖，并令 $X_i=\varphi^{-1}(S_i)$、$Y_i=\psi^{-1}(S_i)$。
若各积 $X_i\times_S Y_i$ 都存在，则 $X\times_S Y$ 存在。}

由（\hyperlink{para:1:3.2.6.3}{3.2.6.3}），只须证明对任意 $i,j$，积 $X_i\times_S Y_j$ 都存在。
令
\[
X_{ij}=X_i\cap X_j=\varphi^{-1}(S_i\cap S_j),\qquad
Y_{ij}=Y_i\cap Y_j=\psi^{-1}(S_i\cap S_j).
\]
由（\hyperlink{para:1:3.2.6.1}{3.2.6.1}），积 $Z_{ij}=X_{ij}\times_S Y_{ij}$ 存在。
若 $T$ 是 $S$-预概形，$g\colon T\to X_i$、$h\colon T\to Y_j$ 为 $S$-态射，
则依 $S$-态射的定义，必有
$\varphi(g(T))=\psi(h(T))\subset S_i\cap S_j$，
故 $g(T)\subset X_{ij}$、$h(T)\subset Y_{ij}$。
由此立即可见，$Z_{ij}$ 是 $X_i,Y_j$ 的积。

\egapar{3.2.6.5}
现在可以完成定理（\hyperlink{para:1:3.2.6}{3.2.6}）的证明。
若 $S$ 为仿射概形，则 $X,Y$ 分别有由仿射开集构成的覆盖 $(U_\alpha)$、$(V_\lambda)$。
由（\hyperlink{para:1:3.2.2}{3.2.2}），$U_\alpha\times_S V_\lambda$ 存在，
故由（\hyperlink{para:1:3.2.6.3}{3.2.6.3}），$X\times_S Y$ 也存在。
若 $S$ 为任意预概形，则 $S$ 有由仿射开集构成的覆盖 $(S_i)$。
设 $\varphi\colon X\to S$、$\psi\colon Y\to S$ 为结构态射，并令
$X_i=\varphi^{-1}(S_i)$、$Y_i=\psi^{-1}(S_i)$。
依上述结果，各积 $X_i\times_{S_i}Y_i$ 存在，
故各积 $X_i\times_S Y_i$ 也存在（\hyperlink{para:1:3.2.5}{3.2.5}）。
由（\hyperlink{para:1:3.2.6.4}{3.2.6.4}），$X\times_S Y$ 存在。

\egastatement{3.2.7}{推论}
设 $Z=X\times_S Y$ 是两个 $S$-预概形的积，$p,q$ 分别为到 $X,Y$ 的投影，
$\varphi,\psi$ 分别为 $X,Y$ 的结构态射。
设 $S'$ 为 $S$ 的开子集，$U,V$ 分别为 $X,Y$ 的开子集，满足
$U\subset\varphi^{-1}(S')$、$V\subset\psi^{-1}(S')$。
则积 $U\times_{S'}V$ 与 $Z$ 在 $p^{-1}(U)\cap q^{-1}(V)$ 上诱导的预概形
（视为 $S'$-预概形）典范地等同。
此外，若 $f\colon T\to X$、$g\colon T\to Y$ 是满足 $f(T)\subset U$、$g(T)\subset V$
的 $S$-态射，则 $S'$-态射 $(f,g)_{S'}$ 与将 $(f,g)_S$ 的目标限制为
$p^{-1}(U)\cap q^{-1}(V)$ 所得的态射等同。

这由（\hyperlink{para:1:3.2.5}{3.2.5}）及（\hyperlink{para:1:3.2.6.1}{3.2.6.1}）得到。

\egapar{3.2.8}
设 $(X_\alpha)$、$(Y_\lambda)$ 为两族 $S$-预概形，
$X,Y$ 分别为这两族的和（\hyperlink{heading:1:3.1}{3.1}）。
则 $X\times_S Y$ 与预概形族 $(X_\alpha\times_S Y_\lambda)_{\alpha,\lambda}$ 的和等同；
这直接由（\hyperlink{para:1:3.2.6.3}{3.2.6.3}）得到。


% END chapters/ch1_s3_2.tex

\par

% BEGIN chapters/ch1_s3_3.tex
\subsection{积的形式性质；基预概形的变换}\hypertarget{heading:1:3.3}{}
\src{108—111}

\egapar{3.3.1}
读者可以注意到，除了（\hyperlink{para:1:3.3.13}{3.3.13}）及（\hyperlink{para:1:3.3.15}{3.3.15}），本小节中所述的一切性质，在任意范畴中，
只要陈述中涉及的积存在，都不加改变地成立。
事实上，显然对于范畴中的任意对象 $S$，都可以完全仿照（\hyperlink{heading:1:2.5}{2.5}）定义
$S$-对象与 $S$-态射的概念。

\egapar{3.3.2}
首先，$X\times_S Y$ 是 $S$-预概形范畴中关于 $X,Y$ 的协变双函子。
只须注意下图交换：
\[
\begin{tikzcd}[column sep=large,row sep=large]
X\times Y\arrow[r,"f\times g"]\arrow[d]
 &X'\times Y'\arrow[r,"f'\times g'"]\arrow[d]
 &X''\times Y''\arrow[d]\\
X\arrow[r,"f"']&X'\arrow[r,"f'"']&X''.
\end{tikzcd}
\]

\egastatement{3.3.3}{命题}
对于任意 $S$-预概形 $X$，$X\times_S S$ 的第一投影
（分别为 $S\times_S X$ 的第二投影）是从该积到 $X$ 的函子性同构，
其逆同构是 $(1_X,\varphi)_S$（分别为 $(\varphi,1_X)_S$），
其中 $\varphi\colon X\to S$ 为结构态射。
因此，在典范同构意义下，可以写成
\[
 X\times_S S=S\times_S X=X.
\]

只须证明三元组 $(X,1_X,\varphi)$ 是 $X,S$ 的积。
若 $T$ 为 $S$-预概形，则从 $T$ 到 $S$ 的唯一 $S$-态射必为结构态射 $\psi\colon T\to S$。
若 $f$ 为从 $T$ 到 $X$ 的 $S$-态射，则必有 $\psi=\varphi\circ f$，断言得证。

\egastatement{3.3.4}{推论}
设 $X,Y$ 为两个 $S$-预概形，$\varphi\colon X\to S$、$\psi\colon Y\to S$ 是其结构态射。
若将 $X$ 与 $X\times_S S$、$Y$ 与 $S\times_S Y$ 典范地等同，
则投影 $X\times_S Y\to X$、$X\times_S Y\to Y$ 分别等同于
$1_X\times\psi$、$\varphi\times1_Y$。

验证是直接的，留给读者。

\egapar{3.3.5}
可以如（\hyperlink{heading:1:3.2}{3.2}）那样定义任意有限个 $n$ 个 $S$-预概形的积。
注意到 $(X_1\times_S X_2\times\cdots\times_S X_{n-1})\times_S X_n$
满足积的定义，对 $n$ 归纳，即可由（\hyperlink{para:1:3.2.6}{3.2.6}）得到这些积的存在性。
与任意范畴中的情形一样，积的唯一性蕴含其交换性与结合性。
例如，若 $p_1,p_2,p_3$ 为 $X_1\times X_2\times X_3$ 的投影，
并将此预概形与 $(X_1\times X_2)\times X_3$ 等同，
则到 $X_1\times X_2$ 的投影与 $(p_1,p_2)_S$ 等同。

\egapar{3.3.6}
设 $S,S'$ 为两个预概形，$\varphi\colon S'\to S$ 是使 $S'$ 成为 $S$-预概形的态射。
对于每个 $S$-预概形 $X$，考虑积 $X\times_S S'$，
并以 $p,\pi'$ 分别表示它到 $X,S'$ 的投影。
赋予 $\pi'$ 后，此积是一个 $S'$-预概形；当将其如此看待时，
记为 $X_{(S')}$ 或 $X_{(\varphi)}$，称其为借助态射 $\varphi$
\emph{将基预概形从 $S$ 扩张到 $S'$ 所得的预概形}，
或 $X$ 在 $\varphi$ 下的\emph{逆像}。
注意，若 $\pi$ 为 $X$ 的结构态射，$\theta$ 为将 $X\times_S S'$ 看成 $S$-预概形时的结构态射，
则下图交换：
\[
\begin{tikzcd}[column sep=large,row sep=large]
X\arrow[d,"\pi"'] & X_{(S')}\arrow[l,"p"']\arrow[d,"\pi'"]\arrow[dl,"\theta" description]\\
S & S'\arrow[l,"\varphi"].
\end{tikzcd}
\]

\egapar{3.3.7}
沿用（\hyperlink{para:1:3.3.6}{3.3.6}）的记号。对于任意 $S$-态射 $f\colon X\to Y$，
以 $f_{(S')}$ 表示 $S'$-态射
$f\times_S1\colon X_{(S')}\to Y_{(S')}$，
称 $f_{(S')}$ 为态射 $f$ 在 $\varphi$ 下的\emph{逆像}。
因此，$X_{(S')}$ 是关于 $X$ 的协变函子，从 $S$-预概形范畴取值于 $S'$-预概形范畴。

\egapar{3.3.8}
预概形 $X_{(S')}$ 也可看成一个泛映射问题的解：
每个 $S'$-预概形 $T$ 都可借助 $\varphi$ 看成 $S$-预概形；
每个 $S$-态射 $g\colon T\to X$ 都唯一地写成 $g=p\circ f$，
其中 $f$ 为一个 $S'$-态射 $T\to X_{(S')}$。
这是将积的定义用于 $S$-态射 $g$ 及 $\psi\colon T\to S'$（$T$ 的结构态射）得到的。
\transnote{本条最后一句原文把与 $\psi$ 配对的态射写为 $f$；
按该句所应用的积的定义及前文的命名，应为 $g\colon T\to X$，此处据此译出。}

\egastatement{3.3.9}{命题（基预概形扩张的传递性）}
设 $S''$ 为预概形，$\varphi'\colon S''\to S'$ 为态射。
则对于每个 $S$-预概形 $X$，存在从 $S''$-预概形
$(X_{(\varphi)})_{(\varphi')}$ 到 $S''$-预概形 $X_{(\varphi\circ\varphi')}$ 的典范函子性同构。

事实上，设 $T$ 为 $S''$-预概形，$\psi$ 为其结构态射，$g$ 为从 $T$ 到 $X$ 的 $S$-态射；
这里将 $T$ 看成结构态射为 $\varphi\circ\varphi'\circ\psi$ 的 $S$-预概形。
由于 $T$ 也可看成结构态射为 $\varphi'\circ\psi$ 的 $S'$-预概形，
可以写 $g=p\circ g'$，其中 $g'$ 为 $S'$-态射 $T\to X_{(\varphi)}$；
再写 $g'=p'\circ g''$，其中 $g''$ 为 $S''$-态射
$T\to(X_{(\varphi)})_{(\varphi')}$：
\[
\begin{tikzcd}[column sep=large,row sep=large]
X\arrow[d,"\pi"'] & X_{(\varphi)}\arrow[l,"p"']\arrow[d,"\pi'"]
 &(X_{(\varphi)})_{(\varphi')}\arrow[l,"p'"']\arrow[d,"\pi''"]\\
S&S'\arrow[l,"\varphi"]&S''\arrow[l,"\varphi'"].
\end{tikzcd}
\]
由泛映射问题的解的唯一性，即得命题。

\begin{sloppypar}
在不致混淆时，此结果也可表述为在典范同构意义下的等式
$(X_{(S')})_{(S'')}=X_{(S'')}$，或
\end{sloppypar}
\begin{equation}
(X\times_S S')\times_{S'}S''=X\times_S S''.\tag{3.3.9.1}\label{eq:1:3.3.9.1}
\end{equation}
（\hyperlink{para:1:3.3.9}{3.3.9}）中所定义的同构的函子性，同样可用态射逆像的传递公式表述为
\begin{equation}
(f_{(S')})_{(S'')}=f_{(S'')}\tag{3.3.9.2}\label{eq:1:3.3.9.2}
\end{equation}
对每个 $S$-态射 $f\colon X\to Y$ 成立。

\egastatement{3.3.10}{推论}
若 $X,Y$ 是两个 $S$-预概形，则从 $S'$-预概形
$X_{(S')}\times_{S'}Y_{(S')}$ 到 $S'$-预概形 $(X\times_S Y)_{(S')}$
存在典范函子性同构。

事实上，在典范同构意义下，由（\hyperref[eq:1:3.3.9.1]{3.3.9.1}）及 $S$-预概形的积的结合性，有
\[
(X\times_S S')\times_{S'}(Y\times_S S')
 =X\times_S(Y\times_S S')=(X\times_S Y)\times_S S'.
\]
（\hyperlink{para:1:3.3.10}{3.3.10}）中同构的函子性表述为公式
\begin{equation}
(u_{(S')},v_{(S')})_{S'}=((u,v)_S)_{(S')}\tag{3.3.10.1}\label{eq:1:3.3.10.1}
\end{equation}
对任意一对 $S$-态射 $u\colon T\to X$、$v\colon T\to Y$ 成立。

换言之，逆像函子 $X_{(S')}$ 与取积运算交换；它也与取和运算交换（\hyperlink{para:1:3.2.8}{3.2.8}）。

\egastatement{3.3.11}{推论}
设 $Y$ 为 $S$-预概形，$f\colon X\to Y$ 为使 $X$ 成为 $Y$-预概形
（从而也成为 $S$-预概形）的态射。
则预概形 $X_{(S')}$ 与积 $X\times_Y Y_{(S')}$ 等同，
且投影 $X\times_Y Y_{(S')}\to Y_{(S')}$ 与 $f_{(S')}$ 等同。

设 $\psi\colon Y\to S$ 为 $Y$ 的结构态射。有交换图
\[
\begin{tikzcd}[column sep=large,row sep=large]
S'\arrow[d] &Y_{(S')}\arrow[l]\arrow[d]&X_{(S')}\arrow[l,"f_{(S')}"']\arrow[d]\\
S&Y\arrow[l,"\psi"]&X\arrow[l,"f"].
\end{tikzcd}
\]
而 $Y_{(S')}$ 与 $S'_{(\psi)}$ 等同，$X_{(S')}$ 与 $S'_{(\psi\circ f)}$ 等同；
结合（\hyperlink{para:1:3.3.9}{3.3.9}）及（\hyperlink{para:1:3.3.4}{3.3.4}），即得推论。

\egapar{3.3.12}
设 $f\colon X\to X'$、$g\colon Y\to Y'$ 是两个 $S$-态射，且都是预概形的单态射
（T，I，1.1）。则 $f\times_S g$ 是单态射。
事实上，设 $p,q$ 为 $X\times_S Y$ 的投影，$p',q'$ 为 $X'\times_S Y'$ 的投影，
$u,v$ 为两个 $S$-态射 $T\to X\times_S Y$。
关系 $(f\times_S g)\circ u=(f\times_S g)\circ v$ 蕴含
\[
p'\circ(f\times_S g)\circ u=p'\circ(f\times_S g)\circ v,
\]
即 $f\circ p\circ u=f\circ p\circ v$。
由于 $f$ 是单态射，有 $p\circ u=p\circ v$；
同样利用 $g$ 是单态射，得 $q\circ u=q\circ v$，故 $u=v$。

因此，对于基预概形的每个扩张 $S'\to S$，
$f_{(S')}\colon X_{(S')}\to Y_{(S')}$ 是单态射。

\egapar{3.3.13}
设 $S,S'$ 是环分别为 $A,A'$ 的两个仿射概形；
态射 $S'\to S$ 对应于环同态 $A\to A'$。
若 $X$ 为 $S$-预概形，则也以 $X_{(A')}$ 或 $X\otimes_A A'$ 表示 $S'$-预概形 $X_{(S')}$。
当 $X$ 本身是环为 $B$ 的仿射概形时，$X_{(A')}$ 为仿射概形，其环为
$B_{(A')}=B\otimes_A A'$，即将 $A$-代数 $B$ 的标量环扩张到 $A'$ 所得的环。

\egapar{3.3.14}
沿用（\hyperlink{para:1:3.3.6}{3.3.6}）的记号。对于每个 $S$-态射 $f\colon S'\to X$，
$f'=(f,1_{S'})_S$ 是一个 $S'$-态射 $S'\to X'=X_{(S')}$，
满足 $p\circ f'=f$、$\pi'\circ f'=1_{S'}$，即 $X'$ 的一个 $S'$-截面。
反过来，若 $f'$ 是这样的 $S'$-截面，则 $f=p\circ f'$ 是 $S$-态射 $S'\to X$。
由此定义了典范的一一对应
\[
\Hom_S(S',X)\xrightarrow{\sim}\Hom_{S'}(S',X').
\]
称 $f'$ 为 $f$ 的\emph{图态射}，记为 $\Gamma_f$。

\egapar{3.3.15}
给定预概形 $X$，它总可看成 $\mathbf Z$-预概形。
由（\hyperlink{para:1:3.3.14}{3.3.14}）特别可知，$X\otimes_{\mathbf Z}\mathbf Z[T]$
（$T$ 为不定元）的 $X$-截面与态射 $X\to\Spec(\mathbf Z[T])$ 一一对应。
\transnote{原文此处写作“morphismes $\mathbf Z[T]\to X$”。
按照所引用的（\hyperlink{para:1:3.3.14}{3.3.14}），应为 $X\to\Spec(\mathbf Z[T])$；
正文明确采用这一方向，而不是把环同态与预概形态射的方向混同。}
我们证明：这些 $X$-截面也与结构层 $\OO_X$ 在 $X$ 上的截面一一对应。
事实上，设 $(U_\alpha)$ 为 $X$ 的仿射开覆盖，
$u\colon X\to X\otimes_{\mathbf Z}\mathbf Z[T]$ 为 $X$-态射，
$u_\alpha$ 为其在 $U_\alpha$ 上的限制。
若 $A_\alpha$ 为仿射概形 $U_\alpha$ 的环，则
$U_\alpha\otimes_{\mathbf Z}\mathbf Z[T]$ 是环为 $A_\alpha[T]$ 的仿射概形（\hyperlink{para:1:3.2.2}{3.2.2}），
而 $u_\alpha$ 典范地对应于一个 $A_\alpha$-同态
$A_\alpha[T]\to A_\alpha$（\hyperlink{para:1:1.7.3}{1.7.3}）。
这样的同态完全由 $T$ 在 $A_\alpha$ 中的像确定，记此像为
$s_\alpha\in A_\alpha=\Gamma(U_\alpha,\OO_X)$。
写出 $u_\alpha,u_\beta$ 在仿射开集 $V\subset U_\alpha\cap U_\beta$ 上的限制相同这一条件，
立即得到 $s_\alpha,s_\beta$ 在 $V$ 上相同。
因此，族 $(s_\alpha)$ 是 $\OO_X$ 的某个全局截面 $s$ 在各 $U_\alpha$ 上的限制。
反过来，显然这样的截面定义了一族态射 $(u_\alpha)$，
它们是某个 $X$-态射 $X\to X\otimes_{\mathbf Z}\mathbf Z[T]$ 在各 $U_\alpha$ 上的限制。
此结果将在 II，1.7.12 中推广。


% END chapters/ch1_s3_3.tex

\par

% BEGIN chapters/ch1_s3_4.tex
\subsection{在预概形中取值的点；几何点}\hypertarget{heading:1:3.4}{}
\src{111—114}

\egapar{3.4.1}
设 $X$ 为预概形。对于任意预概形 $T$，也以 $X(T)$ 表示态射 $T\to X$ 的集合
$\Hom(T,X)$；此集合的元素又称为 $X$ 的\emph{在 $T$ 中取值的点}。
将每个态射 $f\colon T\to T'$ 对应到映射
$X(T')\to X(T)$、$u'\mapsto u'\circ f$，
便可见，对于固定的 $X$，$X(T)$ 是关于 $T$ 的反变函子，
从预概形范畴取值于集合范畴。
此外，每个预概形态射 $g\colon X\to Y$ 都定义一个函子性同态
$X(T)\to Y(T)$，将 $v\in X(T)$ 映到 $g\circ v$。

\egapar{3.4.2}
给定三个集合 $P,Q,R$ 及两个映射 $\varphi\colon P\to R$、$\psi\colon Q\to R$。
积集合 $P\times Q$ 中由满足 $\varphi(p)=\psi(q)$ 的有序对 $(p,q)$ 组成的子集，
称为 $P,Q$ \emph{在 $R$ 上的纤维积}（相对于 $\varphi,\psi$），记为 $P\times_R Q$。
使用（\hyperlink{para:1:3.4.1}{3.4.1}）的记号，$S$-预概形的积的定义（\hyperlink{para:1:3.2.1}{3.2.1}）也可解释为公式
\begin{equation}
(X\times_S Y)(T)=X(T)\times_{S(T)}Y(T),\tag{3.4.2.1}\label{eq:1:3.4.2.1}
\end{equation}
其中映射 $X(T)\to S(T)$、$Y(T)\to S(T)$ 分别对应于结构态射 $X\to S$、$Y\to S$。

\egapar{3.4.3}
若给定预概形 $S$，并只考虑 $S$-预概形与 $S$-态射，则也以 $X(T)_S$ 表示
$S$-态射 $T\to X$ 的集合 $\Hom_S(T,X)$；在不致混淆时，允许省略下标 $S$。
$X(T)_S$ 的元素又称为 $S$-预概形 $X$ 的\emph{在 $S$-预概形 $T$ 中取值的点}
（可能混淆时称为\emph{$S$-点}）。
特别地，$X$ 的一个 $S$-截面就是 $X$ 的一个在 $S$ 中取值的点。
这时（\hyperref[eq:1:3.4.2.1]{3.4.2.1}）也写成
\begin{equation}
(X\times_S Y)(T)_S=X(T)_S\times Y(T)_S.\tag{3.4.3.1}\label{eq:1:3.4.3.1}
\end{equation}
更一般地，若 $Z$ 是 $S$-预概形，$X,Y,T$ 是 $Z$-预概形
（从而也自然是 $S$-预概形），则
\begin{equation}
(X\times_Z Y)(T)_S=X(T)_S\times_{Z(T)}Y(T)_S.\tag{3.4.3.2}\label{eq:1:3.4.3.2}
\end{equation}
注意，若要证明由 $S$-预概形 $W$ 及两个 $S$-态射
$r\colon W\to X$、$s\colon W\to Y$ 组成的三元组 $(W,r,s)$
是 $X,Y$ 在 $Z$ 上的积，则依定义只须验证：对于每个 $S$-预概形 $T$，下图
\[
\begin{tikzcd}[column sep=large,row sep=large]
W(T)_S\arrow[r,"r'"]\arrow[d,"s'"'] &X(T)_S\arrow[d,"\varphi'"]\\
Y(T)_S\arrow[r,"\psi'"']&Z(T)_S
\end{tikzcd}
\]
使 $W(T)_S$ 成为 $X(T)_S,Y(T)_S$ 在 $Z(T)_S$ 上的纤维积。
其中 $r',s'$ 分别对应于 $r,s$，$\varphi',\psi'$ 分别对应于结构态射
$\varphi\colon X\to Z$、$\psi\colon Y\to Z$。

\egapar{3.4.4}
当上述 $T$（分别为 $S$）是环为 $B$（分别为 $A$）的仿射概形时，
在上述记号中将 $T$（分别将 $S$）换成 $B$（分别换成 $A$）。
于是，对于 $X(B)$（分别为 $X(B)_A$）的元素，称之为
$X$ 的\emph{在环 $B$ 中取值的点}（分别为 $A$-预概形 $X$ 的
\emph{在 $A$-代数 $B$ 中取值的点}）。
注意，现在 $X(B)$、$X(B)_A$ 都是关于 $B$ 的协变函子。
同样，以 $X(T)_A$ 表示 $A$-预概形 $X$ 的在 $A$-预概形 $T$ 中取值的点的集合。

\egapar{3.4.5}
特别考虑 $T=\Spec(A)$、$A$ 为局部环的情形。
这时，$X(A)$ 的元素与局部同态 $\OO_x\to A$（$x\in X$）一一对应（\hyperlink{para:1:2.4.4}{2.4.4}）；
称底空间 $X$ 的点 $x$ 为它所对应的 $A$-值点的\emph{落点}。

更特别地，称预概形 $X$ 的在一个域 $K$ 中取值的点为 $X$ 的\emph{几何点}。
因此，给出这样的点，就是给出它在底空间 $X$ 中的落点 $x$，以及 $k(x)$ 的一个扩张域 $K$。
称 $K$ 为相应几何点的\emph{取值域}，也说此几何点\emph{位于} $x$。
由此定义了映射 $X(K)\to X$，将一个 $K$-值几何点映到其落点。
\transnote{这里忠实保留本版的定义：取值域 $K$ 可以是任意域，未附加代数闭等条件。
“落点”对应原文 \textit{localité}，区别于一个给定取值域中的点本身。}

若 $S'=\Spec(K)$ 是 $S$-预概形（即把 $K$ 看成 $S$ 的某个剩余域 $k(s)$ 的扩张，$s\in S$），
且 $X$ 为 $S$-预概形，则 $X(K)_S$ 的一个元素，也称为
\emph{$X$ 的位于 $s$ 之上、在 $K$ 中取值的几何点}，
就是给出一个 $k(s)$-单同态 $k(x)\to K$，
其中 $x$ 是 $X$ 中位于 $s$ 之上的一点，故 $k(x)$ 是 $k(s)$ 的扩张。

特别地，若 $S=\Spec(K)=\{\xi\}$，则 $X$ 的 $K$-值几何点与满足 $k(x)=K$ 的
点 $x\in X$ 等同；这样的点又称为 $K$-预概形 $X$ 的\emph{在 $K$ 上有理的点}。
若 $K'$ 为 $K$ 的扩张，则 $X$ 的 $K'$-值几何点与 $X_{(K')}$ 的
在 $K'$ 上有理的点一一对应（\hyperlink{para:1:3.3.14}{3.3.14}）。

\egastatement{3.4.6}{引理}
设 $X_i$（$1\leq i\leq n$）为 $S$-预概形，$s\in S$，
且对每个 $i$，$x_i$ 是 $X_i$ 中位于 $s$ 之上的一点。
则存在 $k(s)$ 的一个扩张域 $K$，以及积
$Y=X_1\times_S X_2\times_S\cdots\times_S X_n$ 的一个 $K$-值几何点，
使它的各个投影分别位于 $x_i$。

事实上，各域 $k(x_i)$ 存在到同一个 $k(s)$ 的扩张域 $K$ 的
$k(s)$-单同态（Bourbaki，\textit{Alg.}，第 V 章，§4，命题 2）。
各复合 $k(s)\to k(x_i)\to K$ 都相同，故与 $k(x_i)\to K$ 对应的态射
$\Spec(K)\to X_i$ 都是 $S$-态射，从而确定唯一的态射 $\Spec(K)\to Y$。
若 $y$ 为 $Y$ 中相应的点，则它在每个 $X_i$ 中的投影显然为 $x_i$。

\egastatement{3.4.7}{命题}
设 $X_i$（$1\leq i\leq n$）为 $S$-预概形，对每个 $i$ 给定点 $x_i\in X_i$。
则存在 $Y=X_1\times_S X_2\times_S\cdots\times_S X_n$ 的一点 $y$，
使其第 $i$ 个投影为 $x_i$（$1\leq i\leq n$），
当且仅当各 $x_i$ 都位于 $S$ 的同一点 $s$ 之上。

此条件显然必要；引理（\hyperlink{para:1:3.4.6}{3.4.6}）说明它也充分。

换言之，以 $(X)$ 表示 $X$ 的底集合，便有典范满射
$(X\times_S Y)\to(X)\times_{(S)}(Y)$。
必须注意，此映射一般不是单射；也就是说，$X\times_S Y$ 中可能存在多个不同的点，
具有同样的投影 $x\in X$、$y\in Y$。
即使 $S,X,Y$ 是域 $k,K,K'$ 的素谱，也已如此，
因为张量积 $K\otimes_k K'$ 一般具有多个不同的素理想（参见（\hyperlink{para:1:3.4.9}{3.4.9}））。

\egastatement{3.4.8}{推论}
设 $f\colon X\to Y$ 为 $S$-态射，
$f_{(S')}\colon X_{(S')}\to Y_{(S')}$ 为由基预概形扩张 $S'\to S$ 得到的 $S'$-态射。
令 $p\colon X_{(S')}\to X$、$q\colon Y_{(S')}\to Y$ 为投影。
则对于 $X$ 的任意子集 $M$，有
\[
 q^{-1}(f(M))=f_{(S')}(p^{-1}(M)).
\]

事实上，由（\hyperlink{para:1:3.3.11}{3.3.11}），借助交换图
\[
\begin{tikzcd}[column sep=large,row sep=large]
X\arrow[d,"f"']&X_{(S')}\arrow[l,"p"']\arrow[d,"f_{(S')}"]\\
Y&Y_{(S')}\arrow[l,"q"]
\end{tikzcd}
\]
可将 $X_{(S')}$ 与积 $X\times_Y Y_{(S')}$ 等同。
由（\hyperlink{para:1:3.4.7}{3.4.7}），对于 $x\in M$、$y'\in Y_{(S')}$，关系 $q(y')=f(x)$
等价于存在 $x'\in X_{(S')}$，使 $p(x')=x$、$f_{(S')}(x')=y'$。
推论由此得到。

引理（\hyperlink{para:1:3.4.6}{3.4.6}）可进一步精确化如下。

\egastatement{3.4.9}{命题}
设 $X,Y$ 为两个 $S$-预概形，$x\in X$、$y\in Y$ 位于同一点 $s\in S$ 之上。
则 $X\times_S Y$ 中投影分别为 $x,y$ 的点的集合，与把 $k(x),k(y)$
看成 $k(s)$ 的扩张时，二者的复合扩张的类型的集合之间，存在典范的一一对应
（Bourbaki，\textit{Alg.}，第 VIII 章，§8，命题 2）。

令 $p,q$ 分别为 $X\times_S Y$ 到 $X,Y$ 的投影，
$E$ 为底空间中的子空间 $p^{-1}(x)\cap q^{-1}(y)$。
首先注意，态射 $\Spec(k(x))\to S$、$\Spec(k(y))\to S$ 分别分解为
\[
\Spec(k(x))\longrightarrow\Spec(k(s))\longrightarrow S,\qquad
\Spec(k(y))\longrightarrow\Spec(k(s))\longrightarrow S.
\]
由于 $\Spec(k(s))\to S$ 是单态射（\hyperlink{para:1:2.4.7}{2.4.7}），由（\hyperlink{para:1:3.2.4}{3.2.4}）有
\[
\begin{aligned}
P&=\Spec(k(x))\times_S\Spec(k(y))\\
 &=\Spec(k(x))\times_{\Spec(k(s))}\Spec(k(y))
 =\Spec\bigl(k(x)\otimes_{k(s)}k(y)\bigr).
\end{aligned}
\]
以 $P_0$ 表示预概形 $P$ 的底集合。
我们定义两个互逆的映射 $\alpha\colon P_0\to E$、$\beta\colon E\to P_0$。
若 $i\colon\Spec(k(x))\to X$、$j\colon\Spec(k(y))\to Y$ 是典范态射（\hyperlink{para:1:2.4.5}{2.4.5}），
则取 $\alpha$ 为与态射 $i\times_S j$ 对应的底空间映射。
另一方面，依假设，每个 $z\in E$ 都确定两个 $k(s)$-单同态
$k(x)\to k(z)$、$k(y)\to k(z)$，从而确定张量积上的 $k(s)$-代数同态
\[
 k(x)\otimes_{k(s)}k(y)\longrightarrow k(z),
\]
\transnote{此处原文称所得映射为“$k(s)$-monomorphisme produit tensoriel”。
由两个域嵌入所确定的张量积到共同域的映射，一般只是代数同态，并不必单射；
本段构造 $\Spec(k(z))\to P$ 只需要它是同态。正文在此明确译作“代数同态”。}
因而确定态射 $\Spec(k(z))\to P$。
令 $\beta(z)$ 为该态射将 $\Spec(k(z))$ 的唯一点映到 $P_0$ 所得的点。
验证 $\alpha\circ\beta$ 与 $\beta\circ\alpha$ 为恒等映射，
只须使用（\hyperlink{para:1:2.4.5}{2.4.5}）及积的定义（\hyperlink{para:1:3.2.1}{3.2.1}）。
最后，已知 $P_0$ 与 $k(x),k(y)$ 的复合扩张的类型的集合一一对应
（Bourbaki，\textit{Alg.}，第 VIII 章，§8，命题 1）。


% END chapters/ch1_s3_4.tex

\par

% BEGIN chapters/ch1_s3_5.tex
\subsection{满射与单射}\hypertarget{heading:1:3.5}{}
\src{114—117}

\egapar{3.5.1}
一般地，考虑预概形态射的一个性质 $P$，以及以下两个命题：
\begin{enumerate}[label=(\roman*)]
\item 若 $f\colon X\to X'$、$g\colon Y\to Y'$ 是具有性质 $P$ 的两个 $S$-态射，
则 $f\times_S g$ 具有性质 $P$。
\item 若 $f\colon X\to Y$ 是具有性质 $P$ 的 $S$-态射，
则由基预概形扩张得到的每个 $S'$-态射
$f_{(S')}\colon X_{(S')}\to Y_{(S')}$ 都具有性质 $P$。
\end{enumerate}
由于 $f_{(S')}=f\times_S1_{S'}$，可见：若对于每个预概形 $X$，恒等态射 $1_X$
都具有性质 $P$，则 (i) 蕴含 (ii)。
另一方面，$f\times_S g$ 是复合态射
\[
X\times_S Y\xrightarrow{f\times1_Y}X'\times_S Y
\xrightarrow{1_{X'}\times g}X'\times_S Y',
\]
故可见：若两个具有性质 $P$ 的态射的复合仍具有此性质，则 (ii) 蕴含 (i)。
此观察的第一个应用如下。

\egastatement{3.5.2}{命题}
\begin{enumerate}[label=(\roman*)]
\item 若 $f\colon X\to X'$、$g\colon Y\to Y'$ 是满射的 $S$-态射，
则 $f\times_S g$ 满射。
\item 若 $f\colon X\to Y$ 是满射的 $S$-态射，
则对于基预概形的任意扩张 $S'$，$f_{(S')}$ 满射。
\end{enumerate}
两个满射的复合仍为满射，故只须证明 (ii)；
而将（\hyperlink{para:1:3.4.8}{3.4.8}）用于 $M=X$，立即得到此结论。

\egastatement{3.5.3}{命题}
态射 $f\colon X\to Y$ 满射，当且仅当对于每个域 $K$ 及每个态射 $\Spec(K)\to Y$，
存在 $K$ 的一个扩张域 $K'$ 及一个态射 $\Spec(K')\to X$，使下图交换：
\[
\begin{tikzcd}[column sep=large,row sep=large]
X\arrow[d,"f"']&\Spec(K')\arrow[l]\arrow[d]\\
Y&\Spec(K)\arrow[l].
\end{tikzcd}
\]

条件是充分的，因为对于每个 $y\in Y$，只须将它用于与单同态 $k(y)\to K$
对应的态射 $\Spec(K)\to Y$，其中 $K$ 是 $k(y)$ 的一个扩张（\hyperlink{para:1:2.4.6}{2.4.6}）。
反过来，设 $f$ 满射，并令 $y\in Y$ 为 $\Spec(K)$ 的唯一点的像。
存在 $x\in X$，使 $f(x)=y$；考虑对应的单同态 $k(y)\to k(x)$（\hyperlink{para:1:2.2.1}{2.2.1}）。
只须取 $k(y)$ 的一个扩张域 $K'$，使存在 $k(y)$-单同态 $k(x)\to K'$ 与 $K\to K'$
（Bourbaki，\textit{Alg.}，第 V 章，§4，命题 2）。
与 $k(x)\to K'$ 对应的态射 $\Spec(K')\to X$ 即符合要求。

用（\hyperlink{para:1:3.4.5}{3.4.5}）引入的语言，也可说：$Y$ 的每个 $K$-值几何点，
都来自 $X$ 的一个在 $K$ 的某个扩张中取值的几何点。

\egastatement{3.5.4}{定义}
称预概形态射 $f\colon X\to Y$ 为\emph{普遍单射}，或一个\emph{根式态射}，
若对于每个域 $K$，相应的映射 $X(K)\to Y(K)$ 都是单射。

由定义立即得到，每个预概形的单态射（T，1.1）都是根式态射。

\egapar{3.5.5}
为了使态射 $f\colon X\to Y$ 为根式态射，只须定义（\hyperlink{para:1:3.5.4}{3.5.4}）的条件
对每个代数闭域成立。事实上，若 $K$ 是任意域，$K'$ 是 $K$ 的一个代数闭扩张，
则下图交换：
\[
\begin{tikzcd}[column sep=large,row sep=large]
X(K)\arrow[r,"\alpha"]\arrow[d,"\varphi"']&Y(K)\arrow[d,"\varphi'"]\\
X(K')\arrow[r,"\alpha'"']&Y(K').
\end{tikzcd}
\]
其中 $\varphi,\varphi'$ 来自态射 $\Spec(K')\to\Spec(K)$，
$\alpha,\alpha'$ 对应于 $f$。
而 $\varphi$ 单射，$\alpha'$ 依假设也单射，故 $\alpha$ 必单射。

\egastatement{3.5.6}{命题}
设 $f\colon X\to Y$、$g\colon Y\to Z$ 是两个预概形态射。
\begin{enumerate}[label=(\roman*)]
\item 若 $f,g$ 都是根式态射，则 $g\circ f$ 也是根式态射。
\item 反过来，若 $g\circ f$ 是根式态射，则 $f$ 也是根式态射。
\end{enumerate}
结合定义（\hyperlink{para:1:3.5.4}{3.5.4}），此命题归结为对映射
$X(K)\to Y(K)\to Z(K)$ 的相应断言，后者是显然的。

\egastatement{3.5.7}{命题}
\begin{enumerate}[label=(\roman*)]
\item 若 $S$-态射 $f\colon X\to X'$、$g\colon Y\to Y'$ 是根式态射，
则 $f\times_S g$ 也是根式态射。
\item 若 $S$-态射 $f\colon X\to Y$ 是根式态射，
则对于基预概形的每个扩张 $S'\to S$，
$f_{(S')}\colon X_{(S')}\to Y_{(S')}$ 也是根式态射。
\end{enumerate}
由（\hyperlink{para:1:3.5.1}{3.5.1}），只须证明 (i)。我们已知（\hyperref[eq:1:3.4.2.1]{3.4.2.1}）
\[
\begin{gathered}
(X\times_S Y)(K)=X(K)\times_{S(K)}Y(K),\\
(X'\times_S Y')(K)=X'(K)\times_{S(K)}Y'(K).
\end{gathered}
\]
与 $f\times_S g$ 对应的映射
$(X\times_S Y)(K)\to(X'\times_S Y')(K)$
与 $(u,v)\mapsto(f\circ u,g\circ v)$ 等同，命题由此立即得到。

\egastatement{3.5.8}{命题}
态射 $f=(\psi,\theta)\colon X\to Y$ 是根式态射，当且仅当 $\psi$ 单射，
且对于每个 $x\in X$，单同态
$\theta^x\colon k(\psi(x))\to k(x)$ 使 $k(x)$ 成为 $k(\psi(x))$ 的纯不可分扩张。

设 $f$ 为根式态射，先证明 $\psi(x_1)=\psi(x_2)=y$ 必蕴含 $x_1=x_2$。
事实上，存在 $k(y)$ 的扩张域 $K$ 及 $k(y)$-单同态
$k(x_1)\to K$、$k(x_2)\to K$
（Bourbaki，\textit{Alg.}，第 V 章，§4，命题 2）。
对应的态射 $u_1,u_2\colon\Spec(K)\to X$ 满足 $f\circ u_1=f\circ u_2$，
故依假设有 $u_1=u_2$；特别地，$x_1=x_2$。
现在借助 $\theta^x$ 将 $k(x)$ 看成 $k(\psi(x))$ 的扩张。
若它不是纯不可分扩张，则存在到 $k(\psi(x))$ 的某个代数闭扩张 $K$ 的
两个不同的 $k(\psi(x))$-单同态 $k(x)\to K$，
与之对应的两个态射 $\Spec(K)\to X$ 将违背假设。
反过来，结合（\hyperlink{para:1:2.4.6}{2.4.6}），命题中的条件显然足以保证 $f$ 为根式态射。

\egastatement{3.5.9}{推论}
若 $A$ 是环，$S$ 是 $A$ 的乘法子集，则典范态射
$\Spec(S^{-1}A)\to\Spec(A)$ 是根式态射。

事实上，此态射是单态射（\hyperlink{para:1:1.6.2}{1.6.2}）。

\egastatement{3.5.10}{推论}
设 $f\colon X\to Y$ 是根式态射，$g\colon Y'\to Y$ 为态射，
并令 $X'=X_{(Y')}=X\times_Y Y'$。
则根式态射 $f_{(Y')}$（\hyperlink{para:1:3.5.7}{3.5.7}，(ii)）将底空间 $X'$ 双射到 $g^{-1}(f(X))$。
此外，对于每个域 $K$，集合 $X'(K)$ 与以下子集等同：
$Y'(K)$ 在与 $g$ 对应的映射 $Y'(K)\to Y(K)$ 下，对子集 $X(K)\subset Y(K)$ 的逆像。

第一项断言直接由（\hyperlink{para:1:3.5.8}{3.5.8}）及（\hyperlink{para:1:3.4.8}{3.4.8}）得到；第二项由交换图
\[
\begin{tikzcd}[column sep=large,row sep=large]
X'(K)\arrow[r,hook]\arrow[d]&Y'(K)\arrow[d]\\
X(K)\arrow[r,hook]&Y(K)
\end{tikzcd}
\]
及前述纤维积的描述得到。

\egastatement{3.5.11}{注}
称预概形态射 $f=(\psi,\theta)$ 为\emph{单射}，若映射 $\psi$ 是单射。
态射 $f=(\psi,\theta)\colon X\to Y$ 为根式态射，当且仅当对于每个态射 $Y'\to Y$，
$f_{(Y')}\colon X_{(Y')}\to Y'$ 是单射；这就说明了“普遍单射态射”这一术语的缘由。
事实上，必要性由（\hyperlink{para:1:3.5.7}{3.5.7}，(ii)）及（\hyperlink{para:1:3.5.8}{3.5.8}）得到。
反过来，此条件首先蕴含 $\psi$ 单射。
若对于某个 $x\in X$，单同态 $\theta^x\colon k(\psi(x))\to k(x)$ 不是纯不可分的，
则存在 $k(\psi(x))$ 的一个扩张域 $K$，以及对应于同一个态射 $\Spec(K)\to Y$ 的
两个不同的态射 $\Spec(K)\to X$（\hyperlink{para:1:3.5.8}{3.5.8}）。
令 $Y'=\Spec(K)$，便得到 $X_{(Y')}$ 的两个不同的 $Y'$-截面（\hyperlink{para:1:3.3.14}{3.3.14}），
这与 $f_{(Y')}$ 单射的假设矛盾。


% END chapters/ch1_s3_5.tex

\par

% BEGIN chapters/ch1_s3_6.tex
\subsection{纤维}\hypertarget{heading:1:3.6}{}
\src{117—118}

\egastatement{3.6.1}{命题}
设 $f\colon X\to Y$ 为态射，$y\in Y$，
$\mathfrak a_y$ 是 $\OO_y$ 关于 $\mathfrak m_y$-预进制拓扑的一个定义理想。
则投影 $p\colon X\times_Y\Spec(\OO_y/\mathfrak a_y)\to X$
是从预概形 $X\times_Y\Spec(\OO_y/\mathfrak a_y)$ 的底空间
到纤维 $f^{-1}(y)$ 的同胚，后者赋予由 $X$ 的底空间拓扑诱导的拓扑。

由于 $\Spec(\OO_y/\mathfrak a_y)\to Y$ 是根式态射（\hyperlink{para:1:3.5.4}{3.5.4} 及 \hyperlink{para:1:2.4.7}{2.4.7}），
且 $\mathfrak m_y/\mathfrak a_y$ 依假设幂零，故
$\Spec(\OO_y/\mathfrak a_y)$ 只有一点（\hyperlink{para:1:1.1.12}{1.1.12}）。
于是已知（\hyperlink{para:1:3.5.10}{3.5.10} 及 \hyperlink{para:1:3.3.4}{3.3.4}），$p$ 将底集合与 $f^{-1}(y)$ 等同；
只须再证明 $p$ 为同胚。
由（\hyperlink{para:1:3.2.7}{3.2.7}），这是 $X,Y$ 上的局部问题，因而可假设
$X=\Spec(B)$、$Y=\Spec(A)$，其中 $B$ 为 $A$-代数。
态射 $p$ 对应于同态 $1\otimes\varphi\colon B\to B\otimes_A A'$，
这里 $A'=A_y/\mathfrak a_y$，$\varphi$ 为 $A$ 到 $A'$ 的典范映射。
而 $B\otimes_A A'$ 的每个元素都写成
\[
\sum_i b_i\otimes\frac{\varphi(a_i)}{\varphi(s)}
 =\left(\sum_i b_i a_i\otimes1\right)(1\otimes\varphi(s))^{-1},
\qquad s\notin\jj_y,
\]
因此可应用命题（\hyperlink{para:1:1.2.4}{1.2.4}）。

\egapar{3.6.2}
在本专著以下各处，当把态射的一条纤维 $f^{-1}(y)$ 看成赋有 $k(y)$-预概形结构时，
始终是指借助到 $X$ 的投影，将 $X\times_Y\Spec(k(y))$ 的结构搬运到此纤维上所得的预概形。
后一个积也写作 $X\otimes_Y k(y)$ 或 $X\otimes_{\OO_y}k(y)$。
更一般地，若 $B$ 为 $\OO_y$-代数，则以 $X\otimes_Y B$ 或
$X\otimes_{\OO_y}B$ 表示积 $X\times_Y\Spec(B)$。

按上述约定，由（\hyperlink{para:1:3.5.10}{3.5.10}），$X$ 的位于 $y$ 之上、在 $k(y)$ 的一个扩张域 $K$
中取值的点，与 $f^{-1}(y)$ 的 $K$-值点等同。

\egapar{3.6.3}
设 $f\colon X\to Y$、$g\colon Y\to Z$ 是两个态射，$h=g\circ f$ 为其复合。
对于每个 $z\in Z$，纤维 $h^{-1}(z)$ 是一个同构于下式的预概形：
\[
X\times_Z\Spec(k(z))
 =(X\times_Y Y)\times_Z\Spec(k(z))=X\times_Y g^{-1}(z).
\]
特别地，若 $U$ 为 $X$ 的开子集，则由预概形 $f^{-1}(y)$ 在 $U\cap f^{-1}(y)$ 上
诱导的预概形，同构于 $f_U^{-1}(y)$，其中 $f_U$ 是 $f$ 在 $U$ 上的限制。

\egastatement{3.6.4}{命题（纤维的传递性）}
设 $f\colon X\to Y$、$g\colon Y'\to Y$ 为两个态射；
令 $X'=X\times_Y Y'=X_{(Y')}$，$f'=f_{(Y')}\colon X'\to Y'$。
对于每个 $y'\in Y'$，若令 $y=g(y')$，则预概形 $f'^{-1}(y')$
同构于 $f^{-1}(y)\otimes_{k(y)}k(y')$。

事实上，只须注意，两个预概形
\[
(X\otimes_Y k(y))\otimes_{k(y)}k(y'),\qquad
(X\times_Y Y')\otimes_{Y'}k(y')
\]
都由（\hyperref[eq:1:3.3.9.1]{3.3.9.1}）与 $X\times_Y\Spec(k(y'))$ 典范同构。

特别地，若 $V$ 为 $y$ 在 $Y$ 中的开邻域，$f_V$ 表示 $f$ 在 $f^{-1}(V)$
上诱导的预概形上的限制，则预概形 $f^{-1}(y)$ 与 $f_V^{-1}(y)$ 典范地等同。

\egastatement{3.6.5}{命题}
设 $f\colon X\to Y$ 为态射，$y\in Y$，$Z=\Spec(\OO_y)$ 为局部预概形，
$p=(\psi,\theta)\colon X\times_Y Z\to X$ 为投影。
将 $Z$ 的底空间与 $Y$ 的一个子空间等同（\hyperlink{para:1:2.4.2}{2.4.2}）以后，
$p$ 是从 $X\times_Y Z$ 的底空间到 $X$ 的子空间 $f^{-1}(Z)$ 的同胚。
此外，对于每个 $t\in X\times_Y Z$，若令 $z=\psi(t)$，
则 $\theta_t^\sharp\colon\OO_z\to\OO_t$ 是同构。

由于 $Z$（视为 $Y$ 的子空间）包含于每个含 $y$ 的仿射开集（\hyperlink{para:1:2.4.2}{2.4.2}），
可以如（\hyperlink{para:1:3.6.1}{3.6.1}）那样归结到 $X=\Spec(A)$、$Y=\Spec(B)$ 为仿射概形的情形，
其中 $A$ 为 $B$-代数。
这时 $X\times_Y Z$ 是 $A\otimes_B B_y$ 的素谱，
而此环与 $S^{-1}A$ 典范地等同，其中 $S$ 是 $B\setminus\jj_y$ 在 $A$ 中的像（0，\hyperlink{para:0:1.5.2}{1.5.2}）。
由于 $p$ 对应于典范同态 $A\to S^{-1}A$，命题由（\hyperlink{para:1:1.6.2}{1.6.2}）得到。


% END chapters/ch1_s3_6.tex

\par

% BEGIN chapters/ch1_s3_7.tex
\subsection[应用：预概形的模理想约化]{应用：预概形的模 \texorpdfstring{$\mathfrak J$}{J} 约化}
\src{118—119}
\origfootnote{本小节使用第 I 章后文及第 II 章中的概念和结果，
此后不会用到它；本小节只面向熟悉经典代数几何的读者。}

\egapar{3.7.1}
设 $A$ 为环，$X$ 为 $A$-预概形，$\mathfrak J$ 为 $A$ 的理想。
则 $X_0=X\otimes_A(A/\mathfrak J)$ 是一个 $(A/\mathfrak J)$-预概形，
有时称它为由 $X$ 作\emph{模 $\mathfrak J$ 约化}而得。

\egapar{3.7.2}
上述术语尤其用于 $A$ 为局部环、$\mathfrak J$ 为其极大理想的情形。
这时 $X_0$ 是 $A$ 的剩余域 $k=A/\mathfrak J$ 上的预概形。
若 $A$ 还是整环，分式域为 $K$，也可考虑 $K$-预概形 $X'=X\otimes_A K$。
迄今为止，人们通常不甚严格地说 $X_0$ 是由 $X'$ 作模 $\mathfrak J$ 约化而得；
我们将不采用这种说法。
使用此种说法时，$A$ 是一个一维局部环（通常为离散赋值环），
而且或多或少地暗含：给定的 $K$-预概形 $X'$ 是某个 $K$-预概形 $P'$
的闭子预概形；实际上 $P'$ 是典型的射影空间 $\mathbf P_K^r$（参见 II，4.1.1）。
此外，$P'$ 本身具有形式 $P'=P\otimes_A K$，其中 $P$ 是给定的 $A$-预概形；
实际上，使用 II，4.1.1 的记号，$P$ 是 $A$-概形 $\mathbf P_A^r$。
用我们的语言，从 $X'$ 定义 $X_0$ 的过程如下。

考虑仿射概形 $Y=\Spec(A)$。它由两个点组成：唯一闭点 $y=\mathfrak J$，以及泛点 $(0)$。
因此，由泛点组成的单点集 $U$ 是 $Y$ 中的开集 $U=\Spec(K)$。
若 $X$ 是 $A$-预概形，即 $Y$-预概形，且 $\psi\colon X\to Y$ 为结构态射，
则 $X'=X\otimes_A K$ 就是 $X$ 在 $\psi^{-1}(U)$ 上诱导的预概形。
特别地，若 $\varphi\colon P\to Y$ 是结构态射，
则 $P'=\varphi^{-1}(U)$ 的一个闭子预概形 $X'$ 是 $P$ 的一个局部闭子预概形。
若 $P$ 为诺特预概形，例如 $A$ 诺特且 $P$ 为有限型 $A$-预概形，
则存在 $P$ 的一个不小于 $X'$ 的最小闭子预概形 $X=\overline{X'}$（\hyperlink{para:1:9.5.10}{9.5.10}）。
$X'$ 是 $X$ 在开集 $\varphi^{-1}(U)\cap X$ 上诱导的预概形，
故同构于 $X\otimes_A K$（\hyperlink{para:1:9.5.10}{9.5.10}）。
所以，$X'$ 到 $P'=P\otimes_A K$ 的浸入，使我们能典范地把 $X'$ 看成
形如 $X'=X\otimes_A K$ 的预概形，其中 $X$ 为 $A$-预概形。

于是可以考虑模 $\mathfrak J$ 约化的预概形 $X_0=X\otimes_A k$；
它也就是闭点 $y$ 的纤维 $\psi^{-1}(y)$。
迄今由于缺乏适当的术语，人们一直避免明确引入 $A$-预概形 $X$。
但应当注意，通常关于“模 $\mathfrak J$ 约化”预概形 $X_0$ 所作的一切断言，
都应当看成关于 $X$ 本身的更完整断言的推论；
只有按照这种方式解释，才能令人满意地表述和理解这些断言。
而且，所作的假设似乎总归结为关于 $X$ 本身的假设，
与事先给出 $X'$ 到某个 $P'$ 的浸入无关，因而可以得到更加内在的陈述。

\egapar{3.7.3}
最后指出一个很特殊的事实，它无疑曾使这里所讨论情形的概念澄清受到延缓。
若 $A$ 是离散赋值环，且 $X$ 在 $A$ 上固有
（当 $X$ 是某个 $\mathbf P_A^r$ 的闭子预概形时确实如此，参见 II，5.5.4），
则 $X$ 的 $A$-值点与 $X'$ 的 $K$-值点一一对应（II，7.3.8）。
正因为如此，人们常以为自己证明了关于 $X'$ 的结果，
实际却是在证明关于 $X$ 的陈述；
这些陈述以此形式表达时，在不再假设基局部环为一维的情况下仍然成立。


% END chapters/ch1_s3_7.tex

\par

% BEGIN chapters/ch1_s4_1.tex
\section{子预概形与浸入态射}\hypertarget{heading:1:4}{}
\subsection{子预概形}\hypertarget{heading:1:4.1}{}
\src{119—122}

\egapar{4.1.1}
由于拟凝聚层的概念（0，\hyperlink{para:0:5.1.3}{5.1.3}）具有局部性，预概形 $X$ 上的拟凝聚 $\OO_X$-模 $\F$
可以由下述条件刻画：对 $X$ 的每个仿射开集 $V$，$\F|V$ 同构于与某个
$\Gamma(V,\OO_X)$-模相伴的层（\hyperlink{para:1:1.4.1}{1.4.1}）。显然，在预概形 $X$ 上，结构层 $\OO_X$
是拟凝聚的；拟凝聚 $\OO_X$-模之间同态的核、余核、像，以及拟凝聚
$\OO_X$-模的归纳极限和直和，也都是拟凝聚的（\hyperlink{para:1:1.3.7}{1.3.7} 和 \hyperlink{para:1:1.3.9}{1.3.9}）。

\egastatement{4.1.2}{命题}
设 $X$ 为预概形，$\J$ 为 $\OO_X$ 中的拟凝聚理想层。
则层 $\OO_X/\J$ 的支集 $Y$ 是闭集；若以 $\OO_Y$ 表示 $\OO_X/\J$ 在 $Y$ 上的限制，
则 $(Y,\OO_Y)$ 是预概形。

显然，只需考虑 $X$ 为仿射概形的情形（\hyperlink{para:1:2.1.3}{2.1.3}），并证明此时 $Y$ 是 $X$ 的闭子集且为仿射概形。
事实上，若 $X=\Spec(A)$，则 $\OO_X=\wt A$，$\J=\wt{\mathfrak J}$，
其中 $\mathfrak J$ 为 $A$ 的理想（\hyperlink{para:1:1.4.1}{1.4.1}）。此时 $Y=V(\mathfrak J)$，
是 $X$ 的闭子集，并与环 $B=A/\mathfrak J$ 的素谱相认同（\hyperlink{para:1:1.1.11}{1.1.11}）。
此外，若 $\varphi:A\to B=A/\mathfrak J$ 为典范同态，
则直像 ${}^a\varphi_*(\wt B)$ 典范地认同于层
$\wt A/\wt{\mathfrak J}=\OO_X/\J$（\hyperlink{para:1:1.6.3}{1.6.3} 和 \hyperlink{para:1:1.3.9}{1.3.9}），证明完成。

称 $(Y,\OO_Y)$ 为由理想层 $\J$ 定义的 $(X,\OO_X)$ 的\emph{子预概形}。
这是下述一般子预概形概念的一个特例。

\egastatement{4.1.3}{定义}
称赋环空间 $(Y,\OO_Y)$ 为预概形 $(X,\OO_X)$ 的\emph{子预概形}，如果：
\begin{enumerate}[label=\arabic*°]
\item $Y$ 是 $X$ 的局部闭子空间；
\item 以 $U$ 表示 $X$ 中包含 $Y$ 且使 $Y$ 在其中为闭集的最大开集
（即 $X$ 中 $\overline Y\setminus Y$ 的补集），
则 $(Y,\OO_Y)$ 是由 $\OO_X|U$ 的一个拟凝聚理想层定义的 $(U,\OO_X|U)$ 的子预概形。
\end{enumerate}
若 $Y$ 在 $X$ 中为闭集，则称子预概形 $(Y,\OO_Y)$ 是\emph{闭的}；此时 $U=X$。

由此定义及（\hyperlink{para:1:4.1.2}{4.1.2}）立即可知，$X$ 的闭子预概形与 $\OO_X$ 的拟凝聚理想层 $\J$
之间有典范的一一对应。事实上，若两个这样的理想层 $\J,\J'$
所定义的商层具有同一个闭支集 $Y$，且 $\OO_X/\J$ 与 $\OO_X/\J'$ 在 $Y$ 上的限制相同，
则立即得到 $\J'=\J$。\transnote{这里的“相同”保留商结构中的认同，
不是只要求两个商层作为抽象环层同构。}

\egapar{4.1.4}
设 $(Y,\OO_Y)$ 是 $X$ 的子预概形，$U$ 为包含 $Y$ 且使 $Y$ 在其中闭的最大开集，
$V$ 为 $X$ 中包含于 $U$ 的开集。则 $V\cap Y$ 在 $V$ 中闭。
此外，若 $Y$ 由 $\OO_X|U$ 的拟凝聚理想层 $\J$ 定义，
则 $\J|V$ 是 $\OO_X|V$ 的拟凝聚理想层；立即可见，
$Y$ 在 $Y\cap V$ 上诱导的预概形正是由理想层 $\J|V$ 定义的 $V$ 的闭子预概形。反之：

\egastatement{4.1.5}{命题}
设 $(Y,\OO_Y)$ 是赋环空间，$Y$ 为 $X$ 的子空间。
假设存在由 $X$ 的开集组成的 $Y$ 的覆盖 $(V_\alpha)$，使得对每个 $\alpha$，
$Y\cap V_\alpha$ 在 $V_\alpha$ 中闭，且赋环空间
$(Y\cap V_\alpha,\OO_Y|(Y\cap V_\alpha))$ 是 $X$ 在 $V_\alpha$ 上诱导的预概形的闭子预概形。
则 $(Y,\OO_Y)$ 是 $X$ 的子预概形。

假设蕴含 $Y$ 在 $X$ 中局部闭，且包含 $Y$ 并使 $Y$ 在其中闭的最大开集 $U$
包含所有 $V_\alpha$。因而可以归结到 $U=X$、$Y$ 在 $X$ 中闭的情形。
定义 $\OO_X$ 的拟凝聚理想层 $\J$，使 $\J|V_\alpha$ 为定义闭子预概形
$(Y\cap V_\alpha,\OO_Y|(Y\cap V_\alpha))$ 的 $\OO_X|V_\alpha$-理想层，
并对每个不与 $Y$ 相交的开集 $W\subset X$ 令 $\J|W=\OO_X|W$。
由（\hyperlink{para:1:4.1.3}{4.1.3}）和（\hyperlink{para:1:4.1.4}{4.1.4}）立即验证，满足这些条件的理想层 $\J$ 存在且唯一，
并且它定义闭子预概形 $(Y,\OO_Y)$。

特别地，$X$ 在其一个开集上诱导的预概形是 $X$ 的子预概形。

\egastatement{4.1.6}{命题}
$X$ 的一个子预概形（分别地，闭子预概形）的子预概形（分别地，闭子预概形），
典范地认同于 $X$ 的一个子预概形（分别地，闭子预概形）。

$X$ 的局部闭子空间的局部闭子集仍是 $X$ 的局部闭子空间。
因而问题显然具有局部性（\hyperlink{para:1:4.1.5}{4.1.5}），可以假设 $X$ 仿射。
命题遂由下述典范认同得到：若 $\mathfrak J\subset\mathfrak J'$ 是环 $A$ 的两个理想，则
\[
 A/\mathfrak J'\simeq (A/\mathfrak J)/(\mathfrak J'/\mathfrak J).
\]
以下总作上述认同。

\egapar{4.1.7}
设 $Y$ 是预概形 $X$ 的子预概形，以 $\psi$ 表示底空间的典范嵌入 $Y\to X$。
已知逆像 $\psi^*(\OO_X)$ 为限制 $\OO_X|Y$（0，\hyperlink{para:0:3.7.1}{3.7.1}）。
若对每个 $y\in Y$，以 $\omega_y$ 表示典范同态
$(\OO_X)_y\to(\OO_Y)_y$，则这些同态是某个满射环层同态
$\omega:\OO_X|Y\to\OO_Y$ 在各茎上的同态。
只需在 $Y$ 上局部验证此事，因而可假设 $X$ 仿射且 $Y$ 为闭子预概形；
若此时 $\J$ 是定义 $Y$ 的理想层，则 $\omega_y$ 正是
$\OO_X|Y\to(\OO_X/\J)|Y$ 在各茎上的同态。

于是定义了赋环空间的单态射（0，\hyperlink{para:0:4.1.1}{4.1.1}）$j=(\psi,\omega^\flat)$。
它显然是预概形态射 $Y\to X$（\hyperlink{para:1:2.2.1}{2.2.1}），称为\emph{典范嵌入态射}。
若 $f:X\to Z$ 为态射，则也称复合态射 $Y\to X\to Z$ 为 $f$ 在子预概形 $Y$ 上的\emph{限制}。

\egapar{4.1.8}
按照一般定义（T，I，1.1），若预概形态射 $f:Z\to X$ 可分解为
$Z\xrightarrow{g}Y\xrightarrow{j}X$，其中 $Y$ 是 $X$ 的子预概形，
$j$ 为嵌入态射，$g$ 为预概形态射，则称 $f$ \emph{不大于} $j$。
由于 $j$ 是单态射，$g$ 必然唯一。

\egastatement{4.1.9}{命题}
态射 $f:Z\to X$ 不大于嵌入态射 $j:Y\to X$ 的充要条件是：
$f(Z)\subset Y$，且对每个 $z\in Z$，令 $y=f(z)$，由 $f$ 给出的同态
$(\OO_X)_y\to(\OO_Z)_z$ 可分解为
$(\OO_X)_y\to(\OO_Y)_y\to(\OO_Z)_z$；
等价地，前一同态的核包含 $(\OO_X)_y\to(\OO_Y)_y$ 的核。

条件显然必要。为证明充分性，必要时以一个使 $Y$ 在其中闭的开集 $U$ 代替 $X$（\hyperlink{para:1:4.1.3}{4.1.3}），
即可假设 $Y$ 是 $X$ 的闭子预概形。它由 $\OO_X$ 的拟凝聚理想层 $\J$ 定义。
写 $f=(\psi,\theta)$，并以 $\mathscr I$ 表示
$\theta^\sharp:\psi^*(\OO_X)\to\OO_Z$ 的核理想层。
由函子 $\psi^*$ 的性质（0，\hyperlink{para:0:3.7.2}{3.7.2}），假设蕴含
$(\psi^*(\J))_z\subset\mathscr I_z$ 对每个 $z\in Z$ 成立，故 $\psi^*(\J)\subset\mathscr I$。
因此 $\theta^\sharp$ 分解为
\[
 \psi^*(\OO_X)\longrightarrow
 \psi^*(\OO_X)/\psi^*(\J)=\psi^*(\OO_X/\J)
 \xrightarrow{\omega}\OO_Z,
\]
其中第一箭头为典范同态。以 $\psi':Z\to Y$ 表示与 $\psi$ 一致的连续映射，
显然 $\psi'^*(\OO_Y)=\psi^*(\OO_X/\J)$。
另一方面，$\omega$ 显然在茎上是局部同态，故
$g=(\psi',\omega^\flat)$ 是预概形态射 $Z\to Y$（\hyperlink{para:1:2.2.1}{2.2.1}）；
由以上所述 $f=j\circ g$，命题得证。

\egastatement{4.1.10}{推论}
嵌入态射 $Z\to X$ 不大于嵌入态射 $Y\to X$ 的充要条件是 $Z$ 为 $Y$ 的子预概形。

此时记作 $Z\leq Y$；这显然是 $X$ 的子预概形集合上的一个序关系。


% END chapters/ch1_s4_1.tex

\par

% BEGIN chapters/ch1_s4_2.tex
\subsection{浸入态射}\hypertarget{heading:1:4.2}{}
\src{122—124}

\egastatement{4.2.1}{定义}
若态射 $f:Y\to X$ 可分解为 $Y\xrightarrow{g}Z\xrightarrow{j}X$，
其中 $g$ 为同构，$Z$ 为 $X$ 的子预概形（分别地，闭子预概形、在一个开集上诱导的子预概形），
$j$ 为嵌入态射，则称 $f$ 为\emph{浸入}（分别地，\emph{闭浸入}、\emph{开浸入}）。

子预概形 $Z$ 和同构 $g$ 此时被唯一确定。
事实上，若 $Z'$ 为另一个子预概形，$j':Z'\to X$ 为嵌入，$g':Y\to Z'$ 为同构，
且 $j\circ g=j'\circ g'$，则 $j'=j\circ g\circ g'^{-1}$，从而 $Z'\leq Z$（\hyperlink{para:1:4.1.10}{4.1.10}）。
同样得到 $Z\leq Z'$，故 $Z'=Z$；又因 $j$ 是单态射，$g'=g$。
称 $f=j\circ g$ 为浸入 $f$ 的\emph{典范分解}，称 $Z$ 和 $g$ 为与 $f$ \emph{相伴}的子预概形和同构。

显然，浸入是预概形的单态射（\hyperlink{para:1:4.1.7}{4.1.7}），从而是根式态射（\hyperlink{para:1:3.5.4}{3.5.4}）。

\egastatement{4.2.2}{命题}
\begin{enumerate}[label=\alph*)]
\item 态射 $f=(\psi,\theta):Y\to X$ 为开浸入的充要条件是：
$\psi$ 是 $Y$ 到 $X$ 的一个开子集的同胚，且对每个 $y\in Y$，
同态 $\theta_y^\sharp:\OO_{\psi(y)}\to\OO_y$ 为双射。
\item 态射 $f=(\psi,\theta):Y\to X$ 为浸入（分别地，闭浸入）的充要条件是：
$\psi$ 是 $Y$ 到 $X$ 的一个局部闭子集（分别地，闭子集）的同胚，
且对每个 $y\in Y$，同态 $\theta_y^\sharp:\OO_{\psi(y)}\to\OO_y$ 为满射。
\end{enumerate}

 a) 条件显然必要。反之，若条件成立，$\theta^\sharp$ 给出 $\psi^*(\OO_X)$ 与 $\OO_Y$ 的同构。
而 $\psi^*(\OO_X)$ 是从 $\OO_X|\psi(Y)$ 经 $\psi^{-1}$ 搬运结构所得的层，故结论成立。

 b) 条件显然必要，现证明充分性。先设 $X$ 仿射且 $Z=\psi(Y)$ 在 $X$ 中闭。
已知 $\psi_*(\OO_Y)$ 的支集为 $Z$（0，\hyperlink{para:0:3.4.6}{3.4.6}）；若以 $\OO'_Z$ 表示它在 $Z$ 上的限制，
则赋环空间 $(Z,\OO'_Z)$ 是从 $(Y,\OO_Y)$ 经同胚 $\psi:Y\to Z$ 搬运结构所得的。
先证明 $f_*(\OO_Y)=\psi_*(\OO_Y)$ 是拟凝聚 $\OO_X$-模。
对 $x\notin Z$，$\psi_*(\OO_Y)$ 在 $x$ 的适当邻域上为零。
若 $x\in Z$，则有唯一 $y\in Y$ 使 $x=\psi(y)$。
取 $y$ 的仿射开邻域 $V$；$\psi(V)$ 在 $Z$ 中开，故为某个开集 $U\subset X$ 在 $Z$ 上的迹。
$\psi_*(\OO_Y)$ 在 $U$ 上的限制等于直像
$(\psi_V)_*(\OO_Y|V)$ 在 $U$ 上的限制，其中 $\psi_V=\psi|V$。
态射 $(\psi,\theta)$ 在 $(V,\OO_Y|V)$ 上的限制是到 $(X,\OO_X)$ 的态射，
因而形如 $({}^a\varphi,\wt\varphi)$，其中
$\varphi:A=\Gamma(X,\OO_X)\to\Gamma(V,\OO_Y)$ 为环同态（\hyperlink{para:1:1.7.3}{1.7.3}）。
所以 $(\psi_V)_*(\OO_Y|V)$ 是拟凝聚 $\OO_X$-模（\hyperlink{para:1:1.6.3}{1.6.3}）；
拟凝聚性具有局部性，故所断言的拟凝聚性得证。

此外，由 $\psi$ 的同胚性质，对每个 $y\in Y$，典范映射
$\widetilde\psi_y:(\psi_*(\OO_Y))_{\psi(y)}\to\OO_y$ 为同构（0，\hyperlink{para:0:3.4.5}{3.4.5}）。
交换图
\[
\begin{tikzcd}[column sep=large,row sep=large]
\OO_{\psi(y)} \arrow[r,"\theta_{\psi(y)}"] \arrow[d,"\widetilde\psi_y\circ\alpha_{\psi(y)}"'] &
(\psi_*(\OO_Y))_{\psi(y)} \arrow[d,"\widetilde\psi_y"]\\
(\psi^*(\OO_X))_y \arrow[r,"\theta_y^\sharp"'] & \OO_y
\end{tikzcd}
\]
的两个竖箭头都是同构（0，\hyperlink{para:0:3.7.2}{3.7.2}），所以 $\theta_y^\sharp$ 的满射性蕴含
$\theta_{\psi(y)}$ 的满射性。
由于 $\psi_*(\OO_Y)$ 的支集为 $Z=\psi(Y)$，$\theta$ 是从
$\OO_X=\wt A$ 到拟凝聚 $\OO_X$-模 $f_*(\OO_Y)$ 的满同态。
于是存在唯一同构 $\omega:\wt A/\wt{\mathfrak J}\to f_*(\OO_Y)$，
其中 $\mathfrak J$ 为 $A$ 的一个理想，使它与典范同态
$\wt A\to\wt A/\wt{\mathfrak J}$ 的复合等于 $\theta$（\hyperlink{para:1:1.3.8}{1.3.8}）。
若 $\OO_Z$ 是 $\wt A/\wt{\mathfrak J}$ 在 $Z$ 上的限制，则 $(Z,\OO_Z)$ 是 $(X,\OO_X)$ 的子预概形。
$f$ 分解为此子预概形到 $X$ 的典范嵌入与同构 $(\psi_0,\omega_0)$ 的复合，
其中 $\psi_0$ 是把 $\psi$ 看成 $Y\to Z$ 的映射，$\omega_0$ 为 $\omega$ 在 $\OO_Z$ 上的限制。

再考虑一般情形。设 $U$ 为 $X$ 的仿射开集，且 $U\cap\psi(Y)$ 非空并在 $U$ 中闭。
把 $f$ 限制到 $Y$ 在 $\psi^{-1}(U)$ 上诱导的预概形，并看作到 $U$ 的态射，即归结为第一种情形。
故限制是闭浸入 $\psi^{-1}(U)\to U$，其典范分解为 $j_U\circ g_U$，
其中 $g_U$ 是 $\psi^{-1}(U)$ 到 $U$ 的一个子预概形 $Z_U$ 的同构，$j_U:Z_U\to U$ 为嵌入。
若 $V\subset U$ 是另一个仿射开集，则 $Z_U$ 在 $V$ 上的限制 $Z'_V$ 是 $V$ 的子预概形；
$f|\psi^{-1}(V)$ 因而分解为 $j'_V\circ g'_V$。
由典范分解的唯一性（\hyperlink{para:1:4.2.1}{4.2.1}），$Z'_V=Z_V$，$g'_V=g_V$。
所以存在 $X$ 的子预概形 $Z$（\hyperlink{para:1:4.1.5}{4.1.5}），其底空间为 $\psi(Y)$，
在每个 $U\cap\psi(Y)$ 上的限制为 $Z_U$。
各 $g_U$ 是一个同构 $g:Y\to Z$ 在 $\psi^{-1}(U)$ 上的限制，且 $f=j\circ g$，
其中 $j:Z\to X$ 为典范嵌入。

\egastatement{4.2.3}{推论}
设 $X$ 为仿射概形。态射 $f=(\psi,\theta):Y\to X$ 为闭浸入的充要条件是：
$Y$ 为仿射概形，且同态 $\Gamma(\theta):\Gamma(\OO_X)\to\Gamma(\OO_Y)$ 为满射。
\transnote{原版此处写作 $\Gamma(\psi)$；这里按结构层同态 $\theta$ 写作 $\Gamma(\theta)$。}

\egastatement{4.2.4}{推论}
\begin{enumerate}[label=\alph*)]
\item 设 $f:Y\to X$ 为态射，$(V_\lambda)$ 是由 $X$ 的开集组成的 $f(Y)$ 的覆盖。
则 $f$ 为浸入（分别地，开浸入）的充要条件是：
它在每个诱导预概形 $f^{-1}(V_\lambda)$ 上的限制是到 $V_\lambda$ 的浸入（分别地，开浸入）。
\item 设 $f:Y\to X$ 为态射，$(V_\lambda)$ 是 $X$ 的开覆盖。
则 $f$ 为闭浸入的充要条件是：它在每个 $f^{-1}(V_\lambda)$ 上的限制是到 $V_\lambda$ 的闭浸入。
\end{enumerate}

写 $f=(\psi,\theta)$。在 a) 中，各 $\theta_y^\sharp$ 为满射（分别地，双射）；
在 b) 中，它们为满射。故只须分别验证 $\psi$ 是到局部闭子集（分别地，开子集）或到闭子集的同胚。
由假设，$\psi$ 显然为单射，且对每个 $y\in Y$，它把 $y$ 的任一邻域映为
$\psi(y)$ 在 $\psi(Y)$ 中的邻域。
在 a) 中，$\psi(Y)\cap V_\lambda$ 在 $V_\lambda$ 中局部闭（分别地，开），
故 $\psi(Y)$ 在 $\bigcup V_\lambda$ 中局部闭（分别地，开），从而在 $X$ 中也如此。
在 b) 中，$\psi(Y)\cap V_\lambda$ 在 $V_\lambda$ 中闭，而 $X=\bigcup V_\lambda$，
故 $\psi(Y)$ 在 $X$ 中闭。

\egastatement{4.2.5}{命题}
两个浸入（分别地，两个开浸入、两个闭浸入）的复合仍是浸入（分别地，开浸入、闭浸入）。

这由（\hyperlink{para:1:4.1.6}{4.1.6}）直接得到。


% END chapters/ch1_s4_2.tex

\par

% BEGIN chapters/ch1_s4_3.tex
\subsection{浸入的积}\hypertarget{heading:1:4.3}{}
\src{124—125}

\egastatement{4.3.1}{命题}
设 $\alpha:X'\to X$、$\beta:Y'\to Y$ 为两个 $S$-态射。
若 $\alpha,\beta$ 是浸入（分别地，开浸入、闭浸入），则 $\alpha\times_S\beta$
也是浸入（分别地，开浸入、闭浸入）。
此外，若 $\alpha$、$\beta$ 分别把 $X'$、$Y'$ 认同于 $X$、$Y$ 的子预概形 $X''$、$Y''$，
则 $\alpha\times_S\beta$ 把 $X'\times_S Y'$ 的底空间认同于 $X\times_S Y$ 的底空间中的子空间
$p^{-1}(X'')\cap q^{-1}(Y'')$，其中 $p,q$ 为到 $X,Y$ 的投影。

由定义（\hyperlink{para:1:4.2.1}{4.2.1}），可限于 $X',Y'$ 是子预概形而 $\alpha,\beta$ 是嵌入态射的情形。
对于在开集上诱导的子预概形，命题已在（\hyperlink{para:1:3.2.7}{3.2.7}）中证明。
每个子预概形都是某个开集上诱导的预概形的闭子预概形（\hyperlink{para:1:4.1.3}{4.1.3}），
故归结到 $X',Y'$ 为闭子预概形的情形。

首先证明可以假设 $S$ 仿射。取 $S$ 的仿射开覆盖 $(S_\lambda)$，
令 $X_\lambda=\varphi^{-1}(S_\lambda)$、$Y_\lambda=\psi^{-1}(S_\lambda)$，
其中 $\varphi,\psi$ 为结构态射。
$X',Y'$ 分别在 $X_\lambda\cap X'$、$Y_\lambda\cap Y'$ 上的限制 $X'_\lambda,Y'_\lambda$
是 $X_\lambda,Y_\lambda$ 的闭子预概形。
这些预概形都可看作 $S_\lambda$-预概形；相应在 $S$ 和 $S_\lambda$ 上的积相同（\hyperlink{para:1:3.2.5}{3.2.5}）。
若命题在 $S$ 仿射时成立，则 $\alpha\times_S\beta$ 在每个
$X'_\lambda\times_S Y'_\lambda$ 上的限制都是浸入（\hyperlink{para:1:3.2.7}{3.2.7}）。
又由于混合积 $X'_\lambda\times_S Y'_\mu$ 和 $X_\lambda\times_S Y_\mu$
分别认同于
\[
 (X'_\lambda\cap X'_\mu)\times_S(Y'_\lambda\cap Y'_\mu),\qquad
 (X_\lambda\cap X_\mu)\times_S(Y_\lambda\cap Y_\mu)
\]
（\hyperlink{para:1:3.2.6.4}{3.2.6.4}），在每个 $X'_\lambda\times_S Y'_\mu$ 上的限制仍为浸入。
因此由（\hyperlink{para:1:4.2.4}{4.2.4}），$\alpha\times_S\beta$ 也为浸入。

其次证明还可以假设 $X,Y$ 仿射。取它们的仿射开覆盖 $(U_i)$、$(V_j)$，
并以 $X'_i,Y'_j$ 表示 $X',Y'$ 在 $X'\cap U_i,Y'\cap V_j$ 上的限制。
它们分别是 $U_i,V_j$ 的闭子预概形。
$U_i\times_S V_j$ 认同于 $X\times_S Y$ 在 $p^{-1}(U_i)\cap q^{-1}(V_j)$ 上的限制（\hyperlink{para:1:3.2.7}{3.2.7}）；
同样，$X'_i\times_S Y'_j$ 认同于 $X'\times_S Y'$ 在
$p'^{-1}(X'_i)\cap q'^{-1}(Y'_j)$ 上的限制。
置 $\gamma=\alpha\times_S\beta$；由定义，$p\circ\gamma=\alpha\circ p'$、
$q\circ\gamma=\beta\circ q'$。因 $X'_i=\alpha^{-1}(U_i)$、$Y'_j=\beta^{-1}(V_j)$，
\[
 p'^{-1}(X'_i)\cap q'^{-1}(Y'_j)
 =\gamma^{-1}\bigl(p^{-1}(U_i)\cap q^{-1}(V_j)\bigr)
 =\gamma^{-1}(U_i\times_S V_j).
\]
遂可如第一部分那样得出结论。

现设 $X,Y,S$ 仿射，其环分别为 $B,C,A$。
$B,C$ 为 $A$-代数，$X',Y'$ 是环分别为 $B,C$ 的商代数 $B',C'$ 的仿射概形。
此外 $\alpha=({}^a\rho,\wt\rho)$、$\beta=({}^a\sigma,\wt\sigma)$，
其中 $\rho:B\to B'$、$\sigma:C\to C'$ 为典范同态（\hyperlink{para:1:1.7.3}{1.7.3}）。
已知 $X\times_S Y$、$X'\times_S Y'$ 分别为环 $B\otimes_A C$、$B'\otimes_A C'$ 的仿射概形，
且 $\alpha\times_S\beta=({}^a\tau,\wt\tau)$，其中 $\tau=\rho\otimes\sigma$（\hyperlink{para:1:3.2.2}{3.2.2} 和 \hyperlink{para:1:3.2.3}{3.2.3}）。
该同态为满射，故 $\alpha\times_S\beta$ 确为浸入。
若 $\mathfrak b=\Ker\rho$、$\mathfrak c=\Ker\sigma$，且
$u(b)=b\otimes1$、$v(c)=1\otimes c$，则 $\Ker\tau$ 为
$u(\mathfrak b)$ 和 $v(\mathfrak c)$ 所生成的理想之和。
\transnote{原文将此核简写为 $u(\mathfrak b)+v(\mathfrak c)$；
这里按生成理想理解，即
$\Im(\mathfrak b\otimes_A C)+\Im(B\otimes_A\mathfrak c)$，而不是两个集合像的简单和。}
由于 $p=({}^a u,\wt u)$、$q=({}^a v,\wt v)$，
此核在 $\Spec(B\otimes_A C)$ 中对应闭集 $p^{-1}(X')\cap q^{-1}(Y')$
（\hyperref[eq:1:1.2.2.1]{1.2.2.1} 和 \hyperlink{para:1:1.1.2}{1.1.2}，(iii)），证明完成。

\egastatement{4.3.2}{推论}
若 $f:X\to Y$ 是浸入（分别地，开浸入、闭浸入）且为 $S$-态射，
则对任意基预概形扩张 $S'\to S$，$f_{(S')}$ 也是浸入（分别地，开浸入、闭浸入）。


% END chapters/ch1_s4_3.tex

\par

% BEGIN chapters/ch1_s4_4.tex
\subsection{子预概形的逆像}\hypertarget{heading:1:4.4}{}
\src{125—126}

\egastatement{4.4.1}{命题}
设 $f:X\to Y$ 为态射，$Y'$ 为 $Y$ 的子预概形
（分别地，闭子预概形、在一个开集上诱导的预概形），$j:Y'\to Y$ 为嵌入态射。
则投影 $p:X\times_Y Y'\to X$ 是浸入（分别地，闭浸入、开浸入）；
与 $p$ 相伴的 $X$ 的子预概形的底空间为 $f^{-1}(Y')$。
此外，若 $j'$ 是此子预概形到 $X$ 的嵌入态射，
则态射 $h:Z\to X$ 满足 $f\circ h$ 不大于 $j$ 的充要条件是 $h$ 不大于 $j'$。

由于 $p=1_X\times_Y j$（\hyperlink{para:1:3.3.4}{3.3.4}），第一断言由（\hyperlink{para:1:4.3.1}{4.3.1}）得到；
第二断言是（\hyperlink{para:1:3.5.10}{3.5.10}）的特例，交换其中 $X$ 与 $Y'$ 的角色即可。
最后，若 $f\circ h=j\circ h'$，其中 $h':Z\to Y'$，
则由积的定义，$h=p\circ u$，其中 $u:Z\to X\times_Y Y'$，从而得到最后的断言。

称如此定义的 $X$ 的子预概形为 $Y$ 的子预概形 $Y'$ 在态射 $f$ 下的\emph{逆像}，
与（\hyperlink{para:1:3.3.6}{3.3.6}）中更一般的术语一致。
以后把 $f^{-1}(Y')$ 看作 $X$ 的子预概形时，总指此子预概形。

若预概形 $f^{-1}(Y')$ 与 $X$ 相同，则 $j'$ 为恒等态射，故任意 $h:Z\to X$ 均不大于 $j'$；
因此态射 $f:X\to Y$ 分解为 $X\xrightarrow{g}Y'\xrightarrow{j}Y$。

若 $y$ 为 $Y$ 的闭点，且 $Y'=\Spec(k(y))$ 为以 $\{y\}$ 为底空间的最小闭子预概形（\hyperlink{para:1:4.1.9}{4.1.9}），
则闭子预概形 $f^{-1}(Y')$ 典范地同构于（\hyperlink{para:1:3.6.2}{3.6.2}）定义的纤维 $f^{-1}(y)$，以后将它们认同。

\egastatement{4.4.2}{推论}
设 $f:X\to Y$、$g:Y\to Z$ 为态射，$h=g\circ f$。
对 $Z$ 的每个子预概形 $Z'$，$X$ 的子预概形 $f^{-1}(g^{-1}(Z'))$ 与 $h^{-1}(Z')$ 相同。

这由典范同构 $X\times_Y(Y\times_Z Z')\simeq X\times_Z Z'$ 的存在得到（\hyperref[eq:1:3.3.9.1]{3.3.9.1}）。

\egastatement{4.4.3}{推论}
设 $X',X''$ 是 $X$ 的两个子预概形，$j':X'\to X$、$j'':X''\to X$ 为嵌入态射。
则 $j'^{-1}(X'')$ 与 $j''^{-1}(X')$ 均等于子预概形序关系下 $X',X''$ 的下确界
$\inf(X',X'')$，并典范地同构于 $X'\times_X X''$。

这由（\hyperlink{para:1:4.4.1}{4.4.1}）及（\hyperlink{para:1:4.1.10}{4.1.10}）立即得到。

\egastatement{4.4.4}{推论}
设 $f:X\to Y$ 为态射，$Y',Y''$ 是 $Y$ 的子预概形。则
\[
 f^{-1}(\inf(Y',Y''))=\inf(f^{-1}(Y'),f^{-1}(Y'')).
\]
这是因为 $(X\times_Y Y')\times_X(X\times_Y Y'')$ 与
$X\times_Y(Y'\times_Y Y'')$ 之间存在典范同构（\hyperref[eq:1:3.3.9.1]{3.3.9.1}）。

\egastatement{4.4.5}{命题}
设 $f:X\to Y$ 为态射，$Y'$ 是由 $\OO_Y$ 的拟凝聚理想层 $\mathscr K$
定义的闭子预概形（\hyperlink{para:1:4.1.3}{4.1.3}）。则 $X$ 的闭子预概形 $f^{-1}(Y')$
由 $\OO_X$ 的拟凝聚理想层 $f^*(\mathscr K)\OO_X$ 定义。

问题显然在 $X,Y$ 上具有局部性。
只须注意：若 $A$ 为 $B$-代数，$\mathfrak K$ 为 $B$ 的理想，
则 $A\otimes_B(B/\mathfrak K)=A/\mathfrak K A$，再应用（\hyperlink{para:1:1.6.9}{1.6.9}）即可。
\transnote{原文此处写“$B$ 是 $A$-代数”，与紧接的 $A\otimes_B(B/\mathfrak K)$ 不合；
译文按该公式所需的结构写为“$A$ 是 $B$-代数”。}

\egastatement{4.4.6}{推论}
设 $X'$ 是由 $\OO_X$ 的拟凝聚理想层 $\J$ 定义的闭子预概形，$i:X'\to X$ 为嵌入。
$f$ 在 $X'$ 上的限制 $f\circ i$ 不大于嵌入 $j:Y'\to Y$
（即分解为 $j\circ g$，其中 $g:X'\to Y'$）的充要条件是
$f^*(\mathscr K)\OO_X\subset\J$。

只须将（\hyperlink{para:1:4.4.1}{4.4.1}）应用于 $i$，并利用（\hyperlink{para:1:4.4.5}{4.4.5}）。


% END chapters/ch1_s4_4.tex

\par

% BEGIN chapters/ch1_s4_5.tex
\subsection{局部浸入与局部同构}\hypertarget{heading:1:4.5}{}
\src{126—127}

\egastatement{4.5.1}{定义}
设 $f:X\to Y$ 为预概形态射。
若存在 $x\in X$ 的开邻域 $U$ 及 $f(x)$ 在 $Y$ 中的开邻域 $V$，
使 $f$ 在诱导预概形 $U$ 上的限制是到诱导预概形 $V$ 的闭浸入，
则称 $f$ 在 $x$ 处为\emph{局部浸入}。
若 $f$ 在 $X$ 的每一点处均为局部浸入，则称 $f$ 为\emph{局部浸入}。

\egastatement{4.5.2}{定义}
若存在 $x\in X$ 的开邻域 $U$，使态射 $f:X\to Y$ 在诱导预概形 $U$ 上的限制
为到 $Y$ 的开浸入，则称 $f$ 在 $x$ 处为\emph{局部同构}。
若 $f$ 在 $X$ 的每一点处均为局部同构，则称 $f$ 为\emph{局部同构}。

\egapar{4.5.3}
因此，浸入（分别地，闭浸入）$f:X\to Y$ 可以刻画为这样一个局部浸入：
其底空间映射是 $X$ 到 $Y$ 的一个子集（分别地，闭子集）的同胚。
开浸入则可以刻画为单射的局部同构。

\egastatement{4.5.4}{命题}
设 $X$ 为不可约预概形，$f:X\to Y$ 为单射且支配的态射。
若 $f$ 是局部浸入，则 $f$ 是浸入，且 $f(X)$ 在 $Y$ 中开。

取 $x\in X$，并取 $x$ 的开邻域 $U$、$f(x)$ 的开邻域 $V$，
使 $f|U$ 为到 $V$ 的闭浸入。
由于 $U$ 在 $X$ 中稠密，由假设 $f(U)$ 在 $Y$ 中稠密，故 $f(U)=V$，
且 $f$ 是 $U$ 到 $V$ 的同胚。
$f$ 的单射性蕴含 $f^{-1}(V)=U$，命题遂立即得到。

\egastatement{4.5.5}{命题}
\begin{enumerate}[label=(\roman*)]
\item 两个局部浸入（分别地，局部同构）的复合仍是局部浸入（分别地，局部同构）。
\item 设 $f:X\to X'$、$g:Y\to Y'$ 为两个 $S$-态射。
若它们是局部浸入（分别地，局部同构），则 $f\times_S g$ 也如此。
\item 若 $S$-态射 $f$ 是局部浸入（分别地，局部同构），
则任意基预概形扩张 $S'\to S$ 所得的 $f_{(S')}$ 也如此。
\end{enumerate}

根据（\hyperlink{para:1:3.5.1}{3.5.1}），只须证明 (i) 和 (ii)。
(i) 由闭浸入（分别地，开浸入）的传递性，以及下述事实立即得到：
若 $f$ 是 $X$ 到 $Y$ 的一个闭子集的同胚，则对每个开集 $U\subset X$，
$f(U)$ 在 $f(X)$ 中开，故存在开集 $V\subset Y$ 使
$f(U)=V\cap f(X)$，于是 $f(U)$ 在 $V$ 中闭。
\transnote{原文将浸入的传递性引用为（\hyperlink{para:1:4.2.4}{4.2.4}）；直接陈述该性质的是（\hyperlink{para:1:4.2.5}{4.2.5}）。}

为证明 (ii)，以 $p,q$ 表示 $X\times_S Y$ 的投影，$p',q'$ 表示 $X'\times_S Y'$ 的投影。
取 $z\in X\times_S Y$，令 $z'=(f\times_S g)(z)$。
由假设，可分别取 $x=p(z)$、$x'=p'(z')$、$y=q(z)$、$y'=q'(z')$ 的开邻域
$U,U',V,V'$，使 $f|U$、$g|V$ 分别为到 $U',V'$ 的闭浸入（分别地，开浸入）。
$U\times_S V$ 与 $U'\times_S V'$ 的底空间分别认同于 $z,z'$ 的开邻域
$p^{-1}(U)\cap q^{-1}(V)$ 与 $p'^{-1}(U')\cap q'^{-1}(V')$（\hyperlink{para:1:3.2.7}{3.2.7}），
因而命题由（\hyperlink{para:1:4.3.1}{4.3.1}）得到。


% END chapters/ch1_s4_5.tex

\par

% BEGIN chapters/ch1_s5_1.tex
\section{约化预概形；分离条件}\hypertarget{heading:1:5}{}
\subsection{约化预概形}\hypertarget{heading:1:5.1}{}
\src{127—131}

\egastatement{5.1.1}{命题}
设 $(X,\OO_X)$ 为预概形，$\BB$ 为拟凝聚 $\OO_X$-代数。
存在唯一拟凝聚 $\OO_X$-子模 $\N$，使得对每个 $x\in X$，$\N_x$ 为环 $\BB_x$ 的幂零根。
若 $X$ 仿射，因而 $\BB=\wt B$，其中 $B$ 是 $A(X)$-代数，则
$\N=\wt{\mathfrak N}$，其中 $\mathfrak N$ 为 $B$ 的幂零根。

问题具有局部性，故归结到证明最后的断言。
已知 $\wt{\mathfrak N}$ 是拟凝聚 $\OO_X$-模（\hyperlink{para:1:1.4.1}{1.4.1}），
其在 $x$ 处的茎为分式环 $B_x$ 的理想 $\mathfrak N_x$。
只须证明 $B_x$ 的幂零根包含于 $\mathfrak N_x$，反向包含显然。
设 $z/s$ 属于 $B_x$ 的幂零根，其中 $z\in B$，$s\notin\jj_x$。
由假设存在整数 $k$ 使 $(z/s)^k=0$，即存在 $t\notin\jj_x$ 使 $tz^k=0$。
于是 $(tz)^k=0$，故 $z/s=(tz)/(ts)$ 属于 $\mathfrak N_x$。

称如此定义的拟凝聚 $\OO_X$-模 $\N$ 为 $\OO_X$-代数 $\BB$ 的\emph{幂零根}。
特别地，以 $\N_X$ 表示 $\OO_X$ 的幂零根。

\egastatement{5.1.2}{推论}
设 $X$ 为预概形。由理想层 $\N_X$ 定义的 $X$ 的闭子预概形，
是以 $X$ 为底空间的唯一约化子预概形（0，\hyperlink{para:0:4.1.4}{4.1.4}）；
它也是以 $X$ 为底空间的最小子预概形。

由 $\N_X$ 定义的闭子预概形 $Y$ 的结构层为 $\OO_X/\N_X$。
立即可见它是约化的；又因对每个 $x$，$\N_{X,x}\subset\mm_x$，其底空间为 $X$。
为证明其余断言，注意以 $X$ 为底空间的子预概形 $Z$ 由理想层 $\J$
定义（\hyperlink{para:1:4.1.3}{4.1.3}），且每个 $x\in X$ 处 $\J_x\ne\OO_x$。
可以限于仿射情形 $X=\Spec(A)$、$\J=\wt{\mathfrak J}$，其中 $\mathfrak J$ 为 $A$ 的理想。
此时对每个 $x$，$\mathfrak J_x\subset\mm_x$，所以 $\mathfrak J$ 包含于 $A$ 的每个素理想，
从而包含于它们的交，即 $A$ 的幂零根 $\mathfrak N$。
因此 $Y$ 是以 $X$ 为底空间的最小子预概形（\hyperlink{para:1:4.1.9}{4.1.9}）。
若 $Z\ne Y$，则至少存在一个 $x$ 使 $\J_x\ne\N_{X,x}$，
于是由（\hyperlink{para:1:5.1.1}{5.1.1}），$Z$ 不是约化的。

\egastatement{5.1.3}{定义}
称以 $X$ 为底空间的唯一约化子预概形为与 $X$ \emph{相伴的约化预概形}，记为 $X_{\mathrm{red}}$。

所以，预概形 $X$ 约化意味着 $X=X_{\mathrm{red}}$。

\egastatement{5.1.4}{命题}
环 $A$ 的素谱为约化（分别地，整）预概形的充要条件是
$A$ 为约化环（分别地，整环）。

由（\hyperlink{para:1:5.1.1}{5.1.1}）立即可知，$\mathfrak N=(0)$ 是 $X=\Spec(A)$ 约化的充要条件；
关于整环的断言遂由（\hyperlink{para:1:1.1.13}{1.1.13}）得到。
由于整环的每个非零分式环仍为整环，由（\hyperlink{para:1:5.1.4}{5.1.4}）可知，
对每个局部整预概形 $X$ 和每个 $x\in X$，$\OO_x$ 都是整环。
当 $X$ 的底空间局部诺特时，逆命题成立。
事实上，此时 $X$ 约化。若 $U$ 是 $X$ 中作为拓扑空间为诺特的仿射开集，
则 $U$ 只有有限个不可约分支，所以它的环 $A$ 只有有限个极小素理想（\hyperlink{para:1:1.1.14}{1.1.14}）。
若两个这样的分支有公共点 $x$，则 $\OO_x$ 至少有两个不同的极小素理想，因而不是整环。
所以各分支 $U_i$ 是两两不交的开集，并且都是整的。

\egapar{5.1.5}
设 $f=(\psi,\theta):X\to Y$ 为预概形态射。
同态 $\theta_x^\sharp:\OO_{\psi(x)}\to\OO_x$ 将幂零元映为幂零元；
因此取商后得到同态
\[
 \omega:\psi^*(\OO_Y/\N_Y)\longrightarrow\OO_X/\N_X.
\]
显然，各茎上的同态
$\omega_x:\OO_{Y_{\mathrm{red}},\psi(x)}\to\OO_{X_{\mathrm{red}},x}$ 都是局部同态。
故 $(\psi,\omega^\flat)$ 是预概形态射
$X_{\mathrm{red}}\to Y_{\mathrm{red}}$，记为 $f_{\mathrm{red}}$，
称为与 $f$ \emph{相伴的约化态射}。
对两个态射 $f:X\to Y$、$g:Y\to Z$，立即有
$(g\circ f)_{\mathrm{red}}=g_{\mathrm{red}}\circ f_{\mathrm{red}}$。
因此 $X_{\mathrm{red}}$ 是 $X$ 的协变函子。

上述定义说明交换图
\[
\begin{tikzcd}
X_{\mathrm{red}}\arrow[r,"f_{\mathrm{red}}"]\arrow[d] &Y_{\mathrm{red}}\arrow[d]\\
X\arrow[r,"f"']&Y
\end{tikzcd}
\]
成立，竖箭头为嵌入态射。换言之，$X_{\mathrm{red}}\to X$ 为函子态射。
特别地，若 $X$ 约化，则每个态射 $f:X\to Y$ 都分解为
$X\xrightarrow{f_{\mathrm{red}}}Y_{\mathrm{red}}\to Y$；
即 $f$ 不大于嵌入态射 $Y_{\mathrm{red}}\to Y$。

\egastatement{5.1.6}{命题}
若态射 $f:X\to Y$ 为满射、根式态射、浸入、闭浸入、开浸入、局部浸入或局部同构，
则 $f_{\mathrm{red}}$ 具有相应性质。
反之，若 $f_{\mathrm{red}}$ 为满射（分别地，根式态射），则 $f$ 也如此。

对于满射，命题显然。对于根式态射，由 $X$ 和 $X_{\mathrm{red}}$ 在每个点 $x$
处的剩余域 $k(x)$ 相同及（\hyperlink{para:1:3.5.8}{3.5.8}）即可得到。
最后，若 $f=(\psi,\theta)$ 为浸入、闭浸入或局部浸入
（分别地，开浸入或局部同构），则命题由下述事实得到：
若 $\theta_x^\sharp$ 为满射（分别地，双射），
则对两端取幂零根的商后所得的同态也如此（\hyperlink{para:1:5.1.2}{5.1.2} 和 \hyperlink{para:1:4.2.2}{4.2.2}；另参见 \hyperlink{para:1:5.5.12}{5.5.12}）。

\egastatement{5.1.7}{命题}
若 $X,Y$ 为两个 $S$-预概形，则
$X_{\mathrm{red}}\times_{S_{\mathrm{red}}}Y_{\mathrm{red}}$
与 $X_{\mathrm{red}}\times_S Y_{\mathrm{red}}$ 相同，
并典范地认同于 $X\times_S Y$ 的一个具有相同底空间的子预概形。

$X_{\mathrm{red}}\times_S Y_{\mathrm{red}}$ 与 $X\times_S Y$ 的一个具有相同底空间的子预概形的典范认同，
由（\hyperlink{para:1:4.3.1}{4.3.1}）得到。
另一方面，结构态射 $X_{\mathrm{red}}\to S$、$Y_{\mathrm{red}}\to S$
均通过 $S_{\mathrm{red}}$ 分解（\hyperlink{para:1:5.1.5}{5.1.5}）；由于 $S_{\mathrm{red}}\to S$ 为单态射，
第一断言由（\hyperlink{para:1:3.2.4}{3.2.4}）得到。

\egastatement{5.1.8}{推论}
预概形 $(X\times_S Y)_{\mathrm{red}}$ 与
$(X_{\mathrm{red}}\times_{S_{\mathrm{red}}}Y_{\mathrm{red}})_{\mathrm{red}}$ 典范地认同。

这由（\hyperlink{para:1:5.1.2}{5.1.2}）和（\hyperlink{para:1:5.1.7}{5.1.7}）得到。
应当注意，即使 $X,Y$ 约化，$X\times_S Y$ 也未必约化，
因为两个约化代数的张量积可能具有幂零元。

\egastatement{5.1.9}{命题}
设 $X$ 为预概形，$\J$ 为 $\OO_X$ 的拟凝聚理想层，且对某个整数 $n>0$ 有 $\J^n=0$。
以 $X_0$ 表示闭子预概形 $(X,\OO_X/\J)$。
则 $X$ 为仿射概形的充要条件是 $X_0$ 为仿射概形。

条件显然必要，现证明充分性。置 $X_k=(X,\OO_X/\J^{k+1})$，
只须对 $k$ 归纳证明各 $X_k$ 仿射，因而归结到 $\J^2=0$ 的情形。
置
\[
 A=\Gamma(X,\OO_X),\qquad
 A_0=\Gamma(X_0,\OO_{X_0})=\Gamma(X,\OO_X/\J).
\]
典范同态 $\OO_X\to\OO_X/\J$ 诱导环同态 $\varphi:A\to A_0$。
下面将证明 $\varphi$ 为满射，即序列
\begin{equation}
0\longrightarrow\Gamma(X,\J)\longrightarrow\Gamma(X,\OO_X)
\longrightarrow\Gamma(X,\OO_X/\J)\longrightarrow0
\tag{5.1.9.1}\label{eq:1:5.1.9.1}
\end{equation}
正合。先假设这一点成立，并证明它蕴含命题。

注意 $\mathfrak K=\Gamma(X,\J)$ 是 $A$ 中平方为零的理想，
所以是 $A_0=A/\mathfrak K$-模。
由假设 $X_0=\Spec(A_0)$；$X_0$ 与 $X$ 的底空间相同，故
$\mathfrak K=\Gamma(X_0,\J)$。
又因 $\J^2=0$，$\J$ 是拟凝聚 $(\OO_X/\J)$-模，
因而 $\J\simeq\wt{\mathfrak K}$，且对每个 $x\in X_0$ 有 $\J_x=\mathfrak K_x$（\hyperlink{para:1:1.4.1}{1.4.1}）。

置 $X'=\Spec(A)$，考虑与恒等映射 $A\to\Gamma(X,\OO_X)$
相对应的预概形态射 $f=(\psi,\theta):X\to X'$（\hyperlink{para:1:2.2.4}{2.2.4}）。
对 $X$ 的每个仿射开集 $V$，图
\[
\begin{tikzcd}[column sep=large]
A\arrow[r]\arrow[d]&\Gamma(V,\OO_X|V)\arrow[d]\\
A_0=A/\mathfrak K\arrow[r]&\Gamma(V,\OO_{X_0}|V)
\end{tikzcd}
\]
交换，因此图
\[
\begin{tikzcd}
X'&X\arrow[l,"f"']\\
X'_0\arrow[u,"j'"]&X_0\arrow[l,"f_0"]\arrow[u,"j"']
\end{tikzcd}
\]
交换，其中 $X'_0$ 是由拟凝聚理想层 $\wt{\mathfrak K}$ 定义的 $X'$ 的闭子预概形，
$j,j'$ 为典范嵌入。
由于 $X_0$ 仿射，$f_0$ 是同构；又因 $j,j'$ 的底空间映射为恒等映射，
首先得到 $\psi:X\to X'$ 为同胚。
此外，$\J_x=\mathfrak K_x$ 表明
$\theta^\sharp:\psi^*(\OO_{X'})\to\OO_X$ 的限制给出
$\psi^*(\wt{\mathfrak K})\to\J$ 的同构。
另一方面，由 $f_0$ 为同构，取商后 $\theta^\sharp$ 诱导同构
$\psi^*(\OO_{X'}/\wt{\mathfrak K})\to\OO_X/\J$。
由五引理（M，I，1.1）立即得到 $\theta^\sharp$ 本身为同构，故 $f$ 为同构，$X$ 仿射。

所以只剩证明（\hyperref[eq:1:5.1.9.1]{5.1.9.1}）的正合性；它可由 $H^1(X,\J)=0$ 得到。
而 $H^1(X,\J)=H^1(X_0,\J)$，且已知 $\J$ 为拟凝聚 $\OO_{X_0}$-模。
因此所需结论由下述引理得到。

\egastatement{5.1.9.2}{引理}
若 $Y$ 为仿射概形，$\F$ 为拟凝聚 $\OO_Y$-模，则 $H^1(Y,\F)=0$。

此引理将在第 III 章 §1 中，作为更一般的定理
“对每个 $i>0$，$H^i(Y,\F)=0$”的推论予以证明。
为给出独立的证明，注意 $H^1(Y,\F)$ 认同于以 $\F$ 为核的 $\OO_Y$ 的扩张类所构成的模
$\operatorname{Ext}^1_{\OO_Y}(Y;\OO_Y,\F)$（T，4.2.3）。
所以只须证明每个这样的扩张 $\G$ 平凡。
对每个 $y\in Y$，存在其邻域 $V$ 使
$\G|V\simeq\F|V\oplus\OO_Y|V$（0，\hyperlink{para:0:5.4.9}{5.4.9}），故 $\G$ 为拟凝聚 $\OO_Y$-模。
若 $A$ 为 $Y$ 的环，则 $\F=\wt M$、$\G=\wt N$，其中 $M,N$ 为 $A$-模。
由假设，$N$ 是以 $M$ 为核的 $A$-模 $A$ 的扩张（\hyperlink{para:1:1.3.11}{1.3.11}）。
这样的扩张必然平凡，引理以及（\hyperlink{para:1:5.1.9}{5.1.9}）遂得证。

\egastatement{5.1.10}{推论}
设预概形 $X$ 的幂零根 $\N_X$ 为幂零理想层。
则 $X$ 为仿射概形的充要条件是 $X_{\mathrm{red}}$ 为仿射概形。


% END chapters/ch1_s5_1.tex

\par

% BEGIN chapters/ch1_s5_2.tex
\subsection{具有给定底空间的子预概形的存在性}\hypertarget{heading:1:5.2}{}
\src{131—132}

\egastatement{5.2.1}{命题}
对于预概形 $X$ 的底空间中的每个局部闭子空间 $Y$，
存在唯一以 $Y$ 为底空间的约化子预概形。

唯一性由（\hyperlink{para:1:5.1.2}{5.1.2}）得到，故只须证明存在性。
若 $X$ 仿射、环为 $A$，且 $Y$ 在 $X$ 中闭，则命题立即成立：
$\jj(Y)$ 是满足 $V(\mathfrak a)=Y$ 的最大理想 $\mathfrak a\subset A$，
且它等于自身的根（\hyperlink{para:1:1.1.4}{1.1.4}，(i)），所以 $A/\jj(Y)$ 是约化环。

一般情形下，对每个使 $U\cap Y$ 在 $U$ 中闭的仿射开集 $U\subset X$，
考虑由 $A(U)$ 的理想 $\jj(U\cap Y)$ 所相伴的理想层定义的闭子预概形 $Y_U$；它是约化的。
若 $V\subset U$ 是仿射开集，则 $Y_V$ 是 $Y_U$ 在 $V\cap Y$ 上诱导的预概形。
事实上，该诱导预概形是 $V$ 的约化闭子预概形，底空间为 $V\cap Y$，
所以 $Y_V$ 的唯一性给出此断言。

\egastatement{5.2.2}{命题}
设 $X$ 为约化预概形，$f:X\to Y$ 为态射，$Z$ 为 $Y$ 的闭子预概形，且 $f(X)\subset Z$。
则 $f$ 分解为 $X\to Z\xrightarrow{j}Y$，其中 $j$ 为嵌入态射。

由假设，$X$ 的闭子预概形 $f^{-1}(Z)$ 的底空间为整个 $X$（\hyperlink{para:1:4.4.1}{4.4.1}）。
由于 $X$ 约化，该闭子预概形等于 $X$（\hyperlink{para:1:5.1.2}{5.1.2}），故命题由（\hyperlink{para:1:4.4.1}{4.4.1}）得到。

\egastatement{5.2.3}{推论}
设 $X$ 为预概形 $Y$ 的约化子预概形，$Z$ 为以 $\overline X$ 为底空间的 $Y$ 的约化闭子预概形。
则 $X$ 是 $Z$ 在一个开集上诱导的子预概形。

事实上，存在 $Y$ 的开集 $U$ 使 $X=U\cap\overline X$。
由（\hyperlink{para:1:5.2.2}{5.2.2}），$X$ 是 $Z$ 的约化子预概形；
因而由唯一性（\hyperlink{para:1:5.2.1}{5.2.1}），它正是 $Z$ 在开子空间 $X$ 上诱导的子预概形。

\egastatement{5.2.4}{推论}
设 $f:X\to Y$ 为态射，$X'$、$Y'$ 分别为由拟凝聚理想层
$\J\subset\OO_X$、$\mathscr K\subset\OO_Y$ 定义的闭子预概形。
若 $X'$ 约化且 $f(X')\subset Y'$，则 $f^*(\mathscr K)\OO_X\subset\J$。

根据（\hyperlink{para:1:5.2.2}{5.2.2}），$f|X'$ 分解为 $X'\to Y'\to Y$，再应用（\hyperlink{para:1:4.4.6}{4.4.6}）即可。


% END chapters/ch1_s5_2.tex

\par

% BEGIN chapters/ch1_s5_3.tex
\subsection{对角线；态射的图}\hypertarget{heading:1:5.3}{}
\src{132—135}

\egapar{5.3.1}
设 $X$ 为 $S$-预概形。
称 $S$-态射 $(1_X,1_X)_S$ 为 $X$ 到 $X\times_S X$ 的\emph{对角态射}，
记为 $\Delta_{X|S}$、$\Delta_X$，或在无混淆时记为 $\Delta$。
它是满足
\begin{equation}
 p_1\circ\Delta_X=p_2\circ\Delta_X=1_X
 \tag{5.3.1.1}\label{eq:1:5.3.1.1}
\end{equation}
的唯一 $S$-态射，其中 $p_1,p_2$ 为 $X\times_S X$ 的投影（\hyperlink{para:1:3.2.1}{3.2.1}）。
若 $f:T\to X$、$g:T\to Y$ 为两个 $S$-态射，立即可验证
\begin{equation}
 (f,g)_S=(f\times_S g)\circ\Delta_{T|S}.
 \tag{5.3.1.2}\label{eq:1:5.3.1.2}
\end{equation}
读者应当注意，上述定义及（\hyperlink{para:1:5.3.1}{5.3.1}）至（\hyperlink{para:1:5.3.8}{5.3.8}）的结果在任意范畴中都成立，
只要其中出现的积在该范畴中存在。

\egastatement{5.3.2}{命题}
设 $X,Y$ 为两个 $S$-预概形。
若将 $(X\times Y)\times(X\times Y)$ 典范地认同于
$(X\times X)\times(Y\times Y)$，则 $\Delta_{X\times Y}$ 认同于 $\Delta_X\times\Delta_Y$。

事实上，若 $p_1,q_1$ 分别为 $X\times X\to X$、$Y\times Y\to Y$ 的第一投影，
则 $(X\times Y)\times(X\times Y)\to X\times Y$ 的第一投影认同于 $p_1\times q_1$，且
\[
(p_1\times q_1)\circ(\Delta_X\times\Delta_Y)
=(p_1\circ\Delta_X)\times(q_1\circ\Delta_Y)=1_{X\times Y}.
\]
第二投影的论证相同。

\egastatement{5.3.4}{推论}
对任意基预概形扩张 $S'\to S$，$\Delta_{X_{(S')}}$ 典范地认同于 $(\Delta_X)_{(S')}$。
\transnote{原版在（\hyperlink{para:1:5.3.2}{5.3.2}）之后直接编号为（\hyperlink{para:1:5.3.4}{5.3.4}）；此处保留原版编号，未补造（5.3.3）。}

只须注意 $(X\times_S X)_{(S')}$ 典范地认同于
$X_{(S')}\times_{S'}X_{(S')}$（\hyperlink{para:1:3.3.10}{3.3.10}）。

\egastatement{5.3.5}{命题}
设 $X,Y$ 为两个 $S$-预概形，$\varphi:S\to T$ 为态射，因而使每个 $S$-预概形成为 $T$-预概形。
设 $f:X\to S$、$g:Y\to S$ 为结构态射，$p,q$ 为 $X\times_S Y$ 的投影，
$\pi=f\circ p=g\circ q$ 为其到 $S$ 的结构态射。
则图
\begin{equation}
\begin{tikzcd}[column sep=large]
X\times_S Y\arrow[r,"{(p,q)_T}"]\arrow[d,"\pi"']&X\times_T Y\arrow[d,"f\times_T g"]\\
S\arrow[r,"\Delta_{S|T}"']&S\times_T S
\end{tikzcd}
\tag{5.3.5.1}\label{eq:1:5.3.5.1}
\end{equation}
交换，并把 $X\times_S Y$ 认同于两个 $(S\times_T S)$-预概形 $S$ 与 $X\times_T Y$ 的积，
其投影分别为 $\pi$ 与 $(p,q)_T$。

\begin{sloppypar}
根据（\hyperlink{para:1:3.4.3}{3.4.3}），取任意 $T$-预概形 $Z$，将 $X,Y,S$ 分别替换为
$X(Z)_T,Y(Z)_T,S(Z)_T$，即归结到证明集合范畴中对应的命题。
在集合范畴中验证立即可作，留给读者。
\end{sloppypar}


\egastatement{5.3.6}{推论}
令 $P=S\times_T S$，则态射 $(p,q)_T$ 认同于
$1_{X\times_T Y}\times_P\Delta_S$。

这由（\hyperlink{para:1:5.3.5}{5.3.5}）和（\hyperlink{para:1:3.3.4}{3.3.4}）得到。

\egastatement{5.3.7}{推论}
若 $f:X\to Y$ 为 $S$-态射，则图
\[
\begin{tikzcd}[column sep=large]
X\arrow[r,"{(1_X,f)_S}"]\arrow[d,"f"']&X\times_S Y\arrow[d,"f\times_S1_Y"]\\
Y\arrow[r,"\Delta_Y"']&Y\times_S Y
\end{tikzcd}
\]
交换，并把 $X$ 认同于两个 $(Y\times_S Y)$-预概形 $Y$ 与 $X\times_S Y$ 的积。

只须在（\hyperlink{para:1:5.3.5}{5.3.5}）中以 $Y$ 代替 $S$、以 $S$ 代替 $T$，
并注意 $X\times_Y Y=X$（\hyperlink{para:1:3.3.3}{3.3.3}）。

\egastatement{5.3.8}{命题}
态射 $f:X\to Y$ 为预概形单态射的充要条件是
$\Delta_{X|Y}$ 为 $X$ 到 $X\times_Y X$ 的同构。

$f$ 为单态射意味着：对每个 $Y$-预概形 $Z$，对应映射
$f':X(Z)_Y\to Y(Z)_Y$ 为单射。
由于 $Y(Z)_Y$ 只有一个元素，这意味着 $X(Z)_Y$ 至多有一个元素。
\transnote{原文此处将 $X(Z)_Y$ 也说成“只有一个元素”；应允许空集，
因此这里明确写作“至多有一个元素”。}
这又等价于说 $X(Z)_Y\times X(Z)_Y$ 经对角映射典范地同构于 $X(Z)_Y$。
前一集合等于 $(X\times_Y X)(Z)_Y$（\hyperref[eq:1:3.4.3.1]{3.4.3.1}），故此条件正表示 $\Delta_{X|Y}$ 为同构。

\NewNote{\textbf{作者勘误提示：}作者确认此处原证明没有验证对角像局部闭；书末保留原作者的补证。结论与证明缺步应分开判断。 参见\hyperlink{errIII:10}{《勘误与增补（清单二）》第10条}。}
\egastatement{5.3.9}{命题}
对角态射 $\Delta_X$ 是 $X$ 到 $X\times_S X$ 的浸入。

底空间上的连续映射 $p_1$、$\Delta_X$ 满足 $p_1\circ\Delta_X=1_X$，
所以 $\Delta_X$ 是 $X$ 到 $\Delta_X(X)$ 的同胚。
同样，对应于 $p_1$ 与 $\Delta_X$ 的同态的复合
$\OO_x\to\OO_{\Delta_X(x)}\to\OO_x$ 为恒等同态，
所以对应于 $\Delta_X$ 的同态为满射；命题遂由（\hyperlink{para:1:4.2.2}{4.2.2}）得到。
\transnote{原文略去（\hyperlink{para:1:4.2.2}{4.2.2}）所需的像的局部闭性；这一步可用本卷已有结果补足。
令 $f:X\to S$ 为结构态射。对每个 $x\in X$，取 $f(x)$ 的仿射开邻域 $V$，
以及包含于 $f^{-1}(V)$ 的 $x$ 的仿射开邻域 $U$。
在开子预概形 $U\times_S U=U\times_V U$ 内，$\Delta_X(U)$ 是闭子空间，
因为乘法同态 $A(U)\otimes_{A(V)}A(U)\to A(U)$ 为满射（\hyperlink{para:1:3.2.2}{3.2.2} 及 \hyperlink{para:1:4.2.3}{4.2.3}）。
各 $U\times_S U$ 的并是 $\Delta_X(X)$ 的开邻域；对每个这样的开集，
它与 $\Delta_X(X)$ 的交恰为 $\Delta_X(U)$，所以 $\Delta_X(X)$ 在该并内闭，因而局部闭。
本段是译者补证，不属于原著正文。}

称与浸入 $\Delta_X$ 相伴的 $X\times_S X$ 的子预概形（\hyperlink{para:1:4.2.1}{4.2.1}）为 $X\times_S X$ 的\emph{对角线}。

\egastatement{5.3.10}{推论}
在（\hyperlink{para:1:5.3.5}{5.3.5}）的假设下，$(p,q)_T$ 为浸入。

这由（\hyperlink{para:1:5.3.6}{5.3.6}）和（\hyperlink{para:1:4.3.1}{4.3.1}）得到。
称它为 $X\times_S Y$ 到 $X\times_T Y$ 的\emph{典范浸入}。

\egastatement{5.3.11}{推论}
设 $X,Y$ 为两个 $S$-预概形，$f:X\to Y$ 为 $S$-态射。
则 $f$ 的图态射 $\Gamma_f=(1_X,f)_S$（\hyperlink{para:1:3.3.14}{3.3.14}）为 $X$ 到 $X\times_S Y$ 的浸入。

这是（\hyperlink{para:1:5.3.10}{5.3.10}）中以 $Y$ 代替 $S$、以 $S$ 代替 $T$ 的特例（参见 \hyperlink{para:1:5.3.7}{5.3.7}）。
与浸入 $\Gamma_f$ 相伴的子预概形（\hyperlink{para:1:4.2.1}{4.2.1}）称为态射 $f$ 的\emph{图}。
$X\times_S Y$ 中作为态射 $X\to Y$ 的图的子预概形 $G$，由下述性质刻画：
投影 $p_1:X\times_S Y\to X$ 在 $G$ 上的限制是同构 $g:G\to X$。
此时 $G$ 为态射 $p_2\circ g^{-1}$ 的图，其中 $p_2$ 为到 $Y$ 的投影。

特别取 $X=S$，则 $S$-态射 $S\to Y$，即 $Y$ 的 $S$-截面（\hyperlink{para:1:2.5.5}{2.5.5}），等于它们的图态射。
作为 $S$-截面的图的 $Y$ 的子预概形
（即结构态射 $Y\to S$ 的限制将其同构地映到 $S$ 的子预概形），
也称为这些截面的\emph{像}，或略称为 $Y$ 的 $S$-截面。

\egastatement{5.3.12}{推论}
沿用（\hyperlink{para:1:5.3.11}{5.3.11}）的假设和记号。
对态射 $g:S'\to S$，设 $f'$ 为 $f$ 在 $g$ 下的逆像（\hyperlink{para:1:3.3.7}{3.3.7}），
则 $\Gamma_{f'}$ 为 $\Gamma_f$ 在 $g$ 下的逆像。

这是（\hyperref[eq:1:3.3.10.1]{3.3.10.1}）的特例。

\egastatement{5.3.13}{推论}
设 $f:X\to Y$、$g:Y\to Z$ 为两个态射。
若 $g\circ f$ 是浸入（分别地，局部浸入），则 $f$ 也如此。

事实上，$f$ 分解为
$X\xrightarrow{\Gamma_f}X\times_Z Y\xrightarrow{p_2}Y$，
且 $p_2=(g\circ f)\times_Z1_Y$（\hyperlink{para:1:3.3.4}{3.3.4}）。
若 $g\circ f$ 为浸入（分别地，局部浸入），则 $p_2$ 也如此（\hyperlink{para:1:4.3.1}{4.3.1} 和 \hyperlink{para:1:4.5.5}{4.5.5}）；
又因 $\Gamma_f$ 为浸入（\hyperlink{para:1:5.3.11}{5.3.11}），由浸入的传递性
（\hyperlink{para:1:4.2.4}{4.2.4}），分别由（\hyperlink{para:1:4.5.5}{4.5.5}），得到结论。\transnote{原文引用（\hyperlink{para:1:4.2.4}{4.2.4}）；浸入的传递性直接见（\hyperlink{para:1:4.2.5}{4.2.5}）。}

\egastatement{5.3.14}{推论}
设 $j:X\to Y$、$g:X\to Z$ 为两个 $S$-态射。
若 $j$ 是浸入（分别地，局部浸入），则 $(j,g)_S$ 也如此。

若 $p:Y\times_S Z\to Y$ 为第一投影，则 $j=p\circ(j,g)_S$，应用（\hyperlink{para:1:5.3.13}{5.3.13}）即可。

\egastatement{5.3.15}{命题}
若 $f:X\to Y$ 为 $S$-态射，则图
\begin{equation}
\begin{tikzcd}
X\arrow[r,"\Delta_X"]\arrow[d,"f"']&X\times_S X\arrow[d,"f\times_S f"]\\
Y\arrow[r,"\Delta_Y"']&Y\times_S Y
\end{tikzcd}
\tag{5.3.15.1}\label{eq:1:5.3.15.1}
\end{equation}
交换。换言之，$\Delta_X$ 是预概形范畴中的函子态射。

验证立即可作，留给读者。

\egastatement{5.3.16}{推论}
若 $X$ 为 $Y$ 的子预概形，则对角线 $\Delta_X(X)$ 认同于
$\Delta_Y(Y)$ 的一个子预概形，其底空间认同于
\[
 \Delta_Y(Y)\cap p_1^{-1}(X)=\Delta_Y(Y)\cap p_2^{-1}(X),
\]
其中 $p_1,p_2$ 是 $Y\times_S Y$ 的投影。

对嵌入态射 $f:X\to Y$ 应用（\hyperlink{para:1:5.3.15}{5.3.15}）。
已知 $f\times_S f$ 为浸入，并把 $X\times_S X$ 的底空间认同于
$p_1^{-1}(X)\cap p_2^{-1}(X)$（\hyperlink{para:1:4.3.1}{4.3.1}）。
若 $z\in\Delta_Y(Y)\cap p_1^{-1}(X)$，则 $z=\Delta_Y(y)$，
且 $y=p_1(z)\in X$，所以可写 $y=f(y)$。
由图（\hyperref[eq:1:5.3.15.1]{5.3.15.1}）的交换性，$z=\Delta_Y(f(y))$ 属于 $\Delta_X(X)$。

\egastatement{5.3.17}{推论}
设 $f_1,f_2:Y\to X$ 为两个 $S$-态射，$y\in Y$ 满足
$f_1(y)=f_2(y)=x$，且对应的同态 $k(x)\to k(y)$ 相同。
若 $f=(f_1,f_2)_S$，则 $f(y)$ 属于对角线 $\Delta_{X|S}(X)$。

对应于 $f_i$ 的同态 $k(x)\to k(y)$ 定义 $S$-态射
$g_i:\Spec(k(y))\to\Spec(k(x))$，使图
\[
\begin{tikzcd}
\Spec(k(y))\arrow[r,"g_i"]\arrow[d]&\Spec(k(x))\arrow[d]\\
Y\arrow[r,"f_i"']&X
\end{tikzcd}
\]
交换。于是图
\[
\begin{tikzcd}[column sep=large]
\Spec(k(y))\arrow[r,"{(g_1,g_2)_S}"]\arrow[d]&
\Spec(k(x))\times_S\Spec(k(x))\arrow[d]\\
Y\arrow[r,"{(f_1,f_2)_S}"']&X\times_S X
\end{tikzcd}
\]
也交换。由 $g_1=g_2$，$\Spec(k(y))$ 的唯一点在 $(g_1,g_2)_S$ 下的像属于
$\Spec(k(x))\times_S\Spec(k(x))$ 的对角线。
因此结论由（\hyperlink{para:1:5.3.15}{5.3.15}）得到。


% END chapters/ch1_s5_3.tex

\par

% BEGIN chapters/ch1_s5_4.tex
\subsection{分离态射与分离预概形}\hypertarget{heading:1:5.4}{}
\src{135—136}

\egastatement{5.4.1}{定义}
若预概形态射 $f:X\to Y$ 的对角态射 $X\to X\times_Y X$ 是闭浸入，
则称 $f$ 是\emph{分离的}；此时也称 $X$ 是 $Y$ 上的\emph{分离预概形}，或称 $X$ 为 $Y$-\emph{概形}。
若预概形 $X$ 在 $\Spec(\mathbf Z)$ 上分离，则称 $X$ \emph{分离}；
此时也称 $X$ 为\emph{概形}（参见 \hyperlink{para:1:5.5.7}{5.5.7}）。

根据（\hyperlink{para:1:5.3.9}{5.3.9}），$X$ 在 $Y$ 上分离的充要条件是
$\Delta_{X|Y}(X)$ 在 $X\times_Y X$ 的底空间中闭。

\egastatement{5.4.2}{命题}
设 $S\to T$ 为分离态射。若 $X,Y$ 为两个 $S$-预概形，
则典范浸入 $X\times_S Y\to X\times_T Y$（\hyperlink{para:1:5.3.10}{5.3.10}）是闭的。

参照图（\hyperref[eq:1:5.3.5.1]{5.3.5.1}），$(p,q)_T$ 可看作由 $\Delta_{S|T}$ 经基预概形 $S\times_T S$ 的扩张
$f\times_T g:X\times_T Y\to S\times_T S$ 所得的态射。
所以命题由（\hyperlink{para:1:4.3.2}{4.3.2}）得到。

\egastatement{5.4.3}{推论}
设 $Y$ 为 $S$-概形，$f:X\to Y$ 为 $S$-态射。
则图态射 $\Gamma_f:X\to X\times_S Y$（\hyperlink{para:1:5.3.11}{5.3.11}）为闭浸入。

这是（\hyperlink{para:1:5.4.2}{5.4.2}）中以 $Y$ 代替 $S$、以 $S$ 代替 $T$ 的特例。

\egastatement{5.4.4}{推论}
设 $f:X\to Y$、$g:Y\to Z$ 为两个态射，$g$ 分离。
若 $g\circ f$ 为闭浸入，则 $f$ 也为闭浸入。

从（\hyperlink{para:1:5.4.3}{5.4.3}）出发的证明，与从（\hyperlink{para:1:5.3.11}{5.3.11}）出发证明（\hyperlink{para:1:5.3.13}{5.3.13}）相同。

\egastatement{5.4.5}{推论}
设 $Z$ 为 $S$-概形，$j:X\to Y$、$g:X\to Z$ 为两个 $S$-态射。
若 $j$ 为闭浸入，则 $(j,g)_S:X\to Y\times_S Z$ 也为闭浸入。

从（\hyperlink{para:1:5.4.4}{5.4.4}）出发的证明，与从（\hyperlink{para:1:5.3.13}{5.3.13}）出发证明（\hyperlink{para:1:5.3.14}{5.3.14}）相同。

\egastatement{5.4.6}{推论}
若 $X$ 为 $S$-概形，则 $X$ 的每个 $S$-截面（\hyperlink{para:1:2.5.5}{2.5.5}）都是闭浸入。

若 $\varphi:X\to S$ 为结构态射，$\psi:S\to X$ 为 $S$-截面，
只须对 $\varphi\circ\psi=1_S$ 应用（\hyperlink{para:1:5.4.5}{5.4.5}）即可。

\egastatement{5.4.7}{推论}
设 $S$ 为整预概形，$s$ 为其泛点，$X$ 为 $S$-概形。
若 $X$ 的两个 $S$-截面 $f,g$ 满足 $f(s)=g(s)$，则 $f=g$。

事实上，若 $x=f(s)=g(s)$，则对应于 $f,g$ 的同态 $k(x)\to k(s)$ 必然相同。
置 $h=(f,g)_S$，由（\hyperlink{para:1:5.3.17}{5.3.17}）可知 $h(s)$ 属于对角线 $Z=\Delta_X(X)$。
因 $S=\overline{\{s\}}$ 且 $Z$ 由假设为闭，$h(S)\subset Z$。
根据（\hyperlink{para:1:5.2.2}{5.2.2}），$h$ 分解为 $S\to Z\to X\times_S X$，
再由对角线的定义得到 $f=g$。

\egastatement{5.4.8}{注}
反过来，若（\hyperlink{para:1:5.4.3}{5.4.3}）的结论在 $f=1_Y$ 时成立，便得到 $Y$ 在 $S$ 上分离。
同样，若（\hyperlink{para:1:5.4.5}{5.4.5}）的结论可用于两个态射
$Y\xrightarrow{\Delta_Y}Y\times_Z Y\xrightarrow{p_1}Y$，
则得到 $\Delta_Y$ 为闭浸入，因而 $Y$ 在 $Z$ 上分离。
\transnote{原文在后一处引用（\hyperlink{para:1:5.4.5}{5.4.5}）；这里所用的复合态射判据直接见（\hyperlink{para:1:5.4.4}{5.4.4}）。}
最后，若（\hyperlink{para:1:5.4.6}{5.4.6}）的结论对 $Y$-预概形 $Y\times_S Y\to Y$ 的 $Y$-截面 $\Delta_Y$ 成立，
也蕴含 $Y$ 在 $S$ 上分离。


% END chapters/ch1_s5_4.tex

\par

% BEGIN chapters/ch1_s5_5.tex
\subsection{分离性判据}\hypertarget{heading:1:5.5}{}
\src{136—140}

\egastatement{5.5.1}{命题}
\begin{enumerate}[label=(\roman*)]
\item 每个预概形单态射，特别是每个浸入，都是分离态射。
\item 两个分离态射的复合仍分离。
\item 若 $f:X\to X'$、$g:Y\to Y'$ 为两个分离的 $S$-态射，则 $f\times_S g$ 分离。
\item 若 $f:X\to Y$ 为分离的 $S$-态射，则对任意基预概形扩张 $S'\to S$，
$S'$-态射 $f_{(S')}$ 分离。
\item 若两个态射的复合 $g\circ f$ 分离，则 $f$ 分离。
\item $f$ 分离的充要条件是 $f_{\mathrm{red}}$（\hyperlink{para:1:5.1.5}{5.1.5}）分离。
\end{enumerate}

(i) 由（\hyperlink{para:1:5.3.8}{5.3.8}）立即得到。若 $f:X\to Y$、$g:Y\to Z$ 为态射，图
\begin{equation}
\begin{tikzcd}[column sep=large]
X\arrow[r,"\Delta_{X|Z}"]\arrow[dr,"\Delta_{X|Y}"']&X\times_Z X\\
&X\times_Y X\arrow[u,"j"']
\end{tikzcd}
\tag{5.5.1.1}\label{eq:1:5.5.1.1}
\end{equation}
交换，验证立即可作，其中 $j$ 为典范浸入（\hyperlink{para:1:5.3.10}{5.3.10}）。
若 $f,g$ 分离，则 $\Delta_{X|Y}$ 由定义为闭浸入，$j$ 由（\hyperlink{para:1:5.4.2}{5.4.2}）为闭浸入，
所以 $\Delta_{X|Z}$ 也为闭浸入，这证明 (ii)。
\transnote{原文在此引用（\hyperlink{para:1:4.2.4}{4.2.4}）；闭浸入复合的陈述见（\hyperlink{para:1:4.2.5}{4.2.5}）。}
由 (i)、(ii)，(iii)、(iv) 等价（\hyperlink{para:1:3.5.1}{3.5.1}），故只须证明 (iv)。
由（\hyperlink{para:1:3.3.11}{3.3.11}）和（\hyperref[eq:1:3.3.9.1]{3.3.9.1}），有典范认同
\[
 X_{(S')}\times_{Y_{(S')}}X_{(S')}
 \simeq (X\times_Y X)\times_Y Y_{(S')}.
\]
立即可验证，此时对角态射认同于 $\Delta_X\times_Y1_{Y_{(S')}}$，
故命题由（\hyperlink{para:1:4.3.1}{4.3.1}）得到。

为证明 (v)，如（\hyperlink{para:1:5.3.13}{5.3.13}）中那样考虑分解
$X\xrightarrow{\Gamma_f}X\times_Z Y\xrightarrow{p_2}Y$。
由 $p_2=(g\circ f)\times_Z1_Y$，假设及 (iii)、(i) 蕴含 $p_2$ 分离。
$\Gamma_f$ 为浸入，故由 (i) 分离；再由 (ii)，$f$ 分离。
最后，为证明 (vi)，回忆
$X_{\mathrm{red}}\times_{Y_{\mathrm{red}}}X_{\mathrm{red}}$
与 $X_{\mathrm{red}}\times_Y X_{\mathrm{red}}$ 典范地认同（\hyperlink{para:1:5.1.7}{5.1.7}）。
若 $j:X_{\mathrm{red}}\to X$ 为嵌入，则图
\[
\begin{tikzcd}[column sep=large]
X_{\mathrm{red}}\arrow[r,"\Delta_{X_{\mathrm{red}}}"]\arrow[d,"j"']&
X_{\mathrm{red}}\times_Y X_{\mathrm{red}}\arrow[d,"j\times_Y j"]\\
X\arrow[r,"\Delta_{X|Y}"']&X\times_Y X
\end{tikzcd}
\]
交换（\hyperlink{para:1:5.3.15}{5.3.15}）。竖箭头均为底空间的同胚（\hyperlink{para:1:4.3.1}{4.3.1}），故命题成立。

\egastatement{5.5.2}{推论}
若 $f:X\to Y$ 分离，则它在 $X$ 的任意子预概形上的限制分离。

这由（\hyperlink{para:1:5.5.1}{5.5.1}，(i)、(ii)）得到。

\egastatement{5.5.3}{推论}
若 $X,Y$ 为两个 $S$-预概形，且 $Y$ 在 $S$ 上分离，则 $X\times_S Y$ 在 $X$ 上分离。

这是（\hyperlink{para:1:5.5.1}{5.5.1}，(iv)）的特例。

\egastatement{5.5.4}{命题}
设预概形 $X$ 的底空间为有限个闭子集 $X_k$（$1\leq k\leq n$）的并。
仍以 $X_k$ 表示具有该底空间的约化子预概形（\hyperlink{para:1:5.2.1}{5.2.1}）。
设 $f:X\to Y$ 为态射，并为每个 $k$ 取闭子集 $Y_k\subset Y$，使 $f(X_k)\subset Y_k$；
同样以 $Y_k$ 表示相应的约化子预概形。
于是 $f|X_k$ 分解为 $X_k\xrightarrow{f_k}Y_k\to Y$（\hyperlink{para:1:5.2.2}{5.2.2}）。
则 $f$ 分离的充要条件是每个 $f_k$ 分离。

必要性由（\hyperlink{para:1:5.5.1}{5.5.1}，(i)、(ii)、(v)）得到。
反之，若每个 $f_k$ 分离，则 $f$ 的每个限制 $X_k\to Y$ 分离（\hyperlink{para:1:5.5.1}{5.5.1}，(i)、(ii)）。
若 $p_1,p_2$ 为 $X\times_Y X$ 的投影，则子空间 $\Delta_{X_k}(X_k)$
认同于 $\Delta_X(X)\cap p_1^{-1}(X_k)$（\hyperlink{para:1:5.3.16}{5.3.16}）。
这些子空间在 $X\times_Y X$ 中闭，所以它们的并 $\Delta_X(X)$ 也闭。

特别地，若各 $X_k$ 为 $X$ 的不可约分支，则可取 $Y_k$ 为 $Y$ 的不可约分支（0，\hyperlink{para:0:2.1.5}{2.1.5}）。
因而在此情形，（\hyperlink{para:1:5.5.4}{5.5.4}）将分离性的研究归结到整预概形的情形。

\egastatement{5.5.5}{命题}
设 $(Y_\lambda)$ 为预概形 $Y$ 的开覆盖。
态射 $f:X\to Y$ 分离的充要条件是每个限制
$f^{-1}(Y_\lambda)\to Y_\lambda$ 分离。

置 $X_\lambda=f^{-1}(Y_\lambda)$。
由（\hyperlink{para:1:4.2.4}{4.2.4}，b)）及积 $X_\lambda\times_Y X_\lambda$
与 $X_\lambda\times_{Y_\lambda}X_\lambda$ 相同（\hyperlink{para:1:3.2.5}{3.2.5}），
只须证明 $X_\lambda\times_Y X_\lambda$ 覆盖 $X\times_Y X$。
置 $Y_{\lambda\mu}=Y_\lambda\cap Y_\mu$、
$X_{\lambda\mu}=X_\lambda\cap X_\mu=f^{-1}(Y_{\lambda\mu})$。
则 $X_\lambda\times_Y X_\mu$ 认同于
$X_{\lambda\mu}\times_Y X_{\lambda\mu}$（\hyperlink{para:1:3.2.6.4}{3.2.6.4}），
也即 $X_{\lambda\mu}\times_{Y_\lambda}X_{\lambda\mu}$（\hyperlink{para:1:3.2.5}{3.2.5}），
从而认同于 $X_\lambda\times_Y X_\lambda$ 的一个开集（\hyperlink{para:1:3.2.7}{3.2.7}），所断言的覆盖成立。

因此，取 $Y$ 的仿射开覆盖，可将分离态射的研究归结到以仿射概形为靶的分离态射。
\transnote{原文在这一句引用（\hyperlink{para:1:5.5.4}{5.5.4}）；对应的开覆盖判据是刚证明的（\hyperlink{para:1:5.5.5}{5.5.5}）。}

\egastatement{5.5.6}{命题}
设 $Y$ 为仿射概形，$X$ 为预概形，$(U_\alpha)$ 为 $X$ 的仿射开覆盖。
态射 $f:X\to Y$ 分离的充要条件是：对每对指标 $(\alpha,\beta)$，
$U_\alpha\cap U_\beta$ 为仿射开集，且环
$\Gamma(U_\alpha\cap U_\beta,\OO_X)$
由两个环 $\Gamma(U_\alpha,\OO_X)$、$\Gamma(U_\beta,\OO_X)$ 的典范像之并生成。

各 $U_\alpha\times_Y U_\beta$ 构成 $X\times_Y X$ 的开覆盖（\hyperlink{para:1:3.2.7}{3.2.7}）。
以 $p,q$ 表示投影，有
\[
\begin{aligned}
\Delta_X^{-1}(U_\alpha\times_Y U_\beta)
 &=\Delta_X^{-1}\bigl(p^{-1}(U_\alpha)\cap q^{-1}(U_\beta)\bigr)\\
 &=\Delta_X^{-1}(p^{-1}(U_\alpha))\cap
   \Delta_X^{-1}(q^{-1}(U_\beta))=U_\alpha\cap U_\beta.
\end{aligned}
\]
故只须表述 $\Delta_X$ 在 $U_\alpha\cap U_\beta$ 上的限制为到
$U_\alpha\times_Y U_\beta$ 的闭浸入的条件。
由定义，该限制为 $(j_\alpha,j_\beta)_Y$，
其中 $j_\alpha,j_\beta$ 分别是交集到 $U_\alpha,U_\beta$ 的嵌入。
而 $U_\alpha\times_Y U_\beta$ 仿射，其环典范地同构于
\[
 \Gamma(U_\alpha,\OO_X)\otimes_{\Gamma(Y,\OO_Y)}\Gamma(U_\beta,\OO_X)
\]
（\hyperlink{para:1:3.2.2}{3.2.2}）。所以 $U_\alpha\cap U_\beta$ 必须仿射，且将
$h_\alpha\otimes h_\beta$ 映为 $h_\alpha h_\beta$ 的环同态必须为满射（\hyperlink{para:1:4.2.3}{4.2.3}）。证明完成。

\egastatement{5.5.7}{推论}
仿射概形是分离的，因而确实是概形，这说明（\hyperlink{para:1:5.4.1}{5.4.1}）的术语与此前一致。

\egastatement{5.5.8}{推论}
设 $Y$ 为仿射概形。态射 $f:X\to Y$ 分离的充要条件是 $X$ 分离，即 $X$ 为概形。

事实上，（\hyperlink{para:1:5.5.6}{5.5.6}）的判据不依赖于 $f$。

\egastatement{5.5.9}{推论}
态射 $f:X\to Y$ 分离的必要条件是：
对每个使 $Y$ 在其上诱导的预概形分离的开集 $U$，预概形 $f^{-1}(U)$ 分离；
充分条件是：对每个仿射开集 $U\subset Y$，$f^{-1}(U)$ 分离。

必要性由开覆盖判据及（\hyperlink{para:1:5.5.1}{5.5.1}，(ii)）得到；
充分性由开覆盖判据及（\hyperlink{para:1:5.5.8}{5.5.8}）得到，并利用 $Y$ 的仿射开覆盖的存在性。
\transnote{原文在这两处将开覆盖判据引为（\hyperlink{para:1:5.5.4}{5.5.4}）；这里的对应条目是（\hyperlink{para:1:5.5.5}{5.5.5}）。}
特别地，若 $X,Y$ 都是仿射概形，则每个态射 $X\to Y$ 都分离。

\egastatement{5.5.10}{命题}
设 $Y$ 为概形，$f:X\to Y$ 为态射。
对 $X$ 的每个仿射开集 $U$ 及 $Y$ 的每个仿射开集 $V$，$U\cap f^{-1}(V)$ 仿射。

以 $p_1,p_2$ 表示 $X\times_{\mathbf Z}Y$ 的投影。
子空间 $U\cap f^{-1}(V)$ 是
$\Gamma_f(X)\cap p_1^{-1}(U)\cap p_2^{-1}(V)$ 在 $p_1$ 下的像。
而 $p_1^{-1}(U)\cap p_2^{-1}(V)$ 认同于 $U\times_{\mathbf Z}V$ 的底空间（\hyperlink{para:1:3.2.7}{3.2.7}），
故为仿射概形（\hyperlink{para:1:3.2.2}{3.2.2}）。
由于 $\Gamma_f(X)$ 在 $X\times_{\mathbf Z}Y$ 中闭（\hyperlink{para:1:5.4.3}{5.4.3}），
上述交集在 $U\times_{\mathbf Z}V$ 中闭。
于是与 $\Gamma_f$ 相伴的子预概形在该开子空间上诱导的预概形，
是一个仿射概形的闭子预概形，故仿射（\hyperlink{para:1:4.2.3}{4.2.3}）。
最后，由 $\Gamma_f$ 是浸入得到命题。

\egastatement{5.5.11}{例}
例（\hyperlink{para:1:2.3.2}{2.3.2}）的预概形，即“域 $K$ 上的射影直线”，是分离的。
事实上，对于仿射开覆盖 $(X_1,X_2)$，$X_1\cap X_2=U_{12}$ 仿射；
$\Gamma(U_{12},\OO_X)$ 为形如 $f(s)/s^m$、$f\in K[s]$ 的有理分式所成的环，
它由 $K[s]$ 和 $1/s$ 生成，所以（\hyperlink{para:1:5.5.6}{5.5.6}）的条件成立。

仍如例（\hyperlink{para:1:2.3.2}{2.3.2}）选取 $X_1,X_2,U_{12},U_{21}$，但这次取将 $f(s)$ 映为 $f(t)$ 的同构为 $u_{12}$。
粘合后得到一个整的、非分离的预概形 $X$：
（\hyperlink{para:1:5.5.6}{5.5.6}）的第一条件成立，而第二条件不成立。
这里立即可见 $\Gamma(X,\OO_X)\to\Gamma(X_1,\OO_X)=K[s]$ 为同构。
逆同构定义满射态射 $f:X\to\Spec(K[s])$。
若 $y\in\Spec(K[s])$ 满足 $\jj_y\ne(s)$，则 $f^{-1}(y)$ 只有一个点；
若 $\jj_y=(s)$，则 $f^{-1}(y)$ 由两个不同的点组成。
称 $X$ 为“域 $K$ 上将点 $0$ 加倍的仿射直线”。
\transnote{原版在上述两处写作理想 $(0)$，但在本例中加倍的是原点，其理想为 $(s)$；
零理想对应的泛点并未加倍。此处已明确改写为 $(s)$。}

还可给出（\hyperlink{para:1:5.5.6}{5.5.6}）的两个条件都不成立的例子。
首先注意，在域 $K$ 上二元多项式环 $A=K[s,t]$ 的素谱 $Y$ 中，
开集 $U=D(s)\cup D(t)$ 不是仿射开集。
事实上，若 $z$ 为 $\OO_Y$ 在 $U$ 上的截面，则存在整数 $m,n\geq0$，
使 $s^m z$、$t^n z$ 均为关于 $s,t$ 的多项式在 $U$ 上的限制（\hyperlink{para:1:1.4.1}{1.4.1}）。
这显然只有在 $z$ 能延拓到整个 $Y$ 上、并与一个关于 $s,t$ 的多项式相认同时才可能。
若 $U$ 仿射，则嵌入态射 $U\to Y$ 将是同构（\hyperlink{para:1:1.7.3}{1.7.3}），与 $U\ne Y$ 矛盾。

现取仿射概形 $Y_1,Y_2$，分别为环
$A_1=K[s_1,t_1]$、$A_2=K[s_2,t_2]$ 的素谱。
置 $U_{12}=D(s_1)\cup D(t_1)$、$U_{21}=D(s_2)\cup D(t_2)$，
并取 $u_{12}$ 为同构 $Y_2\to Y_1$ 在 $U_{21}$ 上的限制，
该同构对应于将 $f(s_1,t_1)$ 映为 $f(s_2,t_2)$ 的环同构。
于是得到一个（\hyperlink{para:1:5.5.6}{5.5.6}）的两个条件都不满足的例子。
如此得到的整预概形称为“域 $K$ 上将点 $0$ 加倍的仿射平面”。

\egastatement{5.5.12}{注}
给定预概形态射的一个性质 $P$，考虑下述命题：
\begin{enumerate}[label=(\roman*)]
\item 每个闭浸入具有性质 $P$。
\item 两个具有性质 $P$ 的态射的复合具有性质 $P$。
\item 若 $f:X\to X'$、$g:Y\to Y'$ 为具有性质 $P$ 的 $S$-态射，
则 $f\times_S g$ 具有性质 $P$。
\item 若 $f:X\to Y$ 为具有性质 $P$ 的 $S$-态射，
则任意基预概形扩张 $S'\to S$ 所得的 $f_{(S')}$ 具有性质 $P$。
\item 若 $f:X\to Y$、$g:Y\to Z$ 的复合 $g\circ f$ 具有性质 $P$，
且 $g$ 分离，则 $f$ 具有性质 $P$。
\item 若 $f:X\to Y$ 具有性质 $P$，则 $f_{\mathrm{red}}$（\hyperlink{para:1:5.1.5}{5.1.5}）具有性质 $P$。
\end{enumerate}
若 (i)、(ii) 成立，则 (iii) 与 (iv) 等价，且 (v)、(vi) 都是 (i)、(ii)、(iii) 的推论。

第一断言已在（\hyperlink{para:1:3.5.1}{3.5.1}）中证明。
考虑分解（\hyperlink{para:1:5.3.13}{5.3.13}）$X\xrightarrow{\Gamma_f}X\times_Z Y\xrightarrow{p_2}Y$。
等式 $p_2=(g\circ f)\times_Z1_Y$ 表明：
若 $g\circ f$ 具有性质 $P$，则 $p_2$ 由 (iii) 具有性质 $P$。
若 $g$ 分离，则 $\Gamma_f$ 为闭浸入（\hyperlink{para:1:5.4.3}{5.4.3}），因而由 (i) 具有性质 $P$；
最后由 (ii)，$f$ 具有性质 $P$。

再考虑交换图
\[
\begin{tikzcd}
X_{\mathrm{red}}\arrow[r,"f_{\mathrm{red}}"]\arrow[d]&Y_{\mathrm{red}}\arrow[d]\\
X\arrow[r,"f"']&Y
\end{tikzcd}
\]
其中竖箭头为闭浸入（\hyperlink{para:1:5.1.5}{5.1.5}），故由 (i) 具有性质 $P$。
若 $f$ 具有性质 $P$，则由 (ii)，复合态射
$X_{\mathrm{red}}\to Y_{\mathrm{red}}\to Y$ 具有性质 $P$。
由于闭浸入分离（\hyperlink{para:1:5.5.1}{5.5.1}，(i)），由 (v) 得 $f_{\mathrm{red}}$ 具有性质 $P$。

若考虑下述两个命题：
\begin{enumerate}[label=\textup{(\roman*')},start=1]
\item 每个浸入具有性质 $P$；
\end{enumerate}
\noindent(v') 若 $g\circ f$ 具有性质 $P$，则 $f$ 具有性质 $P$；
则以上论证表明，(v') 是 (i')、(ii)、(iii) 的推论。

\egapar{5.5.13}
还应注意，(v)、(vi) 也是 (i)、(iii) 以及下述命题的推论：

\noindent(ii') 若 $j:X\to Y$ 为闭浸入，$g:Y\to Z$ 为具有性质 $P$ 的态射，
则 $g\circ j$ 具有性质 $P$。

同样，(v') 是 (i')、(iii) 以及下述命题的推论：

\noindent(ii'') 若 $j:X\to Y$ 为浸入，$g:Y\to Z$ 为具有性质 $P$ 的态射，
则 $g\circ j$ 具有性质 $P$。

这些断言由（\hyperlink{para:1:5.5.12}{5.5.12}）的论证立即得到。


% END chapters/ch1_s5_5.tex

\par

% BEGIN chapters/ch1_s6_1.tex
\section{有限性条件}\hypertarget{heading:1:6}{}
\subsection{诺特预概形与局部诺特预概形}\hypertarget{heading:1:6.1}{}
\src{140—143}

\egastatement{6.1.1}{定义}
若预概形 $X$ 是有限个仿射开集 $V_i$ 的并（分别地，任意族仿射开集的并），
且每个 $V_i$ 上诱导的仿射概形的环都是诺特环，
则称 $X$ 为\emph{诺特预概形}（分别地，\emph{局部诺特预概形}）。

由（\hyperlink{para:1:1.5.2}{1.5.2}）及问题的局部性立即可知：若 $X$ 局部诺特，
则结构层 $\OO_X$ 是凝聚环层。
此时凝聚 $\OO_X$-模 $\F$ 的每个拟凝聚子模（分别地，拟凝聚商模）都是凝聚的。
问题仍具有局部性；只须应用（\hyperlink{para:1:1.5.1}{1.5.1}）、（\hyperlink{para:1:1.4.1}{1.4.1}）、（\hyperlink{para:1:1.3.10}{1.3.10}），
并利用诺特环上有限型模的子模、商模仍为有限型模的事实。
特别地，$\OO_X$ 的每个拟凝聚理想层都是凝聚的。

若 $X$ 是有限个开集 $W_i$ 的并（分别地，任意族开集的并），
且每个 $W_i$ 上诱导的预概形为诺特（分别地，局部诺特）的，
则显然 $X$ 为诺特（分别地，局部诺特）的。

\egastatement{6.1.2}{命题}
预概形 $X$ 为诺特的充要条件是它局部诺特且底空间拟紧。
此时 $X$ 的底空间是诺特空间。

第一断言由定义及（\hyperlink{para:1:1.1.10}{1.1.10}，(ii)）立即得到。
第二断言由（\hyperlink{para:1:1.1.6}{1.1.6}）及有限个诺特子空间之并为诺特空间的事实（0，\hyperlink{para:0:2.2.3}{2.2.3}）得到。

\egastatement{6.1.3}{命题}
设 $X$ 为环为 $A$ 的仿射概形。
下述条件等价：a) $X$ 诺特；b) $X$ 局部诺特；c) $A$ 诺特。

a)、b) 的等价性由（\hyperlink{para:1:6.1.2}{6.1.2}）及仿射概形的底空间拟紧（\hyperlink{para:1:1.1.10}{1.1.10}）得到；
c) 显然蕴含 a)。
为证明 a) 蕴含 c)，取 $X$ 的有限仿射开覆盖 $(V_i)$，使每个 $A_i=A(V_i)$ 诺特。
设 $(\mathfrak a_n)$ 为 $A$ 的递增理想序列。
它典范地一一对应于 $\wt A=\OO_X$ 中的递增理想层序列
$(\wt{\mathfrak a_n})$（\hyperlink{para:1:1.3.7}{1.3.7}）。
只须证明后一序列稳定。
限制 $\wt{\mathfrak a_n}|V_i$ 是 $\OO_X|V_i$ 的拟凝聚理想层，
因为它是沿典范嵌入 $V_i\to X$ 的逆像（0，\hyperlink{para:0:5.1.4}{5.1.4}）。
故它形如 $\wt{\mathfrak a_{ni}}$，其中 $\mathfrak a_{ni}$ 为 $A_i$ 的理想（\hyperlink{para:1:1.3.7}{1.3.7}）。
由于 $A_i$ 诺特，每个序列 $(\mathfrak a_{ni})_n$ 稳定，从而命题成立。

还应注意，这个论证也证明：若 $X$ 为诺特预概形，
则 $\OO_X$ 的每个递增凝聚理想层序列均稳定。

\egastatement{6.1.4}{命题}
诺特（分别地，局部诺特）预概形的每个子预概形仍为诺特（分别地，局部诺特）的。

只须对诺特预概形 $X$ 证明；又由定义（\hyperlink{para:1:6.1.1}{6.1.1}）立即可归结到 $X$ 仿射的情形。
每个子预概形都是某个开集上诱导的预概形的闭子预概形（\hyperlink{para:1:4.1.3}{4.1.3}），
所以只须分别考虑闭子预概形 $Y$ 与在开集上诱导的子预概形 $Y$。
闭的情形立即成立：若 $X$ 的环为 $A$，则 $Y$ 仿射且其环为 $A/\mathfrak J$（\hyperlink{para:1:4.2.3}{4.2.3}）；
$A$ 诺特（\hyperlink{para:1:6.1.3}{6.1.3}），故 $A/\mathfrak J$ 诺特。
若 $Y$ 在 $X$ 中开，则底空间 $Y$ 诺特（\hyperlink{para:1:6.1.2}{6.1.2}），所以拟紧，
可表示为有限个 $D(f_i)$（$f_i\in A$）的并。
因而归结到 $Y=D(f)$ 的情形。此时 $Y$ 仿射，环同构于 $A_f$（\hyperlink{para:1:1.3.6}{1.3.6}）；
$A$ 诺特，故 $A_f$ 诺特。

\egapar{6.1.5}
两个诺特 $S$-预概形的积未必诺特，即使它们都是仿射概形也是如此；
因为两个诺特代数的张量积未必是诺特环（参见 \hyperlink{para:1:6.3.8}{6.3.8}）。

\egastatement{6.1.6}{命题}
若 $X$ 为诺特预概形，则 $\OO_X$ 的幂零根 $\N_X$ 是幂零的。

可用有限个仿射开集 $U_i$ 覆盖 $X$。
只须证明存在整数 $n_i$ 使 $(\N_X|U_i)^{n_i}=0$；
若 $n$ 为这些 $n_i$ 的最大值，便有 $\N_X^n=0$。
所以归结到 $X=\Spec(A)$、$A$ 诺特的情形。
根据（\hyperlink{para:1:5.1.1}{5.1.1}）和（\hyperlink{para:1:1.3.13}{1.3.13}），只须注意诺特环的幂零根是幂零的（文献 \cite{b11}，127 页，推论 4）。

\egastatement{6.1.7}{推论}
设 $X$ 为诺特预概形。$X$ 为仿射概形的充要条件是 $X_{\mathrm{red}}$ 为仿射概形。

这由（\hyperlink{para:1:6.1.6}{6.1.6}）和（\hyperlink{para:1:5.1.10}{5.1.10}）得到。

\egastatement{6.1.8}{引理}
设 $X$ 为拓扑空间，$x\in X$，$U$ 为 $x$ 的一个只有有限个不可约分支的开邻域。
则存在 $x$ 的邻域 $V$，使包含于 $V$ 的每个 $x$ 的开邻域都连通。

设 $U_i$（$1\leq i\leq m$）为 $U$ 中不包含 $x$ 的不可约分支。
它们的并在 $U$ 中的补集 $V$ 是 $x$ 在 $U$ 中、从而在 $X$ 中的开邻域。
它也是 $U$ 与 $X$ 中所有不包含 $x$ 的不可约分支之并的补集的交（0，\hyperlink{para:0:2.1.6}{2.1.6}）。
\transnote{原文将 $V$ 直接称为上述并在 $X$ 中的补集；
一般还须与 $U$ 相交。例如 $X$ 不可约且 $U\ne X$ 时，原来的直接等式不成立。}
若 $W\subset V$ 为 $x$ 的开邻域，则 $W$ 的不可约分支是 $U$ 中与 $W$ 相交的不可约分支在 $W$ 上的迹（0，\hyperlink{para:0:2.1.6}{2.1.6}）。
所以它们都包含 $x$；由于它们连通，$W$ 也连通。

\egastatement{6.1.9}{推论}
局部诺特拓扑空间局部连通；特别地，其连通分支是开集。

\egastatement{6.1.10}{命题}
设 $X$ 为局部诺特拓扑空间。下述条件等价：
\begin{enumerate}[label=\alph*)]
\item $X$ 的不可约分支是开集；
\item $X$ 的不可约分支与连通分支相同；
\item $X$ 的连通分支不可约；
\item $X$ 的两个不同不可约分支不相交。
\end{enumerate}
若 $X$ 为预概形，则它们还等价于：

\noindent e) 对每个 $x\in X$，$\Spec(\OO_x)$ 不可约，即 $\OO_x$ 的幂零根是素理想。

a) 蕴含 b) 立即成立：不可约空间连通，且 a) 使各不可约分支既开又闭。
b) 显然蕴含 c)。反之，包含一个连通分支 $C$ 且不同于 $C$ 的闭集 $F$ 不可能不可约：
$F$ 不连通，故为两个非空、不交、在 $F$ 中既开又闭、因而在 $X$ 中闭的子集之并。
所以 c) 蕴含 b)。由此立即得到 c) 蕴含 d)，因为不同连通分支不相交。

以上尚未用到局部诺特假设。现利用此假设证明 d) 蕴含 a)。
由（0，\hyperlink{para:0:2.1.6}{2.1.6}），可限于 $X$ 诺特的情形。
此时不可约分支只有有限个，它们闭且两两不交，故也开。

最后，d) 与 e) 的等价性无需假设底空间局部诺特。
由（0，\hyperlink{para:0:2.1.6}{2.1.6}），可归结到 $X=\Spec(A)$ 仿射的情形。
$x$ 只属于一个不可约分支，意味着 $\jj_x$ 只包含一个极小素理想（\hyperlink{para:1:1.1.14}{1.1.14}）；
这又等价于说 $\jj_x\OO_x$ 只包含 $\OO_x$ 的一个极小素理想，从而得到结论。

\egastatement{6.1.11}{推论}
设 $X$ 为局部诺特空间。
$X$ 不可约的充要条件是：$X$ 非空且连通，并且两个不同的不可约分支不相交。
若 $X$ 为预概形，最后一个条件等价于每个 $\Spec(\OO_x)$ 不可约。

最后部分已在（\hyperlink{para:1:6.1.10}{6.1.10}）中证明。
只须证明第一部分条件的充分性：由（\hyperlink{para:1:6.1.10}{6.1.10}），这些条件使不可约分支与连通分支相同；
$X$ 非空且连通，故不可约。

\egastatement{6.1.12}{推论}
设 $X$ 为局部诺特预概形。
$X$ 为整的充要条件是 $X$ 非空且连通，并且对每个 $x\in X$，$\OO_x$ 为整环。
\transnote{原文在本条只写“连通”，未另写“非空”。这里明确补出非空条件，
以与（0，\hyperlink{para:0:2.1.1}{2.1.1}）及（\hyperlink{para:1:6.1.11}{6.1.11}）中的不可约性约定一致。}

\egastatement{6.1.13}{命题}
设 $X$ 为局部诺特预概形，$x\in X$，且 $\OO_x$ 的幂零根 $\N_{X,x}$ 为素理想
（分别地，$\OO_x$ 约化；$\OO_x$ 为整环）。
则存在 $x$ 的一个不可约（分别地，约化；整）的开邻域 $U$。

只须考虑幂零根为素理想和幂零根为零两种情形，第三种假设是前两种的合取。
若 $\N_{X,x}$ 为素理想，则 $x$ 只属于 $X$ 的一个不可约分支 $Y$（\hyperlink{para:1:6.1.10}{6.1.10}）。
不包含 $x$ 的不可约分支之并是闭的，因为这些分支构成局部有限族。
其补集 $U$ 因而开且包含于 $Y$，所以不可约（0，\hyperlink{para:0:2.1.6}{2.1.6}）。
若 $\N_{X,x}=0$，则在 $x$ 的某个邻域内每个 $y$ 都有 $\N_{X,y}=0$：
$\N_X$ 拟凝聚（\hyperlink{para:1:5.1.1}{5.1.1}），因 $X$ 局部诺特而凝聚，故结论由（0，\hyperlink{para:0:5.2.2}{5.2.2}）得到。


% END chapters/ch1_s6_1.tex

\par

% BEGIN chapters/ch1_s6_2.tex
\subsection{阿廷预概形}\hypertarget{heading:1:6.2}{}
\src{143—144}

\egastatement{6.2.1}{定义}
若一个预概形仿射且其环为阿廷环，则称它为\emph{阿廷预概形}。

\egastatement{6.2.2}{命题}
对于预概形 $X$，下述条件等价：
\begin{enumerate}[label=\alph*)]
\item $X$ 为阿廷概形；
\item $X$ 诺特且其底空间离散；
\item $X$ 诺特且其底空间的各点都是闭点，即满足条件 $T_1$。
\end{enumerate}
此时 $X$ 的底空间有限，且 $X$ 的环 $A$ 为各点的局部环（均为阿廷环）的直积。

已知 a) 蕴含最后的断言（文献 \cite{b13}，205 页，定理 3）。
此时 $A$ 的每个素理想都是极大理想，并且是一个局部直积分量的极大理想的逆像，
故 $X$ 有限且离散。于是 a) 蕴含 b)，而 b) 显然蕴含 c)。
为证明 c) 蕴含 a)，先证明 $X$ 有限。
可归结到 $X$ 仿射的情形；一个各素理想均极大的诺特环为阿廷环
（文献 \cite{b13}，203 页），故所断言的有限性成立。
此时 $X$ 的底空间离散，为有限个点 $x_i$ 的拓扑和，局部环 $\OO_{x_i}=A_i$ 为阿廷环。
显然 $X$ 同构于环 $A=\prod_i A_i$ 的素谱所定义的仿射概形（\hyperlink{para:1:1.7.3}{1.7.3}）。


% END chapters/ch1_s6_2.tex

\par

% BEGIN chapters/ch1_s6_3.tex
\subsection{有限型态射}\hypertarget{heading:1:6.3}{}
\src{144—147}

\egastatement{6.3.1}{定义}
若预概形态射 $f:X\to Y$ 满足：$Y$ 是一族具有下述性质的仿射开集 $(V_\alpha)$ 的并，
则称 $f$ 为\emph{有限型态射}：

\noindent(P) $f^{-1}(V_\alpha)$ 是有限个仿射开集 $U_{\alpha\beta}$ 的并，
且每个环 $A(U_{\alpha\beta})$ 都是 $A(V_\alpha)$ 上的有限型代数。

此时也称 $X$ 是 $Y$ 上的\emph{有限型预概形}，或称 $X$ 为\emph{有限型 $Y$-预概形}。

\egastatement{6.3.2}{命题}
若 $f:X\to Y$ 为有限型态射，则 $Y$ 的每个仿射开集 $W$
都具有定义（\hyperlink{para:1:6.3.1}{6.3.1}）中的性质 (P)。

先证明下述引理。

\egastatement{6.3.2.1}{引理}
若仿射开集 $T\subset Y$ 具有性质 (P)，则对每个 $g\in A(T)$，$D(g)$ 也具有性质 (P)。

由假设，$f^{-1}(T)$ 是有限个仿射开集 $Z_j$ 的并，且每个 $A(Z_j)$ 都是有限型 $A(T)$-代数。
设 $\varphi_j:A(T)\to A(Z_j)$ 为对应于 $f|Z_j$ 的环同态（\hyperlink{para:1:2.2.4}{2.2.4}），置 $g_j=\varphi_j(g)$。
则 $f^{-1}(D(g))\cap Z_j=D(g_j)$（\hyperref[eq:1:1.2.2.2]{1.2.2.2}）。
而 $A(D(g_j))=A(Z_j)_{g_j}=A(Z_j)[1/g_j]$ 在 $A(Z_j)$ 上有限型，
由假设在 $A(T)$ 上也有限型，从而在 $A(D(g))=A(T)[1/g]$ 上有限型。引理得证。

由于 $W$ 拟紧（\hyperlink{para:1:1.1.10}{1.1.10}），存在 $W$ 的有限覆盖，成员形如 $D(g_i)$，
每个 $g_i$ 属于某个 $A(V_{\alpha(i)})$。
每个 $D(g_i)$ 拟紧，故为有限个 $D(h_{ik})$ 的并，其中 $h_{ik}\in A(W)$。
若 $\varphi_i:A(W)\to A(D(g_i))$ 为典范映射，则
$D(h_{ik})=D(\varphi_i(h_{ik}))$（\hyperref[eq:1:1.2.2.2]{1.2.2.2}）。
由引理（\hyperlink{para:1:6.3.2.1}{6.3.2.1}），每个 $f^{-1}(D(h_{ik}))$ 都有有限仿射开覆盖 $(U_{ikj})$，
使 $A(U_{ikj})$ 在 $A(D(h_{ik}))=A(W)[1/h_{ik}]$ 上有限型，命题遂得证。

所以可以说，$Y$ 上有限型预概形的概念在 $Y$ 上具有局部性。

\egastatement{6.3.3}{命题}
设 $X,Y$ 为两个仿射概形。
$X$ 在 $Y$ 上有限型的充要条件是 $A(X)$ 为有限型 $A(Y)$-代数。

充分性显然，现证明必要性。置 $A=A(Y)$、$B=A(X)$。
由（\hyperlink{para:1:6.3.2}{6.3.2}），$X$ 有有限仿射开覆盖 $(V_i)$，使各 $A(V_i)$ 为有限型 $A$-代数。
由于 $V_i$ 拟紧，可用有限个 $D(g_{ij})\subset V_i$、$g_{ij}\in B$ 覆盖它。
若 $\varphi_i:B\to A(V_i)$ 对应于典范嵌入 $V_i\to X$，则
\[
 B_{g_{ij}}=A(V_i)_{\varphi_i(g_{ij})}=A(V_i)[1/\varphi_i(g_{ij})],
\]
故 $B_{g_{ij}}$ 为有限型 $A$-代数。
于是可归结到 $V_i=D(g_i)$、$g_i\in B$ 的情形。
由假设存在有限子集 $F_i\subset B$ 及整数 $n_i\geq0$，
使 $B_{g_i}$ 由元素 $b_i/g_i^{n_i}$（$b_i\in F_i$）作为 $A$-代数生成。
因 $g_i$ 只有有限个，可假设各 $n_i$ 相等。
又因 $D(g_i)$ 覆盖 $X$，各 $g_i$ 在 $B$ 中生成单位理想，
故存在 $h_i\in B$ 使 $\sum_i h_i g_i=1$。

令 $F$ 为各 $F_i$ 以及所有 $g_i,h_i$ 的并。证明 $B'=A[F]$ 等于 $B$。
对任意 $b\in B$ 及每个 $i$，$b$ 在 $B_{g_i}$ 中的典范像形如
$b'_i/g_i^{m_i}$，其中 $b'_i\in B'$。
乘以适当的 $g_i$ 的幂后可使各 $m_i$ 同为整数 $m$。
由分式环的定义，存在依赖于 $b$ 的整数 $N>m$，使 $g_i^N b\in B'$ 对每个 $i$ 成立。
在环 $B'$ 中，各 $g_i^N$ 生成单位理想，因为各 $g_i$ 已经如此，且 $h_i\in B'$。
故有 $c_i\in B'$ 使 $\sum_i c_i g_i^N=1$，从而
$b=\sum_i c_i g_i^N b\in B'$。证毕。

\egastatement{6.3.4}{命题}
\begin{enumerate}[label=(\roman*)]
\item 每个闭浸入为有限型态射。
\item 两个有限型态射的复合为有限型态射。
\item 两个有限型 $S$-态射 $f:X\to X'$、$g:Y\to Y'$ 的积 $f\times_S g$ 为有限型态射。
\item 若 $f:X\to Y$ 为有限型 $S$-态射，则任意基预概形扩张 $S'\to S$ 所得的 $f_{(S')}$ 为有限型态射。
\item 若 $g\circ f$ 为有限型态射且 $g$ 分离，则 $f$ 为有限型态射。
\item 若 $f$ 为有限型态射，则 $f_{\mathrm{red}}$ 为有限型态射。
\end{enumerate}

根据（\hyperlink{para:1:5.5.12}{5.5.12}），只须证明 (i)、(ii)、(iv)。
为证明 (i)，可限于闭子预概形的典范嵌入 $X\to Y$，
并可假设 $Y$ 仿射（\hyperlink{para:1:6.3.2}{6.3.2}）。
此时 $X$ 仿射（\hyperlink{para:1:4.2.3}{4.2.3}），其环为 $A/\mathfrak J$，其中 $A=A(Y)$。
$A/\mathfrak J$ 在 $A$ 上有限型，故结论成立。

为证明 (ii)，设 $f:X\to Y$、$g:Y\to Z$ 为有限型态射，$U$ 为 $Z$ 的仿射开集。
$g^{-1}(U)$ 有有限仿射开覆盖 $(V_i)$，使 $A(V_i)$ 在 $A(U)$ 上有限型（\hyperlink{para:1:6.3.2}{6.3.2}）。
同样，每个 $f^{-1}(V_i)$ 有有限仿射开覆盖 $(W_{ij})$，
使 $A(W_{ij})$ 在 $A(V_i)$ 上有限型，故在 $A(U)$ 上有限型。结论成立。

最后证明 (iv)。可限于 $S=Y$ 的情形：
$f_{(S')}$ 也等于把 $f$ 看作 $Y$-态射、沿 $Y_{(S')}\to Y$ 扩张所得的 $f_{(Y_{(S')})}$（\hyperlink{para:1:3.3.9}{3.3.9}）。
以 $p:X_{(S')}\to X$、$q:X_{(S')}\to S'$ 表示投影。
若 $V\subset S$ 为仿射开集，则 $f^{-1}(V)$ 是有限个仿射开集 $W_i$ 的并，
且各 $A(W_i)$ 在 $A(V)$ 上有限型（\hyperlink{para:1:6.3.2}{6.3.2}）。
设 $V'\subset S'$ 为包含于 $g^{-1}(V)$ 的仿射开集。
由 $f\circ p=g\circ q$，$q^{-1}(V')$ 包含于各 $p^{-1}(W_i)$ 的并。
而 $p^{-1}(W_i)\cap q^{-1}(V')$ 认同于 $W_i\times_V V'$（\hyperlink{para:1:3.2.7}{3.2.7}），
它仿射，其环同构于 $A(W_i)\otimes_{A(V)}A(V')$（\hyperlink{para:1:3.2.2}{3.2.2}）。
此环在 $A(V')$ 上有限型，命题得证。

\egastatement{6.3.5}{推论}
设 $f:X\to Y$ 为浸入。
若 $Y$ 的底空间局部诺特，或 $X$ 的底空间诺特，则 $f$ 为有限型态射。

总可假设 $Y$ 仿射（\hyperlink{para:1:6.3.2}{6.3.2}）。
若 $Y$ 的底空间局部诺特，还可假设它诺特；此时子空间 $X$ 也诺特。
所以可以假设 $Y$ 仿射、$X$ 的底空间诺特。
由 $X$ 在 $Y$ 中局部闭（\hyperlink{para:1:4.1.3}{4.1.3}），存在由有限个 $Y$ 的仿射开集
$D(g_i)$（$g_i\in A(Y)$）组成的 $X$ 的覆盖，
使 $X\cap D(g_i)$ 在 $D(g_i)$ 中闭，因而仿射（\hyperlink{para:1:4.2.3}{4.2.3}）。
由（\hyperlink{para:1:6.3.4}{6.3.4}，(i)）和（\hyperlink{para:1:6.3.3}{6.3.3}），$A(X\cap D(g_i))$ 在 $A(D(g_i))$ 上有限型；
而 $A(D(g_i))=A(Y)[1/g_i]$ 在 $A(Y)$ 上有限型，证明完成。

\egastatement{6.3.6}{推论}
设 $f:X\to Y$、$g:Y\to Z$ 为态射。
若 $g\circ f$ 为有限型态射，且 $X$ 诺特或 $X\times_Z Y$ 局部诺特，则 $f$ 为有限型态射。

这由（\hyperlink{para:1:5.5.12}{5.5.12}）的证明以及将（\hyperlink{para:1:6.3.5}{6.3.5}）应用于图浸入 $\Gamma_f$ 立即得到。

\egastatement{6.3.7}{命题}
设 $f:X\to Y$ 为有限型态射。
若 $Y$ 诺特（分别地，局部诺特），则 $X$ 诺特（分别地，局部诺特）。

只须在 $Y$ 诺特时证明。
此时 $Y$ 是有限个仿射开集 $V_i$ 的并，且各 $A(V_i)$ 诺特。
由（\hyperlink{para:1:6.3.2}{6.3.2}），每个 $f^{-1}(V_i)$ 为有限个仿射开集 $W_{ij}$ 的并，
各 $A(W_{ij})$ 为诺特环 $A(V_i)$ 上的有限型代数，故为诺特环。
所以 $X$ 诺特。

\egastatement{6.3.8}{推论}
设 $X$ 为有限型 $S$-预概形。
对任意使 $S'$ 诺特（分别地，局部诺特）的基扩张 $S'\to S$，
$X_{(S')}$ 诺特（分别地，局部诺特）。

由（\hyperlink{para:1:6.3.4}{6.3.4}，(iv)），$X_{(S')}$ 在 $S'$ 上有限型，再用（\hyperlink{para:1:6.3.7}{6.3.7}）即可。
也可表述为：在 $S$-预概形的积 $X\times_S Y$ 中，
若一因子在 $S$ 上有限型，另一因子诺特（分别地，局部诺特），
则该积诺特（分别地，局部诺特）。

\egastatement{6.3.9}{推论}
设 $S$ 为局部诺特预概形，$X$ 为有限型 $S$-预概形。
则每个 $S$-态射 $f:X\to Y$ 都为有限型态射。

可假设 $S$ 诺特。若 $\varphi:X\to S$、$\psi:Y\to S$ 为结构态射，
则 $\varphi=\psi\circ f$；由（\hyperlink{para:1:6.3.7}{6.3.7}），$X$ 诺特，
故 $f$ 由（\hyperlink{para:1:6.3.6}{6.3.6}）为有限型态射。

\egastatement{6.3.10}{命题}
设 $f:X\to Y$ 为有限型态射。
$f$ 为满射的充要条件是：对每个代数闭域 $\Omega$，
由 $f$ 对应的映射 $X(\Omega)\to Y(\Omega)$（\hyperlink{para:1:3.4.1}{3.4.1}）为满射。

条件充分：对每个 $y\in Y$，考虑 $k(y)$ 的代数闭扩张 $\Omega$，
及交换图
\[
\begin{tikzcd}
X\arrow[d,"f"']&\Spec(\Omega)\arrow[l]\arrow[dl]\\
Y&
\end{tikzcd}
\]
即可看出（参见 \hyperlink{para:1:3.5.3}{3.5.3}）。
反之，设 $f$ 满射，$g:\Spec(\Omega)\to Y$ 为态射，$\Omega$ 代数闭。
考虑图
\[
\begin{tikzcd}
X\arrow[d,"f"']&X_{(\Omega)}\arrow[l]\arrow[d,"f_{(\Omega)}"]\\
Y&\Spec(\Omega)\arrow[l,"g"]
\end{tikzcd}
\]
只须证明 $X_{(\Omega)}$ 有一个 $\Omega$-有理点（\hyperlink{para:1:3.3.14}{3.3.14}、\hyperlink{para:1:3.4.3}{3.4.3}、\hyperlink{para:1:3.4.4}{3.4.4}）。
$f$ 满射，故 $X_{(\Omega)}$ 非空；$f$ 有限型，故 $f_{(\Omega)}$ 有限型（\hyperlink{para:1:6.3.4}{6.3.4}，(iv)）。
所以 $X_{(\Omega)}$ 有非空仿射开集 $Z$，其环 $A(Z)$ 为非零有限型 $\Omega$-代数。
由希尔伯特零点定理（文献 \cite{b21}），存在 $\Omega$-同态 $A(Z)\to\Omega$，
从而存在 $X_{(\Omega)}$ 在 $\Spec(\Omega)$ 上的截面，命题得证。


% END chapters/ch1_s6_3.tex

\par

% BEGIN chapters/ch1_s6_4.tex
\subsection{代数预概形}\hypertarget{heading:1:6.4}{}
\src{147—150}

\egastatement{6.4.1}{定义}
给定域 $K$，称 $K$ 上的有限型预概形 $X$ 为\emph{代数 $K$-预概形}；
$K$ 称为 $X$ 的\emph{基域}。
若 $X$ 还是概形，等价地是 $K$-概形（\hyperlink{para:1:5.5.8}{5.5.8}），则称它为\emph{代数 $K$-概形}。

每个代数 $K$-预概形都是诺特的（\hyperlink{para:1:6.3.7}{6.3.7}）。

\egastatement{6.4.2}{命题}
设 $X$ 为代数 $K$-预概形。
点 $x\in X$ 为闭点的充要条件是 $k(x)$ 为 $K$ 的有限次代数扩张。

可假设 $X$ 仿射，其环 $A$ 为有限型 $K$-代数。
事实上，满足 $A(U)$ 为有限型 $K$-代数的仿射开集 $U$ 构成 $X$ 的覆盖（\hyperlink{para:1:6.3.1}{6.3.1}）。
此时 $X$ 的闭点正是对应理想 $\jj_x$ 极大，亦即 $A/\jj_x$ 为域的点，
且该域必然为 $k(x)$。
$A/\jj_x$ 是有限型 $K$-代数，所以若 $x$ 闭，$k(x)$ 是作为 $K$-代数有限型的域，
故必为有限秩 $K$-代数（文献 \cite{b21}）。
反之，若 $k(x)$ 在 $K$ 上秩有限，则 $A/\jj_x\subset k(x)$ 在 $K$ 上也秩有限。
每个作为有限秩 $K$-代数的整环都是域，故 $A/\jj_x=k(x)$，所以 $x$ 闭。

\egastatement{6.4.3}{推论}
设 $K$ 为代数闭域，$X$ 为代数 $K$-预概形。
则 $X$ 的闭点正是 $K$-有理点（\hyperlink{para:1:3.4.4}{3.4.4}），并典范地认同于 $X$ 的 $K$-值点。

\egastatement{6.4.4}{命题}
设 $X$ 为域 $K$ 上的代数预概形。下述性质等价：
\begin{enumerate}[label=\alph*)]
\item $X$ 为阿廷的；
\item $X$ 的底空间离散；
\item $X$ 的底空间只有有限个闭点；
\end{enumerate}
\noindent c') $X$ 的底空间有限；
\begin{enumerate}[label=\alph*),start=4]
\item $X$ 的所有点都是闭点；
\item $X$ 同构于 $\Spec(A)$，其中 $A$ 为有限秩 $K$-代数。
\end{enumerate}

由于 $X$ 诺特，由（\hyperlink{para:1:6.2.2}{6.2.2}），a)、b)、d) 等价，并蕴含 c)、c')；
e) 显然蕴含 a)。剩下证明 c) 蕴含 d)、e)；可限于 $X$ 仿射的情形。
此时 $A(X)$ 是有限型 $K$-代数（\hyperlink{para:1:6.3.3}{6.3.3}），所以是雅各布森环
（文献 \cite{b1}，3-11 至 3-12 页），并且由假设只有有限个极大理想。
有限个素理想的交若仍为素理想，必等于其中一个，
所以 $A(X)$ 的每个素理想都极大，d) 得证。
此外，已知此时 $A(X)$ 是有限型阿廷 $K$-代数（\hyperlink{para:1:6.2.2}{6.2.2}），
所以必为有限秩 $K$-代数（文献 \cite{b21}）。

\egapar{6.4.5}
当（\hyperlink{para:1:6.4.4}{6.4.4}）的条件满足时，称 $X$ 为 $K$ 上的\emph{有限概形}
（参见 II，6.1.1），或\emph{有限 $K$-概形}。
其秩为 $[A:K]$，亦记为 $\operatorname{rg}_K(X)$。
若 $X,Y$ 为两个有限 $K$-概形，则
\begin{equation}
 \operatorname{rg}_K(X\amalg Y)=\operatorname{rg}_K(X)+\operatorname{rg}_K(Y),
 \tag{6.4.5.1}\label{eq:1:6.4.5.1}
\end{equation}
\begin{equation}
 \operatorname{rg}_K(X\times_KY)=\operatorname{rg}_K(X)\operatorname{rg}_K(Y),
 \tag{6.4.5.2}\label{eq:1:6.4.5.2}
\end{equation}
这由（\hyperlink{para:1:3.2.2}{3.2.2}）得到。

\egastatement{6.4.6}{推论}
设 $X$ 为有限 $K$-概形。
对 $K$ 的任意扩张 $K'$，$X\otimes_K K'$ 为有限 $K'$-概形，
其在 $K'$ 上的秩等于 $X$ 在 $K$ 上的秩。

事实上，若 $A=A(X)$，则 $[A\otimes_K K':K']=[A:K]$。

\egastatement{6.4.7}{推论}
设 $X$ 为有限 $K$-概形，置
\[
 n=\sum_{x\in X}[k(x):K]_s.
\]
这里回忆：$[K':K]_s$ 表示扩张 $K'/K$ 的可分次数，
即包含于 $K'$ 的 $K$ 的最大代数可分扩张的次数。
则对 $K$ 的任意代数闭扩张 $\Omega$，$X\otimes_K\Omega$ 的底空间恰有 $n$ 个点，
这些点与 $X$ 的 $\Omega$-值点相认同。

显然可限于 $A=A(X)$ 为局部环的情形（\hyperlink{para:1:6.2.2}{6.2.2}）。
设 $\mathfrak m$ 为其极大理想，$L=A/\mathfrak m$ 为其剩余域，即 $K$ 的一个代数扩张。
$X$ 的 $\Omega$-值点与 $X\otimes_K\Omega$ 的 $\Omega$-截面一一对应（\hyperlink{para:1:3.4.1}{3.4.1} 和 \hyperlink{para:1:3.3.14}{3.3.14}），
也与 $L\to\Omega$ 的 $K$-同态一一对应（\hyperlink{para:1:1.7.3}{1.7.3}）。
因此，由 Bourbaki《代数》第 V 章 §7 第 5 节命题 8，并利用（\hyperlink{para:1:6.4.3}{6.4.3}），得到本推论。

\egapar{6.4.8}
（\hyperlink{para:1:6.4.7}{6.4.7}）中定义的数 $n$ 称为 $A$ 或 $X$ 在 $K$ 上的\emph{可分秩}，
也称为 $X$ 的\emph{几何点数}；它等于 $X(\Omega)_K$ 的元素数。
由此定义立即可知，对 $K$ 的任意扩张 $K'$，$X\otimes_K K'$ 与 $X$ 的几何点数相同。
以 $n(X)$ 表示此数。若 $X,Y$ 为两个有限 $K$-概形，显然
\begin{equation}
 n(X\amalg Y)=n(X)+n(Y).
 \tag{6.4.8.1}\label{eq:1:6.4.8.1}
\end{equation}
同样有
\begin{equation}
 n(X\times_KY)=n(X)n(Y),
 \tag{6.4.8.2}\label{eq:1:6.4.8.2}
\end{equation}
这由将 $n(X)$ 解释为 $X(\Omega)_K$ 的元素数及公式（\hyperref[eq:1:3.4.3.1]{3.4.3.1}）立即得到。

\egastatement{6.4.9}{命题}
设 $K$ 为域，$X,Y$ 为两个代数 $K$-预概形，$f:X\to Y$ 为 $K$-态射，
$\Omega$ 为 $K$ 的一个代数闭扩张，且在 $K$ 上的超越次数无限。
则 $f$ 为满射的充要条件是对应映射
$X(\Omega)_K\to Y(\Omega)_K$（\hyperlink{para:1:3.4.1}{3.4.1}）为满射。

必要性由（\hyperlink{para:1:6.3.10}{6.3.10}）得到，因为 $f$ 必为有限型态射（\hyperlink{para:1:6.3.9}{6.3.9}）。
充分性的论证与（\hyperlink{para:1:6.3.10}{6.3.10}）相同；只须注意，对每个 $y\in Y$，
$k(y)$ 是 $K$ 的有限型扩张，所以 $K$-同构于 $\Omega$ 的一个子域。

\egastatement{6.4.10}{注}
将在第 IV 章看到，（\hyperlink{para:1:6.4.9}{6.4.9}）的结论不加关于 $\Omega/K$ 的超越次数的假设也成立。

\egastatement{6.4.11}{命题}
若 $f:X\to Y$ 为有限型态射，则对每个 $y\in Y$，纤维 $f^{-1}(y)$
是剩余域 $k(y)$ 上的代数预概形；且对每个 $x\in f^{-1}(y)$，
$k(x)$ 是 $k(y)$ 的有限型扩张。

由于 $f^{-1}(y)=X\otimes_Y k(y)$，命题由（\hyperlink{para:1:6.3.4}{6.3.4}，(iv)）和（\hyperlink{para:1:6.3.3}{6.3.3}）得到。

\egastatement{6.4.12}{命题}
设 $f:X\to Y$、$g:Y'\to Y$ 为两个态射，置 $X'=X\times_Y Y'$、$f'=f_{(Y')}$。
取 $y'\in Y'$，令 $y=g(y')$。
若纤维 $f^{-1}(y)$ 为有限 $k(y)$-概形，则纤维 $f'^{-1}(y')$ 为有限 $k(y')$-概形，
且其秩和几何点数分别等于 $f^{-1}(y)$ 的秩和几何点数。

利用纤维的传递性，结论由（\hyperlink{para:1:6.4.6}{6.4.6}）和（\hyperlink{para:1:6.4.8}{6.4.8}）立即得到。
\transnote{原文在此将纤维的传递性引为（\hyperlink{para:1:3.6.5}{3.6.5}）；对应的陈述为（\hyperlink{para:1:3.6.4}{3.6.4}）。}

\egapar{6.4.13}
命题（\hyperlink{para:1:6.4.11}{6.4.11}）表明，有限型态射直观上对应于“代数簇的代数族”，
其中 $Y$ 的点扮演“参数”的角色；这赋予此类态射以“几何”意义。
非有限型态射在下文中主要出现于“基预概形的变换”问题，例如局部化或完备化。


% END chapters/ch1_s6_4.tex

\par

% BEGIN chapters/ch1_s6_5.tex
\subsection{态射的局部确定}\hypertarget{heading:1:6.5}{}
\src{150—152}

\egastatement{6.5.1}{命题}
设 $X,Y$ 为两个 $S$-预概形，$Y$ 在 $S$ 上有限型；
设 $x\in X$、$y\in Y$ 位于同一点 $s\in S$ 之上。
\begin{enumerate}[label=(\roman*)]
\item 若两个 $S$-态射 $f=(\psi,\theta)$、$f'=(\psi',\theta')$ 从 $X$ 到 $Y$，
满足 $\psi(x)=\psi'(x)=y$，且局部 $\OO_s$-同态
$\theta_x^\sharp,\theta_x'^\sharp:\OO_y\to\OO_x$ 相同，
则 $f,f'$ 在 $x$ 的一个开邻域内相同。
\item 再设 $S$ 局部诺特。则对每个局部 $\OO_s$-同态
$\varphi:\OO_y\to\OO_x$，存在 $x$ 的开邻域 $U$
及 $S$-态射 $f=(\psi,\theta):U\to Y$，使
$\psi(x)=y$、$\theta_x^\sharp=\varphi$。
\end{enumerate}

(i) 问题在 $S,X,Y$ 上具有局部性，故可假设它们仿射，环分别为 $A,B,C$。
$f,f'$ 分别形如 $({}^a\varphi,\wt\varphi)$、$({}^a\varphi',\wt\varphi')$，
其中 $\varphi,\varphi':C\to B$ 为 $A$-同态，
$\varphi^{-1}(\jj_x)=\varphi'^{-1}(\jj_x)=\jj_y$，
且由它们诱导的同态 $\varphi_x,\varphi'_x:C_y\to B_x$ 相同。
还可假设 $C$ 为有限型 $A$-代数。
取生成元 $c_i$（$1\leq i\leq n$），置
$b_i=\varphi(c_i)$、$b'_i=\varphi'(c_i)$。
由假设在 $B_x$ 中 $b_i/1=b'_i/1$，
所以有 $s_i\in B\setminus\jj_x$ 使 $s_i(b_i-b'_i)=0$。
可取所有 $s_i$ 相同，记为 $g\in B\setminus\jj_x$。
于是 $B_g$ 中各 $b_i/1=b'_i/1$。
若 $i_g:B\to B_g$ 为典范同态，则 $i_g\circ\varphi=i_g\circ\varphi'$，
故 $f,f'$ 在 $D(g)$ 上的限制相同。

(ii) 可归结到与 (i) 相同的情形，并再假设 $A$ 诺特。
取 $A$-代数 $C$ 的生成元 $c_i$（$1\leq i\leq n$），
设 $\alpha:A[X_1,\ldots,X_n]\to C$ 为将 $X_i$ 映为 $c_i$ 的满同态。
设 $i_y:C\to C_y$ 为典范同态，考虑复合同态
\[
 \beta:A[X_1,\ldots,X_n]\xrightarrow{\alpha}C
 \xrightarrow{i_y}C_y\xrightarrow{\varphi}B_x.
\]
记 $\mathfrak a=\Ker\beta$。
由于 $A$ 诺特，$A[X_1,\ldots,X_n]$ 诺特，故 $\mathfrak a$ 有有限生成系
$Q_j(X_1,\ldots,X_n)$（$1\leq j\leq m$）。
每个 $\varphi(i_y(c_i))$ 可写成 $b_i/s_i$，其中 $b_i\in B$、$s_i\notin\jj_x$；
并可取各分母均为同一个 $g\notin\jj_x$。
由假设，在 $B_x$ 中有 $Q_j(b_1/g,\ldots,b_n/g)=0$。
写
\[
 Q_j(X_1/T,\ldots,X_n/T)=P_j(X_1,\ldots,X_n,T)/T^{k_j},
\]
其中 $P_j$ 为 $k_j$ 次齐次多项式。
置 $z_j=P_j(b_1,\ldots,b_n,g)\in B$。
由假设存在 $h_j\in B\setminus\jj_x$ 使 $h_jz_j=0$，并可使各 $h_j$ 同为 $h$。
于是 $P_j(hb_1,\ldots,hb_n,hg)=0$。
\transnote{此处可在齐次化时取 $k_j\geq1$；若直接取零次多项式的次数为零，
从 $hz_j=0$ 到所写等式的步骤需改为在 $B_{hg}$ 中验证。原文未说明这一选择。}

考虑将 $X_i$ 映为 $hb_i/(hg)$ 的同态
$\rho:A[X_1,\ldots,X_n]\to B_{hg}$。
$\mathfrak a$ 在该同态下的像为零，故 $\Ker\alpha\subset\mathfrak a$ 的像也为零。
所以 $\rho$ 分解为
$A[X_1,\ldots,X_n]\xrightarrow{\alpha}C\xrightarrow{\gamma}B_{hg}$，
且 $\gamma(c_i)=hb_i/(hg)$。
若 $i_x:B_{hg}\to B_x$ 为典范同态，则图
\begin{equation}
\begin{tikzcd}
C\arrow[r,"\gamma"]\arrow[d,"i_y"']&B_{hg}\arrow[d,"i_x"]\\
C_y\arrow[r,"\varphi"']&B_x
\end{tikzcd}
\tag{6.5.1.1}\label{eq:1:6.5.1.1}
\end{equation}
交换。因此 $\varphi=\gamma_x$；因 $\varphi$ 是局部同态，${}^a\gamma(x)=y$。
所以 $f=({}^a\gamma,\wt\gamma)$ 是从邻域 $D(hg)$ 到 $Y$ 的满足要求的 $S$-态射。

\egastatement{6.5.2}{推论}
在（\hyperlink{para:1:6.5.1}{6.5.1}，(ii)）的假设下，若 $X$ 在 $S$ 上也有限型，则可取 $f$ 为有限型态射。

这由（\hyperlink{para:1:6.3.6}{6.3.6}）得到。

\egastatement{6.5.3}{推论}
在（\hyperlink{para:1:6.5.1}{6.5.1}，(ii)）的假设下，再设 $Y$ 为整预概形且 $\varphi$ 为单同态。
则可取 $f=({}^a\gamma,\wt\gamma)$，其中 $\gamma$ 为单同态。

可假设 $C$ 为整环（\hyperlink{para:1:5.1.4}{5.1.4}），于是 $i_y$ 为单同态；
由图（\hyperref[eq:1:6.5.1.1]{6.5.1.1}）得到 $\gamma$ 为单同态。

\egastatement{6.5.4}{命题}
设 $f=(\psi,\theta):X\to Y$ 为有限型态射，$x\in X$、$y=\psi(x)$。
\begin{enumerate}[label=(\roman*)]
\item $f$ 在 $x$ 处为局部浸入（\hyperlink{para:1:4.5.1}{4.5.1}）的充要条件是
$\theta_x^\sharp:\OO_y\to\OO_x$ 为满射。
\item 再设 $Y$ 局部诺特。$f$ 在 $x$ 处为局部同构（\hyperlink{para:1:4.5.2}{4.5.2}）的充要条件是
$\theta_x^\sharp$ 为同构。
\end{enumerate}

(ii) 由（\hyperlink{para:1:6.5.1}{6.5.1}），存在 $y$ 的开邻域 $V$ 及态射 $g:V\to X$，
使 $g\circ f$ 和 $f\circ g$ 分别在 $x$ 和 $y$ 的某个邻域内有定义且等于恒等态射。
由此容易得到 $f$ 为局部同构。

(i) 问题在 $X,Y$ 上局部，可假设它们仿射，环分别为 $A,B$。
$f=({}^a\varphi,\wt\varphi)$，其中 $\varphi:B\to A$ 使 $A$ 成为有限型 $B$-代数；
$\varphi^{-1}(\jj_x)=\jj_y$，且诱导的同态 $\varphi_x:B_y\to A_x$ 为满射。
设 $t_i$（$1\leq i\leq n$）为 $B$-代数 $A$ 的生成元。
由满射性，存在 $b_i\in B$ 及 $c\in B\setminus\jj_y$，
使 $A_x$ 中有 $t_i/1=\varphi(b_i)/\varphi(c)$。
于是存在 $a\in A\setminus\jj_x$，令 $g=a\varphi(c)$，
则在 $A_g$ 中也有 $t_i/1=a\varphi(b_i)/g$（\hyperlink{para:1:1.3.3}{1.3.3}）。

由假设存在系数在 $\varphi(B)$ 中的多项式 $Q$ 使 $a=Q(t_1,\ldots,t_n)$。
写
\[
 Q(X_1/T,\ldots,X_n/T)=P(X_1,\ldots,X_n,T)/T^m,
\]
其中 $P$ 为 $m$ 次齐次多项式。
在 $A_g$ 中有
\[
 a/1=a^mP(\varphi(b_1),\ldots,\varphi(b_n),\varphi(c))/g^m
     =a^m\varphi(d)/g^m
\]
其中 $d\in B$。
由定义 $g/1=(a/1)(\varphi(c)/1)$ 可逆，所以 $a/1$、$\varphi(c)/1$ 均可逆；
从而可写 $a/1=(\varphi(d)/1)(\varphi(c)/1)^{-m}$。
因此 $\varphi(d)/1$ 在 $A_g$ 中也可逆。
置 $h=cd$，则 $\varphi(h)/1$ 在 $A_g$ 中可逆，
所以 $B\xrightarrow{\varphi}A\to A_g$ 分解为
$B\to B_h\xrightarrow{\gamma}A_g$（0，\hyperlink{para:0:1.2.4}{1.2.4}）。

证明 $\gamma$ 为满射，只须证明其像包含各 $t_i/1$ 及 $(g/1)^{-1}$。
事实上，
\[
 (g/1)^{-1}=(\varphi(c)/1)^{m-1}(\varphi(d)/1)^{-1}
          =\gamma(c^m/h),
\]
且 $a/1=\gamma(d^{m+1}/h^m)$，所以
$(a\varphi(b_i))/1=\gamma(b_i d^{m+1}/h^m)$。
又 $t_i/1=(a\varphi(b_i)/1)(g/1)^{-1}$，故满射性得证。
$h$ 的选择使 $\psi(D(g))\subset D(h)$，
且 $f|D(g)=({}^a\gamma,\wt\gamma)$。
因为 $\gamma$ 为满射，该限制是 $D(g)$ 到 $D(h)$ 的闭浸入（\hyperlink{para:1:4.2.3}{4.2.3}）。

\egastatement{6.5.5}{推论}
设 $f=(\psi,\theta):X\to Y$ 为有限型态射，$X$ 不可约，$x$ 为其泛点，$y=\psi(x)$。
\begin{enumerate}[label=(\roman*)]
\item $f$ 在 $X$ 的某一点处为局部浸入的充要条件是
$\theta_x^\sharp:\OO_y\to\OO_x$ 为满射。
\item 再设 $Y$ 不可约且局部诺特。
$f$ 在 $X$ 的某一点处为局部同构的充要条件是：
$y$ 为 $Y$ 的泛点，等价地 $f$ 支配，且 $\theta_x^\sharp$ 为同构；
亦即 $f$ 为双有理态射（\hyperlink{para:1:2.2.9}{2.2.9}）。
\end{enumerate}

由于 $X$ 的每个非空开集都包含 $x$，(i) 显然由（\hyperlink{para:1:6.5.4}{6.5.4}，(i)）得到；
(ii) 同样由（\hyperlink{para:1:6.5.4}{6.5.4}，(ii)）得到。


% END chapters/ch1_s6_5.tex

\par

% BEGIN chapters/ch1_s6_6.tex
\subsection{拟紧态射与局部有限型态射}\hypertarget{heading:1:6.6}{}
\src{152—155}

\egastatement{6.6.1}{定义}
若 $Y$ 的每个拟紧开集在态射 $f:X\to Y$ 下的逆像都是拟紧的，
则称 $f$ 为\emph{拟紧态射}。

设 $\mathfrak B$ 为由拟紧开集组成的 $Y$ 的拓扑基，例如仿射开集基。
$f$ 拟紧的充要条件是：$\mathfrak B$ 中每个成员在 $f$ 下的逆像拟紧，
等价地是有限个仿射开集的并；因为 $Y$ 的每个拟紧开集都是有限个基元素的并。
例如，若 $X$ 拟紧而 $Y$ 仿射，则任意态射 $f:X\to Y$ 均拟紧。
事实上，$X$ 是有限个仿射开集 $U_i$ 的并，
且对 $Y$ 的每个仿射开集 $V$，$U_i\cap f^{-1}(V)$ 仿射（\hyperlink{para:1:5.5.10}{5.5.10}），从而拟紧。

若 $f:X\to Y$ 拟紧，则对每个开集 $V\subset Y$，
限制 $f^{-1}(V)\to V$ 显然拟紧。
反之，若 $(U_\alpha)$ 为 $Y$ 的开覆盖，且各限制
$f^{-1}(U_\alpha)\to U_\alpha$ 拟紧，则 $f$ 拟紧。

\egastatement{6.6.2}{定义}
若对每个 $x\in X$，存在其开邻域 $U$ 及 $y=f(x)$ 的开邻域 $V\supset f(U)$，
使 $f|U:U\to V$ 为有限型态射，
则称态射 $f:X\to Y$ 为\emph{局部有限型态射}。
此时也称 $X$ 为 $Y$ 上的\emph{局部有限型预概形}，或\emph{局部有限型 $Y$-预概形}。

由（\hyperlink{para:1:6.3.2}{6.3.2}）立即可知：若 $f$ 局部有限型，则对每个开集 $W\subset Y$，
限制 $f^{-1}(W)\to W$ 局部有限型。
若 $Y$ 局部诺特，且 $X$ 在 $Y$ 上局部有限型，则 $X$ 局部诺特（\hyperlink{para:1:6.3.7}{6.3.7}）。

\egastatement{6.6.3}{命题}
态射 $f:X\to Y$ 为有限型的充要条件是：它拟紧且局部有限型。

必要性由（\hyperlink{para:1:6.3.1}{6.3.1}）及（\hyperlink{para:1:6.6.1}{6.6.1}）之后的说明立即得到。
反之，设这些条件成立，$U$ 为 $Y$ 的仿射开集，环为 $A$。
对每个 $x\in f^{-1}(U)$，由假设存在开邻域 $V(x)\subset f^{-1}(U)$
及 $y=f(x)$ 的开邻域 $W(x)\subset U$，
使 $f(V(x))\subset W(x)$ 且 $V(x)\to W(x)$ 为有限型态射。
用形如 $D(g)$、$g\in A$ 的较小邻域 $W_1(x)\subset W(x)$ 代替 $W(x)$，
并以 $V(x)\cap f^{-1}(W_1(x))$ 代替 $V(x)$，
可假设 $W(x)=D(g)$。
此时 $W(x)$ 的环为 $A[1/g]$，所以它在 $U$ 上有限型，
从而 $V(x)$ 在 $U$ 上有限型。
又因 $f^{-1}(U)$ 拟紧，它是有限个 $V(x_i)$ 的并，证明完成。

\egastatement{6.6.4}{命题}
\begin{enumerate}[label=(\roman*)]
\item 浸入 $X\to Y$ 在下列任一种情形下拟紧：它是闭浸入；
$Y$ 的底空间局部诺特；$X$ 的底空间诺特。
\item 两个拟紧态射的复合拟紧。
\item 若 $f:X\to Y$ 为拟紧 $S$-态射，则任意基扩张 $g:S'\to S$
所得的 $f_{(S')}:X_{(S')}\to Y_{(S')}$ 拟紧。
\item 两个拟紧 $S$-态射 $f:X\to X'$、$g:Y\to Y'$ 的积 $f\times_S g$ 拟紧。
\item 若 $f:X\to Y$、$g:Y\to Z$ 的复合拟紧，
且 $g$ 分离，或 $X$ 的底空间局部诺特，则 $f$ 拟紧。
\item $f$ 拟紧的充要条件是 $f_{\mathrm{red}}$ 拟紧。
\end{enumerate}

(vi) 显然，因为拟紧性只依赖于底空间的连续映射。
先证明 (v) 中 $X$ 的底空间局部诺特的情形。
置 $h=g\circ f$，设 $U$ 为 $Y$ 的拟紧开集。
$g(U)$ 是 $Z$ 中拟紧但未必开的子空间，故包含于有限个拟紧开集 $V_j$ 的并（\hyperlink{para:1:2.1.3}{2.1.3}）。
所以 $f^{-1}(U)$ 包含于各 $h^{-1}(V_j)$ 的并。
这些是 $X$ 中的拟紧子空间，从而是诺特子空间。
由（0，\hyperlink{para:0:2.2.3}{2.2.3}），$f^{-1}(U)$ 为诺特空间，当然拟紧。

其余断言只须证明 (i)、(ii)、(iii)（\hyperlink{para:1:5.5.12}{5.5.12}）。
(ii) 显然；(i) 在 $Y$ 的底空间局部诺特或 $X$ 的底空间诺特时由（\hyperlink{para:1:6.3.5}{6.3.5}）得到，
对闭浸入则显然。
为证明 (iii)，可限于 $S=Y$（\hyperlink{para:1:3.3.11}{3.3.11}）。
置 $f'=f_{(S')}$，设 $U'$ 为 $S'$ 的拟紧开集。
对每个 $s'\in U'$，取 $g(s')$ 在 $S$ 中的仿射开邻域 $T$，
及 $s'$ 的仿射开邻域 $W\subset U'\cap g^{-1}(T)$。
只须证明 $f'^{-1}(W)$ 拟紧；即归结到 $S,S'$ 都仿射时证明 $X\times_S S'$ 的底空间拟紧。
此时由假设 $X$ 为有限个仿射开集 $V_j$ 的并，
故 $X\times_S S'$ 为仿射概形 $V_j\times_S S'$ 的底空间之并（\hyperlink{para:1:3.2.2}{3.2.2} 和 \hyperlink{para:1:3.2.7}{3.2.7}）。命题得证。

另外，若 $X=X'\amalg X''$ 是两个预概形的和，
则态射 $f:X\to Y$ 拟紧的充要条件是它在 $X',X''$ 上的限制都拟紧。

\egastatement{6.6.5}{命题}
设 $f:X\to Y$ 为拟紧态射。
$f$ 支配的充要条件是：对 $Y$ 的每个不可约分支的泛点 $y$，
$f^{-1}(y)$ 包含 $X$ 的某个不可约分支的泛点。

充分性立即成立，且不必假设 $f$ 拟紧。
为证明必要性，取 $y$ 的仿射开邻域 $U$。
$f^{-1}(U)$ 拟紧，故为有限个仿射开集 $V_i$ 的并。
$f$ 支配的假设蕴含 $y$ 属于某个 $f(V_i)$ 在 $U$ 中的闭包。
显然可假设 $X,Y$ 约化。
由于 $V_i$ 的不可约分支在 $X$ 中的闭包是 $X$ 的不可约分支（0，\hyperlink{para:0:2.1.6}{2.1.6}），
可用 $V_i$ 代替 $X$，用以 $\overline{f(V_i)}\cap U$ 为底空间的 $U$ 的约化闭子预概形代替 $Y$（\hyperlink{para:1:5.2.1}{5.2.1}）。
于是归结到 $X=\Spec(A)$、$Y=\Spec(B)$ 均仿射且约化的情形。
$f$ 支配，所以 $B$ 为 $A$ 的子环（\hyperlink{para:1:1.2.7}{1.2.7}）；
再由 $B$ 的每个极小素理想都是 $B$ 与 $A$ 的某个极小素理想的交（0，\hyperlink{para:0:1.5.8}{1.5.8}），得到命题。

\egastatement{6.6.6}{命题}
\begin{enumerate}[label=(\roman*)]
\item 每个局部浸入都局部有限型。
\item 两个局部有限型态射的复合局部有限型。
\item 局部有限型 $S$-态射沿任意 $S'\to S$ 作基扩张后仍局部有限型。
\item 两个局部有限型 $S$-态射的积局部有限型。
\item 若复合态射 $g\circ f$ 局部有限型，则 $f$ 局部有限型。
\item 若 $f$ 局部有限型，则 $f_{\mathrm{red}}$ 局部有限型。
\end{enumerate}

由（\hyperlink{para:1:5.5.12}{5.5.12}），只须证明 (i)、(ii)、(iii)。
若 $j:X\to Y$ 为局部浸入，则对每个 $x\in X$，
存在 $j(x)$ 的开邻域 $V$ 和 $x$ 的开邻域 $U$，
使 $j|U:U\to V$ 为闭浸入（\hyperlink{para:1:4.5.1}{4.5.1}），所以为有限型态射。

为证明 (ii)，设 $f:X\to Y$、$g:Y\to Z$ 局部有限型，取 $x\in X$。
由假设存在 $g(f(x))$ 的开邻域 $W$ 和 $f(x)$ 的开邻域 $V$，
使 $g(V)\subset W$ 且 $V$ 在 $W$ 上有限型。
$f^{-1}(V)$ 在 $V$ 上局部有限型（\hyperlink{para:1:6.6.2}{6.6.2}），
因而可取 $x$ 的开邻域 $U\subset f^{-1}(V)$，使它在 $V$ 上有限型。
于是 $g(f(U))\subset W$，且 $U$ 在 $W$ 上有限型（\hyperlink{para:1:6.3.4}{6.3.4}，(ii)）。

最后证明 (iii)，可限于 $Y=S$（\hyperlink{para:1:3.3.11}{3.3.11}）。
对每个 $x'\in X'=X_{(S')}$，令 $x$ 为其在 $X$ 中的像，$s$ 为 $x$ 在 $S$ 中的像。
取 $s$ 的开邻域 $T$，以 $T'$ 表示它在 $S'$ 中的逆像；
取 $x$ 的开邻域 $U$，使它的像包含于 $T$ 且在 $T$ 上有限型。
则 $U\times_S T'=U\times_T T'$ 是 $x'$ 的开邻域（\hyperlink{para:1:3.2.7}{3.2.7}），
且在 $T'$ 上有限型（\hyperlink{para:1:6.3.4}{6.3.4}，(iv)）。

\egastatement{6.6.7}{推论}
设 $X,Y$ 为两个在 $S$ 上局部有限型的预概形。
若 $S$ 局部诺特，则 $X\times_S Y$ 局部诺特。

$X$ 在 $S$ 上局部有限型，故局部诺特。
$X\times_S Y$ 在 $X$ 上局部有限型，所以也局部诺特。

\egastatement{6.6.8}{注}
命题（\hyperlink{para:1:6.3.10}{6.3.10}）及其证明立即推广到仅假设 $f$ 局部有限型的情形。
同样，命题（\hyperlink{para:1:6.4.2}{6.4.2}）、（\hyperlink{para:1:6.4.9}{6.4.9}）在其中的 $X,Y$ 仅假设为域 $K$ 上局部有限型预概形时仍成立。


% END chapters/ch1_s6_6.tex

\par

% BEGIN chapters/ch1_s7_1.tex
\section{有理映射}\hypertarget{heading:1:7}{}
\subsection{有理映射与有理函数}\hypertarget{heading:1:7.1}{}
\src{155—158}
\egapar{7.1.1} 设 $X,Y$ 为两个预概形，$U,V$ 为 $X$ 中的两个稠密开集，$f$（分别为 $g$）为 $U$（分别为 $V$）到 $Y$ 的态射。若 $f,g$ 在 $U\cap V$ 的一个稠密开集中重合，就称它们\emph{等价}。由于 $X$ 中有限个稠密开集之交仍是 $X$ 中的稠密开集，这显然是一个等价关系。
\egastatement{7.1.2}{定义} 给定两个预概形 $X,Y$，称从 $X$ 到 $Y$ 的\emph{有理映射}为：从 $X$ 的稠密开子集到 $Y$ 的态射，关于（\hyperlink{para:1:7.1.1}{7.1.1}）所定义关系的一个等价类。若 $X,Y$ 为 $S$-预概形，且该类中存在一个 $S$-态射，则称该类为一个 $S$-有理映射。称从 $S$ 到 $X$ 的任一 $S$-有理映射为 $X$ 的一个\emph{$S$-有理截面}。称 $X$-预概形 $X\otimes_{\mathbf Z}\mathbf Z[T]$（$T$ 为不定元）的任一 $X$-有理截面为预概形 $X$ 上的一个\emph{有理函数}。

当只讨论 $S$-预概形且不致混淆时，将滥用术语，以“有理映射”代替“$S$-有理映射”。

设 $f$ 为从 $X$ 到 $Y$ 的有理映射，$U$ 为 $X$ 的开集。若 $f_1,f_2$ 为属于类 $f$ 的态射，分别定义于 $X$ 的稠密开集 $V,W$，则限制 $f_1|(U\cap V)$ 与 $f_2|(U\cap W)$ 在 $U\cap V\cap W$ 的一个稠密开集中重合；该开集也在 $U$ 中稠密。\transnote{原文此处直接写“在 $U\cap V\cap W$ 中重合”。依（\hyperlink{para:1:7.1.1}{7.1.1}），一般只能断言在其中的一个稠密开集中重合；这里明确写出这一点，不预先使用（\hyperlink{para:1:7.2.2}{7.2.2}）所需的约化性和分离性假设。}因此类 $f$ 定义了从 $U$ 到 $Y$ 的一个有理映射，称为 $f$ 在 $U$ 上的\emph{限制}，记作 $f|U$。

将每个 $S$-态射 $f:X\to Y$ 映到它所属的 $S$-有理映射，就定义了从 $\Hom_S(X,Y)$ 到从 $X$ 到 $Y$ 的 $S$-有理映射集合的典范映射。将 $X$ 的 $Y$-有理截面集合记作 $\Gamma_{\mathrm{rat}}(X/Y)$，于是有典范映射
\[
\Gamma(X/Y)\longrightarrow\Gamma_{\mathrm{rat}}(X/Y).
\]
另外，若 $X,Y$ 为两个 $S$-预概形，从 $X$ 到 $Y$ 的 $S$-有理映射集合显然典范地等同于 $\Gamma_{\mathrm{rat}}((X\times_S Y)/X)$（\hyperlink{para:1:3.3.14}{3.3.14}）。
\egapar{7.1.3} 由（\hyperlink{para:1:7.1.2}{7.1.2}）与（\hyperlink{para:1:3.3.14}{3.3.14}）立即得知，$X$ 上的有理函数典范地等同于结构层 $\OO_X$ 在 $X$ 的稠密开集上的截面的等价类；两个这样的截面等价，是指它们在包含于各自定义集交集中的一个稠密开集上重合。特别地，$X$ 上的有理函数组成一个环 $R(X)$。
\egapar{7.1.4} 当 $X$ 为不可约预概形时，每个非空开集都在 $X$ 中稠密；也可以说，$X$ 的非空开集正是其泛点 $x$ 的开邻域。两个从 $X$ 的非空开子集到 $Y$ 的态射等价，此时就意味着它们在 $x$ 处有相同的芽。换言之，从 $X$ 到 $Y$ 的有理映射（分别为 $S$-有理映射），等同于从 $X$ 的非空开子集到 $Y$ 的态射（分别为 $S$-态射）在 $X$ 的泛点 $x$ 处的芽。特别地：
\egastatement{7.1.5}{命题} 若 $X$ 为不可约预概形，则 $X$ 上的有理函数环 $R(X)$ 典范地等同于 $X$ 的泛点 $x$ 处的局部环 $\OO_x$。这是一个零维局部环，因而在 $X$ 为诺特预概形时是阿廷局部环；在 $X$ 为整预概形时是域，而当 $X$ 还为仿射概形时，它等同于 $A(X)$ 的分式域。

由上述讨论，以及有理函数与 $\OO_X$ 在稠密开集上的截面的等同，第一项断言就是层在一点处的茎的定义。对其余断言，可归结到 $X$ 为环 $A$ 的仿射概形的情形。此时 $\jj_x$ 为 $A$ 的幂零根，因此 $\OO_x$ 为零维；若 $A$ 为整环，则 $\jj_x=(0)$，从而 $\OO_x$ 为 $A$ 的分式域。最后，若 $A$ 为诺特环，则已知 $\jj_x$ 是幂零的（\cite{b11}，第 127 页，推论 4），且 $\OO_x=A_{\jj_x}$ 为阿廷环。

若 $X$ 为整预概形，则对任意 $z\in X$，环 $\OO_z$ 都是整环。任一包含 $z$ 的仿射开集 $U$ 也包含 $x$，而 $R(U)$ 等于 $A(U)$ 的分式域，典范地等同于 $R(X)$；因此 $R(X)$ 也等同于 $\OO_z$ 的分式域。将 $\OO_z$ 典范地等同于 $R(X)$ 的子环，是将每个截面芽 $s\in\OO_z$ 映到唯一的有理函数：该函数为一个在 $z$ 处具有芽 $s$ 的 $\OO_X$-截面的类；此截面必定义于一个稠密开集上。
\egapar{7.1.6} 现在设 $X$ 只有有限个不可约分支 $X_i$（$1\leq i\leq n$）；当 $X$ 的底空间为诺特空间时即属此种情形。令
\[
X_i'=X_i\setminus\bigcup_{j\ne i}(X_i\cap X_j).
\]
这是 $X$ 的一个开集，是不可约的，其泛点 $x_i$ 也是 $X_i$ 的泛点；各 $X_i'$ 两两不交，其并在 $X$ 中稠密（0，\hyperlink{para:0:2.1.6}{2.1.6}）。对 $X$ 的任一稠密开集 $U$，$U_i=U\cap X_i'$ 为 $X_i'$ 的非空稠密开集，各 $U_i$ 两两不交，故 $U'=\bigcup_iU_i$ 在 $X$ 中稠密。给定从 $U'$ 到 $Y$ 的态射，等价于任意给定各 $U_i$ 到 $Y$ 的态射。因此：
\egastatement{7.1.7}{命题} 设 $X,Y$ 为两个预概形（分别为 $S$-预概形），且 $X$ 有有限个不可约分支 $X_i$，泛点为 $x_i$（$1\leq i\leq n$）。若 $R_i$ 为从 $X$ 的开子集到 $Y$ 的态射（分别为 $S$-态射）在 $x_i$ 处的芽的集合，则从 $X$ 到 $Y$ 的有理映射（分别为 $S$-有理映射）集合等同于 $\prod_{i=1}^nR_i$。
\egastatement{7.1.8}{推论} 设 $X$ 为诺特预概形。$X$ 上的有理函数环为阿廷环，其局部分量为 $X$ 的不可约分支的泛点 $x_i$ 处的局部环 $\OO_{x_i}$。
\egastatement{7.1.9}{推论} 设 $A$ 为诺特环，$X=\Spec(A)$。若 $Q$ 为 $A$ 的所有极小素理想之并的补集，则 $X$ 上的有理函数环典范地等同于分式环 $Q^{-1}A$。

这将由下述引理得出。
\egastatement{7.1.9.1}{引理} 元素 $f\in A$ 满足 $D(f)$ 在 $X$ 中稠密，当且仅当 $f\in Q$；$X$ 的每个稠密开集都包含一个形如 $D(f)$ 的开集，其中 $f\in Q$。

先假定此引理已经证明。由于截面环 $\Gamma(D(f),\OO_X)$ 等同于 $A_f$（\hyperlink{para:1:1.3.6}{1.3.6} 与 \hyperlink{para:1:1.3.7}{1.3.7}），且 $f\in Q$ 时的 $D(f)$ 在按 $\supset$ 排序的稠密开集集合中构成一个共尾子集，故由定义（\hyperlink{para:1:7.1.1}{7.1.1}），$X$ 的有理函数环等同于各 $A_f$（$f\in Q$）的归纳极限，其中采用“$g$ 是 $f$ 的倍数”这一预序关系；此极限即 $Q^{-1}A$（0，\hyperlink{para:0:1.4.5}{1.4.5}）。

为证明（\hyperlink{para:1:7.1.9.1}{7.1.9.1}），仍将 $X$ 的不可约分支记作 $X_i$（$1\leq i\leq n$）。$D(f)$ 在 $X$ 中稠密，当且仅当 $D(f)\cap X_i\ne\varnothing$ 对一切 $i$ 成立。令 $\mathfrak p_i=\jj(X_i)$，上述条件就是 $f\notin\mathfrak p_i$ 对所有 $i$ 成立；而各 $\mathfrak p_i$ 正是 $A$ 的极小素理想（\hyperlink{para:1:1.1.14}{1.1.14}），故这些条件等价于 $f\in Q$，得到引理的第一项断言。另一方面，若 $U$ 为 $X$ 的稠密开集，其补集形如 $V(\mathfrak a)$，其中理想 $\mathfrak a$ 不包含于任何一个 $\mathfrak p_i$；因而它也不包含于这些理想之并（\cite{b10}，第 13 页）。故可取 $f\in\mathfrak a\cap Q$，于是 $D(f)\subset U$。证明完毕。
\egapar{7.1.10} 再设 $X$ 不可约，泛点为 $x$。每个非空开集 $U$ 都包含 $x$，因而也包含所有满足 $x\in\overline{\{z\}}$ 的 $z\in X$，所以可将任一态射 $U\to Y$ 与典范态射 $\Spec(\OO_x)\to X$（\hyperlink{para:1:2.4.1}{2.4.1}）复合。两个分别从 $X$ 的非空开子集到 $Y$ 的态射，若在一个非空开子集中重合，经此复合即给出同一态射 $\Spec(\OO_x)\to Y$。换言之，每个从 $X$ 到 $Y$ 的有理映射因此对应一个确定的态射 $\Spec(\OO_x)\to Y$。
\egastatement{7.1.11}{命题} 设 $X,Y$ 为两个 $S$-预概形，$X$ 不可约，泛点为 $x$，且 $Y$ 在 $S$ 上为有限型。若两个从 $X$ 到 $Y$ 的 $S$-有理映射对应同一 $S$-态射 $\Spec(\OO_x)\to Y$，则这两个有理映射相同。若还假设 $S$ 局部诺特，则每个从 $\Spec(\OO_x)$ 到 $Y$ 的 $S$-态射都对应唯一一个从 $X$ 到 $Y$ 的 $S$-有理映射。

考虑到 $X$ 的每个非空开集都是稠密的，这立即由（\hyperlink{para:1:6.5.1}{6.5.1}）得出。
\egastatement{7.1.12}{推论} 假设 $S$ 局部诺特，并满足（\hyperlink{para:1:7.1.11}{7.1.11}）的其余假设。那么，从 $X$ 到 $Y$ 的 $S$-有理映射，等同于 $S$-预概形 $Y$ 的取值于 $S$-预概形 $\Spec(\OO_x)$ 的点。

这就是用（\hyperlink{para:1:3.4.1}{3.4.1}）的术语重述（\hyperlink{para:1:7.1.11}{7.1.11}）。
\egastatement{7.1.13}{推论} 假设满足（\hyperlink{para:1:7.1.12}{7.1.12}）的条件。令 $s$ 为 $x$ 在 $S$ 中的像。给定从 $X$ 到 $Y$ 的一个 $S$-有理映射，等价于给定 $Y$ 中位于 $s$ 之上的一点 $y$，以及一个局部 $\OO_s$-同态 $\OO_y\to\OO_x=R(X)$。

这由（\hyperlink{para:1:7.1.11}{7.1.11}）与（\hyperlink{para:1:2.4.4}{2.4.4}）得出。特别地：
\egastatement{7.1.14}{推论} 在（\hyperlink{para:1:7.1.12}{7.1.12}）的条件下，对给定的 $Y$，从 $X$ 到 $Y$ 的 $S$-有理映射只依赖于 $S$-预概形 $\Spec(\OO_x)$；特别地，对任一 $z\in X$，将 $X$ 换成 $\Spec(\OO_z)$ 后，它们仍保持不变。

事实上，因 $z\in\overline{\{x\}}$，$x$ 是 $Z=\Spec(\OO_z)$ 的泛点，且 $\OO_{X,x}=\OO_{Z,x}$。

当 $X$ 为整预概形时，$R(X)=\OO_x=k(x)$ 是域（\hyperlink{para:1:7.1.5}{7.1.5}），以上推论于是给出：
\egastatement{7.1.15}{推论} 假设满足（\hyperlink{para:1:7.1.12}{7.1.12}）的条件，且 $X$ 为整预概形。令 $s$ 为 $x$ 在 $S$ 中的像。那么，从 $X$ 到 $Y$ 的 $S$-有理映射等同于 $Y\otimes_S k(s)$ 的取值于扩域 $R(X)/k(s)$ 的几何点。换言之，每个这样的有理映射等价于给定 $Y$ 中位于 $s$ 之上的一点 $y$，以及一个 $k(s)$-单同态 $k(y)\to k(x)=R(X)$。

事实上，$Y$ 中位于 $s$ 之上的点等同于 $Y\otimes_S k(s)$ 的点（\hyperlink{para:1:3.6.3}{3.6.3}），而局部 $\OO_s$-同态 $\OO_y\to R(X)$ 等同于 $k(s)$-单同态 $k(y)\to R(X)$。更特别地：
\egastatement{7.1.16}{推论} 设 $k$ 为域，$X,Y$ 为两个 $k$ 上的代数预概形（\hyperlink{para:1:6.4.1}{6.4.1}），且 $X$ 为整预概形。那么，从 $X$ 到 $Y$ 的 $k$-有理映射等同于 $Y$ 的取值于扩域 $R(X)/k$ 的几何点（\hyperlink{para:1:3.4.4}{3.4.4}）。


% END chapters/ch1_s7_1.tex

\par

% BEGIN chapters/ch1_s7_2.tex
\subsection{有理映射的定义域}\hypertarget{heading:1:7.2}{}
\src{158—161}
\egapar{7.2.1} 设 $X,Y$ 为两个预概形，$f$ 为从 $X$ 到 $Y$ 的有理映射。若存在一个包含 $x$ 的稠密开集 $U$，以及属于等价类 $f$ 的态射 $U\to Y$，就称 $f$ \emph{在点 $x\in X$ 处有定义}。所有使 $f$ 有定义的点组成的集合称为 $f$ 的\emph{定义域}；这显然是 $X$ 中的一个稠密开集。
\egastatement{7.2.2}{命题} 设 $X,Y$ 为两个 $S$-预概形，其中 $X$ 约化，$Y$ 在 $S$ 上分离。设 $f$ 为从 $X$ 到 $Y$ 的一个 $S$-有理映射，$U_0$ 为其定义域。那么，存在唯一一个属于类 $f$ 的 $S$-态射 $U_0\to Y$。

由于对属于类 $f$ 的任一态射 $U\to Y$ 都必有 $U\subset U_0$，本命题显然由以下引理得出。
\egastatement{7.2.2.1}{引理} 在（\hyperlink{para:1:7.2.2}{7.2.2}）的假设下，设 $U_1,U_2$ 为 $X$ 中的两个稠密开集，$f_i:U_i\to Y$（$i=1,2$）为两个 $S$-态射。若存在 $V\subset U_1\cap U_2$，它在 $X$ 中开且稠密，并且 $f_1,f_2$ 在 $V$ 中重合，则它们在整个 $U_1\cap U_2$ 中重合。

显然可归结到 $X=U_1=U_2$ 的情形。由于 $X$（因而 $V$）约化，$X$ 是 $X$ 的不小于 $V$ 的最小闭子预概形（\hyperlink{para:1:5.2.2}{5.2.2}）。令 $g=(f_1,f_2)_S:X\to Y\times_S Y$。按假设，对角线 $T=\Delta_Y(Y)$ 为 $Y\times_S Y$ 的闭子预概形，故 $Z=g^{-1}(T)$ 为 $X$ 的闭子预概形（\hyperlink{para:1:4.4.1}{4.4.1}）。若 $h:V\to Y$ 为 $f_1,f_2$ 在 $V$ 上的共同限制，则 $g$ 在 $V$ 上的限制为 $g'=(h,h)_S=\Delta_Y\circ h$。由 $\Delta_Y^{-1}(T)=Y$，得 $g'^{-1}(T)=V$，故 $Z$ 在 $V$ 上诱导 $V$，因而不小于 $V$，于是 $Z=X$。由 $g^{-1}(T)=X$，可知 $g$ 分解为 $\Delta_Y\circ f$，其中 $f:X\to Y$ 为态射（\hyperlink{para:1:4.4.1}{4.4.1}）；按对角态射的定义，$f_1=f_2=f$。

（\hyperlink{para:1:7.2.2}{7.2.2}）所定义的态射 $U_0\to Y$，显然是类 $f$ 中唯一不能延拓到 $X$ 的严格包含 $U_0$ 的开子集上的态射。因此，在（\hyperlink{para:1:7.2.2}{7.2.2}）的假设下，可将从 $X$ 到 $Y$ 的有理映射，等同于从 $X$ 的稠密开子集到 $Y$ 的不能继续延拓至更大开集的态射。在此等同下，（\hyperlink{para:1:7.2.2}{7.2.2}）给出：
\egastatement{7.2.3}{推论} 对 $X,Y$ 作（\hyperlink{para:1:7.2.2}{7.2.2}）的假设，并设 $U$ 为 $X$ 的一个稠密开集。则从 $U$ 到 $Y$ 的 $S$-态射，与从 $X$ 到 $Y$ 的在 $U$ 的每一点都有定义的 $S$-有理映射之间，存在典范双射。

事实上，由（\hyperlink{para:1:7.2.2}{7.2.2}），每个 $S$-态射 $f:U\to Y$ 都有唯一一个延拓它的 $S$-有理映射 $\bar f:X\dashrightarrow Y$。
\egastatement{7.2.4}{推论} 设 $S$ 为概形，$X$ 为约化 $S$-预概形，$Y$ 为 $S$-概形，$f:U\to Y$ 为从 $X$ 的稠密开集 $U$ 到 $Y$ 的 $S$-态射。若 $\bar f$ 为延拓 $f$ 的从 $X$ 到 $Y$ 的 $\mathbf Z$-有理映射，则 $\bar f$ 是一个 $S$-态射（因此也是延拓 $f$ 的 $S$-有理映射）。

事实上，令 $\varphi:X\to S$、$\psi:Y\to S$ 为结构态射，$U_0$ 为 $\bar f$ 的定义域，$j:U_0\to X$ 为嵌入。只须证明 $\psi\circ\bar f=\varphi\circ j$；由于 $f$ 为 $S$-态射，这立即由（\hyperlink{para:1:7.2.2.1}{7.2.2.1}）得出。
\egastatement{7.2.5}{推论} 设 $X,Y$ 为两个 $S$-预概形，$X$ 约化，且 $X,Y$ 均在 $S$ 上分离。设 $p:Y\to X$ 为 $S$-态射，使 $Y$ 成为 $X$-预概形；设 $U$ 为 $X$ 的稠密开集，$f$ 为 $Y$ 的一个 $U$-截面。那么，延拓 $f$ 的从 $X$ 到 $Y$ 的有理映射 $\bar f$ 是 $Y$ 的一个 $X$-有理截面。

须证明 $p\circ\bar f$ 在 $\bar f$ 的定义域上为恒等态射；因 $X$ 在 $S$ 上分离，这仍由（\hyperlink{para:1:7.2.2.1}{7.2.2.1}）得出。
\egastatement{7.2.6}{推论} 设 $X$ 为约化预概形，$U$ 为 $X$ 的一个稠密开集。$\OO_X$ 在 $U$ 上的截面，与 $X$ 上在 $U$ 的每一点都有定义的有理函数 $f$ 之间，存在典范双射。

由（\hyperlink{para:1:7.2.3}{7.2.3}）、（\hyperlink{para:1:7.1.2}{7.1.2}）与（\hyperlink{para:1:7.1.3}{7.1.3}），只须注意到 $X$-预概形 $X\otimes_{\mathbf Z}\mathbf Z[T]$ 在 $X$ 上分离（\hyperlink{para:1:5.5.1}{5.5.1}，(iv)）。
\egastatement{7.2.7}{推论} 设 $Y$ 为约化预概形，$f:X\to Y$ 为分离态射，$U$ 为 $Y$ 的稠密开集，$g:U\to f^{-1}(U)$ 为 $f^{-1}(U)$ 的一个 $U$-截面。令 $Z$ 为 $X$ 的以 $\overline{g(U)}$ 为底空间的约化子预概形（\hyperlink{para:1:5.2.1}{5.2.1}）。则 $g$ 是 $X$ 的一个 $Y$-截面的限制，当且仅当 $f$ 在 $Z$ 上的限制为 $Z$ 到 $Y$ 的同构。换言之，延拓 $g$ 的从 $Y$ 到 $X$ 的有理映射处处有定义，当且仅当上述条件成立（\hyperlink{para:1:7.2.5}{7.2.5}）。

$f$ 在 $f^{-1}(U)$ 上的限制是分离态射（\hyperlink{para:1:5.5.1}{5.5.1}，(i)），故 $g$ 为闭浸入（\hyperlink{para:1:5.4.6}{5.4.6}）。于是 $g(U)=Z\cap f^{-1}(U)$，且 $Z$ 在其开集 $g(U)$ 上诱导的子预概形，与 $f^{-1}(U)$ 中和 $g$ 相伴的闭子预概形相同（\hyperlink{para:1:5.2.1}{5.2.1}）。陈述中的条件显然充分：若成立，令 $f_Z:Z\to Y$ 为 $f$ 的限制，$\bar g:Y\to Z$ 为其逆同构，则 $\bar g$ 延拓 $g$。反之，若 $g$ 为 $X$ 的一个 $Y$-截面 $h$ 在 $U$ 上的限制，则 $h$ 为闭浸入（\hyperlink{para:1:5.4.6}{5.4.6}），故 $h(Y)$ 为闭集；由于它包含于 $Z$，它就等于 $Z$。由（\hyperlink{para:1:5.2.1}{5.2.1}），$h$ 必为 $Y$ 到 $X$ 的闭子预概形 $Z$ 的同构。
\egapar{7.2.8} 设 $X,Y$ 为两个 $S$-预概形，其中 $X$ 约化，$Y$ 在 $S$ 上分离。设 $f$ 为从 $X$ 到 $Y$ 的 $S$-有理映射，$x\in X$。只要 $f$ 的定义域在 $\Spec(\OO_x)$ 上的迹稠密于 $\Spec(\OO_x)$，就可将 $f$ 与典范 $S$-态射 $\Spec(\OO_x)\to X$（\hyperlink{para:1:2.4.1}{2.4.1}）复合；这里将 $\Spec(\OO_x)$ 等同于所有满足 $x\in\overline{\{z\}}$ 的 $z\in X$ 的集合（\hyperlink{para:1:2.4.2}{2.4.2}）。以下情形满足此条件。

$1^\circ$ $X$ 不可约，因而是整预概形。此时 $X$ 的泛点 $\xi$ 也是 $\Spec(\OO_x)$ 的泛点；由于 $f$ 的定义域 $U$ 包含 $\xi$，$U\cap\Spec(\OO_x)$ 也包含 $\xi$，故在 $\Spec(\OO_x)$ 中稠密。

$2^\circ$ $X$ 局部诺特。此时所需断言由以下引理得出。
\egastatement{7.2.8.1}{引理} 设 $X$ 为底空间局部诺特的预概形，$x\in X$。$\Spec(\OO_x)$ 的不可约分支，正是 $X$ 的包含 $x$ 的不可约分支在 $\Spec(\OO_x)$ 上的迹。开集 $U\subset X$ 满足 $U\cap\Spec(\OO_x)$ 稠密于 $\Spec(\OO_x)$，当且仅当它与 $X$ 的每个包含 $x$ 的不可约分支相交；特别地，当 $U$ 在 $X$ 中稠密时，此条件成立。

第二项断言显然由第一项得出，故只须证明后者。$\Spec(\OO_x)$ 包含于每个包含 $x$ 的仿射开集 $U$，而 $U$ 的包含 $x$ 的不可约分支，是 $X$ 的包含 $x$ 的不可约分支在 $U$ 上的迹（0，\hyperlink{para:0:2.1.6}{2.1.6}）。所以可假设 $X$ 为环 $A$ 的仿射概形。$A_x$ 的素理想与 $A$ 中包含于 $\jj_x$ 的素理想一一对应（0，\hyperlink{para:0:1.2.6}{1.2.6}）；因此 $A_x$ 的极小素理想与 $A$ 中包含于 $\jj_x$ 的极小素理想一一对应，从而得到引理（\hyperlink{para:1:1.1.14}{1.1.14}）。

现假设处于以上两种情形之一。若 $U$ 为 $S$-有理映射 $f$ 的定义域，记 $f'$ 为从 $\Spec(\OO_x)$ 到 $Y$ 的、在 $U\cap\Spec(\OO_x)$ 中与 $f$ 重合的有理映射（利用（\hyperlink{para:1:2.4.2}{2.4.2}）的等同）。称 $f'$ 为由 $f$ \emph{诱导}的有理映射。
\egastatement{7.2.9}{命题} 设 $S$ 为局部诺特预概形，$X$ 为约化 $S$-预概形，$Y$ 为有限型 $S$-概形。再假设 $X$ 不可约或局部诺特。设 $f$ 为从 $X$ 到 $Y$ 的 $S$-有理映射，$x\in X$。则 $f$ 在 $x$ 处有定义，当且仅当由 $f$ 诱导的有理映射 $f':\Spec(\OO_x)\dashrightarrow Y$（\hyperlink{para:1:7.2.8}{7.2.8}）是一个态射。

必要性显然，因为 $\Spec(\OO_x)$ 包含于每个包含 $x$ 的开集。证明充分性。由（\hyperlink{para:1:6.5.1}{6.5.1}），存在 $x$ 在 $X$ 中的一个开邻域 $U$，以及一个 $S$-态射 $g:U\to Y$，在 $\Spec(\OO_x)$ 上诱导 $f'$。若 $X$ 不可约，$U$ 在 $X$ 中稠密，故由（\hyperlink{para:1:7.2.3}{7.2.3}），可将 $g$ 看作 $S$-有理映射。此外，$X$ 的泛点属于 $\Spec(\OO_x)$ 及 $f$ 的定义域，所以 $f,g$ 在此点处重合，因而在 $X$ 的一个非空开集中重合（\hyperlink{para:1:6.5.1}{6.5.1}）。由于 $f,g$ 为 $S$-有理映射，它们相同（\hyperlink{para:1:7.2.3}{7.2.3}），故 $f$ 在 $x$ 处有定义。

若 $X$ 局部诺特，可假设 $U$ 诺特。此时 $X$ 的包含 $x$ 的不可约分支 $X_i$ 只有有限个（\hyperlink{para:1:7.2.8.1}{7.2.8.1}），且必要时缩小 $U$，可使它们成为仅有的与 $U$ 相交的不可约分支；这是因为与诺特开集 $U$ 相交的不可约分支只有有限个。如上可知，$f,g$ 在每个 $X_i$ 的一个非空开集中重合。考虑到各 $X_i$ 包含于 $\overline U$，取一个定义于 $U\cup(X\setminus\overline U)$ 的稠密开子集上的态射 $f_1$：它在 $U$ 中等于 $g$，在 $X\setminus\overline U$ 与 $f$ 的定义域之交上等于 $f$。由于 $U\cup(X\setminus\overline U)$ 在 $X$ 中稠密，$f_1,f$ 在 $X$ 的一个稠密开集中重合；又因 $f$ 是有理映射，$f$ 就是 $f_1$ 的延拓（\hyperlink{para:1:7.2.3}{7.2.3}），故在 $x$ 处有定义。


% END chapters/ch1_s7_2.tex

\par

% BEGIN chapters/ch1_s7_3.tex
\subsection{有理函数层}\hypertarget{heading:1:7.3}{}
\src{161—163}
\egapar{7.3.1} 设 $X$ 为预概形。对每个开集 $U\subset X$，记 $R(U)$ 为 $U$ 上的有理函数环（\hyperlink{para:1:7.1.3}{7.1.3}）；这是一个 $\Gamma(U,\OO_X)$-代数。若 $V\subset U$ 为另一开集，则 $\OO_X$ 在 $U$ 的稠密开子集上的任一截面，限制到 $V$ 后，给出 $V$ 的一个稠密开子集上的截面；若两个截面在 $U$ 的一个稠密开子集上重合，它们在 $V$ 上的限制也在 $V$ 的一个稠密开子集上重合。于是定义了代数的双同态 $\rho_{V,U}:R(U)\to R(V)$。显然，对三个开集 $U\supset V\supset W$，有 $\rho_{W,U}=\rho_{W,V}\circ\rho_{V,U}$；因此这些 $R(U)$ 定义了 $X$ 上的一个代数预层。
\egastatement{7.3.2}{定义} 称由这些 $R(U)$ 构成的预层所伴随的 $\OO_X$-代数层为预概形 $X$ 上的\emph{有理函数层}，记作 $\mathscr R(X)$。

对任意预概形 $X$ 及任意开集 $U\subset X$，显然诱导层 $\mathscr R(X)|U$ 就是 $\mathscr R(U)$。
\egastatement{7.3.3}{命题} 设预概形 $X$ 的不可约分支族 $(X_\lambda)$ 局部有限；特别地，当 $X$ 的底空间局部诺特时即如此。那么，$\OO_X$-模 $\mathscr R(X)$ 为拟凝聚层。对 $X$ 的每个只与有限个分支 $X_\lambda$ 相交的开集 $U$，$R(U)=\Gamma(U,\mathscr R(X))$，并典范地等同于那些满足 $U\cap X_\lambda\ne\varnothing$ 的 $X_\lambda$ 的泛点处的局部环的直积。

显然可限于 $X$ 只有有限个不可约分支 $X_i$，其泛点为 $x_i$（$1\leq i\leq n$）的情形。由（\hyperlink{para:1:7.1.7}{7.1.7}），$R(U)$ 典范地等同于满足 $U\cap X_i\ne\varnothing$ 的各 $\OO_{x_i}$ 的直积；原文将这些局部环记为 $R(X_i)$。\transnote{这里的 $R(X_i)$ 旨在表示对应泛点处的有理函数环 $\OO_{X,x_i}$。在尚未给不可约分支指定子预概形结构时，可用（\hyperlink{para:1:7.1.6}{7.1.6}）的不可约开集 $X_i'$ 理解这一记号；不能在非约化情形中擅自换成约化闭分支的函数域。}
再证明预层 $U\mapsto R(U)$ 满足层的公理，于是就有 $R(U)=\Gamma(U,\mathscr R(X))$。由前述表示，它满足（F1）。为验证（F2），设 $(V_\alpha)$ 为开集 $U$ 的开覆盖，$V_\alpha\subset U$，且 $s_\alpha\in R(V_\alpha)$ 对每对指标都在 $V_\alpha\cap V_\beta$ 上有相同限制。于是，对满足 $U\cap X_i\ne\varnothing$ 的每个 $i$，所有满足 $V_\alpha\cap X_i\ne\varnothing$ 的 $s_\alpha$ 在 $R(X_i)$ 中的分量都相同。将此分量记作 $t_i$，则以各 $t_i$ 为分量的 $R(U)$ 中的元素，显然在每个 $V_\alpha$ 上限制为 $s_\alpha$。

最后，为证明 $\mathscr R(X)$ 拟凝聚，可限于 $X=\Spec(A)$ 为仿射概形的情形。取形如 $D(f)$（$f\in A$）的仿射开集作为 $U$，由前面的讨论及定义（\hyperlink{para:1:1.3.4}{1.3.4}）得 $\mathscr R(X)=\wt M$，其中 $M$ 为各 $A$-模 $A_{x_i}$ 的直和。
\egastatement{7.3.4}{推论} 设 $X$ 为只有有限个不可约分支的约化预概形，并令 $X_i$（$1\leq i\leq n$）为以这些不可约分支为底空间的约化闭子预概形（\hyperlink{para:1:5.2.1}{5.2.1}）。若 $h_i:X_i\to X$ 为典范嵌入，则 $\mathscr R(X)$ 为各 $\OO_X$-代数 $(h_i)_*(\mathscr R(X_i))$ 的直积。
\egastatement{7.3.5}{推论} 若 $X$ 不可约，则每个拟凝聚 $\mathscr R(X)$-模 $\F$ 都是常层。

只须证明每个 $x\in X$ 都有一个邻域 $U$，使 $\F|U$ 为常层（0，\hyperlink{para:0:3.6.2}{3.6.2}），故可归结到 $X$ 仿射的情形。还可假设 $\F$ 为某个同态
\[
(\mathscr R(X))^{(I)}\longrightarrow(\mathscr R(X))^{(J)}
\]
的余核（0，\hyperlink{para:0:5.1.3}{5.1.3}），因而只须验证 $\mathscr R(X)$ 为常层。而这显然成立，因为每个非空开集 $U$ 都包含 $X$ 的泛点，故 $\Gamma(U,\mathscr R(X))=R(X)$。
\egastatement{7.3.6}{推论} 若 $X$ 不可约，则对每个拟凝聚 $\OO_X$-模 $\F$，$\F\otimes_{\OO_X}\mathscr R(X)$ 为常层。若 $X$ 还约化，因而是整预概形，则 $\F\otimes_{\OO_X}\mathscr R(X)$ 同构于一个形如 $(\mathscr R(X))^{(I)}$ 的层。

第二项断言来自此时 $R(X)$ 为域这一事实。
\egastatement{7.3.7}{命题} 设预概形 $X$ 局部整或局部诺特。那么 $\mathscr R(X)$ 是拟凝聚 $\OO_X$-代数。若 $X$ 还约化（当 $X$ 局部整时即如此），则典范同态 $\OO_X\to\mathscr R(X)$ 为单射。

问题是局部的，故第一项断言由（\hyperlink{para:1:7.3.3}{7.3.3}）得出；第二项断言立即由（\hyperlink{para:1:7.2.3}{7.2.3}）得出。
\egapar{7.3.8} 设 $X,Y$ 为各有有限个不可约分支的两个预概形，$f:X\to Y$ 为态射，其在 $X$ 的不可约分支的泛点集合上的限制，是到 $Y$ 的不可约分支的泛点集合的满射。原文断言，此时有
\begin{equation}
 f^*(\mathscr R(Y))=\mathscr R(X).
 \tag{7.3.8.1}\label{eq:1:7.3.8.1}
\end{equation}
\transnote{此式及下述论证按原文保留，但在所列假设下，该式一般不成立。譬如 $X=\Spec(k[t])$、$Y=\Spec(k)$，结构态射满足所述泛点条件，却有 $f^*\mathscr R(Y)=\OO_X$，而 $\mathscr R(X)$ 是函数域 $k(t)$ 所给出的常层。仅在泛点处计算茎，不能推出全层同构；（\hyperlink{para:1:7.3.5}{7.3.5}）也不能直接用于这里的任意 $\OO_X$-模。}
事实上，由（\hyperlink{para:1:7.3.3}{7.3.3}），原文将问题归结到 $X,Y$ 不可约、泛点分别为 $x,y$ 且 $f(x)=y$ 的情形。于是
\[
 (f^*(\mathscr R(Y)))_x
 =\OO_y\otimes_{\OO_y}\OO_x=\OO_x
 \qquad\text{（0，\hyperlink{para:0:4.3.1}{4.3.1}）},
\]
并据（\hyperlink{para:1:7.3.5}{7.3.5}）得出（\hyperref[eq:1:7.3.8.1]{7.3.8.1}）。


% END chapters/ch1_s7_3.tex

\par

% BEGIN chapters/ch1_s7_4.tex
\subsection{挠层与无挠层}\hypertarget{heading:1:7.4}{}
\src{163—164}
\egapar{7.4.1} 设 $X$ 为整预概形。对每个 $\OO_X$-模 $\F$，典范同态 $\OO_X\to\mathscr R(X)$ 经张量积定义一个同态（仍称典范同态）
\[
 \F\longrightarrow\F\otimes_{\OO_X}\mathscr R(X).
\]
在每个茎上，它就是 $\F_x\to\F_x\otimes_{\OO_x}R(X)$ 的同态 $z\mapsto z\otimes1$。此同态的核 $\mathscr T$ 是 $\F$ 的子 $\OO_X$-模，称为 $\F$ 的\emph{挠层}；若 $\F$ 拟凝聚，则 $\mathscr T$ 也拟凝聚（\hyperlink{para:1:4.1.1}{4.1.1} 与 \hyperlink{para:1:7.3.6}{7.3.6}）。若 $\mathscr T=0$，称 $\F$ \emph{无挠}；若 $\F=\mathscr T$，称 $\F$ 是\emph{挠层}。对任意 $\OO_X$-模 $\F$，$\F/\mathscr T$ 都无挠。由（\hyperlink{para:1:7.3.5}{7.3.5}）得：
\egastatement{7.4.2}{命题} 若 $X$ 为整预概形，则每个无挠拟凝聚 $\OO_X$-模 $\F$，都同构于一个形如 $(\mathscr R(X))^{(I)}$ 的常层的子层 $\G$，且该常层由 $\G$ 作为 $\mathscr R(X)$-模生成。

$I$ 的基数称为 $\F$ 的\emph{秩}。对 $X$ 的每个非空仿射开集 $U$，$\F$ 的秩等于 $\Gamma(U,\F)$ 作为 $\Gamma(U,\OO_X)$-模的秩；考察包含于 $U$ 的 $X$ 的泛点，即立即得到这一点。特别地：
\egastatement{7.4.3}{推论} 在整预概形 $X$ 上，每个无挠、秩为 $1$ 的拟凝聚 $\OO_X$-模，特别是每个可逆 $\OO_X$-模，都同构于 $\mathscr R(X)$ 的一个子 $\OO_X$-模；反之亦然。\transnote{逆向陈述应就非零拟凝聚子 $\OO_X$-模理解；零子模秩为 $0$，而任意子层也未必拟凝聚。原文逆向措辞未重复这些限制。}
\egastatement{7.4.4}{推论} 设 $X$ 为整预概形，$\LL,\LL'$ 为两个无挠 $\OO_X$-模，$f$（分别为 $f'$）为 $\LL$（分别为 $\LL'$）在 $X$ 上的截面。则 $f\otimes f'=0$，当且仅当 $f,f'$ 中至少一个为零。\transnote{对于所列的任意无挠模层，结论不成立。
取 $X=\Spec(k[t])$、$K=k(t)$，在闭点 $(t)$、$(t-1)$ 分别取从该单点直像得到的层
$\LL=i_{0*}K$、$\LL'=i_{1*}K$，其 $\OO_X$-作用由 $\OO_X\to\mathscr R(X)$ 诱导。
它们的非零茎均为 $K$，按（\hyperlink{para:1:7.4.1}{7.4.1}）都无挠；各自的全局截面 $1$ 非零。
但两者的支集不交，故 $\LL\otimes_{\OO_X}\LL'=0$。
这两个层都不是拟凝聚层。若另外假设两层拟凝聚，则可由（\hyperlink{para:1:7.4.2}{7.4.2}）嵌入函数域上的向量空间所给出的常层，
从而证明结论；这一证明不要求泛点处的茎只有一维。
正文保留原版陈述与证明；以上反例和附加条件分析为译者说明。}

令 $x$ 为 $X$ 的泛点。按假设有 $(f\otimes f')_x=f_x\otimes f'_x=0$。原文指出，$\LL_x,\LL'_x$ 等同于域 $\OO_x$ 的子 $\OO_x$-模，因此该关系推出 $f_x=0$ 或 $f'_x=0$；再由 $\LL,\LL'$ 无挠，得 $f=0$ 或 $f'=0$（\hyperlink{para:1:7.3.5}{7.3.5}）。
\egastatement{7.4.5}{命题} 设 $X,Y$ 为两个整预概形，$f:X\to Y$ 为支配态射。对每个无挠拟凝聚 $\OO_X$-模 $\F$，$f_*(\F)$ 都是无挠 $\OO_Y$-模。

由于 $f_*$ 左正合（0，\hyperlink{para:0:4.2.1}{4.2.1}），由（\hyperlink{para:1:7.4.2}{7.4.2}），只须在 $\F=(\mathscr R(X))^{(I)}$ 时证明命题。$Y$ 的每个非空开集 $U$ 都包含 $Y$ 的泛点，故 $f^{-1}(U)$ 包含 $X$ 的泛点（0，\hyperlink{para:0:2.1.5}{2.1.5}）。因此，此时有
\[
 \Gamma(U,f_*(\F))=\Gamma(f^{-1}(U),\F)=(R(X))^{(I)}.
\]
换言之，$f_*(\F)$ 是茎为 $(R(X))^{(I)}$ 的常层，将其视为 $\mathscr R(Y)$-模，显然无挠。
\egastatement{7.4.6}{命题} 设 $X$ 为整预概形，$x$ 为其泛点。对每个有限型拟凝聚 $\OO_X$-模 $\F$，下列条件等价：
\begin{enumerate}[label=\alph*)]
\item $\F$ 为挠层；
\item $\F_x=0$；
\item $\Supp(\F)\ne X$。
\end{enumerate}
由（\hyperlink{para:1:7.3.5}{7.3.5}）与（\hyperlink{para:1:7.4.1}{7.4.1}），$\F_x=0$ 和 $\F\otimes_{\OO_X}\mathscr R(X)=0$ 等价，故 a) 与 b) 等价。另一方面，$\Supp(\F)$ 在 $X$ 中闭（0，\hyperlink{para:0:5.2.2}{5.2.2}），且 $X$ 的每个非空开集都包含 $x$，故 b) 与 c) 等价。
\egapar{7.4.7} 将（\hyperlink{para:1:7.4.1}{7.4.1}）的定义以滥用术语的方式推广到 $X$ 为只有有限个不可约分支的约化预概形的情形。原文据（\hyperlink{para:1:7.3.4}{7.3.4}）指出，（\hyperlink{para:1:7.4.6}{7.4.6}）中 a) 与 c) 的等价性对这样的预概形仍成立。\transnote{这里原文对 c) 的推广不足以成立：例如两个整仿射概形的不交并上，仅在其中一个分支取结构层、在另一个分支取零的拟凝聚层，其支集不等于全空间，但不是挠层。与 a) 对应的条件应要求支集不包含任何不可约分支，亦即所有分支泛点处的茎都为零；不能只保留 $\Supp(\F)\ne X$。}


% END chapters/ch1_s7_4.tex

\par

% BEGIN chapters/ch1_s8_1.tex
\section{Chevalley 概形}\hypertarget{heading:1:8}{}
\subsection{相关的局部环}\hypertarget{heading:1:8.1}{}
\src{164—165}
对每个局部环 $A$，将其极大理想记作 $\mathfrak m(A)$。
\egastatement{8.1.1}{引理} 设 $A,B$ 为两个局部环，$A\subset B$。下列条件等价：
\begin{enumerate}[label=(\roman*)]
\item $\mathfrak m(B)\cap A=\mathfrak m(A)$；
\item $\mathfrak m(A)\subset\mathfrak m(B)$；
\item $1$ 不属于由 $\mathfrak m(A)$ 在 $B$ 中生成的理想。
\end{enumerate}
显然 (i) 推出 (ii)，而 (ii) 推出 (iii)。最后，若 (iii) 成立，则 $\mathfrak m(B)\cap A$ 包含 $\mathfrak m(A)$ 且不包含 $1$，故等于 $\mathfrak m(A)$。

当（\hyperlink{para:1:8.1.1}{8.1.1}）的等价条件成立时，称 $B$ \emph{支配} $A$；这等价于说嵌入 $A\to B$ 是局部同态。显然，在一个环 $R$ 的所有局部子环的集合中，支配关系是一个序关系。
\egapar{8.1.2} 现在考虑一个域 $R$。对 $R$ 的每个子环 $A$，以 $L(A)$ 表示各局部环 $A_{\mathfrak p}$ 组成的集合，其中 $\mathfrak p$ 遍历 $A$ 的素谱；将这些环等同于 $R$ 的包含 $A$ 的子环。由于 $\mathfrak p=(\mathfrak p A_{\mathfrak p})\cap A$，映射 $\mathfrak p\mapsto A_{\mathfrak p}$ 是 $\Spec(A)$ 到 $L(A)$ 的双射。
\egastatement{8.1.3}{引理} 设 $R$ 为域，$A$ 为其子环。$R$ 的一个局部子环 $M$ 支配某个 $A_{\mathfrak p}\in L(A)$，当且仅当 $A\subset M$。此时，被 $M$ 支配的局部环 $A_{\mathfrak p}$ 唯一，并对应于 $\mathfrak p=\mathfrak m(M)\cap A$。

若 $M$ 支配 $A_{\mathfrak p}$，由（\hyperlink{para:1:8.1.1}{8.1.1}），有 $\mathfrak m(M)\cap A_{\mathfrak p}=\mathfrak p A_{\mathfrak p}$，故 $\mathfrak p$ 唯一。另一方面，若 $A\subset M$，则 $\mathfrak m(M)\cap A=\mathfrak p$ 为 $A$ 的素理想，且 $A\setminus\mathfrak p$ 的元素在 $M$ 中可逆，故 $A_{\mathfrak p}\subset M$ 且 $\mathfrak p A_{\mathfrak p}\subset\mathfrak m(M)$。因此 $M$ 支配 $A_{\mathfrak p}$。
\egastatement{8.1.4}{引理} 设 $R$ 为域，$M,N$ 为 $R$ 的两个局部子环，$P$ 为由 $M\cup N$ 生成的 $R$ 的子环。下列条件等价：
\begin{enumerate}[label=(\roman*)]
\item 存在 $P$ 的素理想 $\mathfrak p$，使 $\mathfrak m(M)=\mathfrak p\cap M$，$\mathfrak m(N)=\mathfrak p\cap N$；
\item 由 $\mathfrak m(M)\cup\mathfrak m(N)$ 在 $P$ 中生成的理想 $\mathfrak a$ 不等于 $P$；
\item 存在 $R$ 的局部子环 $Q$，同时支配 $M,N$。
\end{enumerate}
显然 (i) 推出 (ii)。反之，若 $\mathfrak a\ne P$，则 $\mathfrak a$ 包含于 $P$ 的某个极大理想 $\mathfrak n$；由于 $1\notin\mathfrak n$，$\mathfrak n\cap M$ 包含 $\mathfrak m(M)$ 且不等于 $M$，故 $\mathfrak n\cap M=\mathfrak m(M)$；同理 $\mathfrak n\cap N=\mathfrak m(N)$。若 $Q$ 同时支配 $M,N$，则显然 $P\subset Q$，且
\[
\mathfrak m(M)=(\mathfrak m(Q)\cap P)\cap M,
\qquad
\mathfrak m(N)=(\mathfrak m(Q)\cap P)\cap N,
\]
故 (iii) 推出 (i)；取 $Q=P_{\mathfrak p}$ 即可得到逆向蕴含。

当（\hyperlink{para:1:8.1.4}{8.1.4}）的条件成立时，遵循 C. Chevalley 的术语，称局部环 $M,N$ \emph{相关}。
\egastatement{8.1.5}{命题} 设 $A,B$ 为域 $R$ 的两个子环，$C$ 为由 $A\cup B$ 生成的子环。下列条件等价：
\begin{enumerate}[label=(\roman*)]
\item 对每个包含 $A,B$ 的局部环 $Q$，令 $\mathfrak p=\mathfrak m(Q)\cap A$、$\mathfrak q=\mathfrak m(Q)\cap B$，则有 $A_{\mathfrak p}=B_{\mathfrak q}$；
\item 对 $C$ 的每个素理想 $\mathfrak r$，令 $\mathfrak p=\mathfrak r\cap A$、$\mathfrak q=\mathfrak r\cap B$，则有 $A_{\mathfrak p}=B_{\mathfrak q}$；
\item 若 $M\in L(A)$ 与 $N\in L(B)$ 相关，则它们相同；
\item $L(A)\cap L(B)=L(C)$。
\end{enumerate}
引理（\hyperlink{para:1:8.1.3}{8.1.3}）与（\hyperlink{para:1:8.1.4}{8.1.4}）证明了 (i) 与 (iii) 等价。将 (i) 用于 $Q=C_{\mathfrak r}$，显然得到 (ii)。反之，(ii) 推出 (i)：若 $Q$ 包含 $A\cup B$，则包含 $C$；令 $\mathfrak r=\mathfrak m(Q)\cap C$，由（\hyperlink{para:1:8.1.3}{8.1.3}）有 $\mathfrak p=\mathfrak r\cap A$、$\mathfrak q=\mathfrak r\cap B$。

(iv) 显然推出 (i)。事实上，若 $Q$ 包含 $A\cup B$，由（\hyperlink{para:1:8.1.3}{8.1.3}），它支配一个 $C_{\mathfrak r}\in L(C)$。按假设 $C_{\mathfrak r}\in L(A)\cap L(B)$，而（\hyperlink{para:1:8.1.1}{8.1.1}）与（\hyperlink{para:1:8.1.3}{8.1.3}）给出 $C_{\mathfrak r}=A_{\mathfrak p}=B_{\mathfrak q}$。

最后证明 (iii) 推出 (iv)。取 $Q\in L(C)$。由（\hyperlink{para:1:8.1.3}{8.1.3}），$Q$ 支配某个 $M\in L(A)$ 和某个 $N\in L(B)$；它们相关，故按假设相同。于是 $C\subset M$，所以 $M$ 支配某个 $Q'\in L(C)$（\hyperlink{para:1:8.1.3}{8.1.3}）。因此 $Q$ 支配 $Q'$，由（\hyperlink{para:1:8.1.3}{8.1.3}）必有 $Q=Q'=M$，故 $Q\in L(A)\cap L(B)$。反之，若 $Q\in L(A)\cap L(B)$，则 $C\subset Q$；因而（\hyperlink{para:1:8.1.3}{8.1.3}）$Q$ 支配某个 $Q'\in L(C)\subset L(A)\cap L(B)$。$Q,Q'$ 相关，故相同，于是 $Q=Q'\in L(C)$。证明完毕。


% END chapters/ch1_s8_1.tex

\par

% BEGIN chapters/ch1_s8_2.tex
\subsection{整概形的局部环}\hypertarget{heading:1:8.2}{}
\src{165—168}
\egapar{8.2.1} 设 $X$ 为整预概形，$R$ 为其有理函数域，即 $X$ 的泛点 $a$ 处的局部环。对任意 $x\in X$，已知 $\OO_x$ 典范地等同于 $R$ 的子环（\hyperlink{para:1:7.1.5}{7.1.5}）；对任意有理函数 $f\in R$，其定义域 $\delta(f)$ 为满足 $f\in\OO_x$ 的点 $x\in X$ 组成的开集。由（\hyperlink{para:1:7.2.6}{7.2.6}），对每个开集 $U\subset X$，有
\begin{equation}
\Gamma(U,\OO_X)=\bigcap_{x\in U}\OO_x.
\tag{8.2.1.1}\label{eq:1:8.2.1.1}
\end{equation}
\egastatement{8.2.2}{命题} 设 $X$ 为整预概形，$R$ 为其有理函数域。$X$ 为概形，当且仅当对 $X$ 的任意两点 $x,y$，关系“$\OO_x$ 与 $\OO_y$ 相关”（\hyperlink{para:1:8.1.4}{8.1.4}）蕴含 $x=y$。

假设此条件成立，证明 $X$ 分离。设 $U,V$ 为 $X$ 中的两个不同仿射开集，其环分别为 $A,B$，并等同于 $R$ 的子环。于是 $U$（分别为 $V$）等同于 $L(A)$（分别为 $L(B)$）（\hyperlink{para:1:8.1.2}{8.1.2}）。由假设与（\hyperlink{para:1:8.1.5}{8.1.5}），若 $C$ 为由 $A\cup B$ 生成的子环，则 $W=U\cap V$ 等同于 $L(A)\cap L(B)=L(C)$。已知 $R$ 的任一子环 $E$ 都等于 $L(E)$ 中各局部环的交（\cite{b1}，第 5-03 页，命题 4 bis），所以 $C$ 等同于所有 $z\in W$ 处的 $\OO_z$ 的交，即 $\Gamma(W,\OO_X)$（\hyperref[eq:1:8.2.1.1]{8.2.1.1}）。

考虑 $X$ 在 $W$ 上诱导的预概形。恒等同态 $\varphi:C\to\Gamma(W,\OO_X)$ 对应一个态射 $\Phi=(\psi,\theta):W\to\Spec(C)$（\hyperlink{para:1:2.2.4}{2.2.4}）。我们证明 $\Phi$ 为预概形同构，于是 $W$ 为仿射开集。由 $W$ 与 $L(C)=\Spec(C)$ 的等同，$\psi$ 为双射。另一方面，对每个 $x\in W$，令 $\mathfrak r=\mathfrak m_x\cap C$，则 $\theta_x^\sharp$ 为嵌入 $C_{\mathfrak r}\to\OO_x$；按定义 $C_{\mathfrak r}$ 等同于 $\OO_x$，故该同态为双射。

只须再证明 $\psi$ 为同胚，即 $W$ 的每个闭子集 $F$ 的像 $\psi(F)$ 在 $\Spec(C)$ 中闭。$F$ 是 $U$ 的某个形如 $V(\mathfrak a)$ 的闭集在 $W$ 上的迹，其中 $\mathfrak a$ 为 $A$ 的理想。证明 $\psi(F)=V(\mathfrak a C)$ 即足够。事实上，$C$ 中包含 $\mathfrak a C$ 的素理想就是包含 $\mathfrak a$ 的素理想，亦即形如 $\psi(x)=\mathfrak m_x\cap C$ 且 $\mathfrak a\subset\mathfrak m_x$、$x\in W$ 的理想；而在 $U$ 中，$\mathfrak a\subset\mathfrak m_x$ 等价于 $x\in V(\mathfrak a)$，故有 $\psi(F)=V(\mathfrak a C)$。

因此 $X$ 分离，因为 $U\cap V$ 仿射且其环 $C$ 由 $U,V$ 的环的并 $A\cup B$ 生成（\hyperlink{para:1:5.5.6}{5.5.6}）。

反之，假设 $X$ 分离，并取两点 $x,y\in X$，使 $\OO_x,\OO_y$ 相关。取包含 $x$（分别为 $y$）的仿射开集 $U$（分别为 $V$），环为 $A$（分别为 $B$）。已知 $U\cap V$ 仿射，其环 $C$ 由 $A\cup B$ 生成（\hyperlink{para:1:5.5.6}{5.5.6}）。令 $\mathfrak p=\mathfrak m_x\cap A$、$\mathfrak q=\mathfrak m_y\cap B$，则 $A_{\mathfrak p}=\OO_x$、$B_{\mathfrak q}=\OO_y$。因它们相关，存在 $C$ 的素理想 $\mathfrak r$，使 $\mathfrak p=\mathfrak r\cap A$、$\mathfrak q=\mathfrak r\cap B$（\hyperlink{para:1:8.1.4}{8.1.4}）。由 $U\cap V$ 仿射，存在 $z\in U\cap V$ 使 $\mathfrak r=\mathfrak m_z\cap C$，显然 $x=z$ 且 $y=z$，故 $x=y$。
\egastatement{8.2.3}{推论} 设 $X$ 为整概形，$x,y\in X$。则 $x\in\overline{\{y\}}$，当且仅当 $\OO_x\subset\OO_y$；换言之，当且仅当每个在 $x$ 处有定义的有理函数也在 $y$ 处有定义。

必要性显然，因为有理函数 $f\in R$ 的定义域 $\delta(f)$ 是开集。证明充分性。若 $\OO_x\subset\OO_y$，则存在 $\OO_x$ 的素理想 $\mathfrak p$，使 $\OO_y$ 支配 $(\OO_x)_{\mathfrak p}$（\hyperlink{para:1:8.1.3}{8.1.3}）。又由（\hyperlink{para:1:2.4.2}{2.4.2}），存在 $z\in X$ 使 $x\in\overline{\{z\}}$，且 $\OO_z=(\OO_x)_{\mathfrak p}$。因 $\OO_z,\OO_y$ 相关，由（\hyperlink{para:1:8.2.2}{8.2.2}）得 $z=y$，从而得到推论。
\egastatement{8.2.4}{推论} 若 $X$ 为整概形，则映射 $x\mapsto\OO_x$ 为单射。换言之，若 $x,y$ 为 $X$ 的两个不同点，则存在一个有理函数，在其中一点有定义而在另一点无定义。

这由（\hyperlink{para:1:8.2.3}{8.2.3}）与公理（$T_0$）（\hyperlink{para:1:2.1.4}{2.1.4}）得出。
\egastatement{8.2.5}{推论} 设 $X$ 为底空间诺特的整概形。当 $f$ 遍历 $X$ 的有理函数域 $R$ 时，各集合 $\delta(f)$ 生成 $X$ 的拓扑。

此时 $X$ 的每个闭子集都是有限个不可约闭集的并，也就是有限个形如 $\overline{\{y\}}$ 的集合之并（\hyperlink{para:1:2.1.5}{2.1.5}）。若 $x\notin\overline{\{y\}}$，则存在在 $x$ 处有定义而在 $y$ 处无定义的有理函数 $f$（\hyperlink{para:1:8.2.3}{8.2.3}），亦即 $x\in\delta(f)$ 且 $\delta(f)$ 不与 $\overline{\{y\}}$ 相交。因此 $\overline{\{y\}}$ 的补集为形如 $\delta(f)$ 的集合之并；由第一句话，$X$ 的每个开集都是这类开集的有限交的并。
\egapar{8.2.6} 推论（\hyperlink{para:1:8.2.5}{8.2.5}）表明，给定以 $R$ 为分式域的局部环族 $(\OO_x)_{x\in X}$，就完全确定了 $X$ 的拓扑。也可以等价地如下定义 $X$ 的闭集：给定 $X$ 的有限子集 $\{x_1,\ldots,x_n\}$，考虑所有满足 $\OO_y\subset\OO_{x_i}$ 对至少一个 $i$ 成立的 $y\in X$；遍历所有有限子集，所得集合正是 $X$ 的闭集。此外，一旦知道 $X$ 的拓扑，结构层 $\OO_X$ 也被各 $\OO_x$ 唯一确定，因为（\hyperref[eq:1:8.2.1.1]{8.2.1.1}）给出 $\Gamma(U,\OO_X)=\bigcap_{x\in U}\OO_x$。所以，当 $X$ 为底空间诺特的整概形时，族 $(\OO_x)_{x\in X}$ 完全确定预概形 $X$。
\egastatement{8.2.7}{命题} 设 $X,Y$ 为两个整概形，$f:X\to Y$ 为支配态射（\hyperlink{para:1:2.2.6}{2.2.6}），$K$（分别为 $L$）为 $X$（分别为 $Y$）的有理函数域。则 $L$ 等同于 $K$ 的子域，且对每个 $x\in X$，$\OO_{f(x)}$ 是 $Y$ 的唯一被 $\OO_x$ 支配的局部环。

若 $f=(\psi,\theta)$，$a$ 为 $X$ 的泛点，则 $\psi(a)$ 是 $Y$ 的泛点（0，\hyperlink{para:0:2.1.5}{2.1.5}）。因此 $\theta_a^\sharp$ 为域 $L=\OO_{\psi(a)}$ 到域 $K=\OO_a$ 的单同态。由于 $Y$ 的每个非空仿射开集 $U$ 都包含 $\psi(a)$，由（\hyperlink{para:1:2.2.4}{2.2.4}），与 $f$ 对应的同态
\[
 \Gamma(U,\OO_Y)\longrightarrow\Gamma(\psi^{-1}(U),\OO_X)
\]
就是 $\theta_a^\sharp$ 在 $\Gamma(U,\OO_Y)$ 上的限制。因此，对每个 $x\in X$，$\theta_x^\sharp$ 为 $\theta_a^\sharp$ 在 $\OO_{\psi(x)}$ 上的限制，故也是单同态。又已知 $\theta_x^\sharp$ 为局部同态，所以通过 $\theta_a^\sharp$ 将 $L$ 等同于 $K$ 的子域后，$\OO_{\psi(x)}$ 被 $\OO_x$ 支配（\hyperlink{para:1:8.1.1}{8.1.1}）。它也是 $Y$ 的唯一被 $\OO_x$ 支配的局部环，因为 $Y$ 的两个相关局部环必相同（\hyperlink{para:1:8.2.2}{8.2.2}）。
\egastatement{8.2.8}{命题} 设 $X$ 为不可约预概形，$f:X\to Y$ 为局部浸入（分别为局部同构），并假设态射 $f$ 分离。则 $f$ 为浸入（分别为开浸入）。

令 $f=(\psi,\theta)$。两种情形都只须证明 $\psi$ 为 $X$ 到 $\psi(X)$ 的同胚（\hyperlink{para:1:4.5.3}{4.5.3}）。以 $f_{\mathrm{red}}$ 代替 $f$（\hyperlink{para:1:5.1.6}{5.1.6} 与 \hyperlink{para:1:5.5.1}{5.5.1}，(vi)），可假设 $X,Y$ 约化。令 $Y'$ 为 $Y$ 中以 $\overline{\psi(X)}$ 为底空间的约化闭子预概形，则 $f$ 分解为 $X\xrightarrow{f'}Y'\xrightarrow{j}Y$，其中 $j$ 为典范嵌入（\hyperlink{para:1:5.2.2}{5.2.2}）。由（\hyperlink{para:1:5.5.1}{5.5.1}，(v)），$f'$ 仍分离；它也仍是局部浸入（分别为局部同构）：问题在 $X,Y$ 上都是局部的，故只须考虑 $f$ 为闭浸入（分别为开浸入）的情形，此时断言立即由（\hyperlink{para:1:4.2.2}{4.2.2}）得出。

所以可假设 $f$ 支配，这又使 $Y$ 不可约（0，\hyperlink{para:0:2.1.5}{2.1.5}），因而 $X,Y$ 均为整预概形。问题在 $Y$ 上是局部的，故可假设 $Y$ 仿射；因 $f$ 分离，$X$ 为概形（\hyperlink{para:1:5.5.1}{5.5.1}，(ii)）。于是已处于（\hyperlink{para:1:8.2.7}{8.2.7}）的假设之下。因此对每个 $x\in X$，$\theta_x^\sharp$ 为单射；而 $f$ 为局部浸入，又使 $\theta_x^\sharp$ 满射（\hyperlink{para:1:4.2.2}{4.2.2}），故为双射，即在（\hyperlink{para:1:8.2.7}{8.2.7}）的等同下有 $\OO_{\psi(x)}=\OO_x$。由（\hyperlink{para:1:8.2.4}{8.2.4}），$\psi$ 为单射；当 $f$ 为局部同构时，这就已证明命题（\hyperlink{para:1:4.5.3}{4.5.3}）。

若仅假设 $f$ 为局部浸入，对每个 $x\in X$，存在 $x$ 在 $X$ 中的开邻域 $U$，以及 $\psi(x)$ 在 $Y$ 中的开邻域 $V$，使 $\psi$ 在 $U$ 上的限制为 $U$ 到 $V$ 的一个闭子集的同胚。由于 $U$ 在 $X$ 中稠密，$\psi(U)$ 在 $Y$ 中稠密，因而在 $V$ 中稠密，故 $\psi(U)=V$。由 $\psi$ 单射，$\psi^{-1}(V)=U$，这就证明了 $\psi$ 为 $X$ 到 $\psi(X)$ 的同胚。


% END chapters/ch1_s8_2.tex

\par

% BEGIN chapters/ch1_s8_3.tex
\subsection{Chevalley 概形}\hypertarget{heading:1:8.3}{}
\src{168—169}
\egapar{8.3.1} 设 $X$ 为诺特整概形，$R$ 为其有理函数域；记 $X'$ 为所有局部子环 $\OO_x\subset R$（$x$ 遍历 $X$）组成的集合。$X'$ 满足以下三个条件。

\noindent\textbf{（Sch.1）} 对每个 $M\in X'$，$R$ 是 $M$ 的分式域。

\noindent\textbf{（Sch.2）} 存在 $R$ 的有限个诺特子环 $A_i$，使 $X'=\bigcup_iL(A_i)$，并且对每对指标 $i,j$，由 $A_i\cup A_j$ 生成的 $R$ 的子环 $A_{ij}$ 是有限型 $A_i$-代数。

\noindent\textbf{（Sch.3）} 若 $X'$ 的两个元素 $M,N$ 相关，则它们相同。

事实上，（Sch.1）已见于（\hyperlink{para:1:8.2.1}{8.2.1}），而（Sch.3）由（\hyperlink{para:1:8.2.2}{8.2.2}）得出。为证明（Sch.2），只须用有限个环为诺特环的仿射开集 $U_i$ 覆盖 $X$，并取 $A_i=\Gamma(U_i,\OO_X)$。$X$ 为概形的假设蕴含 $U_i\cap U_j$ 仿射，且 $\Gamma(U_i\cap U_j,\OO_X)=A_{ij}$（\hyperlink{para:1:5.5.6}{5.5.6}）。此外，由于空间 $U_i$ 诺特，浸入 $U_i\cap U_j\to U_i$ 为有限型（\hyperlink{para:1:6.3.5}{6.3.5}），故 $A_{ij}$ 是有限型 $A_i$-代数（\hyperlink{para:1:6.3.3}{6.3.3}）。
\egapar{8.3.2} 以（Sch.1）、（Sch.2）、（Sch.3）为公理的结构，推广了 C. Chevalley 意义下的“概形”。Chevalley 还假设 $R$ 是某个域 $K$ 的有限型扩张，且各 $A_i$ 是有限型 $K$-代数；这使（Sch.2）中的一部分成为多余的\cite{b1}。反之，若在集合 $X'$ 上给定这样的结构，便可利用（\hyperlink{para:1:8.2.6}{8.2.6}）的说明，为它构造一个相伴的整概形 $X$：取 $X$ 的底空间为配备（\hyperlink{para:1:8.2.6}{8.2.6}）所定义拓扑的 $X'$，并取层 $\OO_X$ 使得对每个开集 $U\subset X$，都有
\[
 \Gamma(U,\OO_X)=\bigcap_{x\in U}\OO_x,
\]
限制同态按显然的方式定义。这样确实得到一个整概形，其局部环正是 $X'$ 的元素；我们将此验证留给读者。以后不使用这个结果。


% END chapters/ch1_s8_3.tex

\par

% BEGIN chapters/ch1_s9_1.tex
\section{拟凝聚层的补充知识}\hypertarget{heading:1:9}{}
\subsection{拟凝聚层的张量积}\hypertarget{heading:1:9.1}{}
\src{169—171}
\egastatement{9.1.1}{命题} 设 $X$ 为预概形（分别为局部诺特预概形），$\F,\G$ 为两个拟凝聚（分别为凝聚）$\OO_X$-模。则 $\F\otimes_{\OO_X}\G$ 拟凝聚（分别为凝聚）；若 $\F,\G$ 为有限型，则该张量积也为有限型。若 $\F$ 有有限表示，且 $\G$ 拟凝聚（分别为凝聚），则 $\shHom_{\OO_X}(\F,\G)$ 拟凝聚（分别为凝聚）。

问题是局部的，可假设 $X$ 仿射（分别为诺特仿射）。此外，若 $\F$ 凝聚，可假设它是某个同态 $\OO_X^p\to\OO_X^q$ 的余核。拟凝聚层的断言由（\hyperlink{para:1:1.3.12}{1.3.12}）与（\hyperlink{para:1:1.3.9}{1.3.9}）得出；凝聚层的断言由（\hyperlink{para:1:1.5.1}{1.5.1}）及以下事实得出：若 $M,N$ 为诺特环 $A$ 上的有限型模，则 $M\otimes_A N$ 与 $\Hom_A(M,N)$ 都是有限型 $A$-模。
\egastatement{9.1.2}{定义} 设 $X,Y$ 为两个 $S$-预概形，$p,q$ 为 $X\times_S Y$ 的投影，$\F$（分别为 $\G$）为拟凝聚 $\OO_X$-模（分别为 $\OO_Y$-模）。称预概形 $X\times_S Y$ 上的张量积
\[
 p^*(\F)\otimes_{\OO_{X\times_S Y}}q^*(\G)
\]
为 $\F,\G$ 在 $\OO_S$ 上（或 $S$ 上）的\emph{张量积}，记作 $\F\otimes_{\OO_S}\G$ 或 $\F\otimes_S\G$。

若 $X_i$（$1\leq i\leq n$）为 $S$-预概形，$\F_i$ 为拟凝聚 $\OO_{X_i}$-模，则同样可在 $Z=X_1\times_S\cdots\times_S X_n$ 上定义张量积 $\F_1\otimes_S\cdots\otimes_S\F_n$。由（\hyperlink{para:1:9.1.1}{9.1.1}）与（0，\hyperlink{para:0:5.1.4}{5.1.4}），它为拟凝聚 $\OO_Z$-模；若各 $\F_i$ 凝聚且 $Z$ 局部诺特，则由（\hyperlink{para:1:9.1.1}{9.1.1}）、（0，\hyperlink{para:0:5.3.11}{5.3.11}）与（\hyperlink{para:1:6.1.1}{6.1.1}），它凝聚。

取 $X=Y=S$ 时，定义（\hyperlink{para:1:9.1.2}{9.1.2}）给出通常的 $\OO_S$-模的张量积。此外，由 $q^*(\OO_Y)=\OO_{X\times_S Y}$（0，\hyperlink{para:0:4.3.4}{4.3.4}），张量积 $\F\otimes_S\OO_Y$ 典范地等同于 $p^*(\F)$；同理，$\OO_X\otimes_S\G$ 典范地等同于 $q^*(\G)$。更特别地，取 $Y=S$，并令 $f:X\to Y$ 为结构态射，就有 $\OO_X\otimes_Y\G=f^*(\G)$。因此通常的张量积与逆像都是此一般张量积的特例。

由定义（\hyperlink{para:1:9.1.2}{9.1.2}）立即得知，对固定的 $X,Y$，$\F\otimes_S\G$ 是关于 $\F,\G$ 的协变、可加、右正合双函子。
\egastatement{9.1.3}{命题} 设 $S,X,Y$ 为三个仿射概形，环分别为 $A,B,C$，其中 $B,C$ 为 $A$-代数。设 $M$（分别为 $N$）为 $B$-模（分别为 $C$-模），$\F=\wt M$（分别为 $\G=\wt N$）为相伴的拟凝聚层。则 $\F\otimes_S\G$ 典范地同构于 $(B\otimes_A C)$-模 $M\otimes_A N$ 所伴随的层。

事实上，由（\hyperlink{para:1:1.6.5}{1.6.5}），$\F\otimes_S\G$ 典范地同构于以下 $(B\otimes_A C)$-模所伴随的层：
\[
 \bigl(M\otimes_B(B\otimes_A C)\bigr)
 \otimes_{B\otimes_A C}
 \bigl((B\otimes_A C)\otimes_C N\bigr).
\]
由张量积间的典范同构，该模同构于
\[
 M\otimes_B(B\otimes_A C)\otimes_C N
 =(M\otimes_B B)\otimes_A(C\otimes_C N)=M\otimes_A N.
\]
\egastatement{9.1.4}{命题} 设 $f:T\to X$、$g:T\to Y$ 为两个 $S$-态射，$\F$（分别为 $\G$）为拟凝聚 $\OO_X$-模（分别为 $\OO_Y$-模）。则
\[
 (f,g)_S^*(\F\otimes_S\G)
 =f^*(\F)\otimes_{\OO_T}g^*(\G).
\]
若 $p,q$ 为 $X\times_S Y$ 的投影，则此式由 $(f,g)_S^*\circ p^*=f^*$ 与 $(f,g)_S^*\circ q^*=g^*$（0，\hyperlink{para:0:3.5.5}{3.5.5}），以及代数层的张量积之逆像等于其逆像的张量积（0，\hyperlink{para:0:4.3.3}{4.3.3}）得出。
\egastatement{9.1.5}{推论} 设 $f:X\to X'$、$g:Y\to Y'$ 为两个 $S$-态射，$\F'$（分别为 $\G'$）为拟凝聚 $\OO_{X'}$-模（分别为 $\OO_{Y'}$-模）。则
\[
 (f\times_S g)^*(\F'\otimes_S\G')
 =f^*(\F')\otimes_Sg^*(\G').
\]
这由（\hyperlink{para:1:9.1.4}{9.1.4}）及 $f\times_Sg=(f\circ p,g\circ q)_S$ 得出，其中 $p,q$ 为 $X\times_S Y$ 的投影。
\egastatement{9.1.6}{推论} 设 $X,Y,Z$ 为三个 $S$-预概形，$\F,\G,\HH$ 分别为其上的拟凝聚模。则层 $\F\otimes_S\G\otimes_S\HH$，是 $(\F\otimes_S\G)\otimes_S\HH$ 经典范同构
\[
 X\times_S Y\times_S Z\longrightarrow(X\times_S Y)\times_S Z
\]
的逆像。

事实上，记 $p_1,p_2,p_3$ 为三重积的投影，则此同构写为 $((p_1,p_2)_S,p_3)_S$。

同样，$\G\otimes_S\F$ 经 $X\times_S Y\to Y\times_S X$ 的典范同构的逆像，是 $\F\otimes_S\G$。
\egastatement{9.1.7}{推论} 若 $X$ 为 $S$-预概形，则每个拟凝聚 $\OO_X$-模 $\F$ 都是 $\F\otimes_S\OO_S$ 经典范同构 $X\to X\times_S S$（\hyperlink{para:1:3.3.3}{3.3.3}）的逆像。

此同构为 $(1_X,\varphi)_S$，其中 $\varphi:X\to S$ 为结构态射；推论由（\hyperlink{para:1:9.1.4}{9.1.4}）及 $\varphi^*(\OO_S)=\OO_X$ 得出。
\egapar{9.1.8} 设 $X$ 为 $S$-预概形，$\F$ 为拟凝聚 $\OO_X$-模，$\varphi:S'\to S$ 为态射。将 $X\times_S S'=X_{(\varphi)}=X_{(S')}$ 上的拟凝聚层 $\F\otimes_S\OO_{S'}$ 记作 $\F_{(\varphi)}$ 或 $\F_{(S')}$。因此 $\F_{(S')}=p^*(\F)$，其中 $p:X_{(S')}\to X$ 为投影。
\egastatement{9.1.9}{命题} 设 $\varphi':S''\to S'$ 为态射。对 $S$-预概形 $X$ 上的每个拟凝聚 $\OO_X$-模 $\F$，$(\F_{(\varphi)})_{(\varphi')}$ 是 $\F_{(\varphi\circ\varphi')}$ 经典范同构 $(X_{(\varphi)})_{(\varphi')}\to X_{(\varphi\circ\varphi')}$（\hyperlink{para:1:3.3.9}{3.3.9}）的逆像。

这立即由定义与（\hyperlink{para:1:3.3.9}{3.3.9}）得出，也可写为
\begin{equation}
 (\F\otimes_S\OO_{S'})\otimes_{S'}\OO_{S''}
 =\F\otimes_S\OO_{S''}.
\tag{9.1.9.1}\label{eq:1:9.1.9.1}
\end{equation}
\egastatement{9.1.10}{命题} 设 $Y$ 为 $S$-预概形，$f:X\to Y$ 为 $S$-态射。对每个拟凝聚 $\OO_Y$-模 $\G$ 及每个态射 $S'\to S$，有
\[
 (f_{(S')})^*(\G_{(S')})=(f^*(\G))_{(S')}.
\]
这立即由下图的交换性得出：
\[
\begin{tikzcd}
 X_{(S')}\arrow[r,"f_{(S')}"]\arrow[d]&Y_{(S')}\arrow[d]\\
 X\arrow[r,"f"]&Y.
\end{tikzcd}
\]
\egastatement{9.1.11}{推论} 设 $X,Y$ 为两个 $S$-预概形，$\F$（分别为 $\G$）为拟凝聚 $\OO_X$-模（分别为 $\OO_Y$-模）。层 $\F_{(S')}\otimes_{S'}\G_{(S')}$ 经典范同构
\[
 (X\times_S Y)_{(S')}\longrightarrow X_{(S')}\times_{S'}Y_{(S')}
 \qquad\text{（\hyperlink{para:1:3.3.10}{3.3.10}）}
\]
的逆像，等于 $(\F\otimes_S\G)_{(S')}$。

若 $p,q$ 为 $X\times_S Y$ 的投影，此同构就是 $(p_{(S')},q_{(S')})_{S'}$；推论由（\hyperlink{para:1:9.1.4}{9.1.4}）与（\hyperlink{para:1:9.1.10}{9.1.10}）得出。
\egastatement{9.1.12}{命题} 沿用（\hyperlink{para:1:9.1.2}{9.1.2}）的记号，设 $z\in X\times_S Y$，$x=p(z)$，$y=q(z)$。则茎 $(\F\otimes_S\G)_z$ 同构于
\[
 (\F_x\otimes_{\OO_x}\OO_z)\otimes_{\OO_z}
 (\G_y\otimes_{\OO_y}\OO_z)
 =\F_x\otimes_{\OO_x}\OO_z\otimes_{\OO_y}\G_y.
\]
可归结到仿射情形，所以命题由（\hyperlink{para:1:1.6.5}{1.6.5}）中的茎公式得出。\transnote{原文引用“（\hyperref[eq:1:1.6.5.1]{1.6.5.1}）”，但该处的相应公式在底本中未单独编号。}
\egastatement{9.1.13}{推论} 若 $\F,\G$ 为有限型，则
\[
 \Supp(\F\otimes_S\G)
 =p^{-1}(\Supp(\F))\cap q^{-1}(\Supp(\G)).
\]
由于 $p^*(\F)$ 与 $q^*(\G)$ 都是有限型 $\OO_{X\times_S Y}$-模，利用（\hyperlink{para:1:9.1.12}{9.1.12}）与（0，\hyperlink{para:0:1.7.5}{1.7.5}），可归结到 $\G=\OO_Y$ 的特例，即证明
\begin{equation}
 \Supp(p^*(\F))=p^{-1}(\Supp(\F)).
\tag{9.1.13.1}\label{eq:1:9.1.13.1}
\end{equation}
\transnote{原文左端印作 $\Supp(p^{-1}(\F))$；上下文涉及的是模的拉回 $p^*(\F)$，此处按该含义书写，并保留原编号。}
与（0，\hyperlink{para:0:1.7.5}{1.7.5}）相同的论证，将问题归结到验证：对每个 $z\in X\times_S Y$，令 $x=p(z)$，有 $\OO_z/\mathfrak m_x\OO_z\ne0$。这来自同态 $\OO_x\to\OO_z$ 按假设为局部同态。

我们将本小节中两个因子情形的结果推广到任意多个因子的工作留给读者。


% END chapters/ch1_s9_1.tex

\par

% BEGIN chapters/ch1_s9_2.tex
\subsection{拟凝聚层的直像}\hypertarget{heading:1:9.2}{}
\src{171—172}
\egastatement{9.2.1}{命题} 设 $f:X\to Y$ 为预概形态射。假设 $Y$ 存在一个由仿射开集组成的覆盖 $(Y_\alpha)$，具有下述性质：每个 $f^{-1}(Y_\alpha)$ 都有由包含于它的仿射开集组成的有限覆盖 $(X_{\alpha i})$，并且每个交集 $X_{\alpha i}\cap X_{\alpha j}$ 本身也是有限个仿射开集之并。那么，对每个拟凝聚 $\OO_X$-模 $\F$，$f_*(\F)$ 为拟凝聚 $\OO_Y$-模。

问题在 $Y$ 上是局部的，可假设 $Y$ 等于某个 $Y_\alpha$，并省略下标 $\alpha$。

a) 先设各 $X_i\cap X_j$ 本身就是仿射开集。令 $\F_i=\F|X_i$、$\F_{ij}=\F|(X_i\cap X_j)$，并将它们分别经 $f$ 在 $X_i$ 及 $X_i\cap X_j$ 上的限制所得的直像记作 $\F_i'$ 与 $\F_{ij}'$；已知这些直像拟凝聚（\hyperlink{para:1:1.6.3}{1.6.3}）。令 $\G=\bigoplus_i\F_i'$、$\HH=\bigoplus_{i,j}\F_{ij}'$；它们为拟凝聚 $\OO_Y$-模。定义同态 $u:\G\to\HH$，使 $f_*(\F)$ 为其核，即可推出 $f_*(\F)$ 拟凝聚（\hyperlink{para:1:1.3.9}{1.3.9}）。只须将 $u$ 定义为预层同态。因此，对每个开集 $W\subset Y$，只须定义同态
\[
 u_W:\bigoplus_i\Gamma(f^{-1}(W)\cap X_i,\F)
 \longrightarrow\bigoplus_{i,j}\Gamma(f^{-1}(W)\cap X_i\cap X_j,\F),
\]
使其满足关于 $W$ 的通常相容条件。对截面 $s_i\in\Gamma(f^{-1}(W)\cap X_i,\F)$，以 $s_{ij}$ 表示其在 $f^{-1}(W)\cap X_i\cap X_j$ 上的限制，取
\[
 u_W((s_i)_i)=(s_{ij}-s_{ji})_{i,j}.
\]
相容条件显然成立。令 $\mathscr K=\Ker(u)$，定义 $f_*(\F)\to\mathscr K$：将每个 $s\in\Gamma(f^{-1}(W),\F)$ 映到其在各 $f^{-1}(W)\cap X_i$ 上的限制组成的族 $(s_i)$。层公理（F1）与（F2）（G，II，1.1）说明此同态为双射，从而证明了此种情形。

b) 一般情形，只要证明各 $\F_{ij}'$ 拟凝聚，上面的论证就仍适用。按假设，$X_i\cap X_j$ 为有限个仿射开集 $X_{ijk}$ 之并；这些 $X_{ijk}$ 是一个概形中的仿射开集，故其中任意两个之交仍是仿射开集（\hyperlink{para:1:5.5.6}{5.5.6}）。于是归结到第一种情形，（\hyperlink{para:1:9.2.1}{9.2.1}）得证。
\egastatement{9.2.2}{推论} （\hyperlink{para:1:9.2.1}{9.2.1}）的结论在下列各情形成立：
\begin{enumerate}[label=\alph*)]
\item $f$ 分离且拟紧；
\item $f$ 分离且为有限型；
\item $f$ 拟紧，且 $X$ 的底空间局部诺特。
\end{enumerate}
在 a) 中，各 $X_i\cap X_j$ 仿射（\hyperlink{para:1:5.5.6}{5.5.6}）。b) 是 a) 的特例（\hyperlink{para:1:6.6.3}{6.6.3}）。最后，在 c) 中，可归结到 $Y$ 仿射且 $X$ 的底空间诺特的情形；此时 $X$ 有有限仿射开覆盖 $(X_i)$，而各 $X_i\cap X_j$ 拟紧，故为有限个仿射开集之并（\hyperlink{para:1:2.1.3}{2.1.3}）。


% END chapters/ch1_s9_2.tex

\par

% BEGIN chapters/ch1_s9_3.tex
\subsection{拟凝聚层截面的延拓}\hypertarget{heading:1:9.3}{}
\src{172—173}
\egastatement{9.3.1}{定理} 设 $X$ 为底空间诺特的预概形，或为底空间拟紧的概形。设 $\LL$ 为可逆 $\OO_X$-模（0，\hyperlink{para:0:5.4.1}{5.4.1}），$f$ 为 $\LL$ 在 $X$ 上的截面，$X_f$ 为所有满足 $f(x)\ne0$ 的 $x\in X$ 组成的开集（0，\hyperlink{para:0:5.5.1}{5.5.1}），$\F$ 为拟凝聚 $\OO_X$-模。
\begin{enumerate}[label=(\roman*)]
\item 若 $s\in\Gamma(X,\F)$ 且 $s|X_f=0$，则存在整数 $n>0$ 使 $s\otimes f^{\otimes n}=0$。
\item 对每个截面 $s\in\Gamma(X_f,\F)$，存在整数 $n>0$，使 $s\otimes f^{\otimes n}$ 可延拓为 $\F\otimes\LL^{\otimes n}$ 在 $X$ 上的截面。
\end{enumerate}
(i) 由于 $X$ 的底空间拟紧，它是有限个仿射开集 $U_i$ 之并，可使每个 $\LL|U_i$ 同构于 $\OO_X|U_i$。故归结到 $X$ 仿射且 $\LL=\OO_X$ 的情形。此时 $f$ 等同于 $A(X)$ 的一个元素，且 $X_f=D(f)$；$s$ 等同于某个 $A(X)$-模 $M$ 中的元素，而 $s|X_f$ 等同于 $M_f$ 中的相应元素。由分式模的定义，结论显然。

(ii) 同样可将 $X$ 表为有限个仿射开集 $U_i$（$1\leq i\leq r$）之并，使 $\LL|U_i\simeq\OO_X|U_i$。对每个 $i$，经此同构，$(s\otimes f^{\otimes n})|(U_i\cap X_f)$ 等同于 $(f|(U_i\cap X_f))^n(s|(U_i\cap X_f))$。由（\hyperlink{para:1:1.4.1}{1.4.1}），存在整数 $n>0$，使对每个 $i$，此截面可延拓为 $\F\otimes\LL^{\otimes n}$ 在 $U_i$ 上的截面 $s_i$。记 $s_{ij}$ 为 $s_i$ 在 $U_i\cap U_j$ 上的限制；按定义，$s_{ij}-s_{ji}$ 在 $X_f\cap U_i\cap U_j$ 中为零。

若 $X$ 的底空间诺特，则 $U_i\cap U_j$ 拟紧；若 $X$ 为概形，则 $U_i\cap U_j$ 为仿射开集（\hyperlink{para:1:5.5.6}{5.5.6}），故仍拟紧。由 (i)，存在与 $i,j$ 无关的整数 $m$，使 $(s_{ij}-s_{ji})\otimes f^{\otimes m}=0$。于是立即得到 $\F\otimes\LL^{\otimes(n+m)}$ 在 $X$ 上的一个截面 $s'$，在各 $U_i$ 上诱导 $s_i\otimes f^{\otimes m}$，从而在 $X_f$ 上诱导 $s\otimes f^{\otimes(n+m)}$。

以下推论用更代数化的语言解释定理（\hyperlink{para:1:9.3.1}{9.3.1}）。
\egastatement{9.3.2}{推论} 在（\hyperlink{para:1:9.3.1}{9.3.1}）的假设下，考虑分次环 $A_* =\Gamma_*(\LL)$ 及分次 $A_*$-模 $M_* =\Gamma_*(\LL,\F)$（0，\hyperlink{para:0:5.4.6}{5.4.6}）。若 $f\in A_n$，$n\in\mathbf Z$，则有典范同构
\[
 \Gamma(X_f,\F)\simeq((M_*)_f)_0,
\]
其中右端为分式模 $(M_*)_f$ 的零次齐次元素组成的子群。
\egastatement{9.3.3}{推论} 假设满足（\hyperlink{para:1:9.3.1}{9.3.1}）的假设，且 $\LL=\OO_X$。令 $A=\Gamma(X,\OO_X)$、$M=\Gamma(X,\F)$，则 $A$-模 $\Gamma(X_f,\F)$ 典范地同构于 $M_f$。
\egastatement{9.3.4}{命题} 设 $X$ 为诺特预概形，$\F$ 为凝聚 $\OO_X$-模，$\J$ 为 $\OO_X$ 中的凝聚理想层，且 $\F$ 的支集包含于 $\OO_X/\J$ 的支集中。那么，存在整数 $n>0$，使 $\J^n\F=0$。

由于 $X$ 为有限个环诺特的仿射开集之并，可假设 $X$ 为诺特环 $A$ 的仿射概形。此时 $\F=\wt M$，其中 $M=\Gamma(X,\F)$ 为有限型 $A$-模；$\J=\wt J$，其中 $J=\Gamma(X,\J)$ 为 $A$ 的理想（\hyperlink{para:1:1.4.1}{1.4.1} 与 \hyperlink{para:1:1.5.1}{1.5.1}）。因 $A$ 诺特，$J$ 有有限生成系 $f_i$（$1\leq i\leq m$）。按假设，$\F$ 在 $X$ 上的任一截面在每个 $D(f_i)$ 上都为零。若 $s_j$（$1\leq j\leq q$）为生成 $M$ 的截面，则存在与 $i,j$ 无关的整数 $h$，使 $f_i^h s_j=0$（\hyperlink{para:1:1.4.1}{1.4.1}）；故 $f_i^hs=0$ 对每个 $s\in M$ 成立。取 $n=mh$，即有 $J^nM=0$；相应的 $\OO_X$-模 $\J^n\F=\widetilde{J^nM}$（\hyperlink{para:1:1.3.13}{1.3.13}）也为零。
\egastatement{9.3.5}{推论} 在（\hyperlink{para:1:9.3.4}{9.3.4}）的假设下，存在 $X$ 的闭子预概形 $Y$，其底空间为 $\OO_X/\J$ 的支集，并且若 $j:Y\to X$ 为典范嵌入，则 $\F=j_*(j^*(\F))$。

首先，$\OO_X/\J$ 与 $\OO_X/\J^n$ 的支集相同：若 $\J_x=\OO_x$，则 $\J_x^n=\OO_x$，而另一方面对每个 $x\in X$ 都有 $\J_x^n\subset\J_x$。故利用（\hyperlink{para:1:9.3.4}{9.3.4}），可假设 $\J\F=0$。取 $Y$ 为由 $\J$ 定义的闭子预概形；由于 $\F$ 此时为 $(\OO_X/\J)$-模，结论立即成立。


% END chapters/ch1_s9_3.tex

\par

% BEGIN chapters/ch1_s9_4.tex
\subsection{拟凝聚层的延拓}\hypertarget{heading:1:9.4}{}
\src{174—176}
\egapar{9.4.1} 设 $X$ 为拓扑空间，$\F$ 为 $X$ 上的集合层（分别为群层、环层），$U$ 为 $X$ 的开子集，$\psi:U\to X$ 为典范嵌入，$\G$ 为 $\F|U=\psi^*(\F)$ 的子层。由于 $\psi_*$ 左正合，$\psi_*(\G)$ 为 $\psi_*(\psi^*(\F))$ 的子层。考虑典范同态 $\rho:\F\to\psi_*(\psi^*(\F))$（0，\hyperlink{para:0:3.5.3}{3.5.3}），将 $\F$ 的子层 $\rho^{-1}(\psi_*(\G))$ 记作 $\overline\G$。由定义立即得知，对 $X$ 的任一开集 $V$，$\Gamma(V,\overline\G)$ 由那些在 $V\cap U$ 上的限制属于 $\Gamma(V\cap U,\G)$ 的截面 $s\in\Gamma(V,\F)$ 组成。因此 $\overline\G|U=\psi^*(\overline\G)=\G$；$\overline\G$ 是 $\F$ 的在 $U$ 上诱导 $\G$ 的最大子层。称它为 $\F|U$ 的子层 $\G$ 在 $\F$ 中的\emph{典范延拓}。
\egastatement{9.4.2}{命题} 设 $X$ 为预概形，$U$ 为 $X$ 的开子集，且典范嵌入 $j:U\to X$ 为拟紧态射；当 $X$ 的底空间局部诺特时，对每个 $U$ 都如此（\hyperlink{para:1:6.6.4}{6.6.4}，(i)）。则：
\begin{enumerate}[label=(\roman*)]
\item 对每个拟凝聚 $(\OO_X|U)$-模 $\G$，$j_*(\G)$ 为拟凝聚 $\OO_X$-模，且 $j_*(\G)|U=j^*(j_*(\G))=\G$。
\item 对每个拟凝聚 $\OO_X$-模 $\F$，以及 $\F|U$ 的每个拟凝聚子 $(\OO_X|U)$-模 $\G$，其典范延拓 $\overline\G$（\hyperlink{para:1:9.4.1}{9.4.1}）为 $\F$ 的拟凝聚子 $\OO_X$-模。
\end{enumerate}
若 $j=(\psi,\theta)$，其中 $\psi$ 为底空间的嵌入 $U\to X$，则对每个 $(\OO_X|U)$-模 $\G$，按定义 $j_*(\G)=\psi_*(\G)$。此外，对每个 $\OO_X$-模 $\HH$，由开子集上诱导预概形的定义，有 $j^*(\HH)=\psi^*(\HH)=\HH|U$。因此 (i) 为（\hyperlink{para:1:9.2.2}{9.2.2}，a)）的特例。出于同一理由，$j_*(j^*(\F))$ 拟凝聚；而 $\overline\G$ 为 $j_*(\G)$ 在同态 $\rho:\F\to j_*(j^*(\F))$ 下的原像，故 (ii) 由（\hyperlink{para:1:4.1.1}{4.1.1}）得出。

当 $U$ 拟紧且 $X$ 为概形时，嵌入 $j:U\to X$ 拟紧的假设也成立：此时 $U$ 为有限个仿射开集 $U_i$ 之并，而对 $X$ 的每个仿射开集 $V$，$V\cap U_i$ 仿射（\hyperlink{para:1:5.5.6}{5.5.6}），因而拟紧。
\egastatement{9.4.3}{推论} 设 $X$ 为预概形，$U$ 为 $X$ 的拟紧开集，且嵌入态射 $j:U\to X$ 拟紧。再假设每个拟凝聚 $\OO_X$-模都是其有限型拟凝聚子 $\OO_X$-模的归纳极限；当 $X$ 仿射时，此条件成立。设 $\F$ 为拟凝聚 $\OO_X$-模，$\G$ 为 $\F|U$ 的有限型拟凝聚子 $(\OO_X|U)$-模。那么，存在 $\F$ 的有限型拟凝聚子 $\OO_X$-模 $\G'$，使 $\G'|U=\G$。

事实上，$\G=\overline\G|U$，而由（\hyperlink{para:1:9.4.2}{9.4.2}），$\overline\G$ 拟凝聚，故为其有限型拟凝聚子 $\OO_X$-模 $\HH_\lambda$ 的归纳极限。于是 $\G$ 为各 $\HH_\lambda|U$ 的归纳极限；因其为有限型，$U$ 又拟紧，故它等于其中某个 $\HH_\lambda|U$（0，\hyperlink{para:0:5.2.3}{5.2.3}）。
\egastatement{9.4.4}{注} 假设对每个仿射开集 $U\subset X$，嵌入态射 $U\to X$ 都拟紧。若（\hyperlink{para:1:9.4.3}{9.4.3}）的结论对每个仿射开集 $U$ 及 $\F|U$ 的每个有限型拟凝聚子 $(\OO_X|U)$-模 $\G$ 都成立，则 $\F$ 是其有限型拟凝聚子 $\OO_X$-模的归纳极限。

事实上，对每个仿射开集 $U\subset X$，有 $\F|U=\wt M$，其中 $M$ 为 $A(U)$-模。$M$ 为其有限型子模的归纳极限，因此 $\F|U$ 为其有限型拟凝聚子 $(\OO_X|U)$-模的归纳极限（\hyperlink{para:1:1.3.9}{1.3.9}）。按假设，其中每个子模层都由 $\F$ 的某个有限型拟凝聚子 $\OO_X$-模 $\G_{\lambda,U}$ 在 $U$ 上诱导。这些 $\G_{\lambda,U}$ 的有限和仍为有限型拟凝聚 $\OO_X$-模：问题是局部的，仿射情形已在（\hyperlink{para:1:1.3.10}{1.3.10}）中处理。显然 $\F$ 是这些有限和的归纳极限，断言得证。
\egastatement{9.4.5}{推论} 在（\hyperlink{para:1:9.4.3}{9.4.3}）的假设下，对每个有限型拟凝聚 $(\OO_X|U)$-模 $\G$，存在有限型拟凝聚 $\OO_X$-模 $\G'$，使 $\G'|U=\G$。

因 $\F=j_*(\G)$ 拟凝聚（\hyperlink{para:1:9.4.2}{9.4.2}），且 $\F|U=\G$，只须将（\hyperlink{para:1:9.4.3}{9.4.3}）用于 $\F$。
\egastatement{9.4.6}{引理} 设 $X$ 为预概形，$L$ 为良序集合，$(V_\lambda)_{\lambda\in L}$ 为 $X$ 的仿射开覆盖，$U$ 为 $X$ 的开集；对每个 $\lambda\in L$，令 $W_\lambda=\bigcup_{\mu<\lambda}V_\mu$。假设：$1^\circ$ 对每个 $\lambda$，$V_\lambda\cap W_\lambda$ 拟紧；$2^\circ$ 浸入态射 $U\to X$ 拟紧。那么，对每个拟凝聚 $\OO_X$-模 $\F$，以及 $\F|U$ 的每个有限型拟凝聚子 $(\OO_X|U)$-模 $\G$，存在 $\F$ 的有限型拟凝聚子 $\OO_X$-模 $\G'$，使 $\G'|U=\G$。

令 $U_\lambda=U\cup W_\lambda$。用超限归纳定义族 $(\G_\lambda')$：$\G_\lambda'$ 是 $\F|U_\lambda$ 的有限型拟凝聚子 $(\OO_X|U_\lambda)$-模，满足 $\G_\lambda'|U_\mu=\G_\mu'$（$\mu<\lambda$）及 $\G_\lambda'|U=\G$。将这些层粘合所得的唯一子 $\OO_X$-模 $\G'$（0，\hyperlink{para:0:3.3.1}{3.3.1}），即满足要求。\transnote{在归纳的初始步取 $U_0=U$、$\G_0'=\G$；若原良序指标集有末项，还须做完该末项对应的后继步骤，再取覆盖全体 $X$ 的粘合。原文省略了这两项端点约定。}

设对 $\mu<\lambda$ 已定义 $\G_\mu'$ 并满足上述性质。若 $\lambda$ 为无前驱的非初始指标，取 $\G_\lambda'$ 为 $\F|U_\lambda$ 的唯一子 $(\OO_X|U_\lambda)$-模，使其在每个 $U_\mu$ 上限制为 $\G_\mu'$；这是允许的，因为此时各 $U_\mu$（$\mu<\lambda$）覆盖 $U_\lambda$。若 $\lambda=\mu+1$，则 $U_\lambda=U_\mu\cup V_\mu$，只须定义 $\F|V_\mu$ 的有限型拟凝聚子 $(\OO_X|V_\mu)$-模 $\G_\mu''$，使
\[
 \G_\mu''|(U_\mu\cap V_\mu)=\G_\mu'|(U_\mu\cap V_\mu).
\]
再取 $\G_\lambda'$ 为在 $U_\mu$ 上限制为 $\G_\mu'$、在 $V_\mu$ 上限制为 $\G_\mu''$ 的子模层（0，\hyperlink{para:0:3.3.1}{3.3.1}）。由于 $V_\mu$ 仿射，只要证明 $U_\mu\cap V_\mu$ 拟紧，$\G_\mu''$ 的存在性就由（\hyperlink{para:1:9.4.3}{9.4.3}）保证。而 $U_\mu\cap V_\mu$ 为 $U\cap V_\mu$ 与 $W_\mu\cap V_\mu$ 之并，按假设二者都拟紧。
\egastatement{9.4.7}{定理} 设 $X$ 为预概形，$U$ 为 $X$ 的开集，并假设下列条件之一成立：
\begin{enumerate}[label=\alph*)]
\item $X$ 的底空间局部诺特；
\item $X$ 为拟紧概形，且 $U$ 为拟紧开集。
\end{enumerate}
对每个拟凝聚 $\OO_X$-模 $\F$，以及 $\F|U$ 的每个有限型拟凝聚子 $(\OO_X|U)$-模 $\G$，存在 $\F$ 的有限型拟凝聚子 $\OO_X$-模 $\G'$，使 $\G'|U=\G$。

取 $X$ 的仿射开覆盖 $(V_\lambda)_{\lambda\in L}$，在 b) 中令 $L$ 有限。赋予 $L$ 一个良序结构，只须验证引理（\hyperlink{para:1:9.4.6}{9.4.6}）的条件。在 a) 中这显然成立，因为各 $V_\lambda$ 为诺特空间。在 b) 中，各 $V_\lambda\cap V_\mu$ 仿射（\hyperlink{para:1:5.5.6}{5.5.6}），因而拟紧；由于 $L$ 有限，$V_\lambda\cap W_\lambda$ 拟紧。定理得证。
\egastatement{9.4.8}{推论} 在（\hyperlink{para:1:9.4.7}{9.4.7}）的假设下，对每个有限型拟凝聚 $(\OO_X|U)$-模 $\G$，存在有限型拟凝聚 $\OO_X$-模 $\G'$，使 $\G'|U=\G$。

将（\hyperlink{para:1:9.4.7}{9.4.7}）用于 $\F=j_*(\G)$ 即可；由（\hyperlink{para:1:9.4.2}{9.4.2}），它拟凝聚，且 $\F|U=\G$。
\egastatement{9.4.9}{推论} 设 $X$ 为底空间局部诺特的预概形，或为拟紧概形。那么，每个拟凝聚 $\OO_X$-模都是其有限型拟凝聚子 $\OO_X$-模的归纳极限。

这由（\hyperlink{para:1:9.4.7}{9.4.7}）与注（\hyperlink{para:1:9.4.4}{9.4.4}）得出。
\egastatement{9.4.10}{推论} 在（\hyperlink{para:1:9.4.9}{9.4.9}）的假设下，若拟凝聚 $\OO_X$-模 $\F$ 的每个有限型拟凝聚子 $\OO_X$-模都由其在 $X$ 上的截面生成，则 $\F$ 也由其在 $X$ 上的截面生成。

取点 $x\in X$ 的仿射开邻域 $U$，并取 $\F$ 在 $U$ 上的一个截面 $s$。由 $s$ 生成的 $\F|U$ 的子模层 $\G$ 拟凝聚且为有限型，因此存在 $\F$ 的有限型拟凝聚子 $\OO_X$-模 $\G'$，使 $\G'|U=\G$（\hyperlink{para:1:9.4.7}{9.4.7}）。按假设，存在 $\G'$ 在 $X$ 上的有限个截面 $t_i$，以及 $\OO_X$ 在 $x$ 的某个邻域 $V\subset U$ 上的截面 $a_i$，使 $s|V=\sum_i a_i(t_i|V)$。推论得证。


% END chapters/ch1_s9_4.tex

\par

% BEGIN chapters/ch1_s9_5.tex
\subsection{预概形的闭像；子预概形的闭包}\hypertarget{heading:1:9.5}{}
\src{176—179}
\egastatement{9.5.1}{命题} 设 $f:X\to Y$ 为预概形态射，使 $f_*(\OO_X)$ 为拟凝聚 $\OO_Y$-模；当 $f$ 拟紧，且 $f$ 分离或 $X$ 局部诺特时即如此（\hyperlink{para:1:9.2.2}{9.2.2}）。那么，存在 $Y$ 的最小闭子预概形 $Y'$，使典范嵌入 $j:Y'\to Y$ 不小于 $f$；等价地，使 $X$ 的子预概形 $f^{-1}(Y')$ 与 $X$ 相同（\hyperlink{para:1:4.4.1}{4.4.1}）。\transnote{原文陈述漏写“闭”（\textit{fermé}）；以下证明、推论（\hyperlink{para:1:9.5.2}{9.5.2}）及定义（\hyperlink{para:1:9.5.3}{9.5.3}）均在闭子预概形中讨论最小性。这里补明“闭”，而不把结论解释为对所有子预概形的最小性。}

更精确地：
\egastatement{9.5.2}{推论} 在（\hyperlink{para:1:9.5.1}{9.5.1}）的条件下，设 $f=(\psi,\theta)$，$\J$ 为同态 $\theta:\OO_Y\to f_*(\OO_X)$ 的拟凝聚核。那么，由 $\J$ 定义的 $Y$ 的闭子预概形 $Y'$ 满足（\hyperlink{para:1:9.5.1}{9.5.1}）的条件。

由于函子 $\psi^*$ 正合，典范分解
\[
 \theta:\OO_Y\longrightarrow\OO_Y/\J
 \xrightarrow{\theta'}\psi_*(\OO_X)
\]
给出分解（0，\hyperref[eq:0:3.5.4.3]{3.5.4.3}）
\[
 \theta^\sharp:\psi^*(\OO_Y)
 \longrightarrow\psi^*(\OO_Y)/\psi^*(\J)
 \xrightarrow{\theta'^\sharp}\OO_X.
\]
对每个 $x\in X$，$\theta_x^\sharp$ 为局部同态，故 $\theta_x'^\sharp$ 也为局部同态。将 $\psi$ 视作从 $X$ 到 $Y'$ 的连续映射，记作 $\psi_0$；将限制
\[
 \theta_0=\theta'|Y':(\OO_Y/\J)|Y'
 \longrightarrow\psi_*(\OO_X)|Y'=(\psi_0)_*(\OO_X)
\]
与之配对，便得到预概形态射 $f_0=(\psi_0,\theta_0):X\to Y'$（\hyperlink{para:1:2.2.1}{2.2.1}），满足 $f=j\circ f_0$。\transnote{原文在这段证明中混用了 $X'$、$Y'$，并在下一段使用 $X''$。译文统一用 $Y'$、$Y''$ 表示 $Y$ 的子预概形；不改变它们的角色。}

若 $Y''$ 为 $Y$ 的另一闭子预概形，由拟凝聚理想层 $\J'$ 定义，且嵌入 $j':Y''\to Y$ 不小于 $f$，则首先 $Y''\supset\psi(X)$，因 $Y''$ 闭，故包含 $Y'$ 的底空间。\transnote{此处底空间包含关系使用了 $|Y'|=\overline{\psi(X)}$，该事实紧接着在（\hyperlink{para:1:9.5.4}{9.5.4}）独立证明；也可先用同态分解所得的 $\J'\subset\J$ 直接推出包含关系。}
此外，对每个 $y\in Y''$，$\theta_y$ 必分解为
\[
 \OO_y\longrightarrow\OO_y/\J'_y
 \longrightarrow(\psi_*(\OO_X))_y.
\]
按定义这蕴含 $\J'_y\subset\J_y$，因此 $Y'$ 是 $Y''$ 的闭子预概形（\hyperlink{para:1:4.1.10}{4.1.10}）。
\egastatement{9.5.3}{定义} 若存在 $Y$ 的最小闭子预概形 $Y'$，使典范嵌入 $j:Y'\to Y$ 不小于 $f$，就称 $Y'$ 为 $X$ 经态射 $f$ 所得的\emph{闭像预概形}。
\egastatement{9.5.4}{命题} 若 $f_*(\OO_X)$ 为拟凝聚 $\OO_Y$-模，则 $X$ 经 $f$ 所得闭像的底空间，是 $f(X)$ 在 $Y$ 中的闭包 $\overline{f(X)}$。

由于 $f_*(\OO_X)$ 的支集包含于 $\overline{f(X)}$，沿用（\hyperlink{para:1:9.5.2}{9.5.2}）的记号，对 $y\notin\overline{f(X)}$，有 $\J_y=\OO_y$；因此 $\OO_Y/\J$ 的支集包含于 $\overline{f(X)}$。此外，该支集为闭集，且包含 $f(X)$。事实上，若 $y\in f(X)$，则环 $(\psi_*(\OO_X))_y$ 的单位元不为零，它是截面
\[
 1\in\Gamma(X,\OO_X)=\Gamma(Y,\psi_*(\OO_X))
\]
在 $y$ 处的芽。它是 $\OO_y$ 的单位元经 $\theta_y$ 所得的像，故后者不属于 $\J_y$，从而 $\OO_y/\J_y\ne0$。证明完毕。
\egastatement{9.5.5}{命题（闭像的传递性）} 设 $f:X\to Y$、$g:Y\to Z$ 为两个预概形态射。假设 $X$ 经 $f$ 的闭像 $Y'$ 存在；令 $g'$ 为 $g$ 在 $Y'$ 上的限制，并假设 $Y'$ 经 $g'$ 的闭像 $Z'$ 存在。那么，$X$ 经 $g\circ f$ 的闭像存在，且等于 $Z'$。

只须证明（\hyperlink{para:1:9.5.1}{9.5.1}），$Z'$ 是 $Z$ 的满足以下条件的最小闭子预概形 $Z_1$：$X$ 的闭子预概形 $(g\circ f)^{-1}(Z_1)=f^{-1}(g^{-1}(Z_1))$（\hyperlink{para:1:4.4.2}{4.4.2}）等于 $X$。等价地，须证明 $Z'$ 是使嵌入 $g^{-1}(Z_1)\to Y$ 不小于 $f$ 的最小闭子预概形（\hyperlink{para:1:4.4.1}{4.4.1}）。由闭像 $Y'$ 的存在性，$Z_1$ 具有此性质，当且仅当 $g^{-1}(Z_1)$ 不小于 $Y'$，亦即令 $j:Y'\to Y$ 为嵌入后，有
\[
 j^{-1}(g^{-1}(Z_1))=g'^{-1}(Z_1)=Y'.
\]
按 $Z'$ 的定义，$Z'$ 正是满足此条件的最小闭子预概形。
\egastatement{9.5.6}{推论} 设 $f:X\to Y$ 为 $S$-态射，且 $Y$ 为 $X$ 经 $f$ 的闭像。设 $Z$ 为 $S$-概形。若两个 $S$-态射 $g_1,g_2:Y\to Z$ 满足 $g_1\circ f=g_2\circ f$，则 $g_1=g_2$。

令 $h=(g_1,g_2)_S:Y\to Z\times_S Z$。因对角线 $T=\Delta_Z(Z)$ 为 $Z\times_S Z$ 的闭子预概形，$Y'=h^{-1}(T)$ 为 $Y$ 的闭子预概形（\hyperlink{para:1:4.4.1}{4.4.1}）。令 $u=g_1\circ f=g_2\circ f$，则按积的定义 $h'=h\circ f=(u,u)_S=\Delta_Z\circ u$。因 $\Delta_Z^{-1}(T)=Z$，有 $h'^{-1}(T)=u^{-1}(Z)=X$，故 $f^{-1}(Y')=X$。由（\hyperlink{para:1:4.4.1}{4.4.1}），典范嵌入 $Y'\to Y$ 不小于 $f$，所以按假设 $Y'=Y$。因此 $h$ 分解为 $\Delta_Z\circ v$，其中 $v:Y\to Z$ 为态射（\hyperlink{para:1:4.4.1}{4.4.1}），从而 $g_1=g_2=v$。
\egastatement{9.5.7}{注} 若 $X,Y$ 为 $S$-概形，（\hyperlink{para:1:9.5.6}{9.5.6}）的含义是：当 $Y$ 为 $X$ 经 $f$ 的闭像时，$f$ 是 $S$-概形范畴中的满态射（T，1.1）。我们将在第 V 章证明其逆命题：若闭像 $Y'$ 存在，且 $f$ 为 $S$-概形的满态射，则必有 $Y'=Y$。
\egastatement{9.5.8}{命题} 假设满足（\hyperlink{para:1:9.5.1}{9.5.1}）的假设，$Y'$ 为 $X$ 经 $f$ 的闭像。对 $Y$ 的每个开集 $V$，令 $f_V:f^{-1}(V)\to V$ 为 $f$ 的限制。那么，$f^{-1}(V)$ 经 $f_V$ 在 $V$ 中的闭像存在，且等于 $Y'$ 在其开集 $V\cap Y'$ 上诱导的预概形；换言之，等于 $Y$ 的子预概形 $\inf(V,Y')$（\hyperlink{para:1:4.4.3}{4.4.3}）。

令 $X'=f^{-1}(V)$。由于 $\OO_{X'}$ 经 $f_V$ 的直像就是 $f_*(\OO_X)$ 在 $V$ 上的限制，同态 $\OO_V\to(f_V)_*(\OO_{X'})$ 的核 $\J'$ 显然是 $\J$ 在 $V$ 上的限制，命题随即得出。

此结果可解释为：闭像的形成与作为开浸入的底预概形扩张 $Y_1\to Y$ 可交换。第 IV 章将说明，若 $f$ 分离且拟紧，则对作为平坦态射的扩张 $Y_1\to Y$ 也有相同结论。
\egastatement{9.5.9}{命题} 设态射 $f:X\to Y$ 的闭像 $Y'$ 存在。
\begin{enumerate}[label=(\roman*)]
\item 若 $X$ 约化，则 $Y'$ 也约化。
\item 若满足（\hyperlink{para:1:9.5.1}{9.5.1}）的假设，且 $X$ 不可约（分别为整），则 $Y'$ 也不可约（分别为整）。
\end{enumerate}
按假设，$f$ 分解为 $X\xrightarrow{g}Y'\xrightarrow{j}Y$，其中 $j$ 为典范嵌入。若 $X$ 约化，则 $g$ 分解为 $X\to Y'_{\mathrm{red}}\xrightarrow{j'}Y'$，其中 $j'$ 为典范嵌入（\hyperlink{para:1:5.2.2}{5.2.2}）；由 $Y'$ 的定义，得到 $Y'_{\mathrm{red}}=Y'$。若还满足（\hyperlink{para:1:9.5.1}{9.5.1}）的条件，则由（\hyperlink{para:1:9.5.4}{9.5.4}），$f(X)$ 在 $Y'$ 中稠密；若 $X$ 不可约，则 $Y'$ 也不可约（0，\hyperlink{para:0:2.1.5}{2.1.5}）。关于整预概形的断言是前两项的合取。
\egastatement{9.5.10}{命题} 设 $Y$ 为预概形 $X$ 的子预概形，且典范嵌入 $j:Y\to X$ 拟紧。那么，存在 $X$ 的不小于 $Y$ 的最小闭子预概形 $\overline Y$；其底空间为 $Y$ 的底空间的闭包。$Y$ 的底空间在该闭包中开，且预概形 $Y$ 由 $\overline Y$ 在此开集上诱导。

将（\hyperlink{para:1:9.5.1}{9.5.1}）用于嵌入 $j$ 即可：它分离（\hyperlink{para:1:5.5.1}{5.5.1}），并按假设拟紧。因此（\hyperlink{para:1:9.5.1}{9.5.1}）给出 $\overline Y$ 的存在性，（\hyperlink{para:1:9.5.4}{9.5.4}）说明其底空间为 $Y$ 在 $X$ 中的闭包。因 $Y$ 在 $X$ 中局部闭，它在 $\overline Y$ 中开；最后一项断言由（\hyperlink{para:1:9.5.8}{9.5.8}）用于 $X$ 的一个使 $Y$ 在其中闭的开集 $V$ 得出。

沿用这些记号，若嵌入 $V\to X$ 拟紧，$\J$ 为 $\OO_X|V$ 中定义 $V$ 的闭子预概形 $Y$ 的拟凝聚理想层，则由（\hyperlink{para:1:9.5.1}{9.5.1}），定义 $\overline Y$ 的 $\OO_X$ 的拟凝聚理想层为 $\J$ 的典范延拓 $\overline\J$（\hyperlink{para:1:9.4.1}{9.4.1}），因为它显然是 $\OO_X$ 的在 $V$ 上诱导 $\J$ 的最大拟凝聚理想子层。
\egastatement{9.5.11}{推论} 在（\hyperlink{para:1:9.5.10}{9.5.10}）的假设下，$\OO_{\overline Y}$ 在 $\overline Y$ 的开集 $V$ 上的任一截面，若在 $V\cap Y$ 中为零，则该截面为零。

由（\hyperlink{para:1:9.5.8}{9.5.8}），可归结到 $V=\overline Y$ 的情形。$\OO_{\overline Y}$ 在 $\overline Y$ 上的截面，典范地对应于 $\overline Y\otimes_{\mathbf Z}\mathbf Z[T]$ 的 $\overline Y$-截面（\hyperlink{para:1:3.3.15}{3.3.15}），而后一个预概形在 $\overline Y$ 上分离，故本推论是（\hyperlink{para:1:9.5.6}{9.5.6}）的特例。

若存在 $X$ 的不小于子预概形 $Y$ 的最小闭子预概形 $\overline Y$，则在不致混淆时，称 $\overline Y$ 为 $Y$ 在 $X$ 中的\emph{闭包}。


% END chapters/ch1_s9_5.tex

\par

% BEGIN chapters/ch1_s9_6.tex
\subsection{拟凝聚代数层；结构层的变换}\hypertarget{heading:1:9.6}{}
\src{179—180}
\egastatement{9.6.1}{命题} 设 $X$ 为预概形，$\BB$ 为拟凝聚 $\OO_X$-代数（0，\hyperlink{para:0:5.1.3}{5.1.3}）。一个 $\BB$-模 $\F$ 在赋环空间 $(X,\BB)$ 上拟凝聚，当且仅当 $\F$ 作为 $\OO_X$-模拟凝聚。

问题是局部的，可假设 $X$ 为环 $A$ 的仿射概形；此时 $\BB=\wt B$，其中 $B$ 为 $A$-代数（\hyperlink{para:1:1.4.3}{1.4.3}）。若 $\F$ 在 $(X,\BB)$ 上拟凝聚，还可假设它是某个 $\BB$-同态 $\BB^{(I)}\to\BB^{(J)}$ 的余核。此同态也是 $\OO_X$-模间的 $\OO_X$-同态，而 $\BB^{(I)},\BB^{(J)}$ 为拟凝聚 $\OO_X$-模（\hyperlink{para:1:1.3.9}{1.3.9}，(ii)），故 $\F$ 也是拟凝聚 $\OO_X$-模（\hyperlink{para:1:1.3.9}{1.3.9}，(i)）。

反之，若 $\F$ 为拟凝聚 $\OO_X$-模，则 $\F=\wt M$，其中 $M$ 为 $B$-模（\hyperlink{para:1:1.4.3}{1.4.3}）。$M$ 同构于某个 $B$-同态 $B^{(I)}\to B^{(J)}$ 的余核，故 $\F$ 作为 $\BB$-模，同构于相应同态 $\BB^{(I)}\to\BB^{(J)}$ 的余核（\hyperlink{para:1:1.3.13}{1.3.13}）。证明完毕。

特别地，若 $\F,\G$ 为两个拟凝聚 $\BB$-模，则 $\F\otimes_{\BB}\G$ 为拟凝聚 $\BB$-模；若 $\F$ 还有有限表示，则 $\shHom_{\BB}(\F,\G)$ 也拟凝聚（\hyperlink{para:1:1.3.13}{1.3.13}）。
\egapar{9.6.2} 给定预概形 $X$，若对每个 $x\in X$，存在一个仿射开邻域 $U$，使 $B=\Gamma(U,\BB)$ 是 $A=\Gamma(U,\OO_X)$ 上的有限型代数，就称拟凝聚 $\OO_X$-代数 $\BB$ 为\emph{有限型}。此时 $\BB|U=\wt B$，而对每个 $f\in A$，诱导的 $(\OO_X|D(f))$-代数 $\BB|D(f)$ 为有限型：它同构于 $\widetilde{B_f}$，而 $B_f=B\otimes_A A_f$ 显然是有限型 $A_f$-代数。因各 $D(f)$ 构成 $U$ 的拓扑基，若 $\BB$ 为有限型拟凝聚 $\OO_X$-代数，则对 $X$ 的任一开集 $V$，$\BB|V$ 为有限型拟凝聚 $(\OO_X|V)$-代数。
\egastatement{9.6.3}{命题} 设 $X$ 为局部诺特预概形。则每个有限型拟凝聚 $\OO_X$-代数 $\BB$ 都是凝聚环层（0，\hyperlink{para:0:5.3.7}{5.3.7}）。

仍可限于 $X$ 为诺特环 $A$ 的仿射概形、$\BB=\wt B$ 且 $B$ 为有限型 $A$-代数的情形；于是 $B$ 是诺特环。须证明某个 $\BB$-同态 $\BB^m\to\BB$ 的核 $\N$ 为有限型 $\BB$-模。$\N$ 作为 $\BB$-模同构于 $\wt N$，其中 $N$ 为相应的 $B$-模同态 $B^m\to B$ 的核（\hyperlink{para:1:1.3.13}{1.3.13}）。由于 $B$ 诺特，$B^m$ 的子模 $N$ 为有限型 $B$-模，故存在像为 $N$ 的同态 $B^p\to B^m$。序列 $B^p\to B^m\to B$ 正合，因此相应序列 $\BB^p\to\BB^m\to\BB$ 正合（\hyperlink{para:1:1.3.5}{1.3.5}）。$\N$ 为 $\BB^p\to\BB^m$ 的像（\hyperlink{para:1:1.3.9}{1.3.9}，(i)），命题得证。
\egastatement{9.6.4}{推论} 在（\hyperlink{para:1:9.6.3}{9.6.3}）的假设下，$\BB$-模 $\F$ 凝聚，当且仅当它作为 $\OO_X$-模拟凝聚，且作为 $\BB$-模为有限型。若这些条件成立，且 $\G$ 为 $\F$ 的子 $\BB$-模或商 $\BB$-模，则 $\G$ 为凝聚 $\BB$-模，当且仅当 $\G$ 为拟凝聚 $\OO_X$-模。

由（\hyperlink{para:1:9.6.1}{9.6.1}），对 $\F$ 的条件显然必要。证明充分性。可限于以下情形：$X$ 为诺特环 $A$ 的仿射概形，$\BB=\wt B$，$B$ 为有限型 $A$-代数，$\F=\wt M$，$M$ 为 $B$-模，且存在满射 $\BB^m\to\F\to0$。此时有相应正合列 $B^m\to M\to0$，故 $M$ 为有限型 $B$-模；又因 $B$ 诺特，此同态的核 $P$ 为有限型 $B$-模。因此 $\F$ 为某个 $\BB$-同态 $\BB^p\to\BB^m$ 的余核（\hyperlink{para:1:1.3.13}{1.3.13}），因 $\BB$ 为凝聚环层，$\F$ 凝聚（0，\hyperlink{para:0:5.3.4}{5.3.4}）。同样的论证说明，$\F$ 的拟凝聚子 $\BB$-模（分别为商 $\BB$-模）是有限型，从而得到推论的第二部分。
\egastatement{9.6.5}{命题} 设 $X$ 为拟紧概形，或为底空间诺特的预概形。对每个有限型拟凝聚 $\OO_X$-代数 $\BB$，存在 $\BB$ 的有限型拟凝聚子 $\OO_X$-模 $\E$，使 $\E$ 生成 $\OO_X$-代数 $\BB$。

按假设，存在 $X$ 的有限仿射开覆盖 $(U_i)$，使 $B_i=\Gamma(U_i,\BB)$ 为 $A_i=\Gamma(U_i,\OO_X)$ 上的有限型代数。取生成 $A_i$-代数 $B_i$ 的有限型子 $A_i$-模 $E_i$。由（\hyperlink{para:1:9.4.7}{9.4.7}），存在 $\BB$ 的有限型拟凝聚子 $\OO_X$-模 $\E_i$，使 $\E_i|U_i=\wt E_i$。各 $\E_i$ 的和 $\E$ 显然满足要求。
\egastatement{9.6.6}{命题} 设 $X$ 为底空间局部诺特的预概形，或为拟紧概形。那么，每个拟凝聚 $\OO_X$-代数 $\BB$ 都是其有限型拟凝聚子 $\OO_X$-代数的归纳极限。

由（\hyperlink{para:1:9.4.9}{9.4.9}），$\BB$ 作为 $\OO_X$-模，是其有限型拟凝聚子 $\OO_X$-模的归纳极限。这些子模生成 $\BB$ 的有限型拟凝聚子 $\OO_X$-代数（\hyperlink{para:1:1.3.14}{1.3.14}），而 $\BB$ 当然也是这些子代数的归纳极限。


% END chapters/ch1_s9_6.tex

\par

% BEGIN chapters/ch1_s10_1.tex
\section{形式概形}\hypertarget{heading:1:10}{}
\subsection{仿射形式概形}\hypertarget{heading:1:10.1}{}
\src{180—182}
\egapar{10.1.1} 设 $A$ 为可容许拓扑环（0，\hyperlink{para:0:7.1.2}{7.1.2}）。对 $A$ 的每个定义理想 $J$，$\Spec(A/J)$ 等同于 $\Spec(A)$ 的闭子空间 $V(J)$（\hyperlink{para:1:1.1.11}{1.1.11}），即 $A$ 的开素理想组成的集合。这个拓扑空间不依赖于所取的定义理想 $J$，将它记作 $\frX$。取由定义理想组成的 $A$ 中零点的一个基本邻域系 $(J_\lambda)$。对每个 $\lambda$，令 $\OO_\lambda$ 为 $\Spec(A/J_\lambda)$ 的结构层；它由 $\widetilde{A/J_\lambda}$ 在 $\frX$ 上诱导，后者在 $\frX$ 之外为零。对 $J_\lambda\subset J_\mu$，典范同态 $A/J_\lambda\to A/J_\mu$ 定义环层同态 $u_{\mu\lambda}:\OO_\lambda\to\OO_\mu$（\hyperlink{para:1:1.6.1}{1.6.1}）；各 $\OO_\lambda$ 对这些同态构成一个射影系。

由于 $\frX$ 的拓扑有由拟紧开集组成的基，可为每个 $\OO_\lambda$ 配备相伴的伪离散拓扑环层结构（0，\hyperlink{para:0:3.8.1}{3.8.1}），其底层环层仍为 $\OO_\lambda$，并仍用 $\OO_\lambda$ 表示。它们仍构成拓扑环层的射影系（0，\hyperlink{para:0:3.8.2}{3.8.2}）。将该射影系的极限记作 $\OO_{\frX}$；这是 $\frX$ 上的拓扑环层。对每个拟紧开集 $U\subset\frX$，$\Gamma(U,\OO_{\frX})$ 因而为离散环系 $\Gamma(U,\OO_\lambda)$ 的拓扑射影极限（0，\hyperlink{para:0:3.2.6}{3.2.6}）。
\egastatement{10.1.2}{定义} 给定可容许拓扑环 $A$，称 $\Spec(A)$ 的由 $A$ 的开素理想组成的闭子空间 $\frX$ 为 $A$ 的\emph{形式谱}，记作 $\Spf(A)$。若一个拓扑赋环空间同构于形式谱 $\frX=\Spf(A)$，其拓扑环层 $\OO_{\frX}$ 为伪离散环层 $(\widetilde{A/J_\lambda})|\frX$ 的射影极限，其中 $J_\lambda$ 遍历 $A$ 的定义理想构成的有向集合，就称该拓扑赋环空间为\emph{仿射形式概形}。

谈到形式谱 $\frX=\Spf(A)$ 作为仿射形式概形时，总是指按上述方式定义的拓扑赋环空间 $(\frX,\OO_{\frX})$。

每个仿射概形 $X=\Spec(A)$ 都可通过将 $A$ 赋予离散拓扑，以唯一的方式视作仿射形式概形：此时对拟紧开集 $U$，拓扑环 $\Gamma(U,\OO_X)$ 是离散的；但对 $X$ 的一般开集，通常并非如此。
\egastatement{10.1.3}{命题} 若 $\frX=\Spf(A)$，其中 $A$ 为可容许环，则 $\Gamma(\frX,\OO_{\frX})$ 与 $A$ 拓扑同构。

因 $\frX$ 在 $\Spec(A)$ 中闭，它拟紧，所以 $\Gamma(\frX,\OO_{\frX})$ 拓扑同构于离散环 $\Gamma(\frX,\OO_\lambda)$ 的射影极限。而 $\Gamma(\frX,\OO_\lambda)$ 同构于 $A/J_\lambda$（\hyperlink{para:1:1.3.7}{1.3.7}）；由于 $A$ 分离且完备，它拓扑同构于 $\plim A/J_\lambda$（0，\hyperlink{para:0:7.2.1}{7.2.1}），命题得证。
\egastatement{10.1.4}{命题} 设 $A$ 为可容许环，$\frX=\Spf(A)$。对每个 $f\in A$，令 $\frD(f)=D(f)\cap\frX$；则拓扑赋环空间 $(\frD(f),\OO_{\frX}|\frD(f))$ 同构于仿射形式谱 $\Spf(A_{\{f\}})$（0，\hyperlink{para:0:7.6.15}{7.6.15}）。

对 $A$ 的每个定义理想 $J$，离散环 $S_f^{-1}A/S_f^{-1}J$ 典范地等同于 $A_{\{f\}}/J_{\{f\}}$（0，\hyperlink{para:0:7.6.9}{7.6.9}），所以拓扑空间 $\Spf(A_{\{f\}})$ 典范地等同于 $\frD(f)$（\hyperlink{para:1:1.2.5}{1.2.5} 与 \hyperlink{para:1:1.2.6}{1.2.6}）。此外，对 $\frX$ 的每个包含于 $\frD(f)$ 的拟紧开集 $U$，$\Gamma(U,\OO_\lambda)$ 等同于 $\Spec(S_f^{-1}A/S_f^{-1}J_\lambda)$ 的结构层在 $U$ 上的截面模（\hyperlink{para:1:1.3.6}{1.3.6}）。因此，令 $\frY=\Spf(A_{\{f\}})$，便有 $\Gamma(U,\OO_{\frX})$ 与 $\Gamma(U,\OO_{\frY})$ 的典范等同，命题得证。
\egapar{10.1.5} 将 $\Spf(A)$ 的结构层 $\OO_{\frX}$ 只看作不带拓扑的环层时，它在每个 $x\in\frX$ 处的茎，由（\hyperlink{para:1:10.1.4}{10.1.4}）等同于各 $A_{\{f\}}$（$f\notin\jj_x$）的归纳极限。因此，由（0，\hyperlink{para:0:7.6.17}{7.6.17} 与 \hyperlink{para:0:7.6.18}{7.6.18}）得：
\egastatement{10.1.6}{命题} 对每个 $x\in\frX=\Spf(A)$，茎 $\OO_x$ 为局部环，其剩余域同构于 $k(x)=A_x/\jj_xA_x$。若 $A$ 还为诺特进制环，则 $\OO_x$ 为诺特环。

由于 $k(x)$ 不为零，由此得出环层 $\OO_{\frX}$ 的支集就是 $\frX$。


% END chapters/ch1_s10_1.tex

\par

% BEGIN chapters/ch1_s10_2.tex
\subsection{仿射形式概形的态射}\hypertarget{heading:1:10.2}{}
\src{182—183}
\egapar{10.2.1} 设 $A,B$ 为两个可容许环，$\varphi:B\to A$ 为连续同态。连续映射 ${}^a\varphi:\Spec(A)\to\Spec(B)$（\hyperlink{para:1:1.2.1}{1.2.1}）将 $\frX=\Spf(A)$ 映入 $\frY=\Spf(B)$，因为 $A$ 的开素理想在 $\varphi$ 下的原像是 $B$ 的开素理想。另一方面，对每个 $g\in B$，$\varphi$ 定义一个连续同态
\[
 \Gamma(\frD(g),\OO_{\frY})
 \longrightarrow\Gamma(\frD(\varphi(g)),\OO_{\frX})
\]
（\hyperlink{para:1:10.1.4}{10.1.4}、\hyperlink{para:1:10.1.3}{10.1.3} 与 0，\hyperlink{para:0:7.6.7}{7.6.7}）。这些同态满足由 $g$ 到其倍数所对应的限制相容条件，且 $\frD(\varphi(g))={}^a\varphi^{-1}(\frD(g))$，故它们定义了拓扑环层的连续同态 $\OO_{\frY}\to{}^a\varphi_*(\OO_{\frX})$（0，\hyperlink{para:0:3.2.5}{3.2.5}），仍记作 $\wt\varphi$。于是得到拓扑赋环空间态射 $\Phi=({}^a\varphi,\wt\varphi):\frX\to\frY$。

作为不带拓扑的环层同态，$\wt\varphi$ 对每个 $x\in\frX$ 都定义茎上的同态 $\wt\varphi_x^\sharp:\OO_{{}^a\varphi(x)}\to\OO_x$。
\egastatement{10.2.2}{命题} 设 $A,B$ 为两个可容许拓扑环，$\frX=\Spf(A)$、$\frY=\Spf(B)$。拓扑赋环空间态射 $u=(\psi,\theta):\frX\to\frY$ 形如 $({}^a\varphi,\wt\varphi)$，其中 $\varphi:B\to A$ 为连续环同态，当且仅当对每个 $x\in\frX$，$\theta_x^\sharp:\OO_{\psi(x)}\to\OO_x$ 为局部同态。

必要性：令 $\mathfrak p=\jj_x\in\Spf(A)$，$\mathfrak q=\varphi^{-1}(\mathfrak p)$。若 $g\notin\mathfrak q$，则 $\varphi(g)\notin\mathfrak p$；由 $\varphi$ 导出的同态 $B_{\{g\}}\to A_{\{\varphi(g)\}}$（0，\hyperlink{para:0:7.6.7}{7.6.7}）显然将 $\mathfrak q_{\{g\}}$ 映入 $\mathfrak p_{\{\varphi(g)\}}$。取归纳极限，利用（\hyperlink{para:1:10.1.5}{10.1.5}）与（0，\hyperlink{para:0:7.6.17}{7.6.17}），得 $\wt\varphi_x^\sharp$ 为局部同态。

反之，设 $(\psi,\theta)$ 满足条件。由（\hyperlink{para:1:10.1.3}{10.1.3}），$\theta$ 定义连续环同态
\[
 \varphi=\Gamma(\theta):B=\Gamma(\frY,\OO_{\frY})
 \longrightarrow\Gamma(\frX,\OO_{\frX})=A.
\]
由 $\theta$ 的假设，截面 $\varphi(g)$ 在 $x$ 处的芽可逆，当且仅当 $g$ 在 $\psi(x)$ 处的芽可逆。而由（0，\hyperlink{para:0:7.6.17}{7.6.17}），$\OO_{\frX}$（分别为 $\OO_{\frY}$）的整体截面中，在 $x$（分别为 $\psi(x)$）处具有不可逆芽的，正是 $\jj_x$（分别为 $\jj_{\psi(x)}$）的元素。因此 ${}^a\varphi=\psi$。最后，对每个 $g\in B$，下图交换：
\[
\begin{tikzcd}[column sep=large]
 B=\Gamma(\frY,\OO_{\frY})\arrow[r,"\varphi"]\arrow[d]&\Gamma(\frX,\OO_{\frX})=A\arrow[d]\\
 B_{\{g\}}\arrow[r,"\Gamma(\theta_{\frD(g)})"]&A_{\{\varphi(g)\}}.
\end{tikzcd}
\]
下行的环分别是 $\Gamma(\frD(g),\OO_{\frY})$ 与 $\Gamma(\frD(\varphi(g)),\OO_{\frX})$。由完备分式环的泛性质（0，\hyperlink{para:0:7.6.6}{7.6.6}），对每个 $g$ 有 $\theta_{\frD(g)}=\wt\varphi_{\frD(g)}$，故 $\theta=\wt\varphi$（0，\hyperlink{para:0:3.2.5}{3.2.5}）。

称满足（\hyperlink{para:1:10.2.2}{10.2.2}）条件的拓扑赋环空间态射 $(\psi,\theta)$ 为\emph{仿射形式概形态射}。也可以说，函子 $A\mapsto\Spf(A)$ 与 $\frX\mapsto\Gamma(\frX,\OO_{\frX})$ 给出了可容许环范畴与仿射形式概形范畴的对偶范畴之间的等价（T，I，1.2）。
\egapar{10.2.3} 作为（\hyperlink{para:1:10.2.2}{10.2.2}）的特例，对 $f\in A$，$\frX$ 在 $\frD(f)$ 上诱导的仿射形式概形的典范嵌入，对应于典范连续同态 $A\to A_{\{f\}}$。在（\hyperlink{para:1:10.2.2}{10.2.2}）的假设下，取 $h\in B$，并取 $g\in A$ 为 $\varphi(h)$ 的倍数，则 $\psi(\frD(g))\subset\frD(h)$；$u$ 在 $\frD(g)$ 上的限制，作为 $\frD(g)$ 到 $\frD(h)$ 的态射，是使下图交换的唯一态射 $v$：
\[
\begin{tikzcd}
 \frD(g)\arrow[r,"v"]\arrow[d]&\frD(h)\arrow[d]\\
 \frX\arrow[r,"u"]&\frY.
\end{tikzcd}
\]
它对应于使下图交换的唯一连续同态 $\varphi':B_{\{h\}}\to A_{\{g\}}$（0，\hyperlink{para:0:7.6.7}{7.6.7}）：
\[
\begin{tikzcd}
 A\arrow[d]&B\arrow[l,"\varphi"']\arrow[d]\\
 A_{\{g\}}&B_{\{h\}}\arrow[l,"\varphi'"].
\end{tikzcd}
\]


% END chapters/ch1_s10_2.tex

\par

% BEGIN chapters/ch1_s10_3.tex
\subsection{仿射形式概形的定义理想}\hypertarget{heading:1:10.3}{}
\src{183—185}
\egapar{10.3.1} 设 $A$ 为可容许环，$J$ 为 $A$ 的开理想，$\frX=\Spf(A)$。以 $(J_\lambda)$ 表示 $A$ 中包含于 $J$ 的定义理想所组成的集合。则 $\wt J/\wt J_\lambda$ 是 $\wt A/\wt J_\lambda$ 的理想层。将这些层在 $\frX$ 上的限制的射影极限记作 $J^\Delta$；它等同于 $\OO_{\frX}$ 的理想层（0，\hyperlink{para:0:3.2.6}{3.2.6}）。对每个 $f\in A$，$\Gamma(\frD(f),J^\Delta)$ 为各 $S_f^{-1}J/S_f^{-1}J_\lambda$ 的射影极限，即等同于环 $A_{\{f\}}$ 的开理想 $J_{\{f\}}$（0，\hyperlink{para:0:7.6.9}{7.6.9}）；特别地，$\Gamma(\frX,J^\Delta)=J$。由于各 $\frD(f)$ 构成 $\frX$ 的拓扑基，得
\begin{equation}
 J^\Delta|\frD(f)=(J_{\{f\}})^\Delta.
\tag{10.3.1.1}\label{eq:1:10.3.1.1}
\end{equation}
\egapar{10.3.2} 沿用（\hyperlink{para:1:10.3.1}{10.3.1}）的记号。对每个 $f\in A$，从 $A_{\{f\}}=\Gamma(\frD(f),\OO_{\frX})$ 到
\[
 \Gamma(\frD(f),(\widetilde{A/J})|\frX)
 =S_f^{-1}A/S_f^{-1}J
\]
的典范映射为满射，其核为 $\Gamma(\frD(f),J^\Delta)=J_{\{f\}}$（0，\hyperlink{para:0:7.6.9}{7.6.9}）。因此，这些映射定义了从拓扑环层 $\OO_{\frX}$ 到伪离散环层 $(\widetilde{A/J})|\frX$ 的连续满同态，其核为 $J^\Delta$，称为\emph{典范同态}。这个同态也就是（\hyperlink{para:1:10.2.1}{10.2.1}）中的 $\wt\varphi$，其中 $\varphi:A\to A/J$ 是连续商同态。仿射形式概形态射 $({}^a\varphi,\wt\varphi):\Spec(A/J)\to\frX$ 也称为\emph{典范态射}。原文指出，${}^a\varphi$ 是 $\frX$ 到自身的恒等同胚。\transnote{本小节起初只假设 $J$ 是开理想；在此一般性下，$\Spec(A/J)$ 是 $\frX$ 的闭子空间，不一定等于 $\frX$。这里的“恒等同胚”适用于 $J$ 为定义理想的情形。另，原文此处称目标为“离散环层”，应结合（\hyperlink{para:1:10.1.1}{10.1.1}）的伪离散层约定理解，不能断言任意开集的截面环都离散。}
因此有典范同构
\begin{equation}
 \OO_{\frX}/J^\Delta\simeq(\widetilde{A/J})|\frX.
\tag{10.3.2.1}\label{eq:1:10.3.2.1}
\end{equation}
由 $\Gamma(\frX,J^\Delta)=J$，映射 $J\mapsto J^\Delta$ 显然严格保序。根据上述讨论，对 $J\subset J'$，层 $(J')^\Delta/J^\Delta$ 典范地同构于 $(\widetilde{J'/J})|\frX$，也记作 $(J'/J)^\Delta$。
\egapar{10.3.3} 继续沿用（\hyperlink{para:1:10.3.1}{10.3.1}）的假设与记号。若 $\OO_{\frX}$ 的一个理想层 $\J$ 满足：对每个 $x\in\frX$，存在形如 $\frD(f)$（$f\in A$）的开邻域，使 $\J|\frD(f)=K^\Delta$，其中 $K$ 为 $A_{\{f\}}$ 的定义理想，则称 $\J$ 为 $\frX$ 的\emph{定义理想层}，或简称 $\frX$ 的\emph{定义理想}。
\egastatement{10.3.4}{命题} 对每个 $f\in A$，$\frX$ 的每个定义理想都在 $\frD(f)$ 上诱导一个定义理想。

这由（\hyperref[eq:1:10.3.1.1]{10.3.1.1}）得出。
\egastatement{10.3.5}{命题} 若 $A$ 为可容许环，则 $\frX=\Spf(A)$ 的每个定义理想都形如 $J^\Delta$，其中 $J$ 为 $A$ 的唯一确定的定义理想。

设 $\J$ 为 $\frX$ 的定义理想。按假设及 $\frX$ 的拟紧性，可取有限个 $f_i\in A$，使各 $\frD(f_i)$ 覆盖 $\frX$，并且 $\J|\frD(f_i)=K_i^\Delta$，其中 $K_i$ 为 $A_{\{f_i\}}$ 的定义理想。对每个 $i$，存在 $A$ 的开理想 $H_i$，使 $(H_i)_{\{f_i\}}=K_i$（0，\hyperlink{para:0:7.6.9}{7.6.9}）。取 $A$ 的一个包含于所有 $H_i$ 的定义理想 $H$。

$\J/H^\Delta$ 在 $\Spec(A/H)$ 的结构层 $\widetilde{A/H}$ 中的典范像（\hyperlink{para:1:10.3.2}{10.3.2}），限制到 $\frD(f_i)$ 后与 $\widetilde{H_i/H}$ 的限制相同。因此，此典范像为 $\Spec(A/H)$ 上的拟凝聚层，故形如 $\widetilde{J/H}$，其中 $J$ 为 $A$ 中包含 $H$ 的理想（\hyperlink{para:1:1.4.1}{1.4.1}）。于是 $\J=J^\Delta$（\hyperlink{para:1:10.3.2}{10.3.2}）。此外，对每个 $i$，存在整数 $n_i$ 使 $K_i^{n_i}\subset H_{\{f_i\}}$。令 $n$ 为所有 $n_i$ 的最大值，则 $(\J/H^\Delta)^n=0$，从而（\hyperlink{para:1:10.3.2}{10.3.2}）$\widetilde{(J/H)^n}=0$，最终得到 $(J/H)^n=0$（\hyperlink{para:1:1.3.13}{1.3.13}）。因此 $J$ 为 $A$ 的定义理想（0，\hyperlink{para:0:7.1.4}{7.1.4}）。唯一性来自 $\Gamma(\frX,J^\Delta)=J$。
\egastatement{10.3.6}{命题} 设 $A$ 为进制环，$J$ 为 $A$ 的定义理想，且 $J/J^2$ 为有限型 $A/J$-模。那么，对每个整数 $n>0$，有 $(J^\Delta)^n=(J^n)^\Delta$。

对每个 $f\in A$，由于 $J^n$ 为开理想，利用（\hyperref[eq:1:10.3.1.1]{10.3.1.1}）与（0，\hyperlink{para:0:7.6.12}{7.6.12}），有
\[
 (\Gamma(\frD(f),J^\Delta))^n
 =(J_{\{f\}})^n=(J^n)_{\{f\}}
 =\Gamma(\frD(f),(J^n)^\Delta).
\]
而 $(J^\Delta)^n$ 是预层 $U\mapsto(\Gamma(U,J^\Delta))^n$ 所伴随的层（0，\hyperlink{para:0:4.1.6}{4.1.6}）。因各 $\frD(f)$ 构成拓扑基，命题得证。
\egapar{10.3.7} 若定义理想族 $(\J_\lambda)$ 满足 $\frX$ 的每个定义理想都包含其中某个 $\J_\lambda$，就称该族为\emph{定义理想基本系}。由于 $\J_\lambda=J_\lambda^\Delta$，这等价于各 $J_\lambda$ 构成 $A$ 中零点的基本邻域系。取一族 $f_\alpha\in A$，使各 $\frD(f_\alpha)$ 覆盖 $\frX$。若 $\OO_{\frX}$ 的一个向下有向的理想层族 $(\J_\lambda)$ 满足：对每个 $\alpha$，各 $\J_\lambda|\frD(f_\alpha)$ 构成 $\frD(f_\alpha)$ 的定义理想基本系，那么 $(\J_\lambda)$ 就是 $\frX$ 的定义理想基本系。

事实上，对 $\frX$ 的每个定义理想 $\J$，存在由有限个 $\frD(f_i)$ 组成的覆盖，使对每个 $i$，某个 $\J_{\lambda_i}|\frD(f_i)$ 是包含于 $\J|\frD(f_i)$ 的定义理想。取指标 $\mu$，使 $\J_\mu\subset\J_{\lambda_i}$ 对所有 $i$ 成立。由（\hyperlink{para:1:10.3.3}{10.3.3}），$\J_\mu$ 是 $\frX$ 的定义理想，并显然包含于 $\J$，得到断言。


% END chapters/ch1_s10_3.tex

\par

% BEGIN chapters/ch1_s10_4.tex
\subsection{形式预概形与形式预概形态射}\hypertarget{heading:1:10.4}{}
\src{185—186}
\egapar{10.4.1} 给定拓扑赋环空间 $\frX$，若 $\frX$ 在开集 $U\subset\frX$ 上诱导的拓扑赋环空间为仿射形式概形，就称 $U$ 为\emph{仿射形式开集}。若其环还为进制环（分别为诺特进制环），则称 $U$ 为\emph{进制仿射形式开集}（分别为\emph{诺特仿射形式开集}）。
\egastatement{10.4.2}{定义} 若拓扑赋环空间 $\frX$ 的每一点都具有仿射形式开邻域，就称它为\emph{形式预概形}。若其每一点都具有进制仿射形式开邻域（分别为诺特仿射形式开邻域），则称 $\frX$ 为\emph{进制的}（分别为\emph{局部诺特的}）。若 $\frX$ 局部诺特，且其底空间拟紧，因而诺特，则称 $\frX$ \emph{诺特}。
\egastatement{10.4.3}{命题} 若 $\frX$ 为形式预概形（分别为局部诺特形式预概形），则仿射形式开集（分别为诺特仿射形式开集）构成 $\frX$ 的拓扑基。

这由（\hyperlink{para:1:10.4.2}{10.4.2}）与（\hyperlink{para:1:10.1.4}{10.1.4}）得出，同时用到：若 $A$ 为诺特进制环，则对每个 $f\in A$，$A_{\{f\}}$ 也为诺特进制环（0，\hyperlink{para:0:7.6.11}{7.6.11}）。
\egastatement{10.4.4}{推论} 若 $\frX$ 为形式预概形（分别为局部诺特形式预概形、诺特形式预概形），则它在任意开集上诱导的拓扑赋环空间仍为形式预概形（分别为局部诺特形式预概形、诺特形式预概形）。
\egastatement{10.4.5}{定义} 给定两个形式预概形 $\frX,\frY$，称拓扑赋环空间态射 $(\psi,\theta):\frX\to\frY$ 为一个\emph{形式预概形态射}，若对每个 $x\in\frX$，$\theta_x^\sharp:\OO_{\psi(x)}\to\OO_x$ 为局部同态。

两个形式预概形态射的复合显然仍为这样的态射，因此形式预概形构成一个范畴。将从 $\frX$ 到 $\frY$ 的态射集合记作 $\Hom(\frX,\frY)$。

若 $U$ 为 $\frX$ 的开子集，则在 $U$ 上诱导的形式预概形到 $\frX$ 的典范嵌入，为形式预概形态射，甚至是拓扑赋环空间的单态射（0，\hyperlink{para:0:4.1.1}{4.1.1}）。
\egastatement{10.4.6}{命题} 设 $\frX$ 为形式预概形，$\frS=\Spf(A)$ 为仿射形式概形。从 $\frX$ 到 $\frS$ 的形式预概形态射，与从环 $A$ 到拓扑环 $\Gamma(\frX,\OO_{\frX})$ 的连续同态之间，存在典范双射。

证明与（\hyperlink{para:1:2.2.4}{2.2.4}）相同，只须将“同态”改为“连续同态”，“仿射开集”改为“仿射形式开集”，并使用（\hyperlink{para:1:10.2.2}{10.2.2}）代替（\hyperlink{para:1:1.7.3}{1.7.3}）；细节留给读者。
\egapar{10.4.7} 给定形式预概形 $\frS$，若给定形式预概形 $\frX$ 及态射 $\varphi:\frX\to\frS$，就说它们定义了\emph{$\frS$ 上的形式预概形}，或\emph{$\frS$-形式预概形} $\frX$，并称 $\varphi$ 为其\emph{结构态射}。若 $\frS=\Spf(A)$，其中 $A$ 可容许，则也称 $\frX$ 为\emph{$A$-形式预概形}，或\emph{$A$ 上的形式预概形}。任意形式预概形总可视作离散环 $\mathbf Z$ 上的形式预概形。

若 $\frX,\frY$ 为两个 $\frS$-形式预概形，且态射 $u:\frX\to\frY$ 使下图交换，则称 $u$ 为 $\frS$-态射：
\[
\begin{tikzcd}
 \frX\arrow[rr,"u"]\arrow[dr]&&\frY\arrow[dl]\\
 &\frS&
\end{tikzcd}
\]
其中斜箭头为结构态射。按此定义，对固定的 $\frS$，$\frS$-形式预概形构成范畴。将从 $\frX$ 到 $\frY$ 的 $\frS$-态射集合记作 $\Hom_{\frS}(\frX,\frY)$。当 $\frS=\Spf(A)$ 时，也称 $A$-态射。
\egapar{10.4.8} 由于每个仿射概形可视作仿射形式概形（\hyperlink{para:1:10.1.2}{10.1.2}），每个通常的预概形也可视作形式预概形。此外，由（\hyperlink{para:1:10.4.5}{10.4.5}），对通常的预概形，形式预概形的态射（分别为 $S$-态射）与 §2 定义的态射（分别为 $S$-态射）相同。


% END chapters/ch1_s10_4.tex

\par

% BEGIN chapters/ch1_s10_5.tex
\subsection{形式预概形的定义理想}\hypertarget{heading:1:10.5}{}
\src{186—188}
\egapar{10.5.1}设 $\frX$ 为形式预概形。称 $\OO_{\frX}$ 的理想层 $\J$ 为 $\frX$ 的\emph{定义理想层}（或\emph{定义理想}），如果每个 $x\in\frX$ 都有一个仿射形式开邻域 $U$，使 $\J|U$ 是 $\frX$ 在 $U$ 上诱导的仿射形式概形的定义理想（\hyperlink{para:1:10.3.3}{10.3.3}）。由（\hyperref[eq:1:10.3.1.1]{10.3.1.1}）及（\hyperlink{para:1:10.4.3}{10.4.3}），对于每个开集 $V\subset\frX$，$\J|V$ 都是 $\frX$ 在 $V$ 上诱导的形式预概形的定义理想。

称 $\frX$ 的一族定义理想 $(\J_\lambda)$ 为\emph{定义理想基本系}，如果存在由仿射形式开集组成的覆盖 $(U_\alpha)$，使得对每个 $\alpha$，族 $(\J_\lambda|U_\alpha)$ 是 $\frX$ 在 $U_\alpha$ 上诱导的仿射形式概形的定义理想基本系（\hyperlink{para:1:10.3.7}{10.3.7}）。由（\hyperlink{para:1:10.3.7}{10.3.7}）末尾的说明，当 $\frX$ 是仿射形式概形时，这一定义与该处的定义一致。对于每个开集 $V\subset\frX$，族 $(\J_\lambda|V)$ 也是 $V$ 上诱导的形式预概形的定义理想基本系，见（\hyperref[eq:1:10.3.1.1]{10.3.1.1}）。若 $\frX$ 局部诺特，而 $\J$ 是其定义理想，则由（\hyperlink{para:1:10.3.6}{10.3.6}），各次幂 $\J^n$ 构成 $\frX$ 的定义理想基本系。

\egapar{10.5.2}设 $\frX$ 为形式预概形，$\J$ 为其定义理想。赋环空间 $(\frX,\OO_{\frX}/\J)$ 是一个通常意义的预概形；当 $\frX$ 为仿射形式概形（分别为局部诺特、诺特形式预概形）时，它是仿射的（分别为局部诺特、诺特的）。事实上，立即可归结到仿射情形，而这一情形已在（\hyperlink{para:1:10.3.2}{10.3.2}）中证明。此外，若 $\theta:\OO_{\frX}\to\OO_{\frX}/\J$ 为典范同态，则
\[
u=(1_{\frX},\theta):(\frX,\OO_{\frX}/\J)\longrightarrow(\frX,\OO_{\frX})
\]
是一个形式预概形态射，称为\emph{典范态射}。这也已在仿射情形中证明（\hyperlink{para:1:10.3.2}{10.3.2}），而一般情形立即归结于此。

\egastatement{10.5.3}{命题}设 $\frX$ 为形式预概形，$(\J_\lambda)$ 为其定义理想基本系。则拓扑环层 $\OO_{\frX}$ 是各伪离散环层 $\OO_{\frX}/\J_\lambda$ 的射影极限（0，\hyperlink{para:0:3.8.1}{3.8.1}）。

由于 $\frX$ 的拓扑具有由拟紧仿射形式开集组成的基（\hyperlink{para:1:10.4.3}{10.4.3}），可归结到仿射情形；此时命题是（\hyperlink{para:1:10.3.5}{10.3.5}）、（\hyperlink{para:1:10.3.2}{10.3.2}）和定义（\hyperlink{para:1:10.1.1}{10.1.1}）的结果。

不能确定每个形式预概形都具有定义理想。然而：

\egastatement{10.5.4}{命题}设 $\frX$ 为局部诺特形式预概形。存在 $\frX$ 的一个最大定义理想 $\mathscr T$；它是使预概形 $(\frX,\OO_{\frX}/\J)$ 约化的唯一一个定义理想 $\J$。若 $\J$ 是 $\frX$ 的任一定义理想，则 $\mathscr T$ 是 $\OO_{\frX}/\J$ 的幂零根在 $\OO_{\frX}\to\OO_{\frX}/\J$ 下的逆像。

先设 $\frX=\Spf(A)$，其中 $A$ 为诺特进制环。由 $A$ 的最大定义理想的存在性及其性质（0，\hyperlink{para:0:7.1.6}{7.1.6}及\hyperlink{para:0:7.1.7}{7.1.7}），结合（\hyperlink{para:1:10.3.4}{10.3.4}）与（\hyperlink{para:1:5.1.1}{5.1.1}），立即得到 $\mathscr T$ 的存在性及上述性质。

为证明一般情形，只须证明：若 $U\supset V$ 是 $\frX$ 的两个诺特仿射形式开集，则 $U$ 的最大定义理想 $\mathscr T_U$ 在 $V$ 上诱导 $V$ 的最大定义理想 $\mathscr T_V$。但
\[
\bigl(V,(\OO_{\frX}|V)/(\mathscr T_U|V)\bigr)
\]
是约化的，故此结论由前面的结果得出。

以 $\frX_{\mathrm{red}}$ 表示通常意义的约化预概形 $(\frX,\OO_{\frX}/\mathscr T)$。

\egastatement{10.5.5}{推论}设 $\frX$ 为局部诺特形式预概形，$\mathscr T$ 为其最大定义理想。对于 $\frX$ 的每个开集 $V$，$\mathscr T|V$ 是 $\frX$ 在 $V$ 上诱导的形式预概形的最大定义理想。

\egastatement{10.5.6}{命题}设 $\frX,\frY$ 为两个形式预概形，$\J$（分别为 $\mathscr K$）为 $\frX$（分别为 $\frY$）的定义理想，$f:\frX\to\frY$ 为形式预概形态射。
\begin{enumerate}[label=(\roman*)]
\item 若 $f^*(\mathscr K)\OO_{\frX}\subset\J$，则存在唯一一个通常预概形的态射
\[
f':(\frX,\OO_{\frX}/\J)\longrightarrow(\frY,\OO_{\frY}/\mathscr K),
\]
使下图交换，其中竖直箭头是典范态射：
\begin{equation}\tag{10.5.6.1}\label{eq:1:10.5.6.1}
\begin{tikzcd}
(\frX,\OO_{\frX})\arrow[r,"f"] &(\frY,\OO_{\frY})\\
(\frX,\OO_{\frX}/\J)\arrow[u]\arrow[r,"f'"'] &(\frY,\OO_{\frY}/\mathscr K)\arrow[u]
\end{tikzcd}
\end{equation}
\item 设 $\frX=\Spf(A)$、$\frY=\Spf(B)$ 是仿射形式概形，$\J=J^\Delta$、$\mathscr K=K^\Delta$，其中 $J$（分别为 $K$）为 $A$（分别为 $B$）的定义理想，并设 $f=({}^{a}\varphi,\wt\varphi)$，其中 $\varphi:B\to A$ 为连续同态。则 $f^*(\mathscr K)\OO_{\frX}\subset\J$ 的充要条件是 $\varphi(K)\subset J$；此时 $f'=({}^{a}\varphi',\wt{\varphi'})$，其中 $\varphi':B/K\to A/J$ 是由 $\varphi$ 取商所得的同态。
\end{enumerate}

(i) 写 $f=(\psi,\theta)$。假设说明，$\psi^*(\OO_{\frY})$ 的理想层 $\psi^*(\mathscr K)$ 在 $\theta^\sharp:\psi^*(\OO_{\frY})\to\OO_{\frX}$ 下的像包含于 $\J$（0，\hyperlink{para:0:4.3.5}{4.3.5}）。取商，得到环层同态
\[
\omega:\psi^*(\OO_{\frY}/\mathscr K)
=\psi^*(\OO_{\frY})/\psi^*(\mathscr K)
\longrightarrow\OO_{\frX}/\J.
\]
又因每个 $\theta_x^\sharp$ 是局部同态，$\omega_x$ 也是局部同态。因此赋环空间态射 $(\psi,\omega^\flat)$ 是所求的唯一态射 $f'$（\hyperlink{para:1:2.2.1}{2.2.1}）。

(ii) 仿射形式预概形态射与连续环同态之间的典范函子对应（\hyperlink{para:1:10.2.2}{10.2.2}）表明，所考虑的情形中，关系 $f^*(\mathscr K)\OO_{\frX}\subset\J$ 蕴含 $f'=({}^{a}\varphi',\wt{\varphi'})$，其中 $\varphi'$ 是使下图交换的唯一同态：
\begin{equation}\tag{10.5.6.2}\label{eq:1:10.5.6.2}
\begin{tikzcd}
B\arrow[r,"\varphi"]\arrow[d]&A\arrow[d]\\
B/K\arrow[r,"\varphi'"']&A/J
\end{tikzcd}
\end{equation}
$\varphi'$ 的存在性蕴含 $\varphi(K)\subset J$。反之，若此条件成立，以 $\varphi'$ 表示使（\hyperref[eq:1:10.5.6.2]{10.5.6.2}）交换的唯一同态，并置 $f'=({}^{a}\varphi',\wt{\varphi'})$，则（\hyperref[eq:1:10.5.6.1]{10.5.6.1}）显然交换。考察分别对应于 $f,f'$ 的同态 $\psi^*(\OO_{\frY})\to\OO_{\frX}$ 及 $\psi^*(\OO_{\frY}/\mathscr K)\to\OO_{\frX}/\J$，即知这蕴含 $f^*(\mathscr K)\OO_{\frX}\subset\J$。

上述对应 $f\mapsto f'$ 显然具有函子性。


% END chapters/ch1_s10_5.tex

\par

% BEGIN chapters/ch1_s10_6.tex
\subsection{形式预概形作为预概形的归纳极限}\hypertarget{heading:1:10.6}{}
\src{188—193}
\egapar{10.6.1}设 $\frX$ 为形式预概形，$(\J_\lambda)$ 为其定义理想基本系；对每个 $\lambda$，以 $f_\lambda$ 表示典范态射 $(\frX,\OO_{\frX}/\J_\lambda)\to\frX$（\hyperlink{para:1:10.5.2}{10.5.2}）。对于 $\J_\mu\subset\J_\lambda$，典范同态 $\OO_{\frX}/\J_\mu\to\OO_{\frX}/\J_\lambda$ 定义通常预概形的典范态射
\[
f_{\mu\lambda}:(\frX,\OO_{\frX}/\J_\lambda)\longrightarrow(\frX,\OO_{\frX}/\J_\mu),
\qquad f_\lambda=f_\mu\circ f_{\mu\lambda}.
\]
因此，预概形 $X_\lambda=(\frX,\OO_{\frX}/\J_\lambda)$ 及态射 $f_{\mu\lambda}$ 构成形式预概形范畴中的一个归纳系（\hyperlink{para:1:10.4.8}{10.4.8}）。

\egastatement{10.6.2}{命题}沿用（\hyperlink{para:1:10.6.1}{10.6.1}）的记号，形式预概形 $\frX$ 与各态射 $f_\lambda$ 构成归纳系 $(X_\lambda,f_{\mu\lambda})$ 在形式预概形范畴中的归纳极限（T，I，1.8）。

设 $\frY$ 为形式预概形，对每个 $\lambda$ 给定态射 $g_\lambda=(\psi_\lambda,\theta_\lambda):X_\lambda\to\frY$，且对于 $\J_\mu\subset\J_\lambda$ 有 $g_\lambda=g_\mu\circ f_{\mu\lambda}$。后一条件及 $X_\lambda$ 的定义首先说明，所有 $\psi_\lambda$ 都是底空间之间同一个连续映射 $\psi:\frX\to\frY$。此外，各环层同态
\[
\theta_\lambda^\sharp:\psi^*(\OO_{\frY})\longrightarrow
\OO_{X_\lambda}=\OO_{\frX}/\J_\lambda
\]
构成射影系。取射影极限，得到同态 $\omega:\psi^*(\OO_{\frY})\to\plim\OO_{\frX}/\J_\lambda=\OO_{\frX}$，而赋环空间态射 $g=(\psi,\omega^\flat)$ 显然是使下图交换的唯一态射：
\begin{equation}\tag{10.6.2.1}\label{eq:1:10.6.2.1}
\begin{tikzcd}
X_\lambda\arrow[rr,"g_\lambda"]\arrow[dr,"f_\lambda"']&&\frY\\
&\frX\arrow[ur,"g"']&
\end{tikzcd}
\end{equation}
尚须证明 $g$ 是形式预概形态射。这一问题在 $\frX,\frY$ 上是局部的，因而可设 $\frX=\Spf(A)$、$\frY=\Spf(B)$，其中 $A,B$ 为可容许环，且 $\J_\lambda=J_\lambda^\Delta$，$(J_\lambda)$ 是 $A$ 的定义理想基本系（\hyperlink{para:1:10.3.5}{10.3.5}）。由于 $A=\plim A/J_\lambda$，根据（\hyperlink{para:1:10.2.2}{10.2.2}）的双射对应及射影极限的定义，存在使（\hyperref[eq:1:10.6.2.1]{10.6.2.1}）交换的仿射形式概形态射 $g$。但赋环空间态射的唯一性表明，它与证明开头定义的 $g$ 相同。

下一个命题在某些附加条件下，给出一个给定的通常预概形归纳系在形式预概形范畴中归纳极限的存在性。

\egastatement{10.6.3}{命题}设 $X$ 为拓扑空间，$(\OO_i,u_{ji})$ 为以 $\mathbb N$ 为指标集的环层射影系，其中 $u_{ji}:\OO_i\to\OO_j$（$i\ge j$）。令 $\J_i$ 为 $u_{0i}:\OO_i\to\OO_0$ 的核。假设：
\begin{enumerate}[label=\alph*)]
\item 赋环空间 $(X,\OO_i)$ 是一个预概形 $X_i$；
\item 对每个 $x\in X$ 及每个 $i$，存在 $x$ 的开邻域 $U_i$，使 $\J_i|U_i$ 幂零；
\item 各同态 $u_{ji}$ 都是满射。
\end{enumerate}
令 $\OO_{\frX}$ 为各伪离散环层 $\OO_i$ 的拓扑环层射影极限，$u_i:\OO_{\frX}\to\OO_i$ 为典范同态。则拓扑赋环空间 $\frX=(X,\OO_{\frX})$ 是形式预概形；各 $u_i$ 是满射，其核 $\J^{(i)}$ 构成 $\frX$ 的定义理想基本系，且 $\J^{(0)}$ 是各理想层 $\J_i$ 的射影极限。

首先，在每个茎上 $u_{ji}$ 是满同态，因而是局部同态；故 $v_{ij}=(1_X,u_{ji}):X_j\to X_i$（$i\ge j$）是预概形态射（\hyperlink{para:1:2.2.1}{2.2.1}）。先设每个 $X_i$ 都是环为 $A_i$ 的仿射概形。存在环同态 $\varphi_{ji}:A_i\to A_j$，使 $u_{ji}=\wt{\varphi_{ji}}$（\hyperlink{para:1:1.7.3}{1.7.3}）。由（\hyperlink{para:1:1.6.3}{1.6.3}），$\OO_j$ 是 $X_i$ 上的拟凝聚 $\OO_i$-模层，其外乘由 $u_{ji}$ 定义，它与借助 $\varphi_{ji}$ 视为 $A_i$-模的 $A_j$ 相伴。

对 $f\in A_i$，置 $f'=\varphi_{ji}(f)$。由假设，$D(f)$ 与 $D(f')$ 在 $X$ 中是同一个开集，而对应于 $u_{ji}$ 的截面同态
\[
\Gamma(D(f),\OO_i)=(A_i)_f\longrightarrow
\Gamma(D(f'),\OO_j)=(A_j)_{f'}
\]
就是 $(\varphi_{ji})_f$（\hyperlink{para:1:1.6.1}{1.6.1}）。但将 $A_j$ 视为 $A_i$-模时，$(A_j)_f$ 就是 $(A_i)_f$-模 $(A_j)_{f'}$；所以，将 $\varphi_{ji}$ 视为 $A_i$-模同态时，仍有 $u_{ji}=\wt{\varphi_{ji}}$。由 $u_{ji}$ 满射即得 $\varphi_{ji}$ 满射（\hyperlink{para:1:1.3.9}{1.3.9}）；若 $J_{ji}$ 是其核，则 $u_{ji}$ 的核是拟凝聚 $\OO_i$-模层 $\wt{J_{ji}}$。特别地，$\J_i=\wt{J_i}$，其中 $J_i=\Ker(\varphi_{0i}:A_i\to A_0)$。

假设 b) 蕴含 $\J_i$ 幂零：由于 $X$ 拟紧，可用有限个开集 $U_k$ 覆盖 $X$，且 $(\J_i|U_k)^{n_k}=0$；取 $n=\max n_k$，即有 $\J_i^n=0$。于是 $J_i$ 幂零（\hyperlink{para:1:1.3.13}{1.3.13}）。环 $A=\plim A_i$ 因而可容许（0，\hyperlink{para:0:7.2.2}{7.2.2}），典范同态 $\varphi_i:A\to A_i$ 是满射，其核 $J^{(i)}$ 是 $J_{ik}$（$k\ge i$）的射影极限；各 $J^{(i)}$ 构成 $A$ 的零点邻域基本系。此时（\hyperlink{para:1:10.6.3}{10.6.3}）的各项结论由（\hyperlink{para:1:10.1.1}{10.1.1}）及（\hyperlink{para:1:10.3.2}{10.3.2}）得出，$(X,\OO_{\frX})$ 正是 $\Spf(A)$。

仍在这一特殊情形中，若 $f=(f_i)\in A=\plim A_i$，则所有 $D(f_i)$（$X_i$ 中的仿射开集）都与 $\frX$ 的开集 $\frD(f)$ 相同；$X_i$ 在其上诱导的预概形因而与 $\Spec((A_i)_{f_i})$ 相同。

在一般情形，先注意到，对 $X$ 的每个拟紧开集 $U$，每个 $\J_i|U$ 都幂零，证明与上述相同。现证明每个 $x\in X$ 都有一个同时为所有 $X_i$ 的仿射开集的开邻域 $U$。取 $U$ 为 $X_0$ 的仿射开邻域，并注意 $\OO_{X_0}=\OO_{X_i}/\J_i$。由于 $\J_i|U$ 幂零，由（\hyperlink{para:1:5.1.9}{5.1.9}）知 $U$ 也是每个 $X_i$ 的仿射开集。对于这样的 $U$，上面的仿射情形表明 $(U,\OO_{\frX}|U)$ 为形式预概形，各 $\J^{(i)}|U$ 构成其定义理想基本系，而 $\J^{(0)}|U$ 是各 $\J_i|U$ 的射影极限。结论得证。

\egastatement{10.6.4}{推论}再假设对于 $i\ge j$，$u_{ji}$ 的核是 $\J_i^{j+1}$，且 $\J_1/\J_1^2$ 是 $\OO_0=\OO_1/\J_1$ 上的有限型模层。则 $\frX$ 为进制形式预概形；若 $\J=\J^{(0)}$，$\J^{(n)}$ 为 $\OO_{\frX}\to\OO_n$ 的核，则 $\J^{(n)}=\J^{n+1}$，且 $\J/\J^2$ 与 $\J_1$ 同构。若 $X_0$ 还是局部诺特的（分别为诺特的），则 $\frX$ 也是局部诺特的（分别为诺特的）。

由于 $\frX$ 与 $X_0$ 底空间相同，问题是局部的，可设所有 $X_i$ 仿射。利用 $\J_i=\wt{J_i}$，立即归结到（0，\hyperlink{para:0:7.2.7}{7.2.7}及\hyperlink{para:0:7.2.8}{7.2.8}）的相应断言；此时 $J_1/J_1^2$ 是有限型 $A_0$-模（\hyperlink{para:1:1.3.9}{1.3.9}）。

特别地，\emph{每个局部诺特形式预概形} $\frX$ 都是满足（\hyperlink{para:1:10.6.3}{10.6.3}）与（\hyperlink{para:1:10.6.4}{10.6.4}）条件的一列通常的局部诺特预概形 $(X_n)$ 的归纳极限：只须取 $\frX$ 的一个定义理想 $\J$（\hyperlink{para:1:10.5.4}{10.5.4}），并置 $X_n=(\frX,\OO_{\frX}/\J^{n+1})$，见（\hyperlink{para:1:10.5.1}{10.5.1}）与（\hyperlink{para:1:10.6.2}{10.6.2}）。

\egastatement{10.6.5}{推论}设 $A$ 为可容许环。仿射形式概形 $\frX=\Spf(A)$ 为诺特的充要条件是 $A$ 为诺特进制环。

充分性显然。反之，设 $\frX$ 诺特，$J$ 为 $A$ 的定义理想，$\J=J^\Delta$ 为相应的定义理想。通常预概形 $X_n=(\frX,\OO_{\frX}/\J^{n+1})$ 仿射且诺特，故环 $A_n=A/J^{n+1}$ 诺特（\hyperlink{para:1:6.1.3}{6.1.3}），从而 $J/J^2$ 是有限型 $A/J$-模。由于各 $\J^n$ 构成定义理想基本系（\hyperlink{para:1:10.5.1}{10.5.1}），有 $\OO_{\frX}=\plim\OO_{\frX}/\J^n$（\hyperlink{para:1:10.5.3}{10.5.3}）。据此由（\hyperlink{para:1:10.1.3}{10.1.3}），$A$ 与 $\plim A/J^n$ 拓扑同构，因此是诺特进制环（0，\hyperlink{para:0:7.2.8}{7.2.8}）。

\egastatement{10.6.6}{注}沿用（\hyperlink{para:1:10.6.3}{10.6.3}）的记号，设 $\F_i$ 为 $\OO_i$-模层，并对 $i\ge j$ 给定 $v_{ij}$-态射 $\theta_{ji}:\F_i\to\F_j$，使 $\theta_{kj}\circ\theta_{ji}=\theta_{ki}$（$k\le j\le i$）。由于 $v_{ij}$ 的底连续映射为恒等映射，$\theta_{ji}$ 是空间 $X$ 上的阿贝尔群层同态。若 $\F$ 为阿贝尔群层射影系 $(\F_i)$ 的射影极限，则各 $\theta_{ji}$ 为 $v_{ij}$-态射这一事实，使得通过射影极限在 $\F$ 上定义 $\OO_{\frX}$-模层结构。赋予这一结构后，称 $\F$ 为 $\OO_i$-模层系 $(\F_i)$ 关于 $\theta_{ji}$ 的\emph{射影极限}。在 $v_{ij}^*(\F_i)=\F_j$ 且 $\theta_{ji}$ 为恒等态射的特殊情形中，简言之，称 $\F$ 为满足 $v_{ij}^*(\F_i)=\F_j$（$j\le i$）的系 $(\F_i)$ 的射影极限，而不再提及 $\theta_{ji}$。

\egapar{10.6.7}设 $\frX,\frY$ 为形式预概形，$\J,\mathscr K$ 分别为其定义理想，$f:\frX\to\frY$ 为满足 $f^*(\mathscr K)\OO_{\frX}\subset\J$ 的态射。则对于每个 $n>0$，
\[
f^*(\mathscr K^n)\OO_{\frX}
=\bigl(f^*(\mathscr K)\OO_{\frX}\bigr)^n\subset\J^n.
\]
故由 $f$ 可得通常预概形的态射 $f_n:X_n\to Y_n$（\hyperlink{para:1:10.5.6}{10.5.6}），其中
\[
X_n=(\frX,\OO_{\frX}/\J^{n+1}),\qquad
Y_n=(\frY,\OO_{\frY}/\mathscr K^{n+1}).
\]
由定义立即可知，对 $m\le n$，下图交换；换言之，$(f_n)$ 为态射的归纳系：
\begin{equation}\tag{10.6.7.1}\label{eq:1:10.6.7.1}
\begin{tikzcd}
X_m\arrow[r,"f_m"]\arrow[d]&Y_m\arrow[d]\\
X_n\arrow[r,"f_n"']&Y_n
\end{tikzcd}
\end{equation}

\egapar{10.6.8}反之，设 $(X_n)$（分别为 $(Y_n)$）为满足（\hyperlink{para:1:10.6.3}{10.6.3}）的 b)、c) 条件的通常预概形归纳系，$\frX$（分别为 $\frY$）为其归纳极限。由归纳极限的定义，任一构成归纳系的态射序列 $f_n:X_n\to Y_n$ 都有归纳极限 $f:\frX\to\frY$；它是使下图交换的唯一形式预概形态射：
\[
\begin{tikzcd}
X_n\arrow[r,"f_n"]\arrow[d]&Y_n\arrow[d]\\
\frX\arrow[r,"f"']&\frY
\end{tikzcd}
\]

\egastatement{10.6.9}{命题}设 $\frX,\frY$ 为两个局部诺特形式预概形，$\J,\mathscr K$ 分别为其定义理想。（\hyperlink{para:1:10.6.7}{10.6.7}）所定义的映射 $f\mapsto(f_n)$ 是从满足 $f^*(\mathscr K)\OO_{\frX}\subset\J$ 的态射 $f:\frX\to\frY$ 的集合，到使（\hyperref[eq:1:10.6.7.1]{10.6.7.1}）交换的态射序列 $(f_n)$ 的集合的双射。

若 $f$ 是这样的一个序列的归纳极限，须证明 $f^*(\mathscr K)\OO_{\frX}\subset\J$。问题在 $\frX,\frY$ 上是局部的，故可设 $\frX=\Spf(A)$、$\frY=\Spf(B)$ 仿射，$A,B$ 为诺特进制环，$\J=J^\Delta$、$\mathscr K=K^\Delta$，$J,K$ 分别为 $A,B$ 的定义理想。由（\hyperlink{para:1:10.3.6}{10.3.6}）及（\hyperlink{para:1:10.3.2}{10.3.2}），$X_n=\Spec(A_n)$、$Y_n=\Spec(B_n)$，其中 $A_n=A/J^{n+1}$、$B_n=B/K^{n+1}$。有 $f_n=({}^{a}\varphi_n,\wt{\varphi_n})$，各 $\varphi_n:B_n\to A_n$ 构成射影系，所以 $f=({}^{a}\varphi,\wt\varphi)$，其中 $\varphi=\plim\varphi_n$。（\hyperref[eq:1:10.6.7.1]{10.6.7.1}）在 $m=0$ 时的交换性给出
\[
\varphi_n(K/K^{n+1})\subset J/J^{n+1}
\]
对每个 $n$ 成立；取射影极限得 $\varphi(K)\subset J$，再由（\hyperlink{para:1:10.5.6}{10.5.6}，(ii)）得到所求包含关系。

\egastatement{10.6.10}{推论}设 $\frX,\frY$ 为两个局部诺特形式预概形，$\mathscr T$ 为 $\frX$ 的最大定义理想（\hyperlink{para:1:10.5.4}{10.5.4}）。
\begin{enumerate}[label=(\roman*)]
\item 对 $\frY$ 的每个定义理想 $\mathscr K$ 及每个态射 $f:\frX\to\frY$，都有 $f^*(\mathscr K)\OO_{\frX}\subset\mathscr T$。
\item $\Hom(\frX,\frY)$ 与使（\hyperref[eq:1:10.6.7.1]{10.6.7.1}）交换的态射序列 $(f_n)$ 的集合之间有典范双射，其中
\[
X_n=(\frX,\OO_{\frX}/\mathscr T^{n+1}),\qquad
Y_n=(\frY,\OO_{\frY}/\mathscr K^{n+1}).
\]
\end{enumerate}
(ii) 立即由 (i) 及（\hyperlink{para:1:10.6.9}{10.6.9}）得出。证明 (i) 可归结到 $\frX=\Spf(A)$、$\frY=\Spf(B)$，其中 $A,B$ 为诺特进制环，$\mathscr T=T^\Delta$、$\mathscr K=K^\Delta$，$T$ 为 $A$ 的最大定义理想，$K$ 为 $B$ 的定义理想。写 $f=({}^{a}\varphi,\wt\varphi)$，其中 $\varphi:B\to A$ 连续。$K$ 的元素拓扑幂零（0，\hyperlink{para:0:7.1.4}{7.1.4}，(ii)），故 $\varphi(K)$ 的元素也拓扑幂零；由于 $T$ 是 $A$ 中拓扑幂零元素的集合（0，\hyperlink{para:0:7.1.6}{7.1.6}），有 $\varphi(K)\subset T$。据（\hyperlink{para:1:10.5.6}{10.5.6}，(ii)）即得结论。

\egastatement{10.6.11}{推论}设 $\frS,\frX,\frY$ 为三个局部诺特形式预概形，态射 $f:\frX\to\frS$、$g:\frY\to\frS$ 使 $\frX,\frY$ 成为 $\frS$-形式预概形。设 $\J$、$\mathscr K$、$\LL$ 分别为 $\frS,\frX,\frY$ 的定义理想，并假设
\[
f^*(\J)\OO_{\frX}\subset\mathscr K,\qquad
g^*(\J)\OO_{\frY}=\LL.
\]
置 $S_n=(\frS,\OO_{\frS}/\J^{n+1})$、$X_n=(\frX,\OO_{\frX}/\mathscr K^{n+1})$、$Y_n=(\frY,\OO_{\frY}/\LL^{n+1})$。则 $\Hom_{\frS}(\frX,\frY)$ 与使（\hyperref[eq:1:10.6.7.1]{10.6.7.1}）交换的 $S_n$-态射序列 $u_n:X_n\to Y_n$ 的集合之间有典范双射。

对每个 $\frS$-态射 $u:\frX\to\frY$，由定义 $f=g\circ u$，故
\[
u^*(\LL)\OO_{\frX}
=u^*\bigl(g^*(\J)\OO_{\frY}\bigr)\OO_{\frX}
=f^*(\J)\OO_{\frX}\subset\mathscr K.
\]
推论遂由（\hyperlink{para:1:10.6.9}{10.6.9}）得出。

注意，当 $m\le n$ 时，给定态射 $f_n:X_n\to Y_n$ 就决定使（\hyperref[eq:1:10.6.7.1]{10.6.7.1}）交换的唯一态射 $f_m:X_m\to Y_m$，归结到仿射情形立即可见。因此定义了映射
\[
\varphi_{mn}:\Hom_{S_n}(X_n,Y_n)\longrightarrow\Hom_{S_m}(X_m,Y_m),
\]
各集合 $\Hom_{S_n}(X_n,Y_n)$ 关于这些映射构成射影系。（\hyperlink{para:1:10.6.11}{10.6.11}）也就是说，存在典范双射
\[
\Hom_{\frS}(\frX,\frY)\simeq\plim_n\Hom_{S_n}(X_n,Y_n).
\]


% END chapters/ch1_s10_6.tex

\par

% BEGIN chapters/ch1_s10_7.tex
\subsection{形式预概形的积}\hypertarget{heading:1:10.7}{}
\src{193—194}
\egapar{10.7.1}设 $\frS$ 为形式预概形。由于各 $\frS$-形式预概形构成一个范畴，可以定义 $\frS$-形式预概形的积。

\egastatement{10.7.2}{命题}设 $\frX=\Spf(B)$、$\frY=\Spf(C)$ 为仿射形式概形 $\frS=\Spf(A)$ 上的两个仿射形式概形。令 $\frZ=\Spf(B\ctensor_A C)$，$p_1,p_2$ 为分别对应于典范连续 $A$-同态 $\rho:B\to B\ctensor_A C$、$\sigma:C\to B\ctensor_A C$ 的 $\frS$-态射（\hyperlink{para:1:10.2.2}{10.2.2}）。则 $(\frZ,p_1,p_2)$ 是 $\frS$-仿射形式概形 $\frX,\frY$ 的一个积。

根据（\hyperlink{para:1:10.4.6}{10.4.6}），只须验证：对每个连续 $A$-同态 $\varphi:B\ctensor_A C\to D$，其中 $D$ 为同时是拓扑 $A$-代数的可容许环，使它对应于 $(\varphi\circ\rho,\varphi\circ\sigma)$，便定义了一个双射
\[
\Hom_A(B\ctensor_A C,D)\simeq\Hom_A(B,D)\times\Hom_A(C,D).
\]
这正是完备张量积的泛性质（0，\hyperlink{para:0:7.7.6}{7.7.6}）；这里的同态均为连续同态。

\egastatement{10.7.3}{命题}给定两个 $\frS$-形式预概形 $\frX,\frY$，积 $\frX\times_{\frS}\frY$ 存在。

证明与（\hyperlink{para:1:3.2.6}{3.2.6}）相同，只须将仿射概形（分别为仿射开集）换为仿射形式概形（分别为仿射形式开集），并将命题（\hyperlink{para:1:3.2.2}{3.2.2}）换为（\hyperlink{para:1:10.7.2}{10.7.2}）。预概形积的全部形式性质（\hyperlink{para:1:3.2.7}{3.2.7}、\hyperlink{para:1:3.2.8}{3.2.8}及\hyperlink{para:1:3.3.1}{3.3.1}—\hyperlink{para:1:3.3.12}{3.3.12}），对形式预概形的积均不加修改地成立。

\egapar{10.7.4}设 $\frS,\frX,\frY$ 为三个形式预概形，$f:\frX\to\frS$、$g:\frY\to\frS$ 为两个态射。假设它们分别具有以同一个集合 $I$ 为指标集的定义理想基本系 $(\J_\lambda)$、$(\mathscr K_\lambda)$、$(\LL_\lambda)$，且对每个 $\lambda$ 有
\[
f^*(\J_\lambda)\OO_{\frX}\subset\mathscr K_\lambda,\qquad
g^*(\J_\lambda)\OO_{\frY}\subset\LL_\lambda.
\]
置
\[
S_\lambda=(\frS,\OO_{\frS}/\J_\lambda),\quad
X_\lambda=(\frX,\OO_{\frX}/\mathscr K_\lambda),\quad
Y_\lambda=(\frY,\OO_{\frY}/\LL_\lambda).
\]
当 $\J_\mu\subset\J_\lambda$、$\mathscr K_\mu\subset\mathscr K_\lambda$、$\LL_\mu\subset\LL_\lambda$ 时，$S_\lambda$（分别为 $X_\lambda,Y_\lambda$）是 $S_\mu$（分别为 $X_\mu,Y_\mu$）的具有相同底空间的闭子预概形（\hyperlink{para:1:10.6.1}{10.6.1}）。由于 $S_\lambda\to S_\mu$ 为单态射，积 $X_\lambda\times_{S_\lambda}Y_\lambda$ 与 $X_\lambda\times_{S_\mu}Y_\lambda$ 首先相同（\hyperlink{para:1:3.2.4}{3.2.4}）；后者又与 $X_\mu\times_{S_\mu}Y_\mu$ 的具有相同底空间的一个闭子预概形相同（\hyperlink{para:1:4.3.1}{4.3.1}）。

此时，$\frX\times_{\frS}\frY$ 是各通常预概形 $X_\lambda\times_{S_\lambda}Y_\lambda$ 的归纳极限。事实上，如（\hyperlink{para:1:10.6.2}{10.6.2}）中一样，可归结到 $\frS,\frX,\frY$ 都为仿射形式概形的情形。根据（\hyperlink{para:1:10.5.6}{10.5.6}，(ii)）及关于定义理想基本系的假设，该断言立即由两个代数的完备张量积的定义得出（0，\hyperlink{para:0:7.7.1}{7.7.1}）。

再设 $\frZ$ 为 $\frS$-形式预概形，$(\M_\lambda)$ 为以 $I$ 为指标集的定义理想基本系，$u:\frZ\to\frX$、$v:\frZ\to\frY$ 为两个 $\frS$-态射，满足
\[
u^*(\mathscr K_\lambda)\OO_{\frZ}\subset\M_\lambda,\qquad
v^*(\LL_\lambda)\OO_{\frZ}\subset\M_\lambda.
\]
置 $Z_\lambda=(\frZ,\OO_{\frZ}/\M_\lambda)$，并以 $u_\lambda:Z_\lambda\to X_\lambda$、$v_\lambda:Z_\lambda\to Y_\lambda$ 表示由 $u,v$ 得到的 $S_\lambda$-态射（\hyperlink{para:1:10.5.6}{10.5.6}）。立即可验，$(u,v)_{\frS}$ 为 $(u_\lambda,v_\lambda)_{S_\lambda}$ 的归纳极限。

特别地，当 $\frS,\frX,\frY$ 都局部诺特时，可取某个定义理想的各次幂所组成的基本系（\hyperlink{para:1:10.5.1}{10.5.1}），因而适用上述讨论。但须注意，$\frX\times_{\frS}\frY$ 不一定局部诺特，另见（\hyperlink{para:1:10.13.5}{10.13.5}）。


% END chapters/ch1_s10_7.tex

\par

% BEGIN chapters/ch1_s10_8.tex
\subsection{预概形沿闭子集的形式完备化}\hypertarget{heading:1:10.8}{}
\src{194—198}
\egapar{10.8.1}设 $X$ 为通常意义的局部诺特预概形，$X'$ 为其底空间的闭子集。以 $\Phi$ 表示 $\OO_X$ 中使 $\Supp(\OO_X/\J)=X'$ 的凝聚理想层 $\J$ 的集合。$\Phi$ 非空（\hyperlink{para:1:5.2.1}{5.2.1}、\hyperlink{para:1:4.1.4}{4.1.4}及\hyperlink{para:1:6.1.1}{6.1.1}）；我们用关系 $\supset$ 将它排序。

\egastatement{10.8.2}{引理}有序集 $\Phi$ 是滤过的。若 $X$ 诺特，则对于每个 $\J_0\in\Phi$，其各次幂 $\J_0^n$（$n>0$）的集合在 $\Phi$ 中共尾。

事实上，若 $\J_1,\J_2\in\Phi$，置 $\J=\J_1\cap\J_2$。由于 $\OO_X$ 凝聚，$\J$ 也凝聚（\hyperlink{para:1:6.1.1}{6.1.1}及0，\hyperlink{para:1:5.3.4}{5.3.4}）；对于每个 $x\in X$，$\J_x=(\J_1)_x\cap(\J_2)_x$，故当 $x\notin X'$ 时 $\J_x=\OO_x$，而当 $x\in X'$ 时 $\J_x\ne\OO_x$，即 $\J\in\Phi$。另一方面，若 $X$ 诺特，$\J_0,\J\in\Phi$，则存在 $n>0$ 使 $\J_0^n(\OO_X/\J)=0$（\hyperlink{para:1:9.3.4}{9.3.4}），即 $\J_0^n\subset\J$。

\egapar{10.8.3}设 $\F$ 为凝聚 $\OO_X$-模层。对于每个 $\J\in\Phi$，$\F\otimes_{\OO_X}(\OO_X/\J)$ 为凝聚 $\OO_X$-模层（\hyperlink{para:1:9.1.1}{9.1.1}），其支集包含于 $X'$；通常将它与其在 $X'$ 上的限制等同。当 $\J$ 遍历 $\Phi$ 时，这些层构成阿贝尔群层的射影系。

\egastatement{10.8.4}{定义}给定局部诺特预概形 $X$ 的闭子集 $X'$ 及凝聚 $\OO_X$-模层 $\F$，称
\[
\plim_{\J\in\Phi}\bigl(\F\otimes_{\OO_X}(\OO_X/\J)\bigr)
\]
在 $X'$ 上的限制为 $\F$ \emph{沿 $X'$ 的完备化}，记为 $\F_{/X'}$，无混淆时也记为 $\widehat\F$。它在 $X'$ 上的截面称为 $\F$ 沿 $X'$ 的\emph{形式截面}。

显然，对每个开集 $U\subset X$，有 $(\F|U)_{/(U\cap X')}=\F_{/X'}|(U\cap X')$。通过射影极限，$(\OO_X)_{/X'}$ 成为环层，$\F_{/X'}$ 可视为其上的模层。又因 $X'$ 的拓扑有拟紧开集基，可将 $(\OO_X)_{/X'}$（分别为 $\F_{/X'}$）视为各伪离散环层 $\OO_X/\J$（分别为伪离散群层 $\F\otimes_{\OO_X}(\OO_X/\J)=\F/\J\F$）的射影极限，从而成为拓扑环层（分别为拓扑群层）。通过射影极限，$\F_{/X'}$ 因而成为拓扑 $(\OO_X)_{/X'}$-模层（0，\hyperlink{para:0:3.8.1}{3.8.1}及\hyperlink{para:0:3.8.2}{3.8.2}）。回忆，对每个拟紧开集 $U\subset X$，$\Gamma(U\cap X',(\OO_X)_{/X'})$（分别为 $\Gamma(U\cap X',\F_{/X'})$）是离散环 $\Gamma(U,\OO_X/\J)$（分别为离散群 $\Gamma(U,\F/\J\F)$）的射影极限。

若 $u:\F\to\G$ 为 $\OO_X$-模层同态，则典范地得到
\[
u_\J:\F\otimes_{\OO_X}(\OO_X/\J)\longrightarrow\G\otimes_{\OO_X}(\OO_X/\J)
\]
的射影系。取射影极限并限制到 $X'$，得到连续的 $(\OO_X)_{/X'}$-同态 $u_{/X'}:\F_{/X'}\to\G_{/X'}$，也记为 $\widehat u$，称为 $u$ \emph{沿 $X'$ 的完备化}。若 $v:\G\to\HH$ 为另一个同态，则显然 $(v\circ u)_{/X'}=v_{/X'}\circ u_{/X'}$。因此 $\F_{/X'}$ 是从凝聚 $\OO_X$-模层范畴到拓扑 $(\OO_X)_{/X'}$-模层范畴的加性协变函子。

\egastatement{10.8.5}{命题}$(\OO_X)_{/X'}$ 的支集为 $X'$。拓扑赋环空间 $(X',(\OO_X)_{/X'})$ 为局部诺特形式预概形；若 $\J\in\Phi$，则 $\J_{/X'}$ 为此形式预概形的定义理想。当 $X=\Spec(A)$ 仿射、$A$ 诺特、$\J=\wt J$ 且 $X'=V(J)$ 时，$(X',(\OO_X)_{/X'})$ 典范地与 $\Spf(\widehat A)$ 等同，其中 $\widehat A$ 是 $A$ 关于 $J$-预进制拓扑的分离完备化。

显然只须证明最后一项。由（0，\hyperlink{para:0:7.3.3}{7.3.3}），$J$ 的 $J$-预进制分离完备化 $\widehat J$ 与 $\widehat A$ 中的理想 $J\widehat A$ 等同；$\widehat A$ 为诺特 $\widehat J$-进制环，且 $\widehat A/\widehat J^{\,n}=A/J^n$。后一关系说明，$\widehat A$ 的开素理想是 $\widehat{\mathfrak p}$，其中 $\mathfrak p$ 为 $A$ 中包含 $J$ 的素理想，且其在 $A$ 中的逆像为 $\mathfrak p$，故 $\Spf(\widehat A)=X'$。又有 $\OO_X/\J^n=(A/J^n)^\sim$，所以命题立即由定义得出。

称这样定义的形式预概形为 $X$ \emph{沿 $X'$ 的完备化}，记为 $X_{/X'}$，无混淆时记为 $\widehat X$。取 $X'=X$ 时可以取 $\J=0$，故 $X_{/X}=X$。若 $U$ 为 $X$ 在一个开集上诱导的子预概形，则 $U_{/(U\cap X')}$ 显然典范地与 $X_{/X'}$ 在开集 $U\cap X'\subset X'$ 上诱导的形式子预概形等同。

\egastatement{10.8.6}{推论}通常预概形 $\widehat X_{\mathrm{red}}$ 是 $X$ 的底空间为 $X'$ 的唯一约化子预概形（\hyperlink{para:1:5.2.1}{5.2.1}）。$\widehat X$ 诺特的充要条件是 $\widehat X_{\mathrm{red}}$ 诺特；$X$ 诺特则足以保证这一点。

$\widehat X_{\mathrm{red}}$ 的确定是局部问题（\hyperlink{para:1:10.5.4}{10.5.4}），可再设 $X$ 仿射且环 $A$ 诺特。沿用（\hyperlink{para:1:10.8.5}{10.8.5}）的记号，$\widehat A$ 中拓扑幂零元素的理想 $T$ 是 $A/J$ 的幂零根在 $\widehat A\to\widehat A/\widehat J=A/J$ 下的逆像（0，\hyperlink{para:0:7.1.3}{7.1.3}），故 $\widehat A/T$ 同构于 $A/J$ 对其幂零根的商。第一项断言遂由（\hyperlink{para:1:10.5.4}{10.5.4}）与（\hyperlink{para:1:5.1.1}{5.1.1}）得出。若 $\widehat X_{\mathrm{red}}$ 诺特，则底空间 $X'$ 诺特，故 $X_n=(X',\OO_X/\J^{n+1})$ 诺特（\hyperlink{para:1:6.1.2}{6.1.2}），从而 $\widehat X$ 诺特（\hyperlink{para:1:10.6.4}{10.6.4}）；反向由（\hyperlink{para:1:6.1.2}{6.1.2}）立即得出。

\egapar{10.8.7}典范同态 $\OO_X\to\OO_X/\J$（$\J\in\Phi$）构成射影系；取射影极限，得到环层同态
\[
\theta:\OO_X\longrightarrow\psi_*((\OO_X)_{/X'})
=\plim_{\J\in\Phi}\OO_X/\J,
\]
其中 $\psi:X'\hookrightarrow X$ 为底空间的典范嵌入。以 $i$ 或 $i_X$ 表示典范赋环空间态射 $(\psi,\theta):X_{/X'}\to X$。

通过张量积，对每个凝聚 $\OO_X$-模层 $\F$，典范同态 $\OO_X\to\OO_X/\J$ 给出 $\OO_X$-模层同态 $\F\to\F\otimes_{\OO_X}(\OO_X/\J)$ 的射影系。取射影极限，便得到典范函子同态 $\gamma:\F\to\psi_*(\F_{/X'})$。

\egastatement{10.8.8}{命题}(i) 函子 $\F\mapsto\F_{/X'}$ 是正合的。

(ii) 函子同态 $\gamma^\sharp:i^*(\F)\to\F_{/X'}$ 是 $(\OO_X)_{/X'}$-模层的同构。

(i) 只须证明：若 $0\to\F'\to\F\to\F''\to0$ 为凝聚 $\OO_X$-模层的正合列，$U$ 为 $X$ 的一个仿射开集，环 $A$ 诺特，则
\[
0\to\Gamma(U\cap X',\F'_{/X'})\to\Gamma(U\cap X',\F_{/X'})
\to\Gamma(U\cap X',\F''_{/X'})\to0
\]
正合。此时 $\F|U=\wt M$、$\F'|U=\wt{M'}$、$\F''|U=\wt{M''}$，其中 $M,M',M''$ 为有限型 $A$-模，且 $0\to M'\to M\to M''\to0$ 正合（\hyperlink{para:1:1.5.1}{1.5.1}及\hyperlink{para:1:1.3.11}{1.3.11}）。取 $\J\in\Phi$，并令 $J\subset A$ 使 $\J|U=\wt J$。有
\[
\Gamma\bigl(U\cap X',\F\otimes_{\OO_X}(\OO_X/\J^n)\bigr)
=M\otimes_A(A/J^n)
\]
（\hyperlink{para:1:1.3.12}{1.3.12}），故由射影极限的定义，$\Gamma(U\cap X',\F_{/X'})=\plim_n M\otimes_A(A/J^n)=\widehat M$，即 $M$ 的 $J$-预进制分离完备化。同理，另两项为 $\widehat{M'}$、$\widehat{M''}$。所求断言由诺特环上有限型模的完备化函子的正合性得到（0，\hyperlink{para:0:7.3.3}{7.3.3}）。

(ii) 问题是局部的，可设有正合列 $\OO_X^p\to\OO_X^q\to\F\to0$（0，\hyperlink{para:0:5.3.2}{5.3.2}）。由于 $\gamma^\sharp$ 具有函子性，而 $i^*(\F)$ 与 $\F_{/X'}$ 都为右正合函子，得到具有正合行的交换图
\begin{equation}\tag{10.8.8.1}\label{eq:1:10.8.8.1}
\begin{tikzcd}[column sep=small]
i^*(\OO_X^p)\arrow[r]\arrow[d,"\gamma^\sharp"]&i^*(\OO_X^q)\arrow[r]\arrow[d,"\gamma^\sharp"]&i^*(\F)\arrow[r]\arrow[d,"\gamma^\sharp"]&0\\
(\OO_X^p)_{/X'}\arrow[r]&(\OO_X^q)_{/X'}\arrow[r]&\F_{/X'}\arrow[r]&0
\end{tikzcd}
\end{equation}
又因为这两个函子与有限直和可交换（0，\hyperlink{para:0:3.2.6}{3.2.6}及\hyperlink{para:0:4.3.2}{4.3.2}），可归结到 $\F=\OO_X$。此时 $i^*(\OO_X)=(\OO_X)_{/X'}=\OO_{\widehat X}$（0，\hyperlink{para:0:4.3.4}{4.3.4}），$\gamma^\sharp$ 为 $\OO_{\widehat X}$-模层同态。只须验证它将 $X'$ 的开集上的单位截面映为自身；这是显然的，并说明此时 $\gamma^\sharp$ 为恒等同态。

\egastatement{10.8.9}{推论}赋环空间态射 $i:X_{/X'}\to X$ 是平坦的。

这由（0，\hyperlink{para:0:6.7.3}{6.7.3}）及（\hyperlink{para:1:10.8.8}{10.8.8}，(i)）得出。

\egastatement{10.8.10}{推论}若 $\F,\G$ 为凝聚 $\OO_X$-模层，则存在关于 $\F,\G$ 的典范函子同构
\begin{align}
\F_{/X'}\otimes_{(\OO_X)_{/X'}}\G_{/X'}&\simeq(\F\otimes_{\OO_X}\G)_{/X'},\tag{10.8.10.1}\label{eq:1:10.8.10.1}\\
\bigl(\shHom_{\OO_X}(\F,\G)\bigr)_{/X'}&\simeq
\shHom_{(\OO_X)_{/X'}}(\F_{/X'},\G_{/X'}).\tag{10.8.10.2}\label{eq:1:10.8.10.2}
\end{align}
这由 $i^*(\F)$ 与 $\F_{/X'}$ 的典范等同得出。第一个同构的存在是对任意赋环空间态射成立的结果（0，\hyperref[eq:0:4.3.3.1]{4.3.3.1}）；第二个是对任意平坦态射成立的结果（0，\hyperlink{para:0:6.7.6}{6.7.6}），故由（\hyperlink{para:1:10.8.9}{10.8.9}）得出。

\egastatement{10.8.11}{命题}对于每个凝聚 $\OO_X$-模层 $\F$，由 $\F\to\F_{/X'}$ 得到的典范同态 $\Gamma(X,\F)\to\Gamma(X',\F_{/X'})$ 的核，由在 $X'$ 的某个邻域中为零的截面组成。

由 $\F_{/X'}$ 的定义，这种截面的典范像为零。反之，若 $s\in\Gamma(X,\F)$ 的典范像为零，只须证明每个 $x\in X'$ 在 $X$ 中都有一个使 $s$ 为零的邻域。可因此归结到 $X=\Spec(A)$、$A$ 诺特、$X'=V(J)$、$\F=\wt M$，其中 $M$ 为有限型 $A$-模。此时 $\Gamma(X',\F_{/X'})$ 为 $M$ 的 $J$-预进制分离完备化，截面同态即 $M\to\widehat M$。其核是被 $1+J$ 中的某个元素消去的元素的集合（0，\hyperlink{para:0:7.3.7}{7.3.7}），故存在 $f\in J$ 使 $(1+f)s=0$。对每个 $x\in X'$，$(1_x+f_x)s_x=0$；由于 $J_x\OO_x$ 包含于 $\OO_x$ 的极大理想，$1_x+f_x$ 可逆，所以 $s_x=0$。命题得证。

\egastatement{10.8.12}{推论}$\Supp(\F_{/X'})=\Supp(\F)\cap X'$。

$\F_{/X'}$ 显然为有限型 $(\OO_X)_{/X'}$-模层，见（\hyperlink{para:1:10.8.8}{10.8.8}，(ii)）及（0，\hyperlink{para:0:5.2.4}{5.2.4}）；故其支集闭（0，\hyperlink{para:0:5.2.2}{5.2.2}），并显然包含于 $\Supp(\F)\cap X'$。要证相等，立即归结到证明 $\Gamma(X',\F_{/X'})=0$ 蕴含 $\Supp(\F)\cap X'=\varnothing$；这由（\hyperlink{para:1:10.8.11}{10.8.11}）及（\hyperlink{para:1:1.4.1}{1.4.1}）得出。

\egastatement{10.8.13}{推论}设 $u:\F\to\G$ 为凝聚 $\OO_X$-模层的同态。$u_{/X'}:\F_{/X'}\to\G_{/X'}$ 为零的充要条件是 $u$ 在 $X'$ 的某个邻域中为零。

由（\hyperlink{para:1:10.8.8}{10.8.8}，(ii)），$u_{/X'}$ 与 $i^*(u)$ 等同。将 $u$ 视为 $\HH=\shHom_{\OO_X}(\F,\G)$ 在 $X$ 上的截面，则 $u_{/X'}$ 为 $i^*(\HH)=\HH_{/X'}$ 在 $X'$ 上与之典范对应的截面，见（\hyperref[eq:1:10.8.10.2]{10.8.10.2}）及（0，\hyperlink{para:0:4.4.6}{4.4.6}）。只须对凝聚 $\OO_X$-模层 $\HH$ 应用（\hyperlink{para:1:10.8.11}{10.8.11}）。

\egastatement{10.8.14}{推论}设 $u:\F\to\G$ 为凝聚 $\OO_X$-模层同态。$u_{/X'}$ 为单态射（分别为满态射）的充要条件是 $u$ 在 $X'$ 的某个邻域中为单态射（分别为满态射）。

设 $\mathscr T,\mathscr S$ 分别为 $u$ 的核与余核，故有正合列 $0\to\mathscr T\xrightarrow{v}\F\xrightarrow{u}\G\xrightarrow{w}\mathscr S\to0$。由（\hyperlink{para:1:10.8.8}{10.8.8}，(i)）得到正合列
\[
0\to\mathscr T_{/X'}\xrightarrow{v_{/X'}}\F_{/X'}
\xrightarrow{u_{/X'}}\G_{/X'}\xrightarrow{w_{/X'}}\mathscr S_{/X'}\to0.
\]
若 $u_{/X'}$ 为单态射（分别为满态射），则 $v_{/X'}=0$（分别为 $w_{/X'}=0$），故由（\hyperlink{para:1:10.8.13}{10.8.13}），在 $X'$ 的某个邻域中 $v=0$（分别为 $w=0$）。


% END chapters/ch1_s10_8.tex

\par

% BEGIN chapters/ch1_s10_9.tex
\subsection{将态射延拓到完备化}\hypertarget{heading:1:10.9}{}
\src{198—201}
\egapar{10.9.1}设 $X,Y$ 为两个通常意义的局部诺特预概形，$f:X\to Y$ 为态射，$X',Y'$ 分别为其底空间的闭子集，且 $f(X')\subset Y'$。设 $\J,\mathscr K$ 分别为 $\OO_X,\OO_Y$ 的理想层，使 $\Supp(\OO_X/\J)=X'$、$\Supp(\OO_Y/\mathscr K)=Y'$，且 $f^*(\mathscr K)\OO_X\subset\J$。这样的理想层总存在：例如可取 $\J$ 为定义底空间为 $X'$ 的子预概形的最大理想层（\hyperlink{para:1:5.2.1}{5.2.1}），此时由 $f(X')\subset Y'$ 即得该包含关系（\hyperlink{para:1:5.2.4}{5.2.4}）。

对于每个 $n>0$，有 $f^*(\mathscr K^n)\OO_X\subset\J^n$（0，\hyperlink{para:0:4.3.5}{4.3.5}）。因而置
\[
X_n=(X',(\OO_X/\J^{n+1})|X'),\qquad
Y_n=(Y',(\OO_Y/\mathscr K^{n+1})|Y'),
\]
可由 $f$ 得到态射 $f_n:X_n\to Y_n$；显然各 $f_n$ 构成归纳系。以 $\widehat f:X_{/X'}\to Y_{/Y'}$ 表示其归纳极限（\hyperlink{para:1:10.6.8}{10.6.8}），并略滥用术语，称它为 $f$ 到 $X,Y$ 分别沿 $X',Y'$ 的完备化的\emph{延拓}。

容易验证，这一态射与满足上述条件的 $\J,\mathscr K$ 的选择无关。只须在 $X,Y$ 为环分别是 $A,B$ 的诺特仿射概形时验证。此时 $\J=\wt J$、$\mathscr K=\wt K$，其中 $J\subset A$、$K\subset B$ 为理想，$f$ 对应于满足 $\varphi(K)\subset J$ 的环同态 $\varphi:B\to A$（\hyperlink{para:1:4.4.6}{4.4.6}及\hyperlink{para:1:1.7.4}{1.7.4}）。于是 $\widehat f$ 对应于连续同态 $\widehat\varphi:\widehat B\to\widehat A$（\hyperlink{para:1:10.2.2}{10.2.2}），其中 $\widehat A,\widehat B$ 分别为 $J$-与 $K$-预进制分离完备化（\hyperlink{para:1:10.6.8}{10.6.8}）。若以另一个理想层 $\J'=\wt{J'}$ 代替 $\J$，而仍使 $\Supp(\OO_X/\J')=X'$，则 $A$ 的 $J$-与 $J'$-预进制拓扑相同（\hyperlink{para:1:10.8.2}{10.8.2}）。

注意，按这一定义，对应于 $\widehat f$ 的底空间连续映射 $X'\to Y'$，就是 $f$ 在 $X'$ 上的限制。

\egapar{10.9.2}由上一定义立即可知，赋环空间态射的下图交换，其中竖直箭头为典范态射（\hyperlink{para:1:10.8.7}{10.8.7}）：
\[
\begin{tikzcd}
\widehat X\arrow[r,"\widehat f"]\arrow[d,"i_X"']&\widehat Y\arrow[d,"i_Y"]\\
X\arrow[r,"f"']&Y
\end{tikzcd}
\]

\egapar{10.9.3}设 $Z$ 为第三个预概形，$g:Y\to Z$ 为态射，$Z'\subset Z$ 为满足 $g(Y')\subset Z'$ 的闭子集。若 $\widehat g$ 表示 $g$ 沿 $Y',Z'$ 的完备化，则由（\hyperlink{para:1:10.9.1}{10.9.1}）立即得到 $\widehat{g\circ f}=\widehat g\circ\widehat f$。

\egastatement{10.9.4}{命题}设 $X,Y$ 为两个局部诺特 $S$-预概形，$Y$ 在 $S$ 上为有限型。设 $f,g:X\to Y$ 为满足 $f(X')\subset Y'$、$g(X')\subset Y'$ 的两个 $S$-态射。则 $\widehat f=\widehat g$ 的充要条件是 $f,g$ 在 $X'$ 的某个邻域中相同。

充分性显然，且不需要关于 $Y$ 的有限性假设。为证必要性，先注意到 $\widehat f=\widehat g$ 蕴含对每个 $x\in X'$ 有 $f(x)=g(x)$。问题是局部的，故可设 $X,Y$ 分别为 $x$ 与 $y=f(x)=g(x)$ 的仿射开邻域，环均诺特；又可设 $S$ 仿射，$\Gamma(Y,\OO_Y)$ 为有限型 $\Gamma(S,\OO_S)$-代数（\hyperlink{para:1:6.3.3}{6.3.3}）。于是 $f,g$ 对应于两个 $\Gamma(S,\OO_S)$-同态 $\rho,\sigma:\Gamma(Y,\OO_Y)\to\Gamma(X,\OO_X)$（\hyperlink{para:1:1.7.3}{1.7.3}）；由假设，它们到 $\Gamma(Y,\OO_Y)$ 的分离完备化的连续延拓相同。据（\hyperlink{para:1:10.8.11}{10.8.11}），对每个 $s\in\Gamma(Y,\OO_Y)$，$\rho(s),\sigma(s)$ 在 $X'$ 的某个依赖于 $s$ 的邻域中相同。由有限型假设，可得 $X'$ 的一个邻域 $V$，使得所有这些截面对在 $V$ 中同时相同。若 $h\in\Gamma(X,\OO_X)$ 使 $D(h)$ 为包含于 $V$ 的 $x$ 的邻域，则由前述事实及（\hyperlink{para:1:1.4.1}{1.4.1}，d)）知 $f,g$ 在 $D(h)$ 中相同。\transnote{本句原版将 $h$ 所属的截面环印为 $\Gamma(Y,\OO_X)$；这里按其用途——定义 $X$ 中的主开邻域——写为 $\Gamma(X,\OO_X)$。}

\egastatement{10.9.5}{命题}在（\hyperlink{para:1:10.9.1}{10.9.1}）的假设下，对于每个凝聚 $\OO_Y$-模层 $\G$，存在典范函子同构
\[
(f^*(\G))_{/X'}\simeq\widehat f^{\,*}(\G_{/Y'})
\]
（作为 $(\OO_X)_{/X'}$-模层）。

由（\hyperlink{para:1:10.8.8}{10.8.8}），将左端典范地与 $i_X^*(f^*(\G))$ 等同，右端与 $\widehat f^{\,*}(i_Y^*(\G))$ 等同，命题立即由（\hyperlink{para:1:10.9.2}{10.9.2}）的交换图得出。

\egapar{10.9.6}现在设 $\F$ 为凝聚 $\OO_X$-模层，$\G$ 为凝聚 $\OO_Y$-模层。若 $u:\G\to\F$ 为 $f$-态射，则相应有 $\OO_X$-同态 $u^\sharp:f^*(\G)\to\F$，进而通过完备化得到连续同态
\[
(u^\sharp)_{/X'}:(f^*(\G))_{/X'}\longrightarrow\F_{/X'}.
\]
由（\hyperlink{para:1:10.9.5}{10.9.5}），存在唯一的 $\widehat f$-态射 $v:\G_{/Y'}\to\F_{/X'}$，使 $v^\sharp=(u^\sharp)_{/X'}$。

将三元组 $(\F,X,X')$（$\F$ 为凝聚 $\OO_X$-模层，$X'$ 为 $X$ 的闭子集）视为一个范畴的对象；态射 $(\F,X,X')\to(\G,Y,Y')$ 由满足 $f(X')\subset Y'$ 的预概形态射 $f:X\to Y$ 和 $f$-态射 $u:\G\to\F$ 组成。于是可说，$(X_{/X'},\F_{/X'})$ 是该范畴上的函子，其值为由局部诺特形式预概形 $\frZ$ 与 $\OO_{\frZ}$-模层 $\HH$ 所成的二元组 $(\frZ,\HH)$；后一范畴的态射由形式预概形态射 $g$ 及一个 $g$-态射组成。

\egastatement{10.9.7}{命题}设 $S,X,Y$ 为三个局部诺特预概形，$g:X\to S$、$h:Y\to S$ 为态射，$S'\subset S$ 为闭子集，$X'\subset X$、$Y'\subset Y$ 为满足 $g(X')\subset S'$、$h(Y')\subset S'$ 的闭子集。令 $Z=X\times_S Y$，并假设 $Z$ 局部诺特；令 $Z'=p^{-1}(X')\cap q^{-1}(Y')$，其中 $p,q$ 为积的投影。则 $Z_{/Z'}$ 与 $S_{/S'}$-形式预概形的积
\[
X_{/X'}\times_{S_{/S'}}Y_{/Y'}
\]
相同，其中结构态射为 $\widehat g,\widehat h$，投影为 $\widehat p,\widehat q$。

问题在 $S,X,Y$ 上是局部的，故归结到 $S=\Spec(A)$、$X=\Spec(B)$、$Y=\Spec(C)$，$S'=V(J)$、$X'=V(K)$、$Y'=V(L)$，其中理想满足 $\varphi(J)\subset K$、$\psi(J)\subset L$；$\varphi:A\to B$、$\psi:A\to C$ 分别对应于 $g,h$。于是 $Z=\Spec(B\otimes_A C)$，而 $Z'=V(M)$，其中
\[
M=\Im(K\otimes_A C)+\Im(B\otimes_A L).
\]
若 $\widehat A,\widehat B,\widehat C$ 分别为 $J$-、$K$-、$L$-预进制分离完备化，则 $\widehat B\ctensor_{\widehat A}\widehat C$ 是 $B\otimes_A C$ 关于 $M$-预进制拓扑的分离完备化（0，\hyperlink{para:0:7.7.2}{7.7.2}），故由（\hyperlink{para:1:10.7.2}{10.7.2}）得到结论。

还须注意：若 $T$ 为局部诺特 $S$-预概形，$u:T\to X$、$v:T\to Y$ 为 $S$-态射，$T'\subset T$ 为满足 $u(T')\subset X'$、$v(T')\subset Y'$ 的闭子集，则延拓 $\widehat{(u,v)_S}$ 与 $(\widehat u,\widehat v)_{S_{/S'}}$ 等同。

\egastatement{10.9.8}{推论}设 $X,Y$ 为局部诺特 $S$-预概形，且 $X\times_S Y$ 局部诺特。设 $S'$ 为 $S$ 的闭子集，$X',Y'$ 分别为 $X,Y$ 中其在 $S$ 中的像包含于 $S'$ 的闭子集。对于每个满足 $f(X')\subset Y'$ 的 $S$-态射 $f:X\to Y$，图态射 $\Gamma_{\widehat f}$ 与 $f$ 的图态射的延拓 $\widehat{\Gamma_f}$ 等同。

\egastatement{10.9.9}{推论}设 $X,Y$ 为两个局部诺特预概形，$f:X\to Y$ 为态射，$Y'\subset Y$ 为闭子集，$X'=f^{-1}(Y')$。则 $X_{/X'}$ 借助下列交换图，与形式预概形的积 $X\times_Y Y_{/Y'}$ 等同：
\[
\begin{tikzcd}
X&X_{/X'}\arrow[l,"i_X"']\arrow[d,"\widehat f"]\\
Y\arrow[from=u,"f"']&Y_{/Y'}\arrow[l,"i_Y"]
\end{tikzcd}
\]
只须在（\hyperlink{para:1:10.9.7}{10.9.7}）中将 $S,S'$ 换为 $Y$，并将其中第一因子及其闭子集都换为 $X$。

\egastatement{10.9.10}{注}若 $X=X_1\amalg X_2$ 为和（\hyperlink{heading:1:3.1}{3.1}），$X'=X'_1\cup X'_2$，其中 $X'_i$ 为 $X_i$ 的闭子集（$i=1,2$），则立即可见
\[
X_{/X'}=(X_1)_{/X'_1}\amalg(X_2)_{/X'_2}.
\]


% END chapters/ch1_s10_9.tex

\par

% BEGIN chapters/ch1_s10_10.tex
\subsection{对仿射形式概形上凝聚层的应用}\hypertarget{heading:1:10.10}{}
\src{201—203}
\egapar{10.10.1}本小节中，$A$ 始终表示诺特进制环，$J$ 为 $A$ 的一个定义理想。令 $X=\Spec(A)$、$\frX=\Spf(A)$，后者与 $X$ 的闭子集 $V(J)$ 等同（\hyperlink{para:1:10.1.2}{10.1.2}）。由（\hyperlink{para:1:10.1.2}{10.1.2}）及（\hyperlink{para:1:10.8.4}{10.8.4}），仿射形式概形 $\frX$ 就是仿射概形 $X$ 沿其底空间的闭子集 $\frX$ 的完备化 $X_{/\frX}$。因此，对每个凝聚 $\OO_X$-模层 $\F$，有相应的有限型 $\OO_{\frX}$-模层 $\F_{/\frX}$，它还是拓扑环层 $\OO_{\frX}$ 上的拓扑模层。

每个凝聚 $\OO_X$-模层 $\F$ 都形如 $\wt M$，其中 $M$ 为有限型 $A$-模（\hyperlink{para:1:1.5.1}{1.5.1}）。置 $(\wt M)_{/\frX}=M^\Delta$。若 $u:M\to N$ 为有限型 $A$-模的同态，则相应有 $\wt u:\wt M\to\wt N$，进而有连续同态 $(\wt u)_{/\frX}:(\wt M)_{/\frX}\to(\wt N)_{/\frX}$，记为 $u^\Delta$。显然 $(v\circ u)^\Delta=v^\Delta\circ u^\Delta$。这样定义了从有限型 $A$-模范畴到有限型 $\OO_{\frX}$-模层范畴的加性协变函子 $M\mapsto M^\Delta$。当 $A$ 为离散环时，$M^\Delta=\wt M$。

\egastatement{10.10.2}{命题}(i) $M\mapsto M^\Delta$ 是正合函子，且存在 $A$-模的典范函子同构 $\Gamma(\frX,M^\Delta)\simeq M$。

(ii) 若 $M,N$ 为两个有限型 $A$-模，则存在典范函子同构
\begin{align}
(M\otimes_A N)^\Delta&\simeq M^\Delta\otimes_{\OO_{\frX}}N^\Delta,\tag{10.10.2.1}\label{eq:1:10.10.2.1}\\
(\Hom_A(M,N))^\Delta&\simeq\shHom_{\OO_{\frX}}(M^\Delta,N^\Delta).\tag{10.10.2.2}\label{eq:1:10.10.2.2}
\end{align}
(iii) 映射 $u\mapsto u^\Delta$ 是函子同构
\begin{equation}\tag{10.10.2.3}\label{eq:1:10.10.2.3}
\Hom_A(M,N)\simeq\Hom_{\OO_{\frX}}(M^\Delta,N^\Delta).
\end{equation}

$M^\Delta$ 的正合性来自 $\wt M$（\hyperlink{para:1:1.3.5}{1.3.5}）及 $\F_{/\frX}$（\hyperlink{para:1:10.8.8}{10.8.8}）的正合性。由定义，$\Gamma(\frX,M^\Delta)$ 为 $\Gamma(X,\wt M)=M$ 的 $J$-预进制分离完备化；但由于 $A$ 完备且 $M$ 有限型，$M$ 已分离且完备（0，\hyperlink{para:0:7.3.6}{7.3.6}），从而 (i) 得证。（\hyperref[eq:1:10.10.2.1]{10.10.2.1}）来自（\hyperlink{para:1:1.3.12}{1.3.12}，(i)）与（\hyperref[eq:1:10.8.10.1]{10.8.10.1}）的复合；（\hyperref[eq:1:10.10.2.2]{10.10.2.2}）来自（\hyperlink{para:1:1.3.12}{1.3.12}，(ii)）与（\hyperref[eq:1:10.8.10.2]{10.8.10.2}）的复合。最后，$\Hom_A(M,N)$ 为有限型 $A$-模，可对它应用 (i)，将 $\Gamma(\frX,(\Hom_A(M,N))^\Delta)$ 与 $\Hom_A(M,N)$ 等同，再利用（\hyperref[eq:1:10.10.2.2]{10.10.2.2}），便知（\hyperref[eq:1:10.10.2.3]{10.10.2.3}）为同构。

由（\hyperlink{para:1:10.10.2}{10.10.2}）可得一系列与由（\hyperlink{para:1:1.3.7}{1.3.7}）、（\hyperlink{para:1:1.3.12}{1.3.12}）所得结果类似的推论，我们留给读者表述。还须指出，对正合列 $0\to J\to A\to A/J\to0$ 应用正合函子 $M^\Delta$，结合（\hyperlink{para:1:10.3.2}{10.3.2}），可知此处记作 $J^\Delta$ 的 $\OO_{\frX}$-理想层，与（\hyperlink{para:1:10.3.1}{10.3.1}）中同样记号表示的理想层一致。

\egastatement{10.10.3}{命题}在（\hyperlink{para:1:10.10.1}{10.10.1}）的假设下，$\OO_{\frX}$ 为凝聚环层。

对 $f\in A$，$A_{\{f\}}$ 为诺特进制环（0，\hyperlink{para:0:7.6.11}{7.6.11}）。问题是局部的，故由（\hyperlink{para:1:10.1.4}{10.1.4}），只须证明同态 $v:\OO_{\frX}^n\to\OO_{\frX}$ 的核为有限型 $\OO_{\frX}$-模层。此时 $v=u^\Delta$，其中 $u:A^n\to A$ 为 $A$-同态（\hyperlink{para:1:10.10.2}{10.10.2}）。由于 $A$ 诺特，$\Ker u$ 有限型，故存在同态 $w:A^m\to A^n$ 使 $A^m\xrightarrow{w}A^n\xrightarrow{u}A$ 正合。由（\hyperlink{para:1:10.10.2}{10.10.2}），列 $\OO_{\frX}^m\xrightarrow{w^\Delta}\OO_{\frX}^n\xrightarrow{v}\OO_{\frX}$ 正合，因而 $\Ker v$ 有限型。

\egapar{10.10.4}沿用上述记号，置 $A_n=A/J^{n+1}$，并令
\[
X_n=\Spec(A_n)=(\frX,\OO_{\frX}/\J^{n+1}),\qquad\J=J^\Delta.
\]
对于 $m\le n$，以 $u_{mn}:X_m\to X_n$ 表示对应于典范同态 $A_n\to A_m$ 的预概形态射。形式概形 $\frX$ 是各 $X_n$ 关于 $u_{mn}$ 的归纳极限（\hyperlink{para:1:10.6.3}{10.6.3}）。

\egastatement{10.10.5}{命题}在（\hyperlink{para:1:10.10.1}{10.10.1}）的假设下，设 $\F$ 为 $\OO_{\frX}$-模层。下列条件等价：
\begin{enumerate}[label=\alph*)]
\item $\F$ 凝聚；
\item $\F$ 同构于一列凝聚 $\OO_{X_n}$-模层 $\F_n$ 的射影极限（\hyperlink{para:1:10.6.6}{10.6.6}），且 $u_{mn}^*(\F_n)=\F_m$；
\item 存在有限型 $A$-模 $M$，使 $\F\simeq M^\Delta$；$M$ 由（\hyperlink{para:1:10.10.2}{10.10.2}，(i)）确定到典范同构。
\end{enumerate}

先证 b) 蕴含 c)。有 $\F_n=\wt{M_n}$，其中 $M_n$ 为有限型 $A_n$-模。假设给出 $M_m=M_n\otimes_{A_n}A_m$（$m\le n$）（\hyperlink{para:1:1.6.5}{1.6.5}）。于是各 $M_n$ 关于典范双同态 $M_n\to M_m$ 构成射影系；由 $A_n$ 的定义，它满足（0，\hyperlink{para:0:7.2.9}{7.2.9}）的条件。因此 $M=\plim M_n$ 为有限型 $A$-模，且对每个 $n$，$M_n=M\otimes_A A_n$。故 $\F_n$ 由 $\wt M\otimes_{\OO_X}(\OO_X/\wt{J^{n+1}})$ 在 $X_n$ 上诱导，从而按定义（\hyperlink{para:1:10.8.4}{10.8.4}）有 $\F=M^\Delta$。

反之，c) 蕴含 b)。若 $u_n:X_n\to X$ 为浸入态射，则 $u_n^*(\wt M)=(M\otimes_A A_n)^\sim$ 由 $\wt M\otimes_{\OO_X}(\OO_X/\wt{J^{n+1}})$ 在 $X_n$ 上诱导，并由定义 $M^\Delta=\plim u_n^*(\wt M)$。由于 $u_m=u_n\circ u_{mn}$，层 $\F_n=u_n^*(\wt M)$ 满足 b)。

再证 c) 蕴含 a)。由定义 $\OO_{\frX}=A^\Delta$；由于 $M$ 为某个同态 $A^m\to A^n$ 的余核，据（\hyperlink{para:1:10.10.2}{10.10.2}），$M^\Delta$ 为相应同态 $\OO_{\frX}^m\to\OO_{\frX}^n$ 的余核。环层 $\OO_{\frX}$ 凝聚（\hyperlink{para:1:10.10.3}{10.10.3}），所以 $M^\Delta$ 也凝聚（0，\hyperlink{para:0:5.3.4}{5.3.4}）。

最后证 a) 蕴含 b)。视为 $\OO_{\frX}$-模层，有 $\OO_{X_n}=\OO_{\frX}/\J^{n+1}=A_n^\Delta$。令 $\F_n=\F\otimes_{\OO_{\frX}}\OO_{X_n}$，则它作为 $\OO_{\frX}$-模层凝聚（0，\hyperlink{para:0:5.3.5}{5.3.5}）；又由于 $\J^{n+1}$ 凝聚，它作为 $\OO_{X_n}$-模层也凝聚（0，\hyperlink{para:0:5.3.10}{5.3.10}）。记住 $X_m\to X_n$ 的底连续映射为恒等映射，立即可验 $u_{mn}^*(\F_n)=\F_m$。因此 $\G=\plim\F_n$ 为凝聚 $\OO_{\frX}$-模层，因为已知 b) 蕴含 a)。各典范同态 $\F\to\F_n$ 组成射影系，给出典范同态 $w:\F\to\G$。只须证明 $w$ 双射。

这一问题现在是局部的，可设 $\F$ 为同态 $\OO_{\frX}^m\to\OO_{\frX}^n$ 的余核。由（\hyperlink{para:1:10.10.2}{10.10.2}），此同态形如 $v^\Delta$，其中 $v:A^m\to A^n$，故 $\F\simeq M^\Delta$，$M=\Coker v$。于是
\[
\F_n=M^\Delta\otimes_{\OO_{\frX}}A_n^\Delta=(M\otimes_A A_n)^\Delta.
\]
$M\otimes_A A_n$ 的 $J$-进制拓扑为离散拓扑，所以作为 $\OO_{X_n}$-模层，$(M\otimes_A A_n)^\Delta=(M\otimes_A A_n)^\sim$。前面已知 $M^\Delta=\plim\F_n$，此时 $w$ 就是恒等同态。证毕。

\egastatement{10.10.6}{推论}若 $\F$ 满足（\hyperlink{para:1:10.10.5}{10.10.5}）的条件 b)，则射影系 $(\F_n)$ 同构于系 $(\F\otimes_{\OO_{\frX}}\OO_{X_n})$。

\egapar{10.10.7}现设 $A,B$ 为两个诺特进制环，$\varphi:B\to A$ 为连续同态。以 $J,K$ 分别表示 $A,B$ 的定义理想，且 $\varphi(K)\subset J$。置 $X=\Spec(A)$、$Y=\Spec(B)$、$\frX=\Spf(A)$、$\frY=\Spf(B)$。设 $f:X\to Y$ 为对应于 $\varphi$ 的预概形态射（\hyperlink{para:1:1.6.1}{1.6.1}），$\widehat f:\frX\to\frY$ 为其到完备化的延拓（\hyperlink{para:1:10.9.1}{10.9.1}）；后者也是对应于 $\varphi$ 的形式预概形态射（\hyperlink{para:1:10.2.2}{10.2.2}）。

\egastatement{10.10.8}{命题}对于每个有限型 $B$-模 $N$，存在 $\OO_{\frX}$-模层的典范函子同构
\[
\widehat f^{\,*}(N^\Delta)\simeq(N\otimes_B A)^\Delta.
\]

以 $i_X:\frX\to X$、$i_Y:\frY\to Y$ 表示典范态射，则在典范函子同构意义下，由（\hyperlink{para:1:10.8.8}{10.8.8}）及（\hyperlink{para:1:1.6.5}{1.6.5}）有
\[
N^\Delta=i_Y^*(\wt N),\qquad
(N\otimes_B A)^\Delta=i_X^*((N\otimes_B A)^\sim)=i_X^*(f^*(\wt N)).
\]
命题因而由（\hyperlink{para:1:10.9.2}{10.9.2}）的交换图得出。

\egastatement{10.10.9}{推论}对 $B$ 的每个理想 $\mathfrak b$，有
\[
\widehat f^{\,*}(\mathfrak b^\Delta)\OO_{\frX}=(\mathfrak b A)^\Delta.
\]

设 $j:\mathfrak b\hookrightarrow B$ 为典范嵌入，它对应于 $\OO_{\frY}$-模层的典范嵌入 $j^\Delta:\mathfrak b^\Delta\to\OO_{\frY}$。由定义，左端为同态 $\widehat f^{\,*}(j^\Delta):\widehat f^{\,*}(\mathfrak b^\Delta)\to\OO_{\frX}=\widehat f^{\,*}(\OO_{\frY})$ 的像。由（\hyperlink{para:1:10.10.8}{10.10.8}），这个同态与 $(j\otimes1)^\Delta:(\mathfrak b\otimes_B A)^\Delta\to(B\otimes_B A)^\Delta$ 等同。由于 $j\otimes1$ 的像是理想 $\mathfrak bA$，据（\hyperlink{para:1:10.10.2}{10.10.2}），其相伴层同态的像为 $(\mathfrak bA)^\Delta$，结论得证。


% END chapters/ch1_s10_10.tex

\par

% BEGIN chapters/ch1_s10_11.tex
\subsection{形式预概形上的凝聚层}\hypertarget{heading:1:10.11}{}
\src{204—205}
\egastatement{10.11.1}{命题}若 $\frX$ 为局部诺特形式预概形，则环层 $\OO_{\frX}$ 凝聚，且 $\frX$ 的每个定义理想层都凝聚。

问题是局部的，故归结到诺特仿射形式概形的情形；此时命题由（\hyperlink{para:1:10.10.3}{10.10.3}）及（\hyperlink{para:1:10.10.5}{10.10.5}）得出。

\egapar{10.11.2}设 $\frX$ 为局部诺特形式预概形，$\J$ 为其定义理想层，$X_n=(\frX,\OO_{\frX}/\J^{n+1})$ 为相应的通常局部诺特预概形。因此 $\frX$ 为序列 $(X_n)$ 关于典范态射 $u_{mn}:X_m\to X_n$ 的归纳极限（\hyperlink{para:1:10.6.3}{10.6.3}）。沿用这些记号，有：

\egastatement{10.11.3}{定理}$\OO_{\frX}$-模层 $\F$ 凝聚的充要条件是，它同构于一列凝聚 $\OO_{X_n}$-模层 $\F_n$ 的射影极限，且对 $m\le n$ 有 $u_{mn}^*(\F_n)=\F_m$（\hyperlink{para:1:10.6.6}{10.6.6}）。此时射影系 $(\F_n)$ 同构于系
\[
u_n^*(\F)=\F\otimes_{\OO_{\frX}}\OO_{X_n},
\]
其中 $u_n:X_n\to\frX$ 为典范态射。

问题是局部的，故归结到 $\frX$ 为诺特仿射形式概形的情形，此时定理是（\hyperlink{para:1:10.10.5}{10.10.5}）及（\hyperlink{para:1:10.10.6}{10.10.6}）的结果。因此可以说，\emph{给定一个凝聚 $\OO_{\frX}$-模层，等价于给定一个满足 $u_{mn}^*(\F_n)=\F_m$（$m\le n$）的凝聚 $\OO_{X_n}$-模层射影系。}

\egastatement{10.11.4}{推论}若 $\F,\G$ 为两个凝聚 $\OO_{\frX}$-模层，则沿用（\hyperlink{para:1:10.11.3}{10.11.3}）的记号，可以定义典范函子同构
\begin{equation}\tag{10.11.4.1}\label{eq:1:10.11.4.1}
\Hom_{\OO_{\frX}}(\F,\G)\simeq\plim_n\Hom_{\OO_{X_n}}(\F_n,\G_n).
\end{equation}
右端的射影极限是关于映射 $\theta_n\mapsto u_{mn}^*(\theta_n)$（$m\le n$）而言。同态（\hyperref[eq:1:10.11.4.1]{10.11.4.1}）使 $\theta\in\Hom_{\OO_{\frX}}(\F,\G)$ 对应于序列 $(u_n^*(\theta))$。由（\hyperlink{para:1:10.11.3}{10.11.3}），将一个相容系 $(\theta_n)$ 映至它在 $\Hom_{\OO_{\frX}}(\F,\G)$ 中的射影极限，立即定义了逆同态。

\egastatement{10.11.5}{推论}同态 $\theta:\F\to\G$ 为满射的充要条件是相应的 $\theta_0=u_0^*(\theta):\F_0\to\G_0$ 为满射。

问题是局部的，可设 $\frX=\Spf(A)$、$A$ 为诺特进制环，$\F=M^\Delta$、$\G=N^\Delta$、$\theta=u^\Delta$，其中 $M,N$ 为有限型 $A$-模，$u:M\to N$ 为同态。此时 $\theta_0=\wt{u_0}$，其中 $u_0=u\otimes1:M\otimes_A A/J\to N\otimes_A A/J$。由（\hyperlink{para:1:1.3.9}{1.3.9}）及（\hyperlink{para:1:10.10.2}{10.10.2}），$\theta,u$ 同时满射，$\theta_0,u_0$ 同时满射；而 $u,u_0$ 同时满射（0，\hyperlink{para:0:7.1.14}{7.1.14}），故得结论。

\egapar{10.11.6}定理（\hyperlink{para:1:10.11.3}{10.11.3}）说明，将每个凝聚 $\OO_{\frX}$-模层 $\F$ 视为各伪离散群层 $\F_n$ 的射影极限（0，\hyperlink{para:0:3.8.1}{3.8.1}），可将它视为拓扑 $\OO_{\frX}$-模层。于是由（\hyperlink{para:1:10.11.4}{10.11.4}），任意凝聚 $\OO_{\frX}$-模层的同态 $u:\F\to\G$ 都自动连续（0，\hyperlink{para:0:3.8.2}{3.8.2}）。此外，若 $\HH$ 为凝聚 $\OO_{\frX}$-模层 $\F$ 的凝聚子模层，则对于每个开集 $U\subset\frX$，$\Gamma(U,\HH)$ 是拓扑群 $\Gamma(U,\F)$ 的闭子群。事实上，由 $\Gamma$ 的左正合性，它是 $\Gamma(U,\F)\to\Gamma(U,\F/\HH)$ 的核；$\F/\HH$ 凝聚（0，\hyperlink{para:0:5.3.4}{5.3.4}），所以上述同态连续，而 $\Gamma(U,\F/\HH)$ 为分离拓扑群，即得断言。

\egastatement{10.11.7}{命题}设 $\F,\G$ 为两个凝聚 $\OO_{\frX}$-模层。沿用（\hyperlink{para:1:10.11.3}{10.11.3}）的记号，可以定义拓扑 $\OO_{\frX}$-模层的典范函子同构
\begin{align}
\F\otimes_{\OO_{\frX}}\G&\simeq\plim_n(\F_n\otimes_{\OO_{X_n}}\G_n),\tag{10.11.7.1}\label{eq:1:10.11.7.1}\\
\shHom_{\OO_{\frX}}(\F,\G)&\simeq\plim_n\shHom_{\OO_{X_n}}(\F_n,\G_n).\tag{10.11.7.2}\label{eq:1:10.11.7.2}
\end{align}
这里拓扑如（\hyperlink{para:1:10.11.6}{10.11.6}）所定义。

（\hyperref[eq:1:10.11.7.1]{10.11.7.1}）来自公式
\[
\F_n\otimes_{\OO_{X_n}}\G_n
=(\F\otimes_{\OO_{\frX}}\OO_{X_n})\otimes_{\OO_{X_n}}
(\G\otimes_{\OO_{\frX}}\OO_{X_n})
=(\F\otimes_{\OO_{\frX}}\G)\otimes_{\OO_{\frX}}\OO_{X_n}
\]
及（\hyperlink{para:1:10.11.3}{10.11.3}）。将（\hyperref[eq:1:10.11.7.2]{10.11.7.2}）的两端暂视为不带拓扑的模层，则其同构来自相应同态层的截面的定义，以及对任意诺特仿射形式开集应用（\hyperref[eq:1:10.11.4.1]{10.11.4.1}）。尚须证明它在拟紧开集上的截面同态双连续，因而可归结到 $\frX=\Spf(A)$、$A$ 为诺特进制环。此时 $\F=M^\Delta$、$\G=N^\Delta$，$M,N$ 为有限型 $A$-模（\hyperlink{para:1:10.10.5}{10.10.5}）。结合（\hyperlink{para:1:10.10.2}{10.10.2}）及（\hyperlink{para:1:1.3.12}{1.3.12}，(ii)），只须证明典范同构
\[
\begin{gathered}
\Hom_A(M,N)\simeq\plim_n\Hom_{A_n}(M_n,N_n),\\
M_n=M\otimes_A A_n,\quad N_n=N\otimes_A A_n,
\end{gathered}
\]
是拓扑同构；这已在（0，\hyperlink{para:0:7.8.2}{7.8.2}）中证明。

\egapar{10.11.8}由于 $\Hom_{\OO_{\frX}}(\F,\G)$ 是拓扑群层 $\shHom_{\OO_{\frX}}(\F,\G)$ 的截面群，它自然带有群拓扑。若 $\frX$ 诺特，由（\hyperref[eq:1:10.11.7.2]{10.11.7.2}），各子群 $\Hom_{\OO_{\frX}}(\F,\J^n\G)$ 构成该群的零点邻域基本系。

\egastatement{10.11.9}{命题}设 $\frX$ 为诺特形式预概形，$\F,\G$ 为两个凝聚 $\OO_{\frX}$-模层。在拓扑群 $\Hom_{\OO_{\frX}}(\F,\G)$ 中，满同态（分别为单同态、双射同态）组成一个开子集。

由（\hyperlink{para:1:10.11.5}{10.11.5}），满同态的集合是离散群 $\Hom_{\OO_{X_0}}(\F_0,\G_0)$ 的一个子集在连续映射 $\Hom_{\OO_{\frX}}(\F,\G)\to\Hom_{\OO_{X_0}}(\F_0,\G_0)$ 下的逆像，故第一项成立。为证第二项，用有限个诺特仿射形式开集 $U_i$ 覆盖 $\frX$。$\theta\in\Hom_{\OO_{\frX}}(\F,\G)$ 为单射，当且仅当它在各连续限制映射
\[
\Hom_{\OO_{\frX}}(\F,\G)\longrightarrow
\Hom_{\OO_{\frX}|U_i}(\F|U_i,\G|U_i)
\]
下的像均为单射。于是归结到仿射情形，而这一情形已在（0，\hyperlink{para:0:7.8.3}{7.8.3}）中证明。


% END chapters/ch1_s10_11.tex

\par

% BEGIN chapters/ch1_s10_12.tex
\subsection{形式预概形的进制态射}\hypertarget{heading:1:10.12}{}
\src{206—207}
\egapar{10.12.1}设 $\frX,\frS$ 为两个局部诺特形式预概形。称态射 $f:\frX\to\frS$ 为\emph{进制态射}，如果存在 $\frS$ 的一个定义理想 $\J$，使 $\mathscr K=f^*(\J)\OO_{\frX}$ 为 $\frX$ 的定义理想；也称 $\frX$ 为一个\emph{进制 $\frS$-预概形}（关于 $f$）。此时，对 $\frS$ 的每个定义理想 $\J_1$，$\mathscr K_1=f^*(\J_1)\OO_{\frX}$ 都是 $\frX$ 的定义理想。

事实上，问题是局部的，可设 $\frX,\frS$ 仿射且诺特；于是存在 $n$ 使 $\J^n\subset\J_1$、$\J_1^n\subset\J$（\hyperlink{para:1:10.3.6}{10.3.6}及0，\hyperlink{para:1:7.1.4}{7.1.4}），故 $\mathscr K^n\subset\mathscr K_1$、$\mathscr K_1^n\subset\mathscr K$。第一个关系说明 $\mathscr K_1=K_1^\Delta$，其中 $K_1$ 为 $A=\Gamma(\frX,\OO_{\frX})$ 的开理想；第二个关系说明 $K_1$ 为 $A$ 的定义理想（0，\hyperlink{para:0:7.1.4}{7.1.4}），断言得证。

由此立即可知，若 $\frX,\frY$ 为两个进制 $\frS$-预概形，则每个 $\frS$-态射 $u:\frX\to\frY$ 都是进制的。事实上，若 $f:\frX\to\frS$、$g:\frY\to\frS$ 为结构态射，$\J$ 为 $\frS$ 的定义理想，则 $f=g\circ u$，所以
\[
u^*\bigl(g^*(\J)\OO_{\frY}\bigr)\OO_{\frX}=f^*(\J)\OO_{\frX}
\]
是 $\frX$ 的定义理想，而由假设 $g^*(\J)\OO_{\frY}$ 是 $\frY$ 的定义理想。

\egapar{10.12.2}以下固定一个局部诺特形式预概形 $\frS$ 及其定义理想 $\J$，并置 $S_n=(\frS,\OO_{\frS}/\J^{n+1})$。各局部诺特进制 $\frS$-预概形显然构成范畴。

称通常的局部诺特 $S_n$-预概形的归纳系 $(X_n)$ 为\emph{$(S_n)$-进制归纳系}，如果其结构态射 $f_n:X_n\to S_n$ 使得对 $m\le n$，下图交换，并将 $X_m$ 与 $X_n\times_{S_n}S_m=(X_n)_{(S_m)}$ 等同：
\begin{equation}\tag{10.12.2.1}\label{eq:1:10.12.2.1}
\begin{tikzcd}
X_n\arrow[d,"f_n"']&X_m\arrow[l]\arrow[d,"f_m"]\\
S_n&S_m\arrow[l]
\end{tikzcd}
\end{equation}
进制归纳系构成范畴：只须将态射 $(X_n)\to(Y_n)$ 定义为一个 $S_n$-态射 $u_n:X_n\to Y_n$ 的归纳系，且对 $m\le n$，$u_m$ 与 $(u_n)_{(S_m)}$ 等同。

\egastatement{10.12.3}{定理}进制 $\frS$-预概形的范畴与 $(S_n)$-进制归纳系的范畴典范等价。

这一等价如下得到。若 $\frX$ 为进制 $\frS$-预概形，$f:\frX\to\frS$ 为结构态射，则 $\mathscr K=f^*(\J)\OO_{\frX}$ 为 $\frX$ 的定义理想。使 $\frX$ 对应于 $X_n=(\frX,\OO_{\frX}/\mathscr K^{n+1})$ 的归纳系，结构态射 $f_n:X_n\to S_n$ 由 $f$ 得到（\hyperlink{para:1:10.5.6}{10.5.6}）。

先证 $(X_n)$ 是进制归纳系。若 $f=(\psi,\theta)$，则 $\psi^*(\J)\OO_{\frX}=\mathscr K$，所以对每个 $n$ 有 $\psi^*(\J^n)\OO_{\frX}=\mathscr K^n$；并由 $\psi^*$ 的正合性，对 $m\le n$ 有
\[
\mathscr K^{m+1}/\mathscr K^{n+1}
=\psi^*(\J^{m+1}/\J^{n+1})\,(\OO_{\frX}/\mathscr K^{n+1}).
\]
于是由（\hyperlink{para:1:4.4.5}{4.4.5}）得到结论。此外，立即可验，对进制 $\frS$-预概形的 $\frS$-态射 $u:\frX\to\frY$，相应有 $S_n$-态射 $u_n:X_n\to Y_n$ 的归纳系，且对 $m\le n$，$u_m=(u_n)_{(S_m)}$。

由此确实得到范畴等价，将由下一个更精确的命题得出。

\egastatement{10.12.3.1}{命题}设 $(X_n)$ 为 $S_n$-预概形归纳系，结构态射 $f_n:X_n\to S_n$ 使（\hyperref[eq:1:10.12.2.1]{10.12.2.1}）交换，并对 $m\le n$ 将 $X_m$ 与 $X_n\times_{S_n}S_m$ 等同。则 $(X_n)$ 满足（\hyperlink{para:1:10.6.3}{10.6.3}）的 b)、c) 条件。令 $\frX$ 为其归纳极限，$f:\frX\to\frS$ 为 $(f_n)$ 的归纳极限。若 $X_0$ 局部诺特，则 $\frX$ 局部诺特，且 $f$ 为进制态射。

由于在 $S_n$ 中定义 $S_m$ 的理想层幂零，由（\hyperlink{para:1:4.4.5}{4.4.5}），在 $X_n$ 中定义 $X_m$ 的理想层也幂零，故（\hyperlink{para:1:10.6.3}{10.6.3}）的条件成立。问题因而在 $\frX,\frS$ 上是局部的，可设 $\frS=\Spf(A)$、$\J=J^\Delta$，$A$ 为诺特 $J$-进制环，且 $X_n=\Spec(B_n)$。令 $A_n=A/J^{n+1}$、$J_n=J/J^{n+1}$。假设说明 $B_0$ 诺特，且 $B_m=B_n/J_n^{m+1}B_n$。因此 $B_n\to B_0$ 的核为 $K_n=J_nB_n$，而 $B_n\to B_m$ 的核为 $K_n^{m+1}$。由于 $A_1$ 诺特，$J_1$ 为有限型 $A_1$-模，故 $K_1=K_1/K_1^2$ 为有限型 $B_1$-模，因而也是有限型 $B_0=B_1/K_1$-模。由（\hyperlink{para:1:10.6.4}{10.6.4}），$\frX$ 诺特。

若 $B=\plim B_n$，则 $\frX=\Spf(B)$；令 $K=\Ker(B\to B_0)$，有 $B_n=B/K^{n+1}$。若 $\rho_n:A/J^{n+1}\to B/K^{n+1}$ 为对应于 $f_n$ 的同态，则
\[
K/K^{n+1}=(B/K^{n+1})\rho_n(J/J^{n+1}).
\]
由于对应于 $f$ 的 $\rho:A\to B$ 等于 $\plim\rho_n$，理想 $JB$ 在 $K$ 中稠密；$B$ 的每个理想都闭（0，\hyperlink{para:0:7.3.5}{7.3.5}），故 $K=JB$。若 $\mathscr K=K^\Delta$，则由（\hyperlink{para:1:10.10.9}{10.10.9}）有 $f^*(\J)\OO_{\frX}=\mathscr K$，证明完成。

\egapar{10.12.3.2}对于两个进制 $\frS$-预概形 $\frX,\frY$，上述等价给出典范双射
\[
\Hom_{\frS}(\frX,\frY)\simeq\plim_n\Hom_{S_n}(X_n,Y_n),
\]
其中射影极限关于映射 $u_n\mapsto(u_n)_{(S_m)}$（$m\le n$）而言。


% END chapters/ch1_s10_12.tex

\par

% BEGIN chapters/ch1_s10_13.tex
\subsection{有限型态射}\hypertarget{heading:1:10.13}{}
\src{207—209}
\egastatement{10.13.1}{命题}设 $\frY$ 为局部诺特形式预概形，$\mathscr K$ 为其定义理想，$f:\frX\to\frY$ 为形式预概形态射。下列条件等价：
\begin{enumerate}[label=\alph*)]
\item $\frX$ 局部诺特，$f$ 为进制态射（\hyperlink{para:1:10.12.1}{10.12.1}），且置 $\J=f^*(\mathscr K)\OO_{\frX}$ 后，由 $f$ 得到的态射 $f_0:(\frX,\OO_{\frX}/\J)\to(\frY,\OO_{\frY}/\mathscr K)$ 为有限型；
\item $\frX$ 局部诺特，且为一个 $(Y_n)$-进制归纳系 $(X_n)$ 的归纳极限，其中 $X_0\to Y_0$ 为有限型；
\item $\frY$ 的每一点都有满足下列性质的诺特仿射形式开邻域 $V$：

(Q) $f^{-1}(V)$ 为有限个诺特仿射形式开集 $U_i$ 的并，每个诺特进制环 $\Gamma(U_i,\OO_{\frX})$ 都拓扑同构于 $\Gamma(V,\OO_{\frY})$ 上某个限制形式幂级数代数（0，\hyperlink{para:0:7.5.1}{7.5.1}）对一个必为闭的理想的商。
\end{enumerate}

由（\hyperlink{para:1:10.12.3}{10.12.3}），a) 蕴含 b) 是立即的。为证 b) 蕴含 c)，由于问题在 $\frY$ 上是局部的，可设 $\frY=\Spf(B)$，$B$ 为诺特进制环，$\mathscr K=K^\Delta$。由假设 $X_0$ 在 $Y_0$ 上为有限型，故 $X_0$ 为有限个仿射开集 $U_i$ 的并，且 $X_0$ 在 $U_i$ 上诱导的仿射概形的环 $A_{i0}$ 为 $B/K$ 上的有限型代数（\hyperlink{para:1:6.3.2}{6.3.2}）。由（\hyperlink{para:1:5.1.9}{5.1.9}），$U_i$ 也是每个诺特预概形 $X_n$ 的仿射开集。若 $A_{in}$ 为相应的环，则 b) 蕴含 $A_{im}=A_{in}/K^{m+1}A_{in}$（$m\le n$）。因此 $\frX$ 在 $U_i$ 上诱导的形式预概形同构于 $\Spf(A_i)$，其中 $A_i=\plim_n A_{in}$（\hyperlink{para:1:10.6.4}{10.6.4}）。$A_i$ 为 $KA_i$-进制环，且 $A_i/KA_i\simeq A_{i0}$ 为有限型 $B/K$-代数。由（0，\hyperlink{para:0:7.5.5}{7.5.5}），$A_i$ 拓扑同构于 $B$ 上某个限制形式幂级数代数的商；这样的代数诺特（0，\hyperlink{para:0:7.5.4}{7.5.4}），所取理想必闭。

为证 c) 蕴含 a)，可限于 $\frX=\Spf(A)$ 也仿射，其中 $A$ 是诺特进制环，且拓扑同构于 $B$ 上某个限制形式幂级数代数对闭理想的商。由（0，\hyperlink{para:0:7.5.5}{7.5.5}），$A/KA$ 为有限型 $B/K$-代数，且 $J=KA$ 为 $A$ 的定义理想；所以据（\hyperlink{para:1:10.10.9}{10.10.9}），a) 的条件成立。

注意，当命题（\hyperlink{para:1:10.13.1}{10.13.1}）的条件成立时，由 c) 可知，性质 a) 对 $\frY$ 的任一定义理想 $\mathscr K$ 都成立；因而在 b) 中，所有 $f_n$ 都为有限型态射。

\egastatement{10.13.2}{推论}若（\hyperlink{para:1:10.13.1}{10.13.1}）的条件成立，则 $\frY$ 的每个诺特仿射形式开集 $V$ 都具有性质 (Q)；若 $\frY$ 诺特，则 $\frX$ 也诺特。

这立即由（\hyperlink{para:1:10.13.1}{10.13.1}）及（\hyperlink{para:1:6.3.2}{6.3.2}）得出。

\egastatement{10.13.3}{定义}当（\hyperlink{para:1:10.13.1}{10.13.1}）的等价性质 a)、b)、c) 成立时，称态射 $f$ 为\emph{有限型态射}，或称 $\frX$ 为\emph{有限型 $\frY$-形式预概形}、\emph{$\frY$ 上的有限型形式预概形}。

\egastatement{10.13.4}{推论}设 $\frX=\Spf(A)$、$\frY=\Spf(B)$ 为两个诺特仿射形式概形。$\frX$ 在 $\frY$ 上为有限型的充要条件是，诺特进制环 $A$ 同构于 $B$ 上一个限制形式幂级数代数对闭理想的商。

沿用（\hyperlink{para:1:10.13.1}{10.13.1}）的记号，若 $\frX$ 在 $\frY$ 上为有限型，则 $A/KA$ 为有限型 $B/K$-代数（\hyperlink{para:1:6.3.3}{6.3.3}），而 $KA$ 为 $A$ 的定义理想（\hyperlink{para:1:10.10.9}{10.10.9}）。所以由（0，\hyperlink{para:0:7.5.5}{7.5.5}）得出结论。

\egastatement{10.13.5}{命题}(i) 两个有限型形式预概形态射的复合仍为有限型。

(ii) 设 $\frX,\frS,\frS'$ 为三个局部诺特（分别为诺特）形式预概形，$f:\frX\to\frS$、$g:\frS'\to\frS$ 为两个态射。若 $f$ 为有限型，则 $\frX\times_{\frS}\frS'$ 局部诺特（分别为诺特），并在 $\frS'$ 上为有限型。

(iii) 设 $\frS$ 为局部诺特形式预概形，$\frX',\frY'$ 为两个局部诺特 $\frS$-形式预概形，且 $\frX'\times_{\frS}\frY'$ 局部诺特。若 $\frX,\frY$ 为局部诺特 $\frS$-形式预概形，$f:\frX\to\frX'$、$g:\frY\to\frY'$ 为有限型 $\frS$-态射，则 $\frX\times_{\frS}\frY$ 局部诺特，$f\times_{\frS}g$ 为有限型 $\frS$-态射。

(iii) 由 (i)、(ii) 通过（\hyperlink{para:1:3.5.1}{3.5.1}）的形式论证得出，故只须证明 (i)、(ii)。

设 $f:\frX\to\frY$、$g:\frY\to\frZ$ 为局部诺特形式预概形的有限型态射。若 $\LL$ 为 $\frZ$ 的定义理想，则 $\mathscr K=g^*(\LL)\OO_{\frY}$ 为 $\frY$ 的定义理想，$\J=f^*(g^*(\LL))\OO_{\frX}$ 为 $\frX$ 的定义理想。置 $X_0=(\frX,\OO_{\frX}/\J)$、$Y_0=(\frY,\OO_{\frY}/\mathscr K)$、$Z_0=(\frZ,\OO_{\frZ}/\LL)$，并令 $f_0:X_0\to Y_0$、$g_0:Y_0\to Z_0$ 对应于 $f,g$。由假设 $f_0,g_0$ 为有限型，故对应于 $g\circ f$ 的 $g_0\circ f_0$ 也为有限型（\hyperlink{para:1:6.3.4}{6.3.4}）；据（\hyperlink{para:1:10.13.1}{10.13.1}），$g\circ f$ 为有限型。

在 (ii) 的条件下，$\frS,\frX,\frS'$ 分别为局部诺特预概形序列 $(S_n),(X_n),(S_n')$ 的归纳极限，而且可设 $X_m=X_n\times_{S_n}S_m$（$m\le n$）（\hyperlink{para:1:10.13.1}{10.13.1}）。于是 $\frX\times_{\frS}\frS'$ 为各 $X_n\times_{S_n}S_n'$ 的归纳极限（\hyperlink{para:1:10.7.4}{10.7.4}），并有
\[
X_m\times_{S_m}S_m'
=(X_n\times_{S_n}S_m)\times_{S_m}S_m'
=(X_n\times_{S_n}S_n')\times_{S_n'}S_m'.
\]
由于 $X_0$ 在 $S_0$ 上为有限型，$X_0\times_{S_0}S_0'$ 局部诺特（\hyperlink{para:1:6.3.8}{6.3.8}）。由（\hyperlink{para:1:10.12.3.1}{10.12.3.1}），先知 $\frX\times_{\frS}\frS'$ 局部诺特；又因 $X_0\times_{S_0}S_0'$ 在 $S_0'$ 上为有限型（\hyperlink{para:1:6.3.8}{6.3.8}），由（\hyperlink{para:1:10.12.3.1}{10.12.3.1}）及（\hyperlink{para:1:10.13.1}{10.13.1}），此形式预概形在 $\frS'$ 上为有限型。关于诺特预概形的断言是（\hyperlink{para:1:6.3.8}{6.3.8}）的直接结果，(ii) 得证。

\egastatement{10.13.6}{推论}在（\hyperlink{para:1:10.9.9}{10.9.9}）的假设下，若 $f$ 为有限型态射，则它到完备化的延拓 $\widehat f$ 也是有限型态射。


% END chapters/ch1_s10_13.tex

\par

% BEGIN chapters/ch1_s10_14.tex
\subsection{形式预概形的闭子预概形}\hypertarget{heading:1:10.14}{}
\src{209—212}
\egastatement{10.14.1}{命题}设 $\frX$ 为局部诺特形式预概形，$\A$ 为 $\OO_{\frX}$ 的凝聚理想层。若 $\frY$ 为 $\OO_{\frX}/\A$ 的支集（它是闭集），则拓扑赋环空间 $(\frY,(\OO_{\frX}/\A)|\frY)$ 是局部诺特形式预概形；若 $\frX$ 诺特，则它也诺特。

由（\hyperlink{para:1:10.10.3}{10.10.3}）及（0，\hyperlink{para:0:5.3.4}{5.3.4}），$\OO_{\frX}/\A$ 凝聚，故其支集 $\frY$ 闭（0，\hyperlink{para:0:5.2.2}{5.2.2}）。取 $\frX$ 的定义理想 $\J$，置 $X_n=(\frX,\OO_{\frX}/\J^{n+1})$。环层 $\OO_{\frX}/\A$ 为各
\[
\OO_{\frX}/(\A+\J^{n+1})
=(\OO_{\frX}/\A)\otimes_{\OO_{\frX}}(\OO_{\frX}/\J^{n+1})
\]
的射影极限（\hyperlink{para:1:10.11.3}{10.11.3}），且这些层的支集都为 $\frY$。理想层 $(\A+\J^{n+1})/\J^{n+1}$ 是凝聚 $\OO_{\frX}$-模层，因而也是凝聚 $(\OO_{\frX}/\J^{n+1})$-模层（0，\hyperlink{para:0:5.3.10}{5.3.10}）。令 $Y_n$ 为它在 $X_n$ 中定义的闭子预概形，则 $(\frY,(\OO_{\frX}/\A)|\frY)$ 显然是各 $Y_n$ 的归纳极限。由于（\hyperlink{para:1:10.6.4}{10.6.4}）的条件成立，它局部诺特；若 $\frX$ 诺特，则 $Y_0$ 诺特（\hyperlink{para:1:6.1.4}{6.1.4}），故它也诺特。

\egastatement{10.14.2}{定义}称形如 $(\frY,(\OO_{\frX}/\A)|\frY)$ 的形式预概形为形式预概形 $\frX$ 的\emph{闭子预概形}，其中 $\A$ 为凝聚理想层；称它为\emph{由 $\A$ 定义的闭子预概形}。\transnote{原版此处及下一句写作“凝聚 $\OO_{\frX}$-模层”。根据（\hyperlink{para:1:10.14.1}{10.14.1}）及商环层 $\OO_{\frX}/\A$ 的定义，这里须指 $\OO_{\frX}$ 的凝聚理想层，而非任意凝聚模层；正文据此明确写出。}

凝聚理想层与 $\frX$ 的闭子预概形之间这样定义的对应显然是双射。拓扑赋环空间态射 $j=(\psi,\theta):\frY\to\frX$，其中 $\psi$ 为底空间嵌入、$\theta^\sharp$ 为典范商同态的相应限制，显然是形式预概形态射（\hyperlink{para:1:10.4.5}{10.4.5}），称为 $\frY$ 到 $\frX$ 的\emph{典范嵌入}。

若 $\frX=\Spf(A)$，$A$ 为诺特进制环，则 $\A=\mathfrak a^\Delta$，其中 $\mathfrak a$ 为 $A$ 的理想（\hyperlink{para:1:10.10.5}{10.10.5}）。于是从上面的讨论立即得知，在同构意义下 $\frY=\Spf(A/\mathfrak a)$，且 $j$ 对应于典范同态 $A\to A/\mathfrak a$（\hyperlink{para:1:10.2.2}{10.2.2}）。

称局部诺特形式预概形的态射 $f:\frZ\to\frX$ 为\emph{闭浸入}，如果它可分解为 $\frZ\xrightarrow{g}\frY\xrightarrow{j}\frX$，其中 $g$ 为 $\frZ$ 到 $\frX$ 的一个闭子预概形 $\frY$ 的同构，$j$ 为典范嵌入。由于 $j$ 为赋环空间的单态射，$g,\frY$ 必然唯一。

\egastatement{10.14.3}{命题}闭浸入是有限型态射。

立即归结到 $\frX=\Spf(A)$、$\frY=\Spf(A/\mathfrak a)$ 的情形，再用（\hyperlink{para:1:10.13.1}{10.13.1}，c)）即得命题。

\egastatement{10.14.4}{引理}设 $f:\frY\to\frX$ 为局部诺特形式预概形的态射，$(U_\alpha)$ 为由 $\frX$ 的诺特仿射形式开集组成的 $f(\frY)$ 的覆盖，且各 $f^{-1}(U_\alpha)$ 都是 $\frY$ 的诺特仿射形式开集。$f$ 为闭浸入的充要条件是，$f(\frY)$ 在 $\frX$ 中闭，且对每个 $\alpha$，$f$ 在 $f^{-1}(U_\alpha)$ 上的限制对应于满同态
\[
\Gamma(U_\alpha,\OO_{\frX})\longrightarrow
\Gamma(f^{-1}(U_\alpha),\OO_{\frY})
\]
（\hyperlink{para:1:10.4.6}{10.4.6}）。

必要性显然。反之，若这些条件成立，以 $\mathfrak a_\alpha$ 表示上述截面同态的核。取 $\A|U_\alpha=\mathfrak a_\alpha^\Delta$，并在 $\frX\setminus f(\frY)$ 上取 $\A=\OO_{\frX}$，便可定义凝聚理想层 $\A$。\transnote{原版写成在各 $U_\alpha$ 的并的补集中令 $\A$ 为零，并写 $\mathfrak a_\alpha^\Delta$ 的支集为 $U_\alpha\cap f(\frY)$。对于定义闭子预概形的理想层，零支集应当属于商层 $\OO_{U_\alpha}/\mathfrak a_\alpha^\Delta$；在闭像之外理想层应等于整个结构层。正文明确采用这一与局部商结构一致的表述。}
事实上，由于 $f(\frY)$ 闭，且 $\OO_{U_\alpha}/\mathfrak a_\alpha^\Delta$ 的支集为 $U_\alpha\cap f(\frY)$，只须验证 $\mathfrak a_\alpha^\Delta,\mathfrak a_\beta^\Delta$ 在任意诺特仿射形式开集 $V\subset U_\alpha\cap U_\beta$ 上诱导相同的层。

$f$ 在 $f^{-1}(U_\alpha)$ 上的限制是到 $U_\alpha$ 的闭浸入，因此 $f^{-1}(V)$ 是 $f^{-1}(U_\alpha)$ 的诺特仿射形式开集，且 $f$ 在其上的限制仍为闭浸入。若 $\mathfrak b$ 为相应满同态 $\Gamma(V,\OO_{\frX})\to\Gamma(f^{-1}(V),\OO_{\frY})$ 的核，则由（\hyperlink{para:1:10.10.2}{10.10.2}）立即可知，$\mathfrak a_\alpha^\Delta$ 在 $V$ 上诱导 $\mathfrak b^\Delta$。定义 $\A$ 后，显然 $f=j\circ g$，其中 $j:\frZ\to\frX$ 为 $\A$ 定义的闭子预概形的典范嵌入，$g:\frY\to\frZ$ 为同构。\transnote{原版此处将复合印为 $g\circ j$；按给定的定义域、值域，应为 $j\circ g$。}

\egastatement{10.14.5}{命题}(i) 若 $f:\frZ\to\frY$、$g:\frY\to\frX$ 为局部诺特形式预概形的闭浸入，则 $g\circ f$ 为闭浸入。

(ii) 设 $\frX,\frY,\frS$ 为三个局部诺特形式预概形，$f:\frX\to\frS$ 为闭浸入，$g:\frY\to\frS$ 为态射。则 $\frX\times_{\frS}\frY\to\frY$ 为闭浸入。

(iii) 设 $\frS$ 为局部诺特形式预概形，$\frX',\frY'$ 为局部诺特 $\frS$-形式预概形，且 $\frX'\times_{\frS}\frY'$ 局部诺特。若 $\frX,\frY$ 为局部诺特 $\frS$-形式预概形，$f:\frX\to\frX'$、$g:\frY\to\frY'$ 为闭浸入的 $\frS$-态射，则 $f\times_{\frS}g$ 为闭浸入。

由（\hyperlink{para:1:3.5.1}{3.5.1}），仍只须证明 (i)、(ii)。为证 (i)，可设 $\frY$ 为 $\frX$ 的闭子预概形，由凝聚理想层 $\J\subset\OO_{\frX}$ 定义；$\frZ$ 为 $\frY$ 的闭子预概形，由凝聚理想层 $\mathscr K\subset\OO_{\frY}$ 定义。若 $\psi:\frY\hookrightarrow\frX$ 为底空间嵌入，则 $\psi_*(\mathscr K)$ 为 $\psi_*(\OO_{\frY})=\OO_{\frX}/\J$ 的凝聚理想层（0，\hyperlink{para:0:5.3.12}{5.3.12}），故也是凝聚 $\OO_{\frX}$-模层（0，\hyperlink{para:0:5.3.10}{5.3.10}）。于是 $\OO_{\frX}\to(\OO_{\frX}/\J)/\psi_*(\mathscr K)$ 的核 $\mathscr K_1$ 凝聚（0，\hyperlink{para:0:5.3.4}{5.3.4}），且 $\OO_{\frX}/\mathscr K_1\simeq\psi_*(\OO_{\frY}/\mathscr K)$，这说明 $\frZ$ 同构于 $\frX$ 的闭子预概形。

为证 (ii)，立即可限于 $\frS=\Spf(A)$、$\frX=\Spf(B)$、$\frY=\Spf(C)$，其中 $A$ 为诺特 $J$-进制环，$B=A/\mathfrak a$，$C$ 为拓扑 $A$-代数且为诺特进制环。只须证明 $C\to C\ctensor_A(A/\mathfrak a)$ 为满射。由于 $A/\mathfrak a$ 为有限型 $A$-模，且带 $J$-进制拓扑，根据（0，\hyperlink{para:0:7.7.8}{7.7.8}），$C\ctensor_A(A/\mathfrak a)=C\otimes_A(A/\mathfrak a)=C/\mathfrak aC$，从而得证。

\egastatement{10.14.6}{推论}在（\hyperlink{para:1:10.14.5}{10.14.5}，(ii)）的假设下，设 $p,q$ 为投影，使下图交换：
\[
\begin{tikzcd}
\frX\arrow[d,"f"']&\frX\times_{\frS}\frY\arrow[l,"p"']\arrow[d,"q"]\\
\frS&\frY\arrow[l,"g"]
\end{tikzcd}
\]
则对每个凝聚 $\OO_{\frX}$-模层 $\F$，存在 $\OO_{\frY}$-模层的典范同构
\begin{equation}\tag{10.14.6.1}\label{eq:1:10.14.6.1}
u:g^*(f_*(\F))\simeq q_*(p^*(\F)).
\end{equation}

为定义这样的同态 $u$，由（0，\hyperlink{para:0:4.4.3}{4.4.3}），等价地只须定义
\[
f_*(\F)\longrightarrow g_*(q_*(p^*(\F)))=f_*(p_*(p^*(\F))).
\]
取此同态为 $u^\flat=f_*(\rho)$，其中 $\rho:\F\to p_*(p^*(\F))$ 为典范同态（0，\hyperlink{para:0:4.4.3}{4.4.3}）。为证 $u$ 为同构，立即归结到 $\frS,\frX,\frY$ 分别为诺特进制环 $A,B,C$ 的形式谱，且满足（\hyperlink{para:1:10.14.5}{10.14.5}，(ii)）中上述条件。此时 $\F=M^\Delta$，$M$ 为有限型 $(A/\mathfrak a)$-模（\hyperlink{para:1:10.10.5}{10.10.5}）。由（\hyperlink{para:1:10.10.8}{10.10.8}），（\hyperref[eq:1:10.14.6.1]{10.14.6.1}）的两端分别与 $(C\otimes_A M)^\Delta$ 及 $((C/\mathfrak aC)\otimes_{A/\mathfrak a}M)^\Delta$ 等同。由于
\[
(C/\mathfrak aC)\otimes_{A/\mathfrak a}M
=(C\otimes_A(A/\mathfrak a))\otimes_{A/\mathfrak a}M
\simeq C\otimes_A M,
\]
推论得证。

\egastatement{10.14.7}{推论}设 $X$ 为通常的局部诺特预概形，$Y$ 为其闭子预概形，$j:Y\to X$ 为典范嵌入，$X'$ 为 $X$ 的闭子集，$Y'=Y\cap X'$。则 $\widehat j:Y_{/Y'}\to X_{/X'}$ 为闭浸入；对每个凝聚 $\OO_Y$-模层 $\F$，有
\[
\widehat j_*(\F_{/Y'})=(j_*(\F))_{/X'}.
\]
由于 $Y'=j^{-1}(X')$，只须利用（\hyperlink{para:1:10.9.9}{10.9.9}）并应用（\hyperlink{para:1:10.14.5}{10.14.5}）、（\hyperlink{para:1:10.14.6}{10.14.6}）。


% END chapters/ch1_s10_14.tex

\par

% BEGIN chapters/ch1_s10_15.tex
\subsection{分离的形式预概形}\hypertarget{heading:1:10.15}{}
\src{212—213}
\egastatement{10.15.1}{定义}设 $\frS$ 为形式预概形，$\frX$ 为 $\frS$-形式预概形，$f:\frX\to\frS$ 为结构态射。称态射 $(1_{\frX},1_{\frX})_{\frS}:\frX\to\frX\times_{\frS}\frX$ 为\emph{对角态射}，记为 $\Delta_{\frX/\frS}$，也记为 $\Delta_f$。若底空间 $\frX$ 在 $\Delta_f$ 下的像为 $\frX\times_{\frS}\frX$ 的底空间的闭子集，则称 $\frX$ \emph{在 $\frS$ 上分离}，或称它为一个\emph{$\frS$-形式概形}，也称 $f$ 为\emph{分离态射}。若形式预概形 $\frX$ 在 $\mathbb Z$ 上分离，则称它\emph{分离}，或称它为一个\emph{形式概形}。

\egastatement{10.15.2}{命题}假设形式预概形 $\frS,\frX$ 分别为通常预概形序列 $(S_n),(X_n)$ 的归纳极限，且态射 $f:\frX\to\frS$ 为态射序列 $f_n:X_n\to S_n$ 的归纳极限。则 $f$ 分离的充要条件是 $f_0:X_0\to S_0$ 分离。

此时 $\Delta_f$ 为各 $\Delta_{X_n/S_n}$ 的归纳极限（\hyperlink{para:1:10.7.4}{10.7.4}），而 $\frX$ 的底空间在 $\Delta_{\frX/\frS}$ 下的像（分别为 $\frX\times_{\frS}\frX$ 的底空间）与 $X_0$ 在 $\Delta_{X_0/S_0}$ 下的像（分别为 $X_0\times_{S_0}X_0$ 的底空间）相同，故得结论。

\egastatement{10.15.3}{命题}以下假定所有考虑的形式预概形（分别为形式预概形态射），都是通常预概形序列（分别为通常预概形态射序列）的归纳极限。
\begin{enumerate}[label=(\roman*)]
\item 两个分离态射的复合分离；
\item 若 $f:\frX\to\frX'$、$g:\frY\to\frY'$ 为两个分离的 $\frS$-态射，则 $f\times_{\frS}g$ 分离；
\item 若 $f:\frX\to\frY$ 为分离的 $\frS$-态射，则对于任意基形式预概形扩张 $\frS'\to\frS$，$\frS'$-态射 $f_{(\frS')}$ 分离；
\item 若两个态射的复合 $g\circ f$ 分离，则 $f$ 分离。
\end{enumerate}

这里约定：若同一个形式预概形 $\frZ$ 在同一命题中出现多次，则在各处都将它视为同一个通常预概形序列 $(Z_n)$ 的归纳极限；从 $\frZ$ 出发的态射（分别为到 $\frZ$ 的态射）均视为从 $Z_n$ 出发（分别为到 $Z_n$）的通常预概形态射的归纳极限。

沿用（\hyperlink{para:1:10.15.2}{10.15.2}）的记号，有 $(g\circ f)_0=g_0\circ f_0$，以及 $(f\times_{\frS}g)_0=f_0\times_{S_0}g_0$。因此（\hyperlink{para:1:10.15.3}{10.15.3}）的各项断言立即由（\hyperlink{para:1:10.15.2}{10.15.2}）和通常预概形的相应断言（\hyperlink{para:1:5.5.1}{5.5.1}）得出。

对于同样一类形式预概形及态射，我们将与（\hyperlink{para:1:5.5.5}{5.5.5}）、（\hyperlink{para:1:5.5.9}{5.5.9}）、（\hyperlink{para:1:5.5.10}{5.5.10}）对应的命题留给读者表述，其中应将“仿射开集”换为“满足（\hyperlink{para:1:10.6.3}{10.6.3}）的条件 b) 的仿射形式开集”。类似论证也表明，每个诺特仿射形式概形都分离，这使所用术语得到说明。

\egastatement{10.15.4}{命题}设 $\frS$ 为局部诺特形式预概形，$\frX,\frY$ 为两个局部诺特 $\frS$-形式预概形，且 $\frX$ 或 $\frY$ 在 $\frS$ 上为有限型（故 $\frX\times_{\frS}\frY$ 局部诺特），并且 $\frY$ 在 $\frS$ 上分离。若 $f:\frX\to\frY$ 为 $\frS$-态射，则图态射
\[
\Gamma_f=(1_{\frX},f)_{\frS}:\frX\longrightarrow\frX\times_{\frS}\frY
\]
为闭浸入。

可设 $\frS$ 为局部诺特预概形序列 $(S_n)$ 的归纳极限，$\frX,\frY$ 分别为 $S_n$-预概形序列 $(X_n),(Y_n)$ 的归纳极限，$f$ 为 $S_n$-态射序列 $f_n:X_n\to Y_n$ 的归纳极限。于是 $\frX\times_{\frS}\frY$ 为 $(X_n\times_{S_n}Y_n)$ 的归纳极限，$\Gamma_f$ 为 $(\Gamma_{f_n})$ 的归纳极限（\hyperlink{para:1:10.7.4}{10.7.4}）。由假设，$Y_0$ 在 $S_0$ 上分离（\hyperlink{para:1:10.15.2}{10.15.2}），故 $\Gamma_{f_0}(X_0)$ 为 $X_0\times_{S_0}Y_0$ 的闭子空间。由于这些空间分别与 $\Gamma_f(\frX)$ 及 $\frX\times_{\frS}\frY$ 的底空间相同，首先得到 $\Gamma_f(\frX)$ 闭。

当 $(U,V)$ 遍历满足 $f(U)\subset V$ 的 $\frX,\frY$ 的诺特仿射形式开集对时，各 $U\times_{\frS}V$ 在 $\frX\times_{\frS}\frY$ 中覆盖 $\Gamma_f(\frX)$；若 $f_U:U\to V$ 为限制，则 $\Gamma_{f_U}:U\to U\times_{\frS}V$ 为 $\Gamma_f$ 的相应限制。若证明 $\Gamma_{f_U}$ 为闭浸入，则 $\Gamma_f$ 也为闭浸入（\hyperlink{para:1:10.14.4}{10.14.4}）。换言之，归结到 $\frS=\Spf(A)$、$\frX=\Spf(B)$、$\frY=\Spf(C)$，$A,B,C$ 为诺特进制环，且 $f$ 对应于连续 $A$-同态 $\varphi:C\to B$。此时 $\Gamma_f$ 对应于唯一连续同态 $\omega:B\ctensor_A C\to B$，它与典范同态 $B\to B\ctensor_A C$、$C\to B\ctensor_A C$ 的复合分别为恒等同态及 $\varphi$。显然 $\omega$ 为满射，故断言得证。

\egastatement{10.15.5}{推论}设 $\frS$ 为局部诺特形式预概形，$\frX$ 为有限型 $\frS$-形式预概形。$\frX$ 在 $\frS$ 上分离的充要条件是，对角态射 $\frX\to\frX\times_{\frS}\frX$ 为闭浸入。

\egastatement{10.15.6}{命题}局部诺特形式预概形的闭浸入 $j:\frY\to\frX$ 是分离态射。

沿用（\hyperlink{para:1:10.14.2}{10.14.2}）的记号，$j_0:Y_0\to X_0$ 为闭浸入，故为分离态射；只须应用（\hyperlink{para:1:10.15.2}{10.15.2}）。

\egastatement{10.15.7}{命题}设 $X$ 为通常的局部诺特预概形，$X'$ 为 $X$ 的闭子集，$\widehat X=X_{/X'}$。$\widehat X$ 分离的充要条件是 $\widehat X_{\mathrm{red}}$ 分离；$X$ 分离则足以保证这一点。

沿用（\hyperlink{para:1:10.8.5}{10.8.5}）的记号，$\widehat X$ 分离当且仅当 $X_0$ 分离（\hyperlink{para:1:10.15.2}{10.15.2}）；又因 $\widehat X_{\mathrm{red}}=(X_0)_{\mathrm{red}}$，后者等价于 $\widehat X_{\mathrm{red}}$ 分离（\hyperlink{para:1:5.5.1}{5.5.1}，(vi)）。


% END chapters/ch1_s10_15.tex

\par
